\pdfoutput=1
\documentclass[11pt,a4paper]{article}
\usepackage[T1]{fontenc}
\usepackage[utf8]{inputenc}
\usepackage{amsmath,amssymb,amsthm,mathtools}
\usepackage{bm}
\usepackage{mleftright}\mleftright
\usepackage{jheppub}
\usepackage{microtype}
\usepackage{enumitem}
\usepackage{needspace}
\usepackage{slashed}
\usepackage{simplewick}
\usepackage{tikz}
\usetikzlibrary{arrows.meta,calc,decorations.markings,decorations.pathmorphing,decorations.pathreplacing,patterns,positioning,shapes.geometric,shapes.misc}
\hypersetup{hypertexnames=false,colorlinks=true,linkcolor=black,citecolor=black,urlcolor=blue,
 pdftitle={Quantum Theory of Fields: Lecture Notes},
 pdfsubject={Lecture notes on quantum field theory with the homework problems}}

% Edition of the typed course notes "Quantum Theory of Fields" (Spring 2026),
% retyped from the printed PDFs in originals/fields-notes-qft/lectures. Section and equation numbers
% follow the notes, so references such as "Eqs. (15.3), (15.4)" in the
% homework still point at the same formulas. Every change of substance is
% listed in ../CORRECTIONS.md; ../SOURCE_MAP.md gives the page coverage.

\numberwithin{equation}{section}

% ------------------------------------------------------------ notation
\newcommand{\dd}{\mathop{}\!\mathrm{d}}
\newcommand{\xx}{\mathbf{x}}
\newcommand{\yy}{\mathbf{y}}
\newcommand{\kk}{\mathbf{k}}
\newcommand{\pp}{\mathbf{p}}
\newcommand{\qq}{\mathbf{q}}
\newcommand{\nn}{\mathbf{n}}
\newcommand{\PPv}{\mathbf{P}}
\newcommand{\R}{\mathbf{R}}
\newcommand{\Lag}{\mathcal{L}}
\newcommand{\Ham}{\mathcal{H}}
\newcommand{\Dcal}{\mathcal{D}}
\newcommand{\Tr}{\operatorname{Tr}}
\newcommand{\tr}{\operatorname{tr}}
\newcommand{\Texp}{\operatorname{T}\!\exp}
\newcommand{\Tprod}{\operatorname{T}}
\newcommand{\Const}{\text{Const}}
\newcommand{\bra}[1]{\langle #1|}
\newcommand{\ket}[1]{|#1\rangle}
\newcommand{\braket}[2]{\langle #1|#2\rangle}
\newcommand{\vev}[1]{\langle #1\rangle}
\newcommand{\half}{\tfrac12}
\newcommand{\order}{\mathcal{O}}
\renewcommand{\Re}{\operatorname{Re}}
\renewcommand{\Im}{\operatorname{Im}}

% ------------------------------------------------------- remarks, problems
\theoremstyle{definition}
\newtheorem*{remark}{Remark}
\newtheorem*{exercise}{Exercise}
\newcounter{hwproblem}
\makeatletter
\newenvironment{hwproblem}[1]
 {\par\bigskip\Needspace{6\baselineskip}%
  \gdef\hwcurrent{#1}\phantomsection\def\@currentlabel{#1}%
  \noindent\begin{minipage}{\linewidth}\hrule\smallskip
  \noindent\textbf{Homework Problem #1.}\enspace\ignorespaces}
 {\par\nopagebreak\hfill{\small\hyperref[sol:\hwcurrent]{Solution on page~\pageref*{sol:\hwcurrent}\,$\rightarrow$}}%
  \par\smallskip\hrule\end{minipage}\par\bigskip}
\makeatother
\newlist{hwparts}{enumerate}{1}
\setlist[hwparts]{label=(\alph*),leftmargin=*,topsep=3pt,itemsep=2pt}

% ---------------------------------------------------------- figures
\newcommand{\figfile}[1]{\input{figures/#1.tex}}
\tikzset{
  prop/.style={line width=0.6pt},
  vtx/.style={circle,fill=black,inner sep=0pt,minimum size=3.2pt},
  xvtx/.style={cross out,draw,line width=0.7pt,inner sep=0pt,minimum size=5pt},
  ext/.style={circle,draw,fill=white,inner sep=0pt,minimum size=4pt},
  blob/.style={circle,draw,line width=0.6pt,fill=gray!25,inner sep=0pt},
  pblob/.style={circle,draw,line width=0.6pt,pattern=north east lines,inner sep=0pt},
  wblob/.style={circle,draw,line width=0.6pt,fill=white,inner sep=0pt},
  axis/.style={-{Stealth[length=2.2mm]},line width=0.5pt},
  arr/.style={-{Stealth[length=2mm]},line width=0.6pt},
  otimesv/.style={circle,draw,fill=white,inner sep=0pt,minimum size=7pt,path picture={\draw (path picture bounding box.north west)--(path picture bounding box.south east) (path picture bounding box.south west)--(path picture bounding box.north east);}},
  wavy/.style={decorate,decoration={snake,amplitude=1.2pt,segment length=5pt}},
  midarrow/.style={postaction={decorate},decoration={markings,mark=at position 0.55 with {\arrow{Stealth[length=2mm]}}}},
}
% ---------------------------------------------------------- diagram pieces
% Coordinates are in units of the enclosing picture; \dia scales by 0.5.
\tikzset{
  pics/eight/.style={code={%
    \draw[prop] (0,0.42) circle (0.42); \draw[prop] (0,-0.42) circle (0.42);
    \node[vtx] at (0,0) {};}},
  pics/chain3/.style={code={%
    \draw[prop] (0,0.84) circle (0.42); \draw[prop] (0,0) circle (0.42); \draw[prop] (0,-0.84) circle (0.42);
    \node[vtx] at (0,0.42) {}; \node[vtx] at (0,-0.42) {};}},
  pics/melon/.style={code={%
    \draw[prop] (0,0) ellipse (0.8 and 0.55);
    \draw[prop] (-0.8,0) .. controls (-0.3,0.22) and (0.3,0.22) .. (0.8,0);
    \draw[prop] (-0.8,0) .. controls (-0.3,-0.22) and (0.3,-0.22) .. (0.8,0);
    \node[vtx] at (-0.8,0) {}; \node[vtx] at (0.8,0) {};}},
  pics/tadline/.style={code={%
    \draw[prop] (-1.4,0) -- (1.4,0); \draw[prop] (0,0.42) circle (0.42);
    \node[vtx] at (0,0) {}; \node[vtx] at (-1.4,0) {}; \node[vtx] at (1.4,0) {};}},
  pics/sunline/.style={code={%
    \draw[prop] (-1.7,0) -- (1.7,0); \draw[prop] (0,0) ellipse (0.7 and 0.4);
    \node[vtx] at (-0.7,0) {}; \node[vtx] at (0.7,0) {};
    \node[vtx] at (-1.7,0) {}; \node[vtx] at (1.7,0) {};}},
  pics/line/.style={code={%
    \draw[prop] (-1.4,0) -- (1.4,0); \node[vtx] at (-1.4,0) {}; \node[vtx] at (1.4,0) {};}},
  % blob with legs: \pic{legblob={<radius>}{<angles>}{<leg length>}{<fill style>}}
  pics/legblob/.style n args={4}{code={%
    \foreach \a in {#2} \draw[prop] (\a:#1) -- (\a:#1+#3) node[vtx] {};
    \draw[prop,#4] (0,0) circle (#1);}},
  % proper-vertex blob with dressed legs: \pic{dressed={<radius>}{<angles>}}
  pics/dressed/.style n args={2}{code={%
    \foreach \a in {#2} {\draw[prop] (\a:#1) -- (\a:#1+1.6) node[vtx] {};
      \draw[prop,fill=white] (\a:#1+0.8) circle (0.26);}
    \draw[prop,pattern=north east lines] (0,0) circle (#1);
    \foreach \a in {#2} \node[vtx,minimum size=4pt] at (\a:#1) {};}},
  pics/xvertex/.style={code={%
    \draw[prop] (-0.5,-0.5) -- (0.5,0.5); \draw[prop] (-0.5,0.5) -- (0.5,-0.5); \node[vtx] at (0,0) {};}},
}
\newcommand{\lab}[3][]{\node[#1] at (#2) {\small$#3$};}
% inline diagram: \dia{tikz code}
\newcommand{\dia}[2][0.5]{\begin{tikzpicture}[x=#1cm,y=#1cm,baseline={([yshift=-0.55ex]0,0)}]#2\end{tikzpicture}}
% centred on its bounding box instead of on y=0
\newcommand{\diac}[2][0.5]{\vcenter{\hbox{\begin{tikzpicture}[x=#1cm,y=#1cm]#2\end{tikzpicture}}}}

% ---------------------------------------------------------- layout
\setlength{\emergencystretch}{2em}
\setlist{itemsep=2pt,parsep=0pt,topsep=4pt,partopsep=0pt}
\setlength{\jot}{5pt}
\clubpenalty=10000
\widowpenalty=10000
\displaywidowpenalty=5000
\allowdisplaybreaks[2]
\renewcommand\afterProceedingsSpace{\vspace{20pt}}
\renewcommand\afterTitleSpace{\vspace{14pt}}
\renewcommand\afterRuleSpace{\vspace{16pt}}
\renewcommand\afterEmailSpace{\vspace{24pt}}

\title{Quantum Theory of Fields\\[4pt]\large Lecture notes with homework problems}
\author{Course notes, Spring 2026}
\abstract{These are the course notes on classical and quantum field theory,
retyped from the printed lectures. The figures have been redrawn, errors in
the formulas and text have been corrected, and the nine homework problems
of the course appear at the points in the notes where the material they use
has been covered. Their solutions are in the appendix. Minkowski-space
formulas use the mostly plus metric. A list of corrections is kept with the
source files.}
\setcounter{tocdepth}{2}
\begin{document}
\maketitle
\flushbottom

% Source pp. 1-9 (Lectures1-6.pdf).
\section*{Introduction}
\addcontentsline{toc}{section}{Introduction}

The subject of Field Theory is the dynamical systems (classical or quantum)
with continuously many degrees of freedom. That is, there is some number of
degrees of freedom associated with each point of space $\xx\in\R^3$. Examples
are:
\begin{enumerate}[label=\roman*)]
\item Electromagnetic Theory
\[
\text{Degrees of Freedom:}\qquad \mathbf{E}(\xx,t),\ \mathbf{H}(\xx,t)\,.
\]
\item Gravitation (general Relativity)
\[
\text{Degrees of Freedom} = \text{Spacetime metric}\quad g_{\mu\nu}(\xx,t)\,.
\]
\item More exotic example: Dirac (electron-positron) field
\[
\text{Degrees of Freedom} = \text{Four-component (complex) Lorentz spinor}\quad \psi_a(\xx,t)\,.
\]
\end{enumerate}

Main subject of this course is Quantum Theory of Fields. The states of
quantum fields with local interaction admit interpretation in terms of
propagating and interacting particles. Thus the theory of quantum fields is
the basis of particle theory.

\section{Elements of Classical Field Theory}

We start with a brief review of basic aspects of classical field theory.
Fundamental object in classical mechanics is the \textbf{action},
\[
S=\int L\bigl(q(t),\dot q(t)\bigr)\,dt\,;
\]
the equations of motion are derived from the extremal action principle
$\delta S=0$. In field theory the coordinates are replaced by the field
degrees of freedom,
\[
q(t)=\{q^i(t)\}\to \Phi(\xx,t)\,,\qquad (i\to\xx)
\]
and the action is written as a space-time integral
\[
S=\int \Lag\bigl(\Phi(\xx,t),\dot\Phi(\xx,t),\nabla\Phi(\xx,t)\bigr)\,d^3\xx\,dt\,.
\]
Here $\Phi=\Phi(\xx,t)$ stands for a collection of fields (functions of
$\xx,t$, representing local degrees of freedom), which may include scalar
fields, components of vectors, tensors, spinors, etc. It is conventionally
assumed that $\Phi$ has finitely many individual components. The function
$\Lag$ is called the \textbf{Lagrangian density}; it depends on $\Phi(\xx,t)$
and its derivatives over $\xx$ and $t$, taken at the same space-time point
$(\xx,t)$. Two remarks are in order.

1. The above expression is suitable for imposing conditions of the special
relativity. The integration is over the 4-dimensional Minkowskian space-time,
\[
d^3\xx\,dt=d^4x\,.
\]
Recall that in relativistic theory one deals with 4-dimensional space-time
with the coordinates
\[
x^\mu=(x^0,x^1,x^2,x^3)=(t,\xx)\,,
\]
equipped with the ``pseudo-metric''
\[
ds^2=-c^2\,dt^2+(dx^1)^2+(dx^2)^2+(dx^3)^2\,.
\]
The Lorentz transformations are those linear transformations of the
coordinates $x^\mu$ which preserve $ds^2$. In what follows I will use the
units in which\footnote{These conditions are simply the convention to measure
spatial distances and energy in the units
$[\text{unit of length}]=c\times[\text{unit of time}]$,
$[\text{unit of energy}]=\hbar/[\text{unit of time}]$.}
\[
c=1\,,\qquad \hbar=1\,.
\]
Then
\[
ds^2=\eta_{\mu\nu}\,dx^\mu dx^\nu=dx^\mu\,dx_\mu\,,\qquad\text{where}\qquad x^\mu=(t,\xx)\,,\quad x_\mu=(-t,\xx)\,,
\]
and $\eta_{\mu\nu}=\operatorname{diag}(-1,+1,+1,+1)$ is the Minkowski metric.
Throughout these notes the metric is ``mostly plus'', so a timelike vector
has $v^2=v^\mu v_\mu<0$ and an on-shell momentum has $p^2=-m^2$. In this
convention $px=-p^0x^0+\pp\cdot\xx$.
To ensure relativistic invariance of the classical field theory it suffices
to choose $\Lag$ to be a Lorentz scalar.

2. In writing the above action I already accepted severe restrictions,
namely,
\[
S=\int dt\,L\,,\qquad\text{with}\qquad
L=\int d^3\xx\,\Lag\bigl(\Phi(\xx,t),\dot\Phi(\xx,t),\nabla\Phi(\xx,t)\bigr)\,,
\]
i.e.\ $L$ is a $d^3\xx$ integral of a \emph{spatial density} $\Lag$, which
depends on $\Phi$, $\dot\Phi$, and $\nabla\Phi$, all taken \emph{at the same
point} $\xx$. This means that the field degrees of freedom $\Phi(\xx)$ at
some point $\xx$ do not have a direct interaction with $\Phi(\xx')$ located
at some other point $\xx'$, with $|\xx-\xx'|$ greater than zero. Allowing
such interactions would lead to possible terms in $L$ like
\[
L_{\text{nonlocal}}=\int F\bigl(\Phi(\xx,t),\Phi(\xx',t)\bigr)\,d^3\xx\,d^3\xx'
\]
(``nonlocal interactions''). Nonlocal interactions do not go well along with
the relativistic invariance. For if $S$ includes, say, double integrals over
$\xx$ (as above), relativistic invariance of $S$ forces one to include the
terms which are nonlocal in time as well, like
\[
\int F'\bigl(\Phi(t),\Phi(t')\bigr)\,dt\,dt'\,,
\]
which may violate \emph{causality}: The state in the future affects the
dynamics now. Locality and causality are deeply connected.

The above local form,
$L=\int\Lag\bigl(\Phi(\xx,t),\dot\Phi(\xx,t),\nabla\Phi(\xx,t)\bigr)\,d^3\xx$,
describes the dynamics in which the interaction occurs only between the
degrees of freedom associated with ``infinitely close'' spatial points
(through the $\nabla\Phi$ in $L$).

Also, I did not include higher spatial derivatives of $\Phi$ as the arguments
in $\Lag$. Although such dependence would not violate locality (on a
superficial level), including such dependencies in a relativistic theory
would require to add higher time derivatives, like $\ddot\Phi$, as well. Such
dependence is not allowed in conventional Lagrangian dynamics, which requires
that the time evolution $\Phi(\xx,t)$ is completely determined by the initial
values of $\Phi(\xx)$ and $\dot\Phi(\xx)$ at $t=t_0$.

Anyhow, we start with the above local form of the action. In classical field
theory the field equations of motion follow from the action principle
\[
\delta S=0\,.
\]
Let us combine the spatial and time variable into the 4-coordinate
$x=(x^\mu)=(x^0,\xx)$, with $x^0=t$, and write $\Phi(x)$ instead of
$\Phi(\xx,t)$; the Lagrangian density then is written as
\[
\Lag\bigl(\Phi(x),\partial_\mu\Phi(x)\bigr)\,,
\qquad\text{where}\qquad
\partial_\mu\equiv\frac{\partial}{\partial x^\mu}=\left(\frac{\partial}{\partial t},\nabla\right)\,.
\]
A simple exercise in functional analysis (look up the textbook, page 15 in
PS) yields the classical field equations of motion
\begin{equation}\label{eq:EL}
\partial_\mu\left(\frac{\partial\Lag}{\partial\bigl(\partial_\mu\Phi(x)\bigr)}\right)-\frac{\partial\Lag}{\partial\Phi(x)}=0\,.
\end{equation}

\begin{remark}
Essential step in deriving \eqref{eq:EL} is integrating by parts. The total
derivative term
\[
\int d^4x\,\frac{\partial}{\partial x^\mu}\left(\delta\Phi(x)\,\frac{\partial\Lag}{\partial(\partial_\mu\Phi(x))}\right)
\]
can be reduced to the integral over the boundary (Gauss theorem), which
depends only on the variations $\delta\Phi(x)$ at the boundary; this term does
not affect the ``bulk'' equations of motion \eqref{eq:EL}, which hold at each
point $x$ of the space-time. But if the field dynamics is considered in a
space-time domain $D$, with the boundary $\partial D$, the above term
contributes to the variation of $S$ via the variations of the boundary values
$\delta\Phi(x)\,|_{x\in\partial D}$ of the fields,
\[
\int_D d^4x\,\partial_\mu(\dots)=\int_{\partial D}\delta\Phi(x)\,\frac{\partial\Lag}{\partial\,\partial_\mu\Phi(x)}\,d\Sigma_\mu(x)\,,
\]
where $d\Sigma_\mu(x)$ is the normal\footnote{``Normal'' means that
$d\Sigma_\mu(x)\,T^\mu(x)=0$ for any $T^\mu(x)$ tangent to the boundary
$\partial D$, at any point $x\in\partial D$.} element of the boundary
hypersurface $\partial D$. In this situation the total derivative term
contributes to the boundary conditions (see the homework
Problem~\ref{hw:1}). The default setup is $D=$ full infinite space-time, with
the condition that $\Phi(x)$ decays sufficiently fast at both space and time
infinities (``scattering asymptotic conditions''); in this setup the above
derivative term can be ignored.
\end{remark}

The task of the classical field theory is to find solutions of these
equations with such and such initial or boundary conditions.

One of the most important general results about classical field theory is
the relation between symmetries and conservation laws known as the

\subsection*{Noether's theorem (pp.~17--18 of PS)}

What is a symmetry? Suppose we have a family of continuous transformations of
the field variables
\[
\Phi(x)\to\Phi'(x)=F\bigl(x,\Phi(x)\bigr)
\]
such that the action $S$ does not change
\[
S[\Phi]=S[\Phi']\,.
\]
The term continuous transformations means here that they depend on some
continuous parameter(s), $F=F_s$, in such a way that
\[
F_0\bigl(x,\Phi(x)\bigr)=\Phi(x)\qquad\text{-- no transformation.}
\]
(discrete symmetries is a separate story). Then we can take an infinitesimal
transformation
\[
\Phi(x)\to\Phi'(x)=\Phi(x)+\epsilon\,E\bigl(x,\Phi(x)\bigr)\,,
\qquad\text{where}\qquad
E=\frac{d}{ds}F\Big|_{s=0}\,,
\]
with an infinitesimal $\epsilon$. Substituting in the Lagrangian density, one
finds
\[
\Lag(\Phi',\partial_\mu\Phi')=\Lag(\Phi,\partial_\mu\Phi)
+\epsilon\left[\frac{\partial\Lag}{\partial\Phi}\,E(x,\Phi)+\frac{\partial\Lag}{\partial(\partial_\mu\Phi)}\,\partial_\mu E(x,\Phi)\right]+O(\epsilon^2)\,.
\]

The invariance (i.e.\ the requirement that the action $S=\int\Lag\,d^4x$
does not change) implies
\[
\bigl[\cdots\bigr]=\partial_\mu\tilde J^\mu(\Phi,\partial_\mu\Phi,x)\,,
\]
with some $\tilde J^\mu$ which depend on our field variables in a local way
(it may also have separate dependence on $x$ if $E$ does). Transforming the
left-hand side by parts one finds
\[
\partial_\mu\left(\frac{\partial\Lag}{\partial(\partial_\mu\Phi)}\,E(x,\Phi)\right)
+\left[\frac{\partial\Lag}{\partial\Phi}-\partial_\mu\left(\frac{\partial\Lag}{\partial(\partial_\mu\Phi)}\right)\right]E(x,\Phi)=\partial_\mu\tilde J^\mu\,.
\]
Now, suppose that $\Phi=\Phi_c$, a solution of the classical field equations
of motion. Then the second term vanishes, and we obtain the \emph{continuity
equation}
\[
\partial_\mu J^\mu=0\,,
\qquad\text{where}\qquad
J^\mu(\Phi,\partial\Phi,x)=\frac{\partial\Lag}{\partial(\partial_\mu\Phi)}\,E(x,\Phi)-\tilde J^\mu(\Phi,\partial\Phi,x)\,,
\]
which holds as long as $\Phi$ solves the classical field equations.

The continuity equation then implies a conservation law. Indeed, by the Gauss
theorem
\[
\oint_\Sigma J^\mu\,d\Sigma_\mu=\int_D\partial_\mu J^\mu\,d^4x=0\,.
\]
Here $\Sigma=\partial D$ is any closed 3-dimensional hypersurface in the
4-dimensional space-time, and
$d\Sigma_\mu=\frac{1}{3!}\,\epsilon_{\mu\nu\lambda\rho}\,dx^\nu dx^\lambda dx^\rho$
is the normal element of this hypersurface. Taking the hypersurface to be a
big slab between two equal-time hyperplanes and bringing the spatial boundary
to the spatial infinity (figure~\ref{fig:slab}),
\begin{figure}[ht]
\centering\figfile{slab}
\caption{The closed hypersurface $\Sigma=\partial D$: a slab bounded by the
equal-time hyperplanes $t=t_1$ and $t=t_2$, with the side boundary taken to
spatial infinity. The charges $Q_{t_1}$ and $Q_{t_2}$ are the fluxes of $J^\mu$
through the bottom and top faces.}
\label{fig:slab}
\end{figure}
and also assuming that all fields decay at $\xx\to\infty$, one finds
\[
Q_{t_1}=Q_{t_2}\,,\qquad\text{where}\qquad Q_t=\int J^0\,d^3\xx\,,
\]
that is $Q$ is an integral of motion for the classical field equations. Let
me stress again that both the continuity equation $\partial_\mu J^\mu=0$ and
the conservation law $\frac{d}{dt}Q=0$ are satisfied only if we take
$\Phi=\Phi_c$ (in terminology of QFT they are said to hold ``on shell'').

\subsection*{Energy-Momentum}

The most common symmetry we encounter in field theory is the translational
symmetry. Assume that the Lagrangian density $\Lag(\Phi,\partial\Phi)$ has no
$x$ dependence except for that coming through the field variable $\Phi(x)$.
It means that the dynamics looks the same if we shift the space-time
coordinates
\[
x\to x'=x+a
\]
by a constant 4-vector $a$. This can be viewed as the symmetry with respect
to the following transformation of the field variables
\[
\Phi(x)\to\Phi'(x)=\Phi(x+a)\,,
\]
which leaves the action invariant. This field transformation has the
infinitesimal form
\[
\Phi'(x)=\Phi(x)+da^\mu\,\partial_\mu\Phi(x)\,,
\]
i.e.\ in this example $E(\Phi,\partial\Phi)=\partial_\mu\Phi$. When we do
this transformation the action is going to stay invariant, but the Lagrangian
density changes by the total derivative
\[
\Lag'=\Lag+da^\nu\,\partial_\nu\Lag\,,
\]
where the derivative in the last term can be written as
$\partial_\mu(\delta^\mu_\nu\,\Lag)$, and has the same meaning as
$\partial_\mu\tilde J^\mu$ in the general argument above.

Note that there are four independent symmetries associated with four
components of $a^\mu$. Correspondingly, there are four Noether's currents
\[
T^\mu{}_\nu=-\frac{\partial\Lag}{\partial(\partial_\mu\Phi)}\,\partial_\nu\Phi+\delta^\mu_\nu\,\Lag\,,
\]
which therefore satisfy the continuity equations
\[
\partial_\mu T^\mu{}_\nu=0\qquad\text{on shell.}
\]
This object is known as the \textbf{energy-momentum tensor}. It leads to four
conserved charges
\[
E=-\int T^0{}_0\,d^3\xx\,,
\qquad\text{and}\qquad
P^i=\int T^0{}_i\,d^3\xx\,,\quad\text{with}\quad i=1,2,3\,.
\]
Of these, $E$ is interpreted as the energy of the field configuration
$\Phi(\xx,t)$, and $\mathbf{P}$ is interpreted as its momentum. Therefore one
may take
\begin{align*}
T^{00}&=-T^0{}_0=\mathcal{E}&&\text{for the energy density,}\\
T^{0i}&=T^0{}_i=\mathcal{P}^i&&\text{for the momentum density.}
\end{align*}
The spatial components $T^{ij}$ constitute the (spatial) \emph{stress tensor}.
It is interpreted as usual: Consider any domain $D=D^3$ in the 3-dimensional
space (at a given time). The momentum of the field inside this domain is
\[
P^i_D=\int_D\mathcal{P}^i(\xx)\,d^3\xx\,.
\]
According to the momentum theorem (Newton's second law),
\[
\frac{d}{dt}P^i_D=F^i_D\,,\qquad\text{the net force applied to the field inside }D\,.
\]
Since the interaction is local, this force must come from the surface forces,
\[
F^i_D=-\int_{\partial D}T^{ij}(\xx)\,d\sigma_j=-\int_D\partial_jT^{ij}(\xx)\,d^3\xx\,.
\]
Thus, $dF^i=T^{ij}(\xx)\,d\sigma_j$ is the force exerted by the field on the
surface element $d\sigma_j$.

\subsection*{Klein-Gordon Field}

Let us consider a simple field theory known as the \emph{Klein-Gordon} theory
(it will be our main example for some time). It involves a single-component
scalar field $\Phi(x)=\varphi(\xx,t)$, and the action is
\[
S=\int d^4x\,\left[-\frac12\,\partial_\mu\varphi\,\partial^\mu\varphi-\frac{m^2}{2}\,\varphi^2\right]
\equiv\int dt\,d^3\xx\,\left[\frac12\,\dot\varphi^2-\frac12\,(\nabla\varphi)^2-\frac{m^2}{2}\,\varphi^2\right].
\]
This form of the action shows that the theory is Lorentz-invariant, and the
field $\varphi$ is a scalar (i.e.\ $\varphi(x)$ has definite value at each
space-time point $x$, independent of the choice of the inertial frame).
Variation yields the Klein-Gordon field equation
\[
-\partial_\mu\partial^\mu\varphi+m^2\,\varphi=\partial_t^2\varphi-\nabla^2\varphi+m^2\,\varphi=0\,.
\]
Applying the above equations one finds for the energy-momentum tensor
\[
T_{\mu\nu}=\partial_\mu\varphi\,\partial_\nu\varphi-\frac12\,\eta_{\mu\nu}\,\bigl(\partial_\lambda\varphi\,\partial^\lambda\varphi+m^2\,\varphi^2\bigr)\,,
\]
with $\eta_{\mu\nu}=\operatorname{diag}(-1,1,1,1)$, so that
$\mathcal{E}=T^{00}=-T^0{}_0$ and $\mathcal{P}^i=T^{0i}=T^0{}_i$. Thus,
\[
\mathcal{E}=\frac12\Bigl(\dot\varphi^2+(\nabla\varphi)^2+m^2\,\varphi^2\Bigr)\,,
\qquad\text{and}\qquad
\boldsymbol{\mathcal{P}}=-\dot\varphi\,\nabla\varphi\,.
\]
We see in this example that the energy-momentum tensor $T_{\mu\nu}$ is
symmetric,
\[
T_{\mu\nu}=T_{\nu\mu}\,.
\]

\begin{exercise}
Show that in any Lorentz-invariant theory involving only a scalar field
$\varphi(x)$ the energy-momentum tensor always comes out symmetric,
$T_{\mu\nu}=T_{\nu\mu}$. Check that in such theory all components of the
antisymmetric tensor
\[
M^{\mu\nu}=\int\bigl[x^\mu\,T^{\nu0}-x^\nu\,T^{\mu0}\bigr]\,d^3\xx
\]
are conserved on-shell.
\end{exercise}

The equations of motion of the Klein-Gordon theory
\[
-\partial_\mu\partial^\mu\varphi+m^2\,\varphi=\partial_t^2\varphi-\nabla^2\varphi+m^2\,\varphi=0
\]
are \emph{linear}. That is, any linear combination of solutions is again a
solution. (In general, this property defines the \emph{free field theory}).
Therefore, general solution is a superposition of ``elementary solutions''.
In the infinite space-time, one can take the plane waves
\[
e^{ipx}=e^{-ip^0x^0+i\pp\cdot\xx}\qquad\text{with}\qquad p^0=\omega_\pp:=\sqrt{\pp^2+m^2}
\]
as well as the complex-conjugate waves
\begin{equation}\label{eq:KGgeneral}
\varphi(x)=\int\frac{d^3\pp}{2\omega_\pp\,(2\pi)^3}\,\bigl(A_\pp\,e^{ipx}+A^*_\pp\,e^{-ipx}\bigr)\,.
\end{equation}
This form represents general solution of the Klein-Gordon equation, which is
bounded at the spatial infinity.

\begin{hwproblem}{1}\label{hw:1}
Consider a classical field theory in the semi-infinite space,
$x^1,x^2\in(-\infty,\infty)$ but $x^3\in[0,\infty)$, with the boundary at
$x^3=0$. Assume that there is a single scalar field $\varphi(\xx,t)$, and the
Lagrangian has the form
\[
L=\int_0^\infty dz\int d^2x_\parallel\,\left[\frac12\,\dot\varphi^2-\frac12\,\bigl(\nabla\varphi\bigr)^2-\frac{m^2}{2}\,\varphi^2\right]
-\int d^2x_\parallel\,V\bigl(\varphi_B(x_\parallel,t)\bigr)\,,
\]
with some function $V(\varphi_B)$. Here $z\equiv x^3$, and
$x_\parallel=(x^1,x^2)$ are the coordinates along the boundary $z=0$, and
$\varphi_B(x_\parallel,t)$ stands for the boundary values of $\varphi$,
\[
\varphi_B(x_\parallel,t)=\varphi(\xx,t)\big|_{z=0}\,.
\]
\begin{hwparts}
\item Find the field equations of motion, including the boundary conditions
at $z=0$.
\item Which of the quantities $E$, $P^1$, $P^2$, $P^3$ are still conserved?
\item Assuming that $V(\varphi_B)=M\varphi_B^2$ with $M>0$, find general
solutions of the field equations, bounded at $z\to\infty$.
\item Consider now $V(\varphi_B)=M\varphi_B^2$ with $M<0$. Find additional
solutions which appear in this case.
\end{hwparts}
\end{hwproblem}

% Source pp. 10-19 (Lectures1-6.pdf).
\section{Quantization}

Before going to quantum theory, let us recall the Hamiltonian formalism of
classical mechanics and see how it looks for fields.

\subsection*{Hamiltonian Formalism}

In Lagrangian dynamics with the generalized coordinates $q^i$ one passes to
the Hamiltonian formalism through introducing the conjugate momenta
\[
p_i=\frac{\partial L(q,\dot q)}{\partial\dot q^i}\,,
\]
and then excluding $\dot q$ in favor of $p$, by solving these equations. The
Hamiltonian function $H$ appears as the Legendre transform
\[
H(p,q)=\sum_i p_i\,\dot q^i-L\,.
\]
The equations of motion then take the canonical form
\[
\dot p_i=\{p_i,H\}=-\frac{\partial H}{\partial q^i}\,,\qquad
\dot q^i=\{q^i,H\}=\frac{\partial H}{\partial p_i}\,,
\]
in terms of the Poisson brackets
\[
\{f,g\}:=\sum_i\frac{\partial f}{\partial q^i}\,\frac{\partial g}{\partial p_i}-\frac{\partial f}{\partial p_i}\,\frac{\partial g}{\partial q^i}\,.
\]
In a field theory we have continuously many generalized coordinates
$\varphi(\xx)$ (I assume here a scalar field, for simplicity), labeled by the
spatial coordinates $\xx$. In this case the conjugated momenta are defined as
\[
\pi(\xx)=\frac{\delta L}{\delta\dot\varphi(\xx)}\,,
\]
with the basic Poisson brackets
\[
\{\varphi(\xx),\pi(\xx')\}=\delta(\xx-\xx')\,.
\]
In the Klein-Gordon theory
\[
\pi(\xx)=\dot\varphi(\xx)\,,
\]
and the Hamiltonian is
\begin{equation}\label{eq:HKG}
H=\int d^3\xx\,\left[\frac12\,\pi^2(\xx)+\frac12\,\bigl(\nabla\varphi\bigr)^2(\xx)+\frac{m^2}{2}\,\varphi^2(\xx)\right].
\end{equation}
Of course, it coincides with the above expression for the energy,
$E=\int T^{00}\,d^3\xx$, expressed through $\pi$ rather than $\dot\varphi$.

\subsection*{Quantum Klein-Gordon field}

We can try now to develop quantum theory of the Klein-Gordon field, following
the usual rules of canonical quantization. One starts with the correspondence
\begin{center}
\begin{tabular}{ll}
Classical Theory & Quantum Theory\\[2pt]
Phase space & Hilbert space $\Ham$\\
$p_i,\,q^i$ & Hermitian operators $\hat p_i,\hat q^i:\ \Ham\to\Ham$\\
$H(p,q)$ & $\hat H(\hat p,\hat q):\ \Ham\to\Ham$
\end{tabular}
\end{center}
where the Hermitian operators $\hat q,\hat p$ must obey the canonical
commutation relations
\[
[\hat q^i,\hat p_j]=i\,\delta^i_j\,.
\]

\begin{remark}
In QM with finitely many degrees of freedom, in the Schroedinger's
quantization one can take the space of square-integrable functions $\Psi(q)$
for $\Ham$, and define the corresponding representation
\[
\hat q^i\,\Psi(q)=q^i\,\Psi(q)\,,\qquad
\hat p_i\,\Psi(q)=-i\,\frac{\partial}{\partial q^i}\,\Psi(q)\,,
\]
and $\hat H=H(\hat p,\hat q)$. The time evolution of the state $\Psi(q)$ is
governed by the Schroedinger equation
\[
i\,\frac{\partial}{\partial t}\,\Psi(q)=\hat H\,\Psi(q)\,.
\]
\end{remark}

In principle, one can follow this prescription and introduce the operators
$\hat\varphi(\xx)$, $\hat\pi(\xx)$ and demand that they satisfy the canonical
commutators in the form
\[
[\hat\varphi(\xx),\hat\pi(\xx')]=i\,\delta(\xx-\xx')\,.
\]
Furthermore, one can take the space of functionals $\Psi[\varphi(\xx)]$ as
the Hilbert space $\Ham$; it provides a formal representation of the above
canonical commutation relations through the identifications
\[
\hat\varphi(\xx)\,\Psi[\varphi]=\varphi(\xx)\,\Psi[\varphi]\,,\qquad
\hat\pi(\xx)\,\Psi[\varphi]=-i\,\frac{\delta}{\delta\varphi(\xx)}\,\Psi[\varphi]\,.
\]
The Hamiltonian operator then is
\[
\hat H=\int d^3\xx\,\left[-\frac12\,\frac{\delta^2}{\delta\varphi(\xx)^2}+\frac12\,\bigl(\nabla\varphi\bigr)^2+\frac{m^2}{2}\,\varphi^2\right],
\]
and it remains to find all solutions of the stationary Schroedinger equation
\[
\hat H\,\Psi[\varphi]=E\,\Psi[\varphi]\,.
\]

\begin{exercise}
Show that this Schroedinger equation admits formal solution of the form
\[
\Psi_0[\varphi]=\Const\,\exp\Bigl\{\int\varphi(\xx)\,G(\xx-\xx')\,\varphi(\xx')\,d^3\xx\,d^3\xx'\Bigr\}\,.
\]
Find the kernel $G(\xx-\xx')$. (At this point, do not worry when you discover
divergent integrals). Try to look for further solutions in the form
\[
\Psi_n[\varphi]\sim\int F(\xx_1,\dots,\xx_n)\,\varphi(\xx_1)\cdots\varphi(\xx_n)\,
\exp\Bigl\{\int\varphi(\xx)\,G(\xx-\xx')\,\varphi(\xx')\,d^3\xx\,d^3\xx'\Bigr\}\,.
\]
\end{exercise}

Although this straightforward approach, with certain refinements, would work
in the Klein-Gordon theory, it has a number of unpleasant features. One is
that the scalar product in $\Ham$ has to be defined through the functional
integral
\[
(\Psi_1,\Psi_2)\sim\int\Psi_1^*[\varphi]\,\Psi_2[\varphi]\,D[\varphi]\,,
\]
and certain care must be taken to give it a satisfactory definition (we are
going to use functional integrals, in somewhat different context, though).
This and other features of the above representation in terms of the
variational derivatives makes this approach somewhat cumbersome. Much shorter
route is based on explicit separation of variables, as we discuss below.

One observes that the Hamiltonian functional $H[\varphi]$ of the
Klein-Gordon theory is quadratic in the field variables $\varphi(\xx)$ and
their conjugate momenta $\pi(\xx)$; this means that the theory can be
understood as the \emph{collection of harmonic oscillators} (this observation
actually suggests the Gaussian form of $\Psi_0[\varphi]$ in the above
Exercise). It is then straightforward to make linear transformation of the
canonical variables $\varphi(\xx),\pi(\xx)$ which separates the variables into
a system of non-interacting harmonic oscillators, and then quantize the
individual oscillators. This transformation is described in detail in PS,
and I refer to \textbf{Sect.~2.3} there for the detailed calculation.

Instead, here I will follow equivalent routine based on special symmetries of
the KG theory. Recall that the Lagrangian density is
\[
\Lag=-\frac12\,\partial_\mu\varphi\,\partial^\mu\varphi-\frac{m^2}{2}\,\varphi^2\,,
\]
and observe that the transformation of the field
\[
\varphi(x)\to\varphi'(x)=\varphi(x)+f(x)\,,
\]
where $f(x)$ is an arbitrary solution of the KG equation
\[
\partial_\mu\partial^\mu f(x)-m^2\,f(x)=0\,,
\]
changes the Lagrangian density by a total derivative term
\[
\Lag(\varphi',\partial\varphi')=\Lag(\varphi,\partial\varphi)+\partial_\mu\tilde J^\mu_f\,,
\qquad\text{where}\qquad
\tilde J^\mu_f(x)=-\varphi(x)\,\partial^\mu f(x)
\]
(in writing this equation I have assumed that $f$ is infinitesimal and
neglected the terms $\sim f^2$). The situation meets the conditions of the
Noether's theorem. It follows that every solution $f$ of the KG equation
generates a conserved current $J^\mu_f$, such that
\[
\partial_\mu J^\mu_f=0\qquad\text{on shell.}
\]
Explicitly
\[
J^\mu_f=\frac{\partial\Lag}{\partial(\partial_\mu\varphi)}\,f-\tilde J^\mu_f=\varphi\,\partial^\mu f-\partial^\mu\varphi\,f\,.
\]
The corresponding Integrals of Motion are (at a time $t$)
\[
A_f=\int d^3\xx\,\Bigl(\dot\varphi(\xx,t)\,f(\xx,t)-\varphi(\xx,t)\,\dot f(\xx,t)\Bigr)
=\int d^3\xx\,\Bigl(\pi(\xx)\,f(\xx)-\varphi(\xx)\,\dot f(\xx)\Bigr)\,,
\]
where $\pi(\xx)=\dot\varphi(\xx)$ is the canonical momentum conjugate to
$\varphi(\xx)$. The integrals $A_f$ are Lorentz scalars. This becomes explicit
if one rewrites their definition in covariant form
\[
A_f=\int_\Sigma d\Sigma_\mu\,\bigl(\varphi\,\partial^\mu f-\partial^\mu\varphi\,f\bigr)\,,
\]
where $\Sigma$ is any space-like 3-surface (figure~\ref{fig:spacelike});
because of the continuity equation the currents $J^\mu_f$ obey, the above
integrals do not depend on specific choice of $\Sigma$.
\begin{figure}[ht]
\centering\figfile{spacelike}
\caption{A space-like 3-surface $\Sigma$. The integral $A_f$ over $\Sigma$
does not depend on the choice of $\Sigma$.}
\label{fig:spacelike}
\end{figure}

In classical theory one straightforwardly evaluates the Poisson brackets
\[
\{A_f,A_g\}=\int d^3\xx\,\Bigl(f(\xx,t)\,\dot g(\xx,t)-\dot f(\xx,t)\,g(\xx,t)\Bigr)
\]
(this expression also can be written in the covariant form, in terms of the
integral over arbitrary $\Sigma$). In quantum theory we should replace, as
usual,
\[
\{A,B\}\to-i\,[\hat A,\hat B]\,,
\]
so that
\[
[\hat A_f,\hat A_g]=-i\int d^3\xx\,\Bigl(\dot f(\xx)\,g(\xx)-f(\xx)\,\dot g(\xx)\Bigr)\,.
\]
It is convenient to use standard plane-wave solutions of the Klein-Gordon
equation
\[
f_\pp=e^{i\omega_\pp t-i\pp\cdot\xx}\,,\qquad f^*_\pp=e^{-i\omega_\pp t+i\pp\cdot\xx}\,,
\qquad\text{where}\qquad \omega_\pp=\sqrt{\pp^2+m^2}\,.
\]
Denoting
\[
\hat A_\pp=\hat A_{f_\pp}\,,\qquad \hat A^\dagger_\pp=\hat A_{f^*_\pp}\,,
\]
we find the commutation relations
\[
[\hat A_\pp,\hat A^\dagger_{\pp'}]=(2\pi)^3\,2\omega_\pp\,\delta^{(3)}(\pp-\pp')\,;\qquad
[\hat A_\pp,\hat A_{\pp'}]=0\,;\qquad
[\hat A^\dagger_\pp,\hat A^\dagger_{\pp'}]=0\,.
\]
These are the \emph{creation and annihilation operators}. The above operators
$\hat A_\pp$ and $\hat A^\dagger_\pp$ differ in normalization from $a_\pp$ and
$a^\dagger_\pp$ used in PS (see \textbf{Sect.~2.3})
\[
\hat A_\pp=\sqrt{2\omega_\pp}\,a_\pp\,;\qquad \hat A^\dagger_\pp=\sqrt{2\omega_\pp}\,a^\dagger_\pp\,.
\]
Advantage of $A,A^\dagger$ is their Lorentz invariance.

The Hamiltonian can be expressed through the creation and annihilation
operators as
\begin{equation}\label{eq:HAA}
\hat H=\int\frac{d^3\pp}{(2\pi)^3}\,\frac14\,\Bigl(\hat A^\dagger_\pp\,\hat A_\pp+\hat A_\pp\,\hat A^\dagger_\pp\Bigr)\,,
\end{equation}
which indeed is the Hamiltonian of a collection of harmonic oscillators, one
oscillator per each wave-vector $\pp$, the associated frequency being
$\omega_\pp$. It is straightforward to check that $\hat A^\dagger_\pp$ raises
(and $\hat A_\pp$ lowers) the energy by the amount $\omega_\pp$, namely
\[
[\hat H,\hat A^\dagger_\pp]=\omega_\pp\,\hat A^\dagger_\pp\,,\qquad
[\hat H,\hat A_\pp]=-\,\omega_\pp\,\hat A_\pp\,.
\]
Similarly, for the momentum
\[
\hat{\PPv}=-\int d^3\xx\,\hat\pi(\xx)\,\nabla\hat\varphi(\xx)
\]
we have
\[
\hat{\PPv}=\int\frac{d^3\pp}{(2\pi)^3}\,\frac{\pp}{4\omega_\pp}\,\Bigl(\hat A^\dagger_\pp\,\hat A_\pp+\hat A_\pp\,\hat A^\dagger_\pp\Bigr)\,.
\]
Hence
\[
[\hat{\PPv},\hat A^\dagger_\pp]=\pp\,\hat A^\dagger_\pp\,,\qquad
[\hat{\PPv},\hat A_\pp]=-\,\pp\,\hat A_\pp\,,
\]
i.e.\ $\hat A^\dagger_\pp$ adds $\pp$ to the total momentum.

To find a physically acceptable solution of quantum KG theory it remains to
find appropriate representation of the above commutation relations in a
Hilbert space $\Ham$. Classically, the Hamiltonian \eqref{eq:HKG} is
non-negative definite. Correspondingly, under physically meaningful
quantization the quantum Hamiltonian \eqref{eq:HAA} must be bounded from
below. Recall from the theory of harmonic oscillator that in order to satisfy
this property one has to assume that $\Ham$ contains the ground state
$\ket{0}$ such that
\[
\hat A_\pp\ket{0}=0\qquad\text{for all}\quad\pp\,.
\]
Just as in the harmonic oscillator, by applying the operators
$A^\dagger_\pp$ to $\ket{0}$ one generates the complete space of normalizable
states. In QFT this is referred as the \emph{Fock space},
\[
\Ham_{\mathrm{KG}}=\text{Fock space}=\operatorname{Span}\Bigl\{\hat A^\dagger_{\pp_1}\hat A^\dagger_{\pp_2}\cdots\hat A^\dagger_{\pp_N}\ket{0}\Bigr\}\,.
\]
The state $\ket{0}$ is called \emph{vacuum state}, or \emph{Fock vacuum}. Note
that the above basis already solves the diagonalization problem for $\hat H$:
these vectors are the eigenstates of both $\hat H$ and $\hat{\PPv}$, with the
eigenvalues
\begin{align*}
E&=\omega_{\pp_1}+\omega_{\pp_2}+\ldots+\omega_{\pp_N}+E_0\,,\\
\PPv&=\pp_1+\pp_2+\ldots+\pp_N+\PPv_0\,,
\end{align*}
where $E_0$ and $\PPv_0$ are the eigenvalues associated with the vacuum state
$\ket{0}$,
\[
\hat H\ket{0}=E_0\ket{0}\,,\qquad \hat{\PPv}\ket{0}=\PPv_0\ket{0}\,.
\]
The basic vectors $\hat A^\dagger_{\pp_1}\hat A^\dagger_{\pp_2}\cdots\hat A^\dagger_{\pp_N}\ket{0}$
are interpreted as the states of $N$ identical free (Bose) particles, with
the momenta $\pp_1,\dots,\pp_N$.

The above expression for the Hamiltonian can be rewritten as
\[
\hat H=E_0+\int\frac{d^3\pp}{(2\pi)^3}\,\frac12\,\hat A^\dagger_\pp\,\hat A_\pp\,,
\qquad\text{with}\qquad
E_0=\int\frac{d^3\pp}{(2\pi)^3}\,\frac14\,[\hat A_\pp,\hat A^\dagger_\pp]\,,
\]
which of course is the sum of the zero-point energies of all the constituent
oscillators.

There are two evident problems with the last expression:

(i) It is not zero, and moreover it contains seemingly meaningless factor
$\delta^{(3)}(\mathbf{0})$ (recall that
$[\hat A_\pp,\hat A^\dagger_{\pp'}]=2\omega_\pp\,(2\pi)^3\,\delta^{(3)}(\pp-\pp')$).

(ii) The remaining integral over $\pp$ diverges at large $\pp$.

Note that these problems don't seem to be as bad for $\PPv_0$, where we have
\[
\PPv_0=\int\frac{d^3\pp}{(2\pi)^3}\,\frac{\pp}{4\omega_\pp}\,[\hat A_\pp,\hat A^\dagger_\pp]=0\,,
\]
since the contributions of $\pp$ and $-\pp$ cancel pairwise. This is good, as
we do not want the vacuum to have any net momentum. But this seems to create
another problem, if one thinks that the energy and momentum should be the
components of 4-vector $(E_0,\PPv_0)$. If $\PPv_0=0$ and $E_0\neq0$, it seems
the Lorentz invariance is broken.

To sort this part of the problem out, let us recall that we are dealing with
a system in which the degrees of freedom $\varphi(\xx)$ are attached to all
points of the space. For such system one expects that its vacuum energy
$E_0$, whatever it is, must be proportional to the volume of the space,
\[
E_0=\varepsilon\,V^{(3)}\,,
\]
where $\varepsilon$ is the vacuum energy density. This answers the question
how to interpret the $\delta^{(3)}(\mathbf 0)$ factor: Recall that
\[
(2\pi)^3\,\delta^{(3)}(\pp-\pp')=\int d^3\xx\,e^{i(\pp-\pp')\cdot\xx}\,,
\]
and so the above nonsense factor in fact is
\[
(2\pi)^3\,\delta^{(3)}(\mathbf 0)=\int d^3\xx\ \to\ V^{(3)}\,.
\]
This also solves the problem of Lorentz invariance. The energy density
$\varepsilon$ transforms as $00$ component of 4-tensor (the energy-momentum
tensor),
\[
\varepsilon=\bra{0}T^{00}\ket{0}\,.
\]
The Lorentz-invariance of the vacuum state $\ket{0}$ only requires that
\[
\bra{0}T^{\mu\nu}\ket{0}=-\varepsilon_0\,\eta^{\mu\nu}\,,\qquad
\eta^{\mu\nu}=\operatorname{diag}(-1,+1,+1,+1)\,,
\]
and we soon will have techniques to check that this equation indeed holds in
the KG theory. It is OK for a Lorentz-invariant state to have nonzero energy
density.

Still, something has to be done about the badly divergent integral for this
energy density,
\[
\varepsilon_0=\frac12\int\frac{d^3\pp}{(2\pi)^3}\,\omega_\pp\,.
\]
The integral diverges at large momenta, i.e.\ the divergent contribution comes
from small length scales. This is not too surprising, since our idea of the
field theory was to associate a degree of freedom with every point of the
space, and there are continuously many points in each finite part of it. One
could assume that perhaps some microscopic physics takes care of the problem,
replacing the above mathematically meaningless expression by the ``cutoff''
integral
\[
\varepsilon_0=\frac12\int\frac{d^3\pp}{(2\pi)^3}\,\omega_\pp\,\Phi\left(\frac{\pp^2}{\Lambda^2}\right)\,,
\]
where the ``cutoff'' factor $\Phi(\xi)$ is close to 1 when $\xi\ll1$ and
decays fast at large $\xi$, and $\Lambda$ represents the energy scale beyond
which the unknown microscopic physics becomes essential.

Whereas introducing the cutoff rids us of mathematically meaningless
expression, needless to say it does not solve the problem of vacuum energy
density, which now is finite but completely undetermined. In many practical
cases one can simply ignore the problem. Large class of physical quantities
are only sensitive to differences of energies, and in such problems the
vacuum energy is invisible. But in fact the vacuum energy is measurable, in
principle. Generally, if we want to be able to apply QFT to physics, we need
to know what happens when we couple it to gravity, and gravity is very
sensitive to the energy density. Even without gravity, suppose we can confine
our system to finite part of the space, say inside a box. Then we want to know
what forces it exerts upon the wall of that box. The last situation can
actually be realized in laboratory, and it leads to a measurable effect, the
so called Casimir effect. We are going to discuss it next, but let me make one
more remark, in order to clarify the nature of the problem.

When a divergent integral appears in otherwise sensible calculations, it
usually signals some ambiguity in the theory. Indeed, there is an ambiguity
in the definition of the KG theory. One can add a constant term to the
Lagrangian density,
\[
\Lag\to\Lag+\Const\,.
\]
This extra term is fully consistent with both relativistic symmetry and
locality. Of course this modification is irrelevant in the classical theory
(again, as long as gravity is ignored), since it does not affect the
equations of motion. But the Lagrangian density will pass this term to the
energy-momentum tensor, and finally to the Hamiltonian,
\[
\hat H\to\hat H-\Const\int d^3\xx\,.
\]
One can now ``absorb'' the vacuum energy divergence into this constant. We
write
\[
\Const=\varepsilon_0-\varepsilon\,,
\]
where $\varepsilon_0$ is the above ugly integral, and $\varepsilon$ is a
finite constant (of the dimension $[\text{Mass}]^4$). The result is that the
divergent integral disappears, but at the same time we learn that the quantum
KG theory in fact has not one but two parameters, $m^2$ and $\varepsilon$.
This extra term we added to the action is the simplest example of what we
will call the \emph{counterterms}.

% Source pp. 20-26 (Lectures1-6.pdf).
\section{Casimir Effect}

To understand possible manifestations of the vacuum energy, let us now
briefly discuss the \emph{Casimir Effect}. Originally, the Casimir effect
refers to the situation where we have two parallel conducting plates, of large
size $L\times L$, separated by a distance $a$ (figure~\ref{fig:plates}).
\begin{figure}[ht]
\centering\figfile{plates}
\caption{Two parallel conducting plates of size $L\times L$ at separation $a$.}
\label{fig:plates}
\end{figure}
Due to the quantum nature of the electromagnetic field in between the plates,
more precisely, due to the effect of the plates on the zero-point energy of
quantum electromagnetic field, the plates are attracted with the force
\[
\frac{\text{Force}}{L^2}=-\hbar c\,\frac{\pi^2}{240\,a^4}\,,
\]
where it is assumed that
\[
L\gg a\gg\text{atomic distances.}
\]

\subsection*{Casimir effect in KG theory}

Let us study the nature of this effect using our scalar field theory, the KG
theory, instead of the electrodynamics; the result is not going to be very
different. Since the real photons are massless (and to make calculations
simpler), we will assume that $m^2=0$. In macroscopic electromagnetic theory
(valid at the scales $\gg$ atomic distances) the presence of the conducting
plates is taken into account by imposing the boundary conditions
$\mathbf{E}_\parallel=0$, $\mathbf{B}_\perp=0$ at the surface of the conductor.
To mimic this situation we consider instead the KG field in the presence of
two plates with the boundary condition
\[
\varphi(\xx)=0\qquad\text{at the surface.}
\]
Let us concentrate attention on the field between the plates. In order to
take into account the boundary conditions, in this case we must use, instead
of the functions $f_\pp(x)$ and $f^*_\pp(x)$ above, the following complete set
of solutions
\[
f_{\pp_\parallel,n}(\xx,t)=\sqrt{\frac2a}\;e^{i\omega_{\pp_\parallel,n}t-i\pp_\parallel\cdot\xx_\parallel}\,
\sin\left(\frac{\pi n}{a}\,x_\perp\right)\,,
\]
and the corresponding complex-conjugated functions $f^*_{\pp_\parallel,n}$.
Here $\xx_\parallel$ and $x_\perp$ are the components of $\xx$ parallel and
perpendicular to the plates, $\xx=(\xx_\parallel,x_\perp)$, and
\[
\omega_{\pp_\parallel,n}=\sqrt{\pp_\parallel^2+\frac{\pi^2n^2}{a^2}}
\]
with positive integer $n$,
\[
n=1,2,3,\dots\,.
\]
These functions solve the KG equation with $m^2=0$ as well as the boundary
conditions
\[
f_{\pp_\parallel,n}\big|_{x_\perp=0}=0\,,\qquad\text{and}\qquad
f_{\pp_\parallel,n}\big|_{x_\perp=a}=0\,.
\]
Exactly as we did in the infinite-volume case, we now introduce the operators
\[
\hat A_{\pp_\parallel,n}=\int_0^a \dd x_\perp\int \dd^2\xx_\parallel\,
\bigl[\hat\pi\,f_{\pp_\parallel,n}-\hat\varphi\,\dot f_{\pp_\parallel,n}\bigr]\,,
\]
and similarly for $\hat A^\dagger_{\pp_\parallel,n}$, with
$f^*_{\pp_\parallel,n}$ replacing $f_{\pp_\parallel,n}$. It is not difficult
then to show that
\[
[\hat A_{\pp_\parallel,n},\hat A^\dagger_{\pp'_\parallel,n'}]
=2\omega_{\pp_\parallel,n}\,\delta_{nn'}\,(2\pi)^2\,\delta(\pp_\parallel-\pp'_\parallel)\,,
\]
and to compute the Hamiltonian of the part of the system residing between the
plates,
\[
\hat H=\sum_{n=1}^\infty\int\frac{\dd^2\pp_\parallel}{(2\pi)^2}\,\frac14\,
\Bigl(\hat A^\dagger_{\pp_\parallel,n}\hat A_{\pp_\parallel,n}+\hat A_{\pp_\parallel,n}\hat A^\dagger_{\pp_\parallel,n}\Bigr)
+a\,L^2\,(\varepsilon-\varepsilon_0)\,,
\]
where the last term comes from vacuum energy counterterm
\[
(\varepsilon-\varepsilon_0)\int_0^a \dd x_\perp\int \dd^2\xx_\parallel=a\,L^2\,(\varepsilon-\varepsilon_0)\,.
\]
We still have to demand that all the operators $\hat A_{\pp_\parallel,n}$
nullify the ground state, since otherwise $\hat H$ is unbounded from below.
Therefore the vacuum eigenvalue of this Hamiltonian is
\[
E_{\text{vac}}=\frac{L^2}{2}\sum_{n=1}^\infty\int\frac{\dd^2\pp_\parallel}{(2\pi)^2}\,
\sqrt{\pp_\parallel^2+\frac{\pi^2n^2}{a^2}}\;\Phi\left(\frac{\pp_\parallel^2+\pi^2n^2/a^2}{\Lambda^2}\right)
-a\,L^2\,\varepsilon_0+a\,L^2\,\varepsilon\,,
\]
where I have introduced again the cutoff factor $\Phi$; without it this
expression is as divergent as it was in the absence of the boundaries.

Now, the term $-a\,L^2\,\varepsilon_0$ is the divergent part of the vacuum
energy counterterm. It is
\[
-a\,L^2\,\varepsilon_0=-\frac{L^2}{2}\,a\int\frac{\dd p_\perp}{(2\pi)}\int\frac{\dd^2\pp_\parallel}{(2\pi)^2}\,
\sqrt{\pp_\parallel^2+p_\perp^2}\;\Phi\left(\frac{\pp_\parallel^2+p_\perp^2}{\Lambda^2}\right)\,,
\]
where I have separated the $\dd p_\perp$ and $\dd^2\pp_\parallel$ parts of the
integral over $\dd^3\pp=\dd p_\perp\dd^2\pp_\parallel$. After the change of
variables
\[
p_\perp=\pi\,\tau/a
\]
it takes the form
\[
-\frac{L^2}{2}\int_0^\infty \dd\tau\int\frac{\dd^2\pp_\parallel}{(2\pi)^2}\,
\sqrt{\pp_\parallel^2+\frac{\pi^2\tau^2}{a^2}}\;\Phi\left(\frac{\pp_\parallel^2+\pi^2\tau^2/a^2}{\Lambda^2}\right)\,.
\]
We see that the interesting part of the ground-state energy (rather, of the
energy per unit area of the plates) can be written as
\[
\frac{E}{L^2}-a\,\varepsilon=\int\frac{\dd^2\pp_\parallel}{(2\pi)^2}\,
\left[\sum_{n=1}^\infty F_{\pp_\parallel}(n)-\int_0^\infty F_{\pp_\parallel}(\tau)\dd\tau\right]\,,
\]
where
\[
F_{\pp_\parallel}(\tau)=\frac12\,\sqrt{\pp_\parallel^2+\pi^2\tau^2/a^2}\;\Phi\left(\frac{\pp_\parallel^2+\pi^2\tau^2/a^2}{\Lambda^2}\right)\,.
\]
The $\sum F(n)-\int F(\tau)$ expression above can be handled with the help of
the Abel-Plana summation formula (see the Appendix to this section)
\begin{equation}\label{eq:AbelPlana}
\sum_{n=1}^\infty F(n)-\int_0^\infty F(\tau)\dd\tau
=-\frac12\,F(0)+i\int_0^\infty\frac{F(it+0)-F(-it+0)}{e^{2\pi t}-1}\dd t\,,
\end{equation}
which is valid provided $F(\tau)$ is analytic in the right half-plane of
complex $\tau$. The first term
\[
-\frac12\,F_{\pp_\parallel}(0)\ \to\ -\frac14\int\frac{\dd^2\pp_\parallel}{(2\pi)^2}\,
\sqrt{\pp_\parallel^2}\;\Phi(\pp_\parallel^2/\Lambda^2)
\]
does not depend on the separation $a$; it is interpreted as the surface
energy. It diverges, again signaling ambiguity. When there is a surface, a
constant surface term can be added to the action. We ignore this term since,
being independent of $a$, it does not contribute to the force.

The interesting contribution is extracted from the second term, which involves
\[
i\int_0^\infty\frac{F_{\pp_\parallel}(it+0)-F_{\pp_\parallel}(-it+0)}{e^{2\pi t}-1}\dd t\,.
\]
The shift ``$+0$'' in the arguments in \eqref{eq:AbelPlana} is important
since the function $F_{\pp_\parallel}(\tau)$ has square-root branch points at
the imaginary axis. Note that if it was regular, the integral would vanish in
view of the obvious $\tau\to-\tau$ symmetry of $F_{\pp_\parallel}(\tau)$. In
fact, one can take advantage of this symmetry to write the integrand as
\[
i\int_0^\infty\frac{F(it+0)-F(it-0)}{e^{2\pi t}-1}\dd t
=-2\int_{|\pp_\parallel|a/\pi}^\infty\frac{\Im F(it+0)}{e^{2\pi t}-1}\dd t\,,
\]
where the last form reflects the fact that the function
$F_{\pp_\parallel}(\tau)$ has the branch cut from $i|\pp_\parallel|a/\pi$ to
$i\infty$ (and from $-i|\pp_\parallel|a/\pi$ to $-i\infty$). We have, at
$|t|>a|\pp_\parallel|/\pi$,
\[
\Im F_{\pp_\parallel}(it+0)=\frac12\,\sqrt{\frac{\pi^2t^2}{a^2}-\pp_\parallel^2}\;
\Phi\left(\frac{\pp_\parallel^2-\pi^2t^2/a^2}{\Lambda^2}\right):=\frac12\,G_{\pp_\parallel}(t)\,.
\]
Note that the factor $(e^{2\pi t}-1)^{-1}$ makes this integral fast
convergent at large $t$, even without the cutoff factor $\Phi$. With large
$\Lambda\gg a^{-1}$ the cutoff factor has no effect, and can be dropped,
\[
G_{\pp_\parallel}(t)\to\sqrt{\frac{\pi^2t^2}{a^2}-\pp_\parallel^2}\,.
\]
The interesting contribution to the vacuum energy then is
\[
-\int\frac{\dd t}{e^{2\pi t}-1}\int_{|\pp_\parallel|<\pi t/a}\frac{\dd^2\pp_\parallel}{(2\pi)^2}\,
\sqrt{\frac{\pi^2t^2}{a^2}-\pp_\parallel^2}\,,
\]
where I have interchanged the order in which the integrals are taken. The
integral over $\pp_\parallel$ is evaluated in closed form
\[
\int_0^{\pi^2t^2/a^2}\frac{\dd p^2}{4\pi}\,\sqrt{\frac{\pi^2t^2}{a^2}-p^2}=\frac1{6\pi}\left(\frac{\pi t}{a}\right)^3\,.
\]
The contribution to the energy is
\[
-\frac{\pi^2}{6\,a^3}\int_0^\infty\frac{t^3\dd t}{e^{2\pi t}-1}=-\frac{\pi^2}{6\,a^3}\,\frac{\Gamma(4)\,\zeta(4)}{(2\pi)^4}\,,
\]
where I have used
\[
\int_0^\infty\frac{x^{s-1}\dd x}{e^x-1}=\Gamma(s)\,\zeta(s)
\]
($\zeta(4)=\pi^4/90$). Thus we find
\[
\frac{E}{L^2}=a\,\varepsilon-\frac12\,F(0)-\frac{\pi^2}{1440\,a^3}\,,
\]
where $F(0)$ stands for the surface term $\int\frac{\dd^2\pp_\parallel}{(2\pi)^2}F_{\pp_\parallel}(0)$.

The contribution $a\,\varepsilon$ is not too interesting: there is the same
term in the vacuum energy density outside the space between the plates, and
hence this term does not lead to any force. The interesting term
$-\pi^2/1440\,a^3$ shows that the plates attract; the force (per unit plate
area) is
\[
\frac{\text{Force}}{L^2}=-\frac{\partial}{\partial a}\left(-\frac{\pi^2}{1440\,a^3}\right)=-\frac{\pi^2\,\hbar c}{480\,a^4}\,.
\]
Let me stress that this result is completely independent of the form of
auxiliary cutoff factor $\Phi$; this is exactly what we want from quantum
field theory. Another useful observation is that in order to obtain this
result we don't really need to send $\Lambda$ to Planck's energies or
something, it suffices to have $a^2\gg1/\Lambda^2$.

In the real case of electromagnetic field the calculations can be done in very
similar way. The resulting force is twice as large, due to the fact that
photon field has two degrees of freedom for each $\pp$, as the photon has two
polarization states.

Few words about the $a$-independent term
\[
-\frac12\,F(0)=-\frac1{16\pi}\int_0^\infty \dd\xi\,\sqrt\xi\;\Phi(\xi/\Lambda^2)\,.
\]
It does not contribute to the force, but still needs some interpretation. It
is divergent, and it is easy to design special counter-term to absorb this
divergence. But it is instructive to think about real case of electromagnetic
field and real conducting plates made of atoms, conducting electrons, and all
that. In this case the cutoff represents the limit of validity of the
macroscopic electrodynamics, and the cutoff energy $\Lambda$ in this
expression should be taken around atomic energy scale. Having sufficient
knowledge about the microscopic structure of the conductor we could in
principle calculate the cutoff factor and hence this contribution in terms of
the atomic parameters. The nice feature of the above result is in that it
predicts some attractive force between the plates which is completely
independent of all these microscopic details. This points to another area
where quantum field theory applies, when it is not that we can't know the
short-distance physics, rather we don't want to know. In this area the
quantum field theory describes universal, i.e.\ not too sensitive to
microscopic details, properties of matter.

\subsection*{Appendix to Sect.~3: Abel-Plana summation formula}

Let $F(\tau)$ be a function analytic in the right half-plane, and decaying at
$\Re\tau\to+\infty$ sufficiently fast, so that the integral
$\int_0^\infty F(\tau)\dd\tau$ converges. Start with the integral
\[
I=\int_{C_0}\frac{F(\tau)\dd\tau}{e^{2\pi i\tau}-1}\,,
\]
where the integration is over the contour $C_0$ which encloses all positive
integer points $\tau=1,2,3,\dots$, as shown in the left part of
figure~\ref{fig:abelplana}.
\begin{figure}[ht]
\centering\figfile{abelplana}
\caption{Left: the contour $C_0$ encircles the poles at $\tau=1,2,3,\dots$
counterclockwise and passes to the right of $\tau=0$. Right: the deformed
contour. The part $C_+$ above the real axis is turned onto the positive
imaginary axis, and $C_-$ onto the negative imaginary axis. The small
quarter-circles $C'_\pm$ go around $\tau=0$.}
\label{fig:abelplana}
\end{figure}
The integrand has poles at the integer points $\tau=n$, with the residues
$F(n)/2\pi i$. Therefore
\[
I=\sum_{n=1}^\infty F(n)\,.
\]
On the other hand, the integral can be split into two parts
\[
I=\int_{C_0}=\int_{C_+}+\int_{C_-}=I_++I_-\,,
\]
where $C_+$ is the part of the contour $C_0$ that lies above the real axis,
and $C_-$ is the part below the real axis. Note that the factor
\[
\frac1{e^{2\pi i\tau}-1}
\]
decays fast into the lower half-plane, hence the contour $C_-$ can be rotated
into the lower half-plane to go along the contour $\tilde C_-$,
\[
I_-=\int_{\tilde C_-}\frac{F(\tau)\dd\tau}{e^{2\pi i\tau}-1}
=-i\int_\epsilon^\infty\frac{F(-it+0)\dd t}{e^{2\pi t}-1}+\int_{C'_-}\frac{F(\tau)\dd\tau}{e^{2\pi i\tau}-1}\,,
\]
where for the straight part of $\tilde C_-$ I have changed
\[
\tau=-it\,,
\]
and $C'_-$ is the small quarter-circle below the real axis, seen in the right
part of the figure. It is easy to check that in the limit $\epsilon\to0$ this
part of the integral yields $-F(0)/4$.

On the other hand the factor $(e^{2\pi i\tau}-1)^{-1}$ does not decay into the
upper half plane, but one can write
\[
\frac1{e^{2\pi i\tau}-1}=-1-\frac1{e^{-2\pi i\tau}-1}\,,
\]
where the second term does decay into the upper half-plane, so that
\[
I_+=-\int_{C_+}F(\tau)\dd\tau-\int_{\tilde C_+}\frac{F(\tau)\dd\tau}{e^{-2\pi i\tau}-1}
=\int_\epsilon^\infty F(\tau)\dd\tau+i\int_\epsilon^\infty\frac{F(it+0)\dd t}{e^{2\pi t}-1}+\int_{C'_+}\,,
\]
where the second term in the last form represents the contribution of the
straight part of $\tilde C_+$ (with $\tau=it+0$), and the last term is the
contribution of the small quarter-circle above the real axis; again, it is
easy to check that this last contribution is $-F(0)/4$ in the limit
$\epsilon\to0$. Adding $I_+$ to $I_-$, and taking the limit $\epsilon=0$,
yields \eqref{eq:AbelPlana}.

\begin{hwproblem}{2}\label{hw:2}
Derive Casimir force generated by massless scalar field theory in
$d+1$-dimensional space-time ($d$ spatial dimensions plus time). Assume that
the plates have $d-1$ spatial dimensions, i.e.\
$\xx=(\xx_\parallel,x_\perp)$, $\xx_\parallel\in\R^{d-1}$.
\end{hwproblem}

% Source pp. 27-37 (Lectures1-6.pdf).
\section{More on Quantum Klein-Gordon Theory}

The space of states $\Ham$ of the quantum KG theory is the Fock space, defined
as the linear envelope of the basic states
\[
\ket{\pp_1,\dots,\pp_N}=\hat A^\dagger_{\pp_1}\cdots\hat A^\dagger_{\pp_N}\ket{0}\,,
\]
where the vacuum $\ket0$ satisfies
\[
\hat A_\pp\ket0=0\qquad\text{for all}\quad\pp\,.
\]
Here I have returned to the default case of the infinite Minkowski space-time,
so that the 3-momenta $\pp_i$ may take continuous unrestricted values.

The states $\ket{\pp_1,\dots,\pp_N}$ are interpreted as $N$-particle states.
Their scalar products follow from the commutation relations
$[\hat A_\pp,\hat A^\dagger_{\pp'}]=2\omega_\pp\,(2\pi)^3\,\delta^{(3)}(\pp-\pp')$.
In particular, for the one-particle states
\[
\braket{\pp}{\pp'}=2\omega_\pp\,(2\pi)^3\,\delta^{(3)}(\pp-\pp')\,.
\]
Since the operators $\hat A^\dagger_\pp$ are Lorentz invariant (only the
momentum $\pp$ transforms under the Lorentz transformations), the particles
have zero spin. Since the creation operators $\hat A^\dagger_\pp$ commute, the
$N$-particle states are symmetric with respect to permutations of the momenta
$\pp_i$: the particles are bosons.

The spectrum of the energy-momentum operator $(\hat H-E_0,\hat{\PPv})$ is the
composition of spectra of the $N$-particle states. In the space of vectors
$(E-E_0,\PPv)$ the vacuum resides at the origin of coordinates, the one
particle states form the mass-shell hyperboloid $(E-E_0)^2-\PPv^2=m^2$, the
two-particle states fill the inside of the hyperboloid
$(E-E_0)^2-\PPv^2=4m^2$, etc.\ (figure~\ref{fig:spectrum}).
\begin{figure}[ht]
\centering\figfile{spectrum}
\caption{Spectrum of $(\hat H-E_0,\hat{\PPv})$ in the quantum KG theory. The
vacuum sits at the origin and the one-particle states $\ket{\pp}$ on the mass
shell $(E-E_0)^2-\PPv^2=m^2$. The two-particle states $\ket{\pp_1\pp_2}$ fill
the region $(E-E_0)^2-\PPv^2\ge4m^2$, the three-particle states the region
$(E-E_0)^2-\PPv^2\ge9m^2$, and so on.}
\label{fig:spectrum}
\end{figure}
Correspondingly, the spectrum of the operator
$\hat M^2=(\hat H-E_0)^2-\hat{\PPv}^2$ has two discrete components, $0$ and
$m^2$, and continuous part from $4m^2$ to infinity.

\subsection*{Klein-Gordon propagator}

By their definition, the creation and annihilation operators
$\hat A^\dagger_\pp,\hat A_\pp$ are the Fourier components of the local field
operators $\hat\varphi(\xx)$ and $\hat\pi(\xx)$,
\begin{align*}
\hat\varphi(\xx)&=\int d\mu(\pp)\,\bigl[\hat A_\pp+\hat A^\dagger_{-\pp}\bigr]\,e^{i\pp\cdot\xx}\,,\\
\hat\pi(\xx)&=\int d\mu(\pp)\,i\omega_\pp\,\bigl[-\hat A_\pp+\hat A^\dagger_{-\pp}\bigr]\,e^{i\pp\cdot\xx}\,,
\end{align*}
where I have introduced special notation for the Lorentz-invariant measure on
the mass shell,
\[
d\mu(\pp)=\frac1{2\omega_\pp}\,\frac{d^3\pp}{(2\pi)^3}\,.
\]

\begin{exercise}
Check that the measure $d\mu(\pp)$ is Lorentz-invariant. An instructive way to
do that is to check the identity
\[
d\mu(\pp)=\int_{p^0>0}\frac{dp^0}{2\pi}\;2\pi\,\delta(p^\mu p_\mu-m^2)\,\frac{d^3\pp}{(2\pi)^3}\,,
\]
i.e.\ we have
\[
\int d\mu(\pp)\,F(\pp,\omega_\pp)=\int_{p^0>0}\frac{d^4p}{(2\pi)^4}\,2\pi\,\delta(p^\mu p_\mu-m^2)\,F(\pp,p^0)\,.
\]
\end{exercise}

The operators $\hat\varphi(\xx)$ and $\hat\pi(\xx)$ are the
``Schroedinger'', or the ``equal time'' field operators. As defined, they obey
the canonical commutation relations
\[
\bigl[\hat\pi(\xx),\hat\varphi(\xx')\bigr]=-i\,\delta^{(3)}(\xx-\xx')\,.
\]
As is common in quantum mechanics, it is often useful to trade the time
evolution of the states for the time evolution of the operators. This is known
as the Heisenberg picture. In our case, introduce the Heisenberg field
operator
\[
\hat\varphi(x)=\hat\varphi(\xx,t)=e^{i\hat Ht}\,\hat\varphi(\xx)\,e^{-i\hat Ht}\,,\qquad x=(\xx,t)\,.
\]
In terms of the creation and annihilation operators it is
\[
\hat\varphi(x)=\int d\mu(\pp)\,\bigl[\hat A_\pp\,f^*_\pp(x)+\hat A^\dagger_\pp\,f_\pp(x)\bigr]\,,
\]
where as before $f_\pp(x)=e^{i\omega_\pp t-i\pp\cdot\xx}$. Then
\[
\hat\varphi(\xx)=\hat\varphi(\xx,t=0)\,,\qquad \hat\pi(\xx)=\partial_t\hat\varphi(\xx,t)\big|_{t=0}\,.
\]

\subsection*{Commutators}

It is interesting and instructive to see how the canonical commutators of the
``Schroedinger'' operators extend to the Heisenberg field operators
$\hat\varphi(x)$. For this calculation, it is useful to write
\[
\hat\varphi(x)=\hat\varphi_-(x)+\hat\varphi_+(x)\,,
\]
where
\[
\hat\varphi_-(x)=\int d\mu(\pp)\,\hat A_\pp\,f^*_\pp(x)\,,\qquad
\hat\varphi_+(x)=\int d\mu(\pp)\,\hat A^\dagger_\pp\,f_\pp(x)\,.
\]
The operator $\hat\varphi_+(x)$ contains only the creation operators
$\hat A^\dagger_\pp$; it can be interpreted as an operator creating particle
at the space-time point $x$; likewise, $\hat\varphi_-(x)$ absorbs the particle
at the space-time point $x$.\footnote{A subtlety in this interpretation will
be discussed later in Lecture~6.} Since $\hat\varphi_-(x)$ obviously commutes
with $\hat\varphi_-(x')$, and the same is true for the $\hat\varphi_+$
component, we have
\[
\bigl[\hat\varphi(x),\hat\varphi(x')\bigr]=D_-(x-x')-D_+(x-x')\,,
\]
where
\begin{align*}
D_-(x-x')&=\bigl[\hat\varphi_-(x),\hat\varphi_+(x')\bigr]\,,\\
D_+(x-x')&=\bigl[\hat\varphi_-(x'),\hat\varphi_+(x)\bigr]\equiv D_-(x'-x)\,.
\end{align*}
Explicitly
\begin{align*}
D_-(\xx,t)&=\int d\mu(\pp)\,e^{-i\omega_\pp t+i\pp\cdot\xx}\,,\\
D_+(\xx,t)&=\int d\mu(\pp)\,e^{i\omega_\pp t-i\pp\cdot\xx}\,.
\end{align*}
Some properties of these integrals are worth mentioning.

i) The integrals for $D_-(x)$ and $D_+(x)$ are Lorentz-invariant, i.e.\ if two
4-vectors $x$ and $x'$ are related by a Lorentz transformation, $x'=\Lambda x$,
then $D(x)=D(x')$. Indeed, the integrals can be written in explicitly
covariant form, say
\[
D_-(x)=\int_{p^0>0}\frac{d^4p}{(2\pi)^4}\,2\pi\,\delta(p^\mu p_\mu-m^2)\,e^{-ip_\mu x^\mu}\,.
\]
Then the change of the integration variables $p=\Lambda^{-1}p'$ verifies that
$D_-(x)=D_-(\Lambda x)$. Similar representation exists for $D_+(x)$.

ii) The integrals are not absolutely convergent. It is usually convenient to
compute such integrals as limiting values of absolutely convergent integrals.
In our case it can be done by considering the integrals at complex values of
the time variable $t$. For instance, the above integral defining
$D_-(\xx,t)$ converges absolutely if $t$ has \textbf{negative imaginary part}
(remember that $\omega_\pp$ is positive), therefore this integral defines the
function $D_-(\xx,t)$ of complex variable $t$ which is analytic in the lower
half-plane of the complex $t$-plane. Then, for real $t$ we need to take the
limit $\lim_{\epsilon\to-0}D_-(\xx,t+i\epsilon)$ (this is usually written as
$D_-(\xx,t-i0)$). Likewise, the integral for $D_+$ defines the function
$D_+(\xx,t)$ of complex $t$, analytic in the upper half-plane, and for real
$t$ we take $D_+(\xx,t+i0)$.

Let us consider the case of real $t$ such that the 4-vector $x^\mu=(t,\xx)$ is
space-like,
\[
x^\mu x_\mu=t^2-\xx^2<0
\]
(in this consideration we assume that $\xx$ is real). In the complex $t$-plane
we are looking at the segment
\[
-|\xx|<\Re t<|\xx|
\]
of the real axis. It is easy to show that at this segment we have
\[
D_+(\xx,t)=D_-(\xx,t)\qquad\text{for}\quad -|\xx|<t<|\xx|\,.
\]
It suffices to verify this identity at $t=0$, since one can always transform
to the frame where $t=0$, as long as $x^\mu$ is space-like. For $t=0$ we have
\[
D_-(\xx,0)=\lim_{\epsilon\to+0}\int e^{-\omega_\pp\epsilon+i\pp\cdot\xx}\,d\mu(\pp)
\qquad\text{and}\qquad
D_+(\xx,0)=\lim_{\epsilon\to+0}\int e^{-\omega_\pp\epsilon-i\pp\cdot\xx}\,d\mu(\pp)\,,
\]
and the identity is established by the change of the integration variables
$\pp\to-\pp$. Note that this argument also shows that $D_\pm(\xx,t)$ return
real values when $t$ takes values within the above segment of the real axis.

It follows that the function $D_-(\xx,t)$, which was originally defined (by
our integral) in the lower half-plane of $t$, in fact can be analytically
continued, through the above segment of the real axis, to the upper
half-plane, and there it coincides with $D_+(\xx,t)$ (defined, again, through
its integral representation written above); see figure~\ref{fig:tplane}.
\begin{figure}[ht]
\centering\figfile{tplane}
\caption{The complex $t$-plane at fixed $\xx$. The function $D(\xx,t)$ is
analytic away from the two cuts $t\le-|\xx|$ and $t\ge|\xx|$ along the real
axis. Its values in the upper half-plane are given by the integral for
$D_+(\xx,t)$, and in the lower half-plane by the integral for $D_-(\xx,t)$.
The two continue into each other through the segment $-|\xx|<t<|\xx|$.}
\label{fig:tplane}
\end{figure}

In other words, in fact there is a single function of complex $t$ (I will
denote it as $D(\xx,t)$), analytic in both the upper and lower half-planes,
including the segment $-|\xx|<t<|\xx|$ of the real axis, and $D_+(\xx,t)$
($D_-(\xx,t)$) is given by its values in the upper (lower) half-plane. In
particular, for real $t$
\[
D_+(\xx,t)=D(\xx,t+i0)\,,\qquad D_-(\xx,t)=D(\xx,t-i0)\,.
\]
As we have seen, these two limits coincide when $-|\xx|<t<|\xx|$, but outside
this segment they don't. That means the function $D(\xx,t)$ has two branch
cuts in the complex $t$-plane, from $-\infty$ to $-|\xx|$, and from $|\xx|$ to
$+\infty$.

It is not difficult to identify the geometric meaning of the branching points
$t=\pm|\xx|$. They represent points on the future and past components of the
light cone associated with the space-time point $0$, that is the geometric
place of all points which can be connected to $0$ by light rays
(figure~\ref{fig:lightcone}). At fixed $\xx$ the real $t$ line intersects the
light cone at two points $\pm|\xx|$. When $t$ is regarded as the complex
variable, the intersection points become exactly the branching points in the
complex $t$ plane. At real $t$, when crossing the light cone one needs to
choose the branch of the analytic function (similar to choosing a branch of
$\sqrt z$ at negative $z$). The functions $D_+$ and $D_-$ are the two branches
corresponding to two ways of going around the branching points in the complex
$t$-plane.
\begin{figure}[ht]
\centering\figfile{lightcone}
\caption{The light cone of the origin. At fixed $\xx$ the line of real $t$
(thick) meets the future and past sheets of the cone at $t=|\xx|$ and
$t=-|\xx|$. Between these points the separation from the origin is
space-like.}
\label{fig:lightcone}
\end{figure}

Let us come back to the commutator
$[\hat\varphi(x),\hat\varphi(0)]=D_-(x)-D_+(x)$. Here I have set $x'=0$; in
general case $x$ has to be replaced by the difference $x-x'$. We see that the
commutator is equal to the discontinuity
\[
\bigl[\hat\varphi(\xx,t),\hat\varphi(\mathbf0,0)\bigr]=D(\xx,t-i0)-D(\xx,t+i0)\,.
\]
When the separation $x=(t,\xx)$ is space-like the discontinuity vanishes, and
we conclude that
\[
\bigl[\hat\varphi(x),\hat\varphi(x')\bigr]=0\qquad\text{for all space-like separations}\quad x-x'\,.
\]
This is important conclusion. Recall that in quantum mechanics commutator of
two operators allows one to determine how, in a given state, the measurement
of one observable can affect results of the measurements of another
observable. Our calculation shows that the measurement at some space-time
point $x$ can not affect measurements at the point $x'$ as long as the
separation $x-x'$ is space-like. It expresses the \textbf{causality} of the
quantum field theory. This is general requirement of quantum field theory
called the \textbf{local commutativity}. It states that for any local fields
$\mathcal O_1(x)$, $\mathcal O_2(x')$ we must have
\[
\bigl[\hat{\mathcal O}_1(x),\hat{\mathcal O}_2(x')\bigr]=0\qquad\text{for all space-like separations}\quad x-x'\,.
\]
The discontinuities across the branch cuts $|t|>|\xx|$ do not vanish, so when
the separation is time-like the commutator takes non-zero values. We see that
the commutator has its support inside the light cone.

Consider now the expectation value $\bra0\hat\varphi(x)\,\hat\varphi(x')\ket0$
of the product of two field operators. In the KG theory this expectation
value can be expressed through the same analytic function $D(\xx,t)$. Indeed,
since $\hat\varphi_-(x)\ket0=0$ and $\bra0\hat\varphi_+(x)=0$, we have
\[
\bra0\hat\varphi(x)\,\hat\varphi(x')\ket0=\bra0\hat\varphi_-(x)\,\hat\varphi_+(x')\ket0
=D_-(x-x')=D(\xx-\xx',t-t'-i0)\,.
\]
If we interchange the positions of the operators, by similar calculation
\[
\bra0\hat\varphi(x')\,\hat\varphi(x)\ket0=D_+(x-x')=D(\xx-\xx',t-t'+i0)\,.
\]
Thus, these expectation values are described by the function $D(\xx,t)$ taken
at the upper or lower edges of the light-cone branch cuts
(figure~\ref{fig:tcontours}).
\begin{figure}[ht]
\centering\figfile{tcontours}
\caption{Contours in the complex $t$-plane. Along $C_+$, just above both cuts,
$D$ gives $\bra0\hat\varphi(\mathbf0,0)\,\hat\varphi(\xx,t)\ket0$; along
$C_-$, just below both cuts, it gives
$\bra0\hat\varphi(\xx,t)\,\hat\varphi(\mathbf0,0)\ket0$. The Feynman contour
$C_F$ runs above the left cut and below the right cut.}
\label{fig:tcontours}
\end{figure}

\subsection*{Feynman Propagator}

Important role in our future studies will belong to the \textbf{time ordered}
expectation values
\[
D_F(x-x')=D_F(\xx-\xx',t-t')=\bra0\Tprod\bigl(\hat\varphi(\xx,t)\,\hat\varphi(\xx',t')\bigr)\ket0\,,
\]
where the symbol $\Tprod$ signifies the time ordered product of the operators,
\[
\Tprod\bigl(\hat\varphi(\xx,t)\,\hat\varphi(\xx',t')\bigr)=
\begin{cases}
\hat\varphi(\xx,t)\,\hat\varphi(\xx',t') & \text{for}\quad t>t'\\
\hat\varphi(\xx',t')\,\hat\varphi(\xx,t) & \text{for}\quad t<t'\,.
\end{cases}
\]
This expectation value is also called the \textbf{Feynman propagator}. One can
visualize it in figure~\ref{fig:tcontours} by first considering complex values
of $t$ along the contour $C_F$, and then taking the limit when this contour is
brought to the real axis; $D_F$ corresponds to the values of $D$ along $C_F$.

The Feynman propagator can be written as
\[
D_F(\xx,t)=
\begin{cases}
D_-(\xx,t)=\int d\mu(\pp)\,e^{-i\omega_\pp t+i\pp\cdot\xx} & \text{for}\quad t>0\\[2pt]
D_+(\xx,t)=\int d\mu(\pp)\,e^{i\omega_\pp t-i\pp\cdot\xx} & \text{for}\quad t<0\,,
\end{cases}
\]
but the most convenient representation is given by the 4-dimensional integral
\[
D_F(\xx,t)=\int\frac{d^3\pp}{(2\pi)^3}\int_{\tilde C_F}\frac{d\omega}{2\pi}\,
\frac{i\,e^{-i\omega t+i\pp\cdot\xx}}{\omega^2-\pp^2-m^2}\,.
\]
The integrand has two poles in $\omega$, at $\omega=\pm\omega_\pp$, and the
integration contour $\tilde C_F$ for $\omega$ goes around these poles as shown
in figure~\ref{fig:omegaF}.
\begin{figure}[ht]
\centering\figfile{omegaF}
\caption{The contour $\tilde C_F$ in the complex $\omega$-plane passes below
the pole at $\omega=-\omega_\pp$ and above the pole at $\omega=\omega_\pp$.}
\label{fig:omegaF}
\end{figure}

If $t>0$ the integrand decays exponentially into the lower half-plane of the
variable $\omega$, and one can perform the integration over this variable by
closing the contour around $\omega=\omega_\pp$. The integral reduces to the
residue at this pole, which yields $D_-(\xx,t)$. Similarly, when $t<0$ the
integration contour can be closed around the pole $\omega=-\omega_\pp$, and the
residue calculation yields $D_+(\xx,t)$.

This integration prescription is usually expressed as
\[
D_F(x)=\int\frac{d\omega}{2\pi}\,\frac{d^3\pp}{(2\pi)^3}\,\frac{i\,e^{-i\omega t+i\pp\cdot\xx}}{\omega^2-\pp^2-m^2+i0}
=\int\frac{d^4p}{(2\pi)^4}\,\frac{i\,e^{-ip_\mu x^\mu}}{p_\mu p^\mu-m^2+i0}\,.
\]
This form of the Feynman propagator is explicitly covariant. It also shows that
$D_F(x)$ is the Green's function of the KG equation,
\[
\bigl(\partial_\mu\partial^\mu+m^2\bigr)\,D_F(x)=-i\,\delta^{(4)}(x)\,.
\]
Note that the functions $D_\pm(x)$ satisfy homogeneous KG equation
\[
\bigl(\partial_\mu\partial^\mu+m^2\bigr)\,D_\pm(x)=0\,,
\]
since in the KG theory the Heisenberg operator $\hat\varphi(x)$ itself
satisfies homogeneous KG equation.

\begin{exercise}
The Feynman propagator is defined as
\[
D_F(\xx,t)=\theta(t)\,\bra0\hat\varphi(\xx,t)\,\hat\varphi(\mathbf0,0)\ket0+\theta(-t)\,\bra0\hat\varphi(\mathbf0,0)\,\hat\varphi(\xx,t)\ket0\,.
\]
Applying the differential operator $\partial_\mu\partial^\mu+m^2$ and using the
equal-time commutation relations
\[
\bigl[\hat\pi(\xx,t),\hat\varphi(\xx',t)\bigr]=-i\,\delta^{(3)}(\xx-\xx')\,,
\]
confirm that
\[
\bigl(\partial_\mu\partial^\mu+m^2\bigr)\,D_F(x)=-i\,\delta^{(4)}(x)\,.
\]
\end{exercise}

Finally, let us consider the ``euclidean'', or ``imaginary time'' correlation
function. It is defined in terms of the function $D(\xx,t)$ taken at pure
imaginary values of the time variable. Namely, let
\[
t=-i\,x_4\,,
\]
with real $x_4$. Then
\[
D_E(\xx,x_4)=D(\xx,-ix_4)\,,
\]
i.e.\ $D_E$ is the function $D$ with $t$ taken along the contour $C_E$ in the
complex $t$-plane (figure~\ref{fig:tE}).
\begin{figure}[ht]
\centering\figfile{tE}
\caption{The Euclidean contour $C_E$ runs along the imaginary $t$ axis,
between the two light-cone cuts.}
\label{fig:tE}
\end{figure}
Note that this contour can be obtained from $C_F$ by $90^\circ$ rotation.

As is well known, by taking pure imaginary values of $t$ one converts the
Minkowski space-time into 4-dimensional Euclidean space, since
\[
-d\tau^2\to ds^2=dx_4^2+d\xx^2\,,
\]
and the Lorentz symmetry becomes the orthogonal group $O(4)$ of rotations of
the Euclidean 4-vectors
\[
x_E=(\xx,x_4)\,.
\]

The function $D_E(x_E)$ is called the Euclidean correlation function of the KG
theory. We will discuss its significance later. For now, let us observe that
since $D_E$ is related to $D_F$ simply by rotating the time contour, it
satisfies the Euclidean-space version of the inhomogeneous KG equation
\[
\Bigl(m^2-\frac{\partial^2}{\partial x_4^2}-\nabla^2\Bigr)\,D_E(\xx,x_4)=\delta(x_4)\,\delta^{(3)}(\xx)\,,
\]
or in symmetric form
\[
\bigl(m^2-\Delta_E\bigr)\,D_E(x_E)=\delta^{(4)}(x_E)\,,
\]
where $\Delta_E$ stands for the 4-dimensional Laplacian. This equation is
easily solved by Fourier transformation
\[
D_E(x_E)=\int\frac{d^4p_E}{(2\pi)^4}\,\frac{e^{ip_Ex_E}}{p_E^2+m^2}\,,
\]
where the integration is over the 4-dimensional Euclidean momentum space,
$p_E=(\pp,p_4)$, and $p_E^2=\pp^2+p_4^2$. Note that now the integrand does not
have singularities in the integration domain. Also, unlike the Minkowski space
version of this equation, the solution is unique if we demand that $D_E(x_E)$
decays at the Euclidean infinity $x_E\to\infty$. By introducing auxiliary
integration, one can write this function as
\begin{equation}\label{eq:Schwinger}
D_E(x_E)=\int_0^\infty d\tau\int\frac{d^4p_E}{(2\pi)^4}\,e^{-\tau(p_E^2+m^2)}\,e^{ip_Ex_E}\,.
\end{equation}
This form is known as the Schwinger's proper-time representation.

Clearly, the above Euclidean momentum space integral for $D_E(x_E)$ is related
to the covariant 4-momentum integral for the Feynman propagator. The
correlation function $D_E(\xx,x_4)$ is the analytic continuation of
$D_F(\xx,t)$ obtained by rotating the time contour $C_F$ to $C_E$.
Correspondingly, the integration contour $\tilde C_F$ in the complex
$\omega$-plane is rotated by $90^\circ$ to the contour $\tilde C_E$
(figure~\ref{fig:omegaE}).
\begin{figure}[ht]
\centering\figfile{omegaE}
\caption{Wick rotation in the $\omega$-plane. The contour $\tilde C_F$ can be
turned counterclockwise by $90^\circ$ into $\tilde C_E$, along the imaginary
axis, without crossing the poles at $\omega=\pm\omega_\pp$.}
\label{fig:omegaE}
\end{figure}

This example suggests that some analytic characteristics of the KG theory,
and perhaps more general QFT, may take simpler and more symmetric form when
continued to pure imaginary time. This idea can be refined with the help of
the path integral formulation of quantum mechanics, which we consider next.

\begin{hwproblem}{7}\label{hw:7}
\begin{hwparts}
\item Show that the Euclidean Klein-Gordon propagator
\[
D_E(x_E):=\int\frac{d^4p}{(2\pi)^4}\,\frac{e^{ipx_E}}{p^2+m^2}
\]
behaves as
\[
D_E(x_E)\to\Const\,\frac{m^{1/2}}{R^{3/2}}\,e^{-mR}\qquad\text{for}\quad mR\to\infty\,,
\]
where $R$ stands for $|x_E|$. Evaluate the $\Const$.
\item Using the Schwinger's representation, find $D_E$ explicitly in terms of
the Bessel's functions. (Consult e.g.\ A.~Erd\'elyi, H.~Bateman, ``Higher
Transcendental Functions'', Volume~II,
\url{https://authors.library.caltech.edu/43491/}).
\item Analyze the singularity at $x_E^2=0$. For $x_E=(\xx,x_4)$, consider
analytic continuation to complex $x_4=it$. Describe analytic structure of
$D(\xx,t)$ as the function of complex $t$. Find the discontinuity across the
branch cuts associated with the future and past components of the light cone
$t^2=\xx^2$.
\end{hwparts}
\end{hwproblem}

% Source pp. 38-48 (Lectures1-6.pdf).
\section{Path Integrals}

Let me briefly review the method of path integral in ordinary quantum
mechanics. Here I will combine it with the idea of analytic continuation in
the time variable. We will start with the ``time evolution'' in pure imaginary
values of the time parameter. (Besides technical advantage of improving
convergence of the integral, the imaginary-time path integral makes evident
remarkable relation to classical statistical mechanics, see below.) The
real-time transition amplitudes can be recovered by analyticity. For
derivation directly in real time see \textbf{Sect.~9.1 of PS}.

\subsection*{Path integral in Quantum Mechanics}

Given a quantum mechanical system, with the space of states $\Ham$ and the
time-independent Hamiltonian $\hat H$, many problems can be reduced to the
calculations of appropriate matrix elements of its time-evolution operator,
\[
\bra{f}e^{-\frac i\hbar\hat Ht}\ket{i}\,.
\]
In terms of the stationary states $\ket n$ this can be written as
\[
\sum_n\braket{f}{n}\braket{n}{i}\,e^{-\frac i\hbar E_nt}\,.
\]
In a typical problem, the energy spectrum $E_n$ is bounded from below,
$E_n\ge E_0$, and $E_n$ grow to infinity with $n$. The above expression is a
sum of oscillating terms, and we would like to have an efficient mathematical
tool to sort them out. It is useful to consider this sum at complex values of
$t$. If the number of states does not grow exponentially with $E$ (and in
ordinary quantum system it never does), this sum converges absolutely if $t$
has negative imaginary part; thus it defines an analytic function of $t$ in
the lower half-plane. The real-time matrix element is the limiting value of
this analytic function. One can start with pure imaginary $t$,
\[
t=-i\tau\,,
\]
and try to evaluate the matrix elements of the corresponding bounded Hermitian
operator
\[
\bra{f}e^{-\frac1\hbar\hat H\tau}\ket{i}\,,
\]
anticipating that the real-time answer will be then obtained by analytic
continuation.

Consider a system with one degree of freedom $q$, and assume that the
Hamiltonian operator is ``local'' in the basis which diagonalizes $\hat q$;
this means that the matrix elements
\[
\bra{q}\hat H\ket{q'}
\]
have a support restricted to $q=q'$. Most typical Hamiltonians belong to this
class, for instance for $H=\frac12p^2+V(q)$ the above matrix element is
($\hbar=1$ from now on)
\[
-\frac12\,\delta''(q-q')+V(q)\,\delta(q-q')\,.
\]
In such cases it is possible to develop a path-integral representation for the
matrix element
\[
\bra{q_f}e^{-\hat H\tau}\ket{q_i}\,.
\]
The composition property (which follows from $\mathbb{I}=\int \dd q\,\ket q\bra q$)
\[
\bra{q_f}e^{-\hat H(\tau_1+\tau_2)}\ket{q_i}=\int\bra{q_f}e^{-\hat H\tau_1}\ket{q}\,\bra{q}e^{-\hat H\tau_2}\ket{q_i}\dd q
\]
can be applied repeatedly to reduce the above matrix element to the
composition of the matrix elements with arbitrarily small ``time'' intervals
$\Delta\tau$,
\[
\bra{q_f}e^{-\hat H\tau}\ket{q_i}=\int\prod_{k=1}^{n-1}\dd q_k\,\prod_{k=1}^n\bra{q_k}e^{-\hat H\Delta\tau}\ket{q_{k-1}}\,,
\]
where $q_0=q_i$, $q_n=q_f$, and $\Delta\tau=\tau/n$. In the limit of small
$\Delta\tau$, as the consequence of the ``locality'' of $\hat H$ in the
$q$-space, only the matrix elements $\bra{q'}\exp(-\hat H\Delta\tau)\ket{q}$
with $q'$ close to $q$ will bring significant contributions to the integral.

To see this, consider again the Hamiltonian of the form
\[
\hat H=\frac12\,\hat p^2+V(\hat q)\,.
\]
It is straightforward to show that in this case the matrix element
\[
G(q,q'|\tau)=\bra{q}e^{-\hat H\tau}\ket{q'}
\]
satisfies the differential equation of the diffusion type
\[
-\frac{\partial}{\partial\tau}\,G(q,q'|\tau)=\left(-\frac12\,\frac{\partial^2}{\partial q^2}+V(q)\right)\,G(q,q'|\tau)\,,
\]
with the initial condition
\[
G(q,q'|0)=\delta(q-q')\,.
\]
For $V(q)=0$ the problem is easily solved,
\[
G(q,q'|\tau)=\frac1{\sqrt{2\pi\tau}}\,e^{-\frac{(q-q')^2}{2\tau}}\,.
\]
This function indeed decays fast at $|q-q'|\gg\sqrt\tau$. In the case of
nonzero potential $V$ the problem still can be solved in the limit of small
$\tau$. We write
\[
G(q,q'|\tau)=e^{-\sigma(q,q'|\tau)}
\]
and look for $\sigma$ of the form of the following small-$\tau$ expansion
\[
\sigma(q,q'|\tau)=\frac{\sigma_0(q,q')}{\tau}+a\,\log\tau+b+\tau\,\sigma_1(q,q')+\dots\,.
\]
Substituting this ansatz into the above diffusion equation one finds (I
suggest doing it as an \textbf{Exercise})
\[
\sigma_0(q,q')=\frac{(q-q')^2}{2}\,,\qquad a=1/2\,,\qquad
\sigma_1(q,q')=\frac1{q-q'}\int_{q'}^qV(y)\dd y\,,\qquad\text{etc.}
\]
The leading terms do not depend on $V$, and are the same as for the free case,
and $\sigma_1$ is the potential $V$ averaged over the interval $[q',q]$. The
constant $b=\frac12\log2\pi$ is determined by the initial condition at
$\tau=0$.

Because of the factor $\exp(-\sigma_0(q,q')/\tau)$ the function
$G(q,q'|\tau)$ at small $\tau$ is sharply peaked around $|q-q'|=0$ with the
width $\sqrt\tau$. This is typical for the \emph{Brownian motion}, where
$\bar R\sim\sqrt t$ (our equation \emph{is} the diffusion equation). If the
potential $V$ is differentiable (which we assume), we have
\[
\sigma_1(q,q')=\frac12\,V(q)+\frac12\,V(q')+O(\tau)
\]
(I have used the fact that $(q-q')^2\sim\tau$). Therefore, for finite $\tau$
but large $n$ (i.e.\ small $\Delta\tau=\tau/n$) we can write
\[
\bra{q_f}e^{-\hat H\tau}\ket{q_i}=\lim_{n\to\infty}\left(\frac1{2\pi\Delta\tau}\right)^{n/2}\int\left(\prod_{k=1}^{n-1}\dd q_k\right)\,e^{-\mathcal A\{q_k\}}\,,
\]
where
\[
\mathcal A\{q_k\}=\sum_{k=1}^n\left[\frac{(q_k-q_{k-1})^2}{2\,\Delta\tau}+\Delta\tau\,\frac{V(q_k)+V(q_{k-1})}{2}\right]\,,
\]
and $\Delta\tau=\tau/n$, and it is understood that
\[
q_0=q_i\,,\qquad\text{and}\qquad q_n=q_f\,.
\]
Before the limit $n\to\infty$ is taken, the integration in this formula can be
understood as going over the piecewise linear ``trajectories'' $q(\tau)$
running from $q_i$ to $q_f$ (figure~\ref{fig:piecewise}).
\begin{figure}[ht]
\centering\figfile{piecewise}
\caption{A piecewise linear trajectory from $q_i$ at $\tau_0$ to $q_f$ at
$\tau_n$. The intermediate values $q_k=q(\tau_k)$ are integrated over.}
\label{fig:piecewise}
\end{figure}

When $n\to\infty$, $\Delta\tau\to0$, the function $\mathcal A\{q_k\}$ in the
exponential has formal limit
\[
\mathcal A[q(\tau)]=\int_0^\tau \dd\tau\,\left[\frac12\,\bigl(\dd q(\tau)/\dd\tau\bigr)^2+V\bigl(q(\tau)\bigr)\right]\,.
\]
It coincides with the ``imaginary time action'', i.e.\ the action associated
with our system, with $t$ replaced by $-i\tau$,
\[
e^{\frac i\hbar S[q(t)]}\ \xrightarrow{\ t\to-i\tau\ }\ e^{-\mathcal A[q(\tau)]}\,.
\]
It is conventional to write the $n\to\infty$ limit of the above $n$-fold
integral as
\[
\bra{q_f}e^{-\hat H(\tau_f-\tau_i)}\ket{q_i}=\int\bigl[Dq(\tau)\bigr]_{\substack{q(\tau_i)=q_i\\q(\tau_f)=q_f}}\,e^{-\mathcal A[q(\tau)]}\,,
\]
and to call it the (imaginary time) path integral, since the r.h.s.\ involves
the continuous path $q(\tau)$ from $q_i$ to $q_f$.

Few remarks.
\begin{enumerate}
\item The path integral method makes the most explicit use of the
\emph{superposition principle} of quantum theory. The latter states that if
some process can go several distinct ways, the full transition amplitude is
the sum of the amplitudes associated with each possible way. Here (when
returning back to real time) we sum the amplitudes associated with all
distinct paths from $q_i$ to $q_f$.

\item What kind of paths enter the path integral in the limit $n\to\infty$?
Although it is correct to think of the limiting paths as continuous functions
$q(\tau)$, absolute majority of the paths entering the integral are not smooth
curves. This can be observed directly from the fact that for
$\tau_1-\tau_2\to0$ we have $|q(\tau_1)-q(\tau_2)|\sim\sqrt{|\tau_1-\tau_2|}$,
not $\sim|\tau_1-\tau_2|$ as should be true for a continuously differentiable
function. From this point of view the two terms in the action play
significantly different role. The ``kinetic'' term $\frac12\int(\dd q/\dd\tau)^2\dd\tau$
selects the class of the paths entering the integral, i.e.\ those paths for
which $\bigl(q(\tau+\Delta\tau)-q(\tau)\bigr)^2/\Delta\tau$ remains finite as
$\Delta\tau\to0$; in this sense the factor
$\exp\bigl(-\frac12\int(\dd q/\dd\tau)^2\dd\tau\bigr)$ should be regarded as a part
of the functional measure. The factor $\exp\bigl(-\int V(q(\tau))\dd\tau\bigr)$
weights the paths according to their average potential energy.

\item Note that the symbol
\[
\bigl[Dq(\tau)\bigr]=\lim_{n\to\infty}\left(\frac1{2\pi\Delta\tau}\right)^{n/2}\prod_{k=1}^{n-1}\dd q_k
\]
contains divergent (in the limit) factor
\[
\left(\frac1{2\pi\Delta\tau}\right)^{n/2}=e^{-\frac{\log(2\pi\Delta\tau)}{2\,\Delta\tau}\,L}\,,
\]
where $L=(\tau_f-\tau_i)=n\,\Delta\tau$ is the ``volume'' of the imaginary
time interval. This factor is reminiscent of the irritating infinity we had
for the vacuum energy in the KG theory, and absorbing it into $[Dq(\tau)]$ is
analogous to subtracting $E_0$ from $\hat H$.

\item We have derived the above path integral representation of the evolution
operator in the imaginary time $\tau$. But one can repeat the above
calculation step by step, considering the paths in the real time $t$ instead
of $\tau$ (this is the form in which the path integral is usually introduced
in quantum mechanics, see for instance \S9.1 of the textbook). In this way one
arrives at the expression
\[
\bra{q_f}e^{-\frac i\hbar HT}\ket{q_i}=\int\bigl[Dq(t)\bigr]\,e^{\frac i\hbar S[q(t)]}\,,
\]
with the phase $e^{\frac i\hbar S}$ replacing $e^{-\mathcal A}$. This of
course is equivalent approach, although, unlike the imaginary time scheme, the
integrals involved are not absolutely convergent. Usually, the most efficient
way to handle such integrals is through analytic continuation (in variables or
parameters) which makes the integral absolutely convergent, with subsequent
taking the limiting values. The imaginary-time approach simply makes this
procedure explicit.

\item Generalization of the above construction to the systems with many degrees
of freedom $q_a$ is straightforward, if the Hamiltonian is of the form
\[
H=\sum_a\frac12\,p_a^2+V(q_a)\,;
\]
one obtains
\[
\int\bigl[Dq_a\bigr]\,e^{-\mathcal A[q_a(\tau)]}
\]
for the evolution operator. In more complicated cases the path integral in the
phase space can be useful; this is briefly discussed in Sect.~9.1 of PS.
\end{enumerate}

\subsection*{Relation to Classical Statistical Mechanics}

At this point, let me deviate to discuss remarkable relation which exists
between the \textbf{Path Integral in Quantum Mechanics} (in its imaginary time
version) and the \textbf{Classical Statistical Mechanics}. Let us recall some
generalities of the classical statistical mechanics.

Suppose we have a classical dynamical system with the phase-space coordinates
$\{P_i,Q_i\}$ (I will use capitals here because the relation between these
$Q_i$ and $q_a$ appearing in the path integral is not going to be
straightforward), and the Hamiltonian $\Ham(P_i,Q_i)$. In a thermal
equilibrium state at a temperature $T$ the probability distribution of
microscopic states is given by the Gibbs formula
\[
\mathcal P(P_i,Q_i)\prod_i \dd P_i\dd Q_i=Z^{-1}\,e^{-\beta\Ham(P_i,Q_i)}\,\frac1{N!}\prod_i \dd P_i\dd Q_i\,,
\]
where $\beta=1/kT$, and $Z$ is the partition function
\[
Z=\int\frac{\prod_i \dd P_i\dd Q_i}{N!}\,e^{-\beta\Ham(P_i,Q_i)}\,.
\]
Calculating the partition function is one of the main problems in statistical
mechanics because the thermodynamic quantities are expressed through $Z$, say
the free energy is $\beta F=-\log Z$. Again, a typical form of the Hamiltonian
is
\[
\Ham(P_i,Q_i)=\sum_i\frac12\,P_i^2+U(\{Q_i\})\,.
\]
In this case the integration over $P_i$ is easy to perform,
\[
Z=\frac1{N!}\left(\frac{2\pi}{\beta}\right)^{\frac N2}Z_{\text{conf}}\,,
\]
where $N$ is the number of the degrees of freedom, $i=1,2,\dots,N$, and
$Z_{\text{conf}}$ is the configuration-space integral
\[
Z_{\text{conf}}=\int\prod_i \dd Q_i\,e^{-\beta U(\{Q_i\})}
\]
which contains most of interesting physics.

As a special example of such system, consider an elastic string whose
configuration is described by a function $q(\tau)$, $\tau\in[0,L]$; the values
$q(\tau)$ for all $\tau$ within this interval play the role of the
configuration-space coordinates (figure~\ref{fig:string}),
\[
\{Q_i\}\to q(\tau)\,.
\]
\begin{figure}[ht]
\centering\figfile{string}
\caption{A configuration $q(\tau)$ of the elastic string on $[0,L]$ with the
ends held at $q(0)=q_i$ and $q(L)=q_f$.}
\label{fig:string}
\end{figure}
We also assume that the string lies in a potential well $V(q)$, so that the
potential energy of the string is
\[
U[q(\tau)]=\int_0^L\dd\tau\,\left[V\bigl(q(\tau)\bigr)+\frac12\,\left(\frac{\dd q(\tau)}{\dd\tau}\right)^2\right]\,,
\]
where the last term accounts for the elastic energy; to simplify the equations
I have assumed that the elastic string tension equals 1.

To evaluate the configuration part of the partition function, we must
integrate over the functions $q(\tau)$ satisfying given boundary conditions at
the ends of the interval $[0,L]$, say
\[
q(0)=q_i\,,\qquad q(L)=q_f\,,
\]
with the statistical weights $\exp(-\beta\,U[q(\tau)])$. Thus
\[
Z_{\text{conf}}=\int\bigl[Dq(\tau)\bigr]\,e^{-\beta U[q(\tau)]}\,.
\]
With the above form of $U[q(\tau)]=\mathcal A[q(\tau)]$, this expression is
identical to the path integral for the imaginary time transition amplitude,
\[
Z_{\text{conf}}\sim\bra{q_f}e^{-\frac1\hbar\hat HL}\ket{q_i}\,,
\]
provided we also identify
\[
\beta=1/\hbar\,.
\]
We see that quantum mechanics with a single degree of freedom $q$, and the
Hamiltonian operator
\[
\hat H=\frac12\,\hat p^2+V(\hat q)\,,
\]
is related in this way to classical statistical mechanics of a system with
continuously many degrees of freedom $q(\tau)$.

\subsection*{Correlation Functions}

In statistical mechanics one may also be interested in \emph{correlation
functions}. In our example one may ask how the string positions $q(\tau)$ at
different points, say at $\tau_1,\tau_2,\dots,\tau_N$, are correlated in the
thermal equilibrium state. In figure~\ref{fig:stringcorr} I made a trivial
shift of the variable $\tau$, so that from now on $\tau\in[-L/2,L/2]$.
\begin{figure}[ht]
\centering\figfile{stringcorr}
\caption{The string on $[-L/2,L/2]$, with its positions $q(\tau_1),\dots,q(\tau_N)$
marked at $\tau_1\le\tau_2\le\dots\le\tau_N$.}
\label{fig:stringcorr}
\end{figure}

The correlations are described by the averaged products
$q(\tau_1)q(\tau_2)\cdots q(\tau_N)$, which in turn are given by the integrals
\[
\bigl\langle q(\tau_1)\,q(\tau_2)\cdots q(\tau_N)\bigr\rangle
=Z^{-1}\int q(\tau_1)q(\tau_2)\cdots q(\tau_N)\,e^{-\mathcal A[q(\tau)]}\,\bigl[Dq(\tau)\bigr]\,,
\]
because $Z^{-1}e^{-\mathcal A[q(\tau)]}\,[Dq(\tau)]$ is the statistical
weight of the microstate $q(\tau)$. It is instructive to find out how these
correlation functions are interpreted in associated quantum mechanical
problem.

Without loss of generality we can assume that
\[
-L/2\le\tau_1\le\tau_2\le\dots\le\tau_N\le L/2\,.
\]
The integration over all microstates $q(\tau)$ can be performed in two steps:
First, we fix the values of $q(\tau)$ at the points $\tau_1,\tau_2,\dots,\tau_N$
to be
\[
q(\tau_1)=q_1\,,\quad q(\tau_2)=q_2\,,\quad\dots\quad q(\tau_N)=q_N\,,
\]
and integrate over all $q(\tau)$ subject to these constraints. After that, the
integration over $q_1,q_2,\dots,q_N$ is performed to complete the integral.
Comparing this with the path integral representation for the evolution
operator, say
\[
\bra{q_2}e^{-\hat H(\tau_2-\tau_1)}\ket{q_1}=\int\bigl[Dq(\tau)\bigr]_{\substack{q(\tau_1)=q_1\\q(\tau_2)=q_2}}\,e^{-\mathcal A[q(\tau)]}\,,
\]
we find
\begin{multline*}
\bigl\langle q(\tau_1)q(\tau_2)\cdots q(\tau_N)\bigr\rangle=Z^{-1}\int\left[\prod_{k=1}^N\dd q_k\right]\times q_1q_2\cdots q_N\\
\times\bra{q_f}e^{-\hat H(L/2-\tau_N)}\ket{q_N}\bra{q_N}e^{-\hat H(\tau_N-\tau_{N-1})}\ket{q_{N-1}}\cdots\bra{q_1}e^{-\hat H(\tau_1+L/2)}\ket{q_i}\,.
\end{multline*}
Recall that in the coordinate representation the operator $\hat q$ acts by
multiplication
\[
\hat q\,\Psi(q)=q\,\Psi(q)\,,
\]
so that in this expression we can replace
\[
\int \dd q\,\ket{q}\,q\,\bra{q}=\hat q\,.
\]
Then the above expression can be written as
\[
\bra{q_f}e^{-\hat H(L/2-\tau_N)}\,\hat q\,e^{-\hat H(\tau_N-\tau_{N-1})}\,\hat q\cdots\hat q\,e^{-\hat H(\tau_1+L/2)}\ket{q_i}
\]
which is best expressed in terms of the ``Euclidean'' (i.e.\ imaginary time)
Heisenberg operators
\[
\hat q_E(\tau)=e^{\hat H\tau}\,\hat q\,e^{-\hat H\tau}
\]
($\tau=it$), allowing one to write the above formula as
\[
\bra{q_f}e^{-\hat H\frac L2}\,\hat q_E(\tau_N)\,\hat q_E(\tau_{N-1})\cdots\hat q_E(\tau_1)\,e^{-\hat H\frac L2}\ket{q_i}\,.
\]
To simplify things further, let's take the limit $L\to\infty$, thus turning to
the correlation functions of the infinite string. In this limit the exact
boundary conditions, e.g.\ the values $q_i$ and $q_f$, become unimportant.
Indeed, the states $\ket{q_i}$ and $\ket{q_f}$ can be expanded in the
stationary states $\ket n$ of the Hamiltonian $\hat H$, say
\[
\ket{q_i}=\sum_n\ket n\braket{n}{q_i}=\sum_n\ket n\,\Psi_n^*(q_i)\,,
\]
where it is assumed that
\[
\braket{n}{m}=\delta_{n,m}\,.
\]
In the limit $L\to\infty$ the contribution of the ground state $\ket0$
dominates, since all other states are suppressed by the factor
$\sim e^{-(E_n-E_0)\frac L2}$. Thus we have
\[
\bigl\langle q(\tau_1)q(\tau_2)\cdots q(\tau_N)\bigr\rangle=\frac{\braket{q_f}{0}\braket{0}{q_i}}{Z}\,e^{-E_0L}\,
\bra0\hat q_E(\tau_N)\,\hat q_E(\tau_{N-1})\cdots\hat q_E(\tau_1)\ket0\,.
\]
Furthermore, note that
\[
Z=\bra{q_f}e^{-\hat HL}\ket{q_i}\ \to\ \braket{q_f}{0}\braket{0}{q_i}\,e^{-E_0L}\,,
\]
hence the first factor in the above equation equals 1 (in the limit
$L=\infty$), and we have
\[
\bigl\langle q(\tau_1)\cdots q(\tau_N)\bigr\rangle=\bra0\hat q_E(\tau_N)\,\hat q_E(\tau_{N-1})\cdots\hat q_E(\tau_1)\ket0\,.
\]
Note that the operators here are arranged in the order of $\tau$ increasing
from the right to the left. Let me stress that the Euclidean Heisenberg
operators can not be arranged any other way, as otherwise the product of the
operators (the sum over the intermediate states) does not converge. More
generally, if we consider complex values of $\tau_i$ (which we will), the above
product of the operators makes direct sense in the domain
\[
\Re\tau_1\le\Re\tau_2\le\dots\le\Re\tau_N
\]
where the sums over the intermediate states converge (hence the sums yield
analytic function).

To relate the above Euclidean correlation function to the real-time ones one
analytically continues to pure imaginary values of $\tau$,
\[
\tau_k=it_k+0\,,
\]
where ``$+0$'' shift indicates that if there is a branch cut along the
imaginary axis the values at the right edge are to be taken.

% Source pp. 49-55 (Lectures1-6.pdf).
\section{Path integral for a relativistic particle}

So far, we have discussed the path integral representation in ordinary quantum
mechanics. The real-time transition amplitude is represented as a suitably
defined integral
\[
\bra{q_f}e^{-\frac i\hbar HT}\ket{q_i}=\int\bigl[Dq(t)\bigr]_{\substack{q(0)=q_i\\q(T)=q_f}}\,e^{\frac i\hbar S[q(t)]}\,,
\]
over all ``paths'' $q(t)$ from $q_i$ to $q_f$. Now I want to discuss how the
idea of the path integral can be applied to a relativistic particle.

As we have seen (Lecture~4), the KG theory is the theory of spinless Bose
particles, and we expect that the Feynman propagator $D_F(x_f-x_i)$ (and
related Euclidean correlation function $D_E$) admit an interpretation in terms
of amplitude associated with the particle propagation from $x_i$ to $x_f$. Can
we write such amplitude as an integral over the ``paths''?

In classical mechanics the action of a relativistic particle travelling from
$x_i$ to $x_f$ is
\[
S=-m\int_{x_i}^{x_f}\sqrt{dx_\mu\,dx^\mu}=-m\int_{(t_i,\xx_i)}^{(t_f,\xx_f)}\sqrt{dt^2-d\xx^2}\,,
\]
where the integration is performed along a given path (generally described by
a parametric curve $x^\mu(s)$) from $x_i$ to $x_f$. In classical theory it is
assumed that $(dx^0)^2\ge d\xx^2$ along the whole path (the particle velocity
never exceeds the speed of light). Then the time component $t=x^0$ of $x^\mu$
can serve as the parameter along the path, and we can write
\[
S[\xx(t)]=-m\int_{t_i}^{t_f}dt\,\sqrt{1-\dot\xx^2}\,.
\]
In quantum theory one may attempt to write for the transition amplitude
\[
\int_{\text{paths}(i\to f)}e^{\frac i\hbar S}\,[D\xx(t)]\,,
\]
where the integration is over ``all paths'' from $x_i$ to $x_f$
(figure~\ref{fig:relpath}).
\begin{figure}[ht]
\centering\figfile{relpath}
\caption{Two of the paths from $x_i$ to $x_f$ that enter the sum. The
element $dx^\mu$ of a path may be time-like or space-like.}
\label{fig:relpath}
\end{figure}

The problems appear when we want to give a precise meaning to this expression.
On the formal side, we observe that the Hamiltonian associated with the
problem,
\[
H=\sqrt{m^2+\pp^2}
\]
(where $\pp$ is the momentum conjugated to $\xx$), is not ``local'' in the
sense we discussed the last time; the matrix elements $\bra{\xx}H\ket{\xx'}$
do not vanish when $\xx\neq\xx'$ (check it against the Exercise in
Lecture~2). More fundamentally, the problem is about the class of the paths
which we need to include into the path integral. The expression
$dx_\mu dx^\mu$ in the square root in $S$ can be either positive or negative,
depending on whether $dx^\mu$ is time-like or space-like. One can try to
exclude from the integration the paths having space-like $dx^\mu$ anywhere,
but it does not seem to be right idea if we want to describe the propagator
$D_F$ this way. Indeed, unlike the commutator, the function $D_F(\xx,t)$ does
not vanish outside the light cone (at $|t|<|\xx|$). It is exponentially small
if you go far outside the light cone, but it is not zero. This is similar to a
non-relativistic particle in the classically inaccessible region of its
configuration space: the wave function is exponentially small but does not
vanish (``quantum tunneling''). Thus we have to admit the paths having
space-like $dx^\mu$ somewhere along them. Once such paths are admitted, the
Lorentz invariance requires that paths in which $dx^0$ is allowed to take
negative values (i.e.\ the paths with the folds going ``backward in time'')
are admitted as well (figure~\ref{fig:backpath}).
\begin{figure}[ht]
\centering\figfile{backpath}
\caption{A path with a fold going backward in time.}
\label{fig:backpath}
\end{figure}

For such paths $\xx(t)$ is not a function, and the integral over $[D\xx(t)]$
does not apply. Besides, whereas for the case of space-like $dx^\mu$ there is
a natural choice of the branch of the square root $\sqrt{dx_\mu dx^\mu}$ in the
action ($e^{iS}$ must receive exponentially small contributions from the
space-like parts of the path), we have no clear intuition which sign to choose
in the case of the time-like ``backward in time'' paths.

All these ambiguities are resolved by adopting the Euclidean ($=$ imaginary
time) formulation of the path integral. Replacing as before
\[
t=-i\,x_4\,,
\]
we have $iS\to-\mathcal A$, where the Euclidean action is
\[
\mathcal A=m\int\sqrt{(dx_E)^2}\,,
\]
with explicitly positive expression
\[
dx_E^2=dx_4^2+d\xx^2
\]
under the square root. In the subsequent discussion I will often omit the
subscript $E$ for the Euclidean 4-vectors. With this, one can define the
Euclidean amplitude as the integral
\[
\int_{\text{paths}(x_i\to x_f)}e^{-\mathcal A[\text{path}]}
\]
over all paths connecting two Euclidean points $x_i$ and $x_f$; one expects
that the real-time propagator will be accessible through analytic
continuation, as we previously discussed.

Note that
\[
\mathcal A[\text{path}]=m_0\,L[\text{path}]\,,
\]
where $L[\text{path}]$ is just the Euclidean length of the path. I have
replaced $m$ by the ``bare mass'' $m_0$; we will see that this parameter does
not necessarily coincide with the physical mass but requires
``renormalization''.

By analogy with our definition of the path integral in ordinary quantum
mechanics, let us define this path integral as a limit of finite-dimensional
integrals. Namely, for given (large) number $n$, let us consider the piecewise
linear paths from $x_i=0$ to $x_f=x$, each linear piece having fixed length
$\Delta$ (figure~\ref{fig:steps}).
\begin{figure}[ht]
\centering\figfile{steps}
\caption{A piecewise linear path from $0$ to $x$ made of $n$ pieces of equal
length $\Delta$.}
\label{fig:steps}
\end{figure}
The overall length of such path is
\[
L_n=n\,\Delta\,.
\]
Let $\nu_k$ be the unit 4-vector in the direction of the $k$-th piece;
$|\nu_k|=1$. In this discrete approximation the path integral can be written
as
\[
\sum_{n=0}^\infty e^{-m_0\Delta n}\int\Bigl(\prod_{k=1}^nd\mu(\nu_k)\Bigr)\,\delta^{(4)}\Bigl(\Delta\sum_{k=1}^n\nu_k-x\Bigr)\,.
\]
Here $d\mu(\nu_k)$ stands for the usual measure on $S^3$, and the
delta-function ensures that the vectors $\Delta\,\nu_k$ add up to $x$, i.e.\
the path indeed is from $0$ to $x=x_f-x_i$. Finally, the sum over $n$ must be
performed since we want to include all paths, with arbitrary length $L$.

The delta-function can be represented as
\[
\delta^{(4)}\Bigl(\Delta\sum\nu_k-x\Bigr)=\int\frac{d^4P}{(2\pi)^4}\,e^{iP(x-\Delta\sum\nu_k)}\,.
\]
For fixed $P$ the integration over $\prod d\mu(\nu_k)$ splits into $n$
identical integrals $I(P\Delta)=\int d\mu(\nu)\,e^{-i\Delta\nu P}$,
\[
\prod_{k=1}^n\int d\mu(\nu_k)\,e^{-i\Delta\nu_kP}=\bigl[I(P\Delta)\bigr]^n\,.
\]
Certainly, the integral $I$ can be evaluated explicitly in terms of the
Bessel's function. However, there is no need in its explicit form. We are
going to send $\Delta\to0$ (in which limit a significant contribution comes
from large $n$). Therefore only the behavior of $I$ at small values of its
argument is important. We find
\[
I(P\Delta)=A-A\,\xi\,P^2\Delta^2+O(\Delta^4)\,,
\]
where exact values of the constants $A=\int_{S^3}d\mu(\nu)=2\pi^2$ and
$\xi=1/8$ are not important. We have
\[
I^n\sim A^n\,e^{-\Delta^2n\,\xi P^2}\,.
\]
Now the integration over $P$ can be performed, and we obtain for given $n$
\[
\bigl(4\pi\,n\Delta^2\xi\bigr)^{-2}\,e^{-\frac{x^2}{4\xi n\Delta^2}}\,.
\]
Note that $\Delta$ enters here in the combination $n\Delta^2$, not $n\Delta$
as one could naively expect. Attempting to take the naive limit
\[
\Delta\to0\,,\qquad n\to\infty\qquad\text{with}\quad L=n\Delta=\text{finite}
\]
would lead to the propagation amplitude which vanishes at any finite $|x|$.
This result is not surprising. Our calculation was nearly identical to the
well known probability calculation for a Brownian particle. We know that a
Brownian particle with typical microscopic velocity $\bar v$ will be found,
after the elapsed time $T$, at the distance
\[
\bar x=\bar v\,\sqrt{T\,\Delta t}
\]
away from its initial position. Here $\Delta t$ is the typical time between
collisions. In this analogy
\[
\Delta\sim\bar v\,\Delta t\,,\qquad n\sim T/\Delta t\,.
\]
Thus we expect to have
\[
\bar x^2=n\,\Delta^2\,,
\]
in agreement with our calculation.

We conclude that the right continuous limit is achieved by taking
\[
\Delta\to0\,,\qquad n\to\infty\qquad\text{with}\quad s=\xi\,n\,\Delta^2=\text{finite}\,,
\]
and it is the parameter $s$ that characterizes the geometry of the typical
path; the microscopic length $L=n\Delta$ becomes infinitely large in this
limit. The absolute majority of the paths contributing to the functional
integral are extremely crumpled at all scales, as figure~\ref{fig:crumpled}
illustrates.
\begin{figure}[ht]
\centering\figfile{crumpled}
\caption{A typical path is crumpled at all scales: a small piece of it,
magnified (right), looks as crumpled as the whole.}
\label{fig:crumpled}
\end{figure}
The parameter $s$ has the dimension of the $[\text{length}]^2$; it is said that
the typical path has the fractal dimension 2.

Finally, there are two unpleasant factors to be dealt with before we can take
the limit $\Delta\to0$. These are
\[
A^n\,e^{-m_0n\Delta}=e^{-n(m_0\Delta-\log A)}\,.
\]
Consider $m_0$ as a parameter which can carry some dependence on $\Delta$, and
is free to be adjusted (in the limit $\Delta\to0$) so that
\[
\Delta\,m_0(\Delta)-\log A\ \to\ \xi\,m^2\,\Delta^2\qquad\text{as}\quad\Delta\to0\,,
\]
where it is $m^2$, not $m_0$, which remains constant in this limit. Then the
above two factors combine into
\[
e^{-m^2\xi n\Delta^2}\to e^{-m^2s}\,.
\]
For large $n$ the sum is replaced by the integral over $s$, and one obtains
\[
\int_0^\infty ds\,e^{-m^2s}\int\frac{d^4P}{(2\pi)^4}\,e^{-sP^2}\,e^{iPx}\,,
\]
which is exactly the Schwinger's proper-time representation \eqref{eq:Schwinger}
of the Euclidean correlation function $D_E(x)$ of the KG theory. Thus this
correlation function admits interpretation as the integral over the paths in
Euclidean space-time. As we know, the Feynman propagator $D_F$ coincides with
the analytic continuation of this correlation function back to the real time,
$x_4=it$.

\begin{remark}
In our calculation above we have assumed that the linear pieces of the
discretized paths have equal lengths $\Delta$. This is not essential, and is
done for simplicity only. One can change details of the finite-dimensional
approximations, but in the limit the same correlation function $D_E$ comes
out (though the relation between the ``bare'' and ``renormalized'' mass
parameters $m_0$ and $m$ does depend on the details of the approximation).
This is an elementary example of \emph{universality}.
\end{remark}

In principle, it is possible to develop a theory of interacting relativistic
particles based on the idea of integrating over all the particles paths,
including the paths of many particles, with interaction events (``vertices'').
But much more universal and powerful approach is based on quantization of the
field variables. In the associated path integral, one integrates not over the
paths of the particles, but over the paths of the field variables, i.e.\ over
all the ``trajectories'' $\Phi(\xx,t)$ of the fields. In what follows we
elaborate this idea in the KG theory.

% Source pp. 56-60 (Lectures7-12.pdf).
\section{Path Integral in Klein-Gordon Theory}

Let us now apply the idea of the (Euclidean) path integral to the KG field
theory. In this case the role of the generalized coordinates $q$ is taken by
the field degrees of freedom $\varphi(\xx)$, thus we have to replace
\[
q\to\varphi(\xx)\,,\qquad q(\tau)\to\varphi(\xx,\tau)=\varphi(x_E)\,,
\]
where $\tau$ is the ``imaginary time'' which I will denote $x_4$ below, and
$x_E=(\xx,x_4)$. The Euclidean action is obtained by substituting $t=-ix_4$ in
the KG action, i.e.
\[
\mathcal A=\int dx_4\,d\xx\,\frac12\Bigl[(\partial_4\varphi)^2+(\nabla\varphi)^2+m^2\varphi^2\Bigr]
=\int d^4x_E\,\left[\frac12\,(\partial_\mu\varphi)^2+\frac{m^2}{2}\,\varphi^2\right]\,,
\]
where now $\mu=1,2,3,4$. We are interested in the path integral
\[
Z=\int\bigl[D\varphi(x)\bigr]\,e^{-\mathcal A[\varphi(x)]}\,.
\]
More generally, we will be interested in the correlation functions defined as
the ratios
\begin{equation}\label{eq:corrdef}
\bigl\langle\varphi(x_1)\varphi(x_2)\cdots\varphi(x_N)\bigr\rangle
=\frac1Z\int\bigl[D\varphi\bigr]\,\varphi(x_1)\varphi(x_2)\cdots\varphi(x_N)\,e^{-\mathcal A[\varphi]}\,.
\end{equation}
Here and below I drop the subscript $E$ and use the notation $x_i$ for the
\emph{Euclidean} vectors, unless stated otherwise. As we discussed before, when
continued to complex values of $(x_i)_4$ and specialized to $(x_i)_4=it_i$,
these correlation functions yield the time-ordered expectation values
\[
\bra0\Tprod\bigl(\hat\varphi(\xx_1,t_1)\cdots\hat\varphi(\xx_N,t_N)\bigr)\ket0\,,
\]
where $\hat\varphi(\xx,t)$ denotes the real-time Heisenberg field operator.

Few words about the definition of the path integral in the case of fields. For
the path integral in quantum mechanics we had a definition in terms of a limit
$\Delta\tau\to0$ of a finite-dimensional integral with discretized ``time''
$\tau$. In principle, there are many ways to build finite approximations in
quantum field theory. For example, the 4-dimensional space can be replaced by,
say, the hypercubic lattice, in which $x^\mu$ take discrete values
\[
x^\mu=\Delta\,n^\mu\,,\qquad n^\mu\in\mathbb Z\,.
\]
The derivatives and the integration in the action can be replaced by finite
differences and summations, i.e.
\[
\partial_\mu\varphi(x)\to\frac1\Delta\,\bigl(\varphi(x+e^\mu\Delta)-\varphi(x)\bigr)\,;\qquad
\int d^4x\to\Delta^4\sum_x
\]
($e^\mu$ is the unit vector in the direction $\mu$); with this one can set
\begin{equation}\label{eq:latticemeasure}
\bigl[D\varphi\bigr]=\lim_{\Delta\to0}\prod_xd\varphi(x)\,.
\end{equation}
In the Klein-Gordon theory such definition works more or less
straightforwardly. In more interesting cases with an interaction defining the
integral over $[D\varphi]$ can become highly nontrivial problem, which is
closely related to the problem of \emph{critical behavior}; many subsequent
lectures will be devoted to this subject. Let me also note here that in the
context of field theory the integrals like those above are usually referred to
as \emph{functional integrals} (the term ``path integral'' is used
interchangeably).

The action $\mathcal A$ of the KG theory is quadratic in $\varphi$. This makes it
possible to evaluate the above path integrals explicitly. I am not going to
explain how these integrals are actually evaluated; this calculation is
presented in \textbf{PS, \S\,9.2}. Instead I will derive the above correlation
functions using formal properties of the functional integral. The resulting
formula is known as the

\subsection*{Wick's Theorem}

\begin{itemize}
\item The correlation function of any odd number of fields vanishes,
\[
\bigl\langle\varphi(x_1)\cdots\varphi(x_{2N+1})\bigr\rangle=0\,.
\]
This follows from the fact that the action $\mathcal A$ is an even functional
of $\varphi$, while in this case the product of an odd number of factors
$\varphi$ in the integrand is odd.
\item The two-point correlation function is
\[
\bigl\langle\varphi(x_1)\varphi(x_2)\bigr\rangle=D_E(x_1-x_2)\,,
\]
where $D_E(x)$ is the familiar function
\[
D_E(x)=\int\frac{d^4p}{(2\pi)^4}\,\frac1{p^2+m^2}\,e^{ipx}\,.
\]
From now on I will usually drop the subscript $E$, so that $D(x)$ will stand
for this function. Also, $x$ and $p$ now stand for the Euclidean vectors
$x^\mu=(\xx,x_4)$, $p^\mu=(\pp,p_4)$, with $xp:=\xx\cdot\pp+x_4p_4$.
\item The $2N$-point correlation function
\[
\bigl\langle\varphi(x_1)\cdots\varphi(x_{2N})\bigr\rangle
\]
is a sum of $(2N-1)!!\equiv1\cdot3\cdot\ldots\cdot(2N-1)$ terms, each term
corresponding to one of $(2N-1)!!$ distinct ways of ``pairing'' (or
``contracting'') between $2N$ field insertions
$\varphi(x_1)\cdots\varphi(x_{2N})$. A pairing of the insertions
$\varphi(x_i)$ with $\varphi(x_j)$ is usually denoted by an overbrace
\[
\dots\contraction{}{\varphi(x_i)}{\cdots}{\varphi(x_j)}\varphi(x_i)\cdots\varphi(x_j)\dots\,.
\]
Every contraction gives rise to the factor $D(x-x')$,
\[
\contraction{}{\varphi(x)}{\,}{\varphi(x')}\varphi(x)\,\varphi(x')=D(x-x')\,.
\]
\end{itemize}
The last rule of pairing can be expressed in a compact form, in terms of the
following

\subsection*{(Wick's) Recursion Relation}
\begin{equation}\label{eq:WickRec}
\bigl\langle\varphi(x_1)\varphi(x_2)\cdots\varphi(x_N)\bigr\rangle
=\sum_{k=2}^ND(x_1-x_k)\,\bigl\langle\varphi(x_2)\cdots\widehat{\varphi(x_k)}\cdots\varphi(x_N)\bigr\rangle\,,
\end{equation}
where the hat indicates that this insertion is dropped. This relation
expresses $N$-point correlation function in terms of the $N-2$-point ones.
When supplemented with the condition $\langle\varphi(x)\rangle=0$, it contains
all the statements of the Wick's theorem ($\langle1\rangle=1$ by definition).

Let us derive this relation. The main assumption will be the invariance of the
functional measure $D[\varphi]=D[\varphi']$ under the functional shifts
\[
\varphi(x)\to\varphi'(x):=\varphi(x)+c(x)
\]
with arbitrary fixed function $c(x)$. This invariance is an evident property
of the discrete approximation \eqref{eq:latticemeasure}, so we assume it.

Consider the functional integral representing the $N-1$-point correlation
function,
\[
Z^{-1}\int\bigl[D\varphi\bigr]\,\varphi(x_2)\cdots\varphi(x_N)\,e^{-\mathcal A[\varphi]}\,.
\]
Note that the product of the insertions starts with $\varphi(x_2)$, and
involves $N-1$ terms. Make the following change of the integration variables
\[
\varphi(x)\to\varphi'(x)=\varphi(x)+\epsilon(x)\,,
\]
where $\epsilon(x)$ is an arbitrary infinitesimal function. The action changes
as follows
\[
\mathcal A[\varphi+\epsilon]=\mathcal A[\varphi]+\int d^4x_1\,\epsilon(x_1)\,\bigl[m^2-\Delta_{x_1}\bigr]\varphi(x_1)+O(\epsilon^2)\,,
\]
where I have assumed that the function $\epsilon(x)$ decays sufficiently fast
making the integration by parts possible. Here
\[
\Delta_x:=\sum_{\mu=1}^4\left(\frac{\partial}{\partial x^\mu}\right)^2
\]
is 4D Laplace operator. Since the change of integration variables does not
change the value of the integral, we have the identity
\begin{multline*}
0=\sum_{k=2}^N\epsilon(x_k)\int\bigl[D\varphi\bigr]\,\varphi(x_2)\cdots\widehat{\varphi(x_k)}\cdots\varphi(x_N)\,e^{-\mathcal A}\\
-\int d^4x_1\,\epsilon(x_1)\,\bigl(m^2-\Delta_{x_1}\bigr)\int\bigl[D\varphi\bigr]\,\varphi(x_1)\varphi(x_2)\cdots\varphi(x_N)\,e^{-\mathcal A}\,.
\end{multline*}
As the function $\epsilon(x)$ is arbitrary, we must have
\begin{equation}\label{eq:WickSD}
\bigl(m^2-\Delta_{x_1}\bigr)\bigl\langle\varphi(x_1)\varphi(x_2)\cdots\varphi(x_N)\bigr\rangle
=\sum_{k=2}^N\delta^{(4)}(x_1-x_k)\,\bigl\langle\varphi(x_2)\cdots\widehat{\varphi(x_k)}\cdots\varphi(x_N)\bigr\rangle\,.
\end{equation}
Assuming that the $N-2$-point correlation functions in the right hand side are
known, we can think of this identity as of the equation for the $N$-point
correlation function in the left hand side. The solution has the form of our
recursion relation \eqref{eq:WickRec}, where the function $D(x-x')$ solves the
equation
\[
\bigl(m^2-\Delta_x\bigr)\,D(x-x')=\delta^{(4)}(x-x')\,.
\]

The Wick's theorem can be derived in many other ways. Direct evaluation of the
Gaussian functional integral is performed in the \textbf{Section 9.2 of PS}.
Derivation of equivalent rules in the operator formalism is explained in
\textbf{Sect.~4.3 of PS}.

The rules of the Wick's theorem are conveniently represented in terms of
diagrams. With every factor $D(x-x')$ one associates a line connecting the
points $x$ and $x'$,
\[
D(x-x')=\contraction{}{\varphi(x)}{}{\varphi(x')}\varphi(x)\varphi(x')
=\ x\;\dia{\draw[prop] (0,0) node[vtx]{} -- (3,0) node[vtx]{};}\;x'\,.
\]
The diagram consists of the lines representing the contractions. For example
\begin{align*}
\bigl\langle\varphi(x_1)\varphi(x_2)\varphi(x_3)\varphi(x_4)\bigr\rangle
&=\contraction{}{\varphi(x_1)}{}{\varphi(x_2)}\varphi(x_1)\varphi(x_2)\,
\contraction{}{\varphi(x_3)}{}{\varphi(x_4)}\varphi(x_3)\varphi(x_4)
+\ \text{two other ways of contractions}\\[4pt]
&=\figfile{wick4}
\end{align*}
Likewise, $2N$-point correlation function is represented as the sum of
$(2N-1)!!$ diagrams involving $N$ lines.

% Source pp. 61-72 (Lectures7-12.pdf).
\section{Interaction. \texorpdfstring{$\varphi^4$}{phi**4} theory}

The KG theory describes free particles. Its action is quadratic in the field
variables. We want to add interactions. This is done by adding new terms to the
action. It should be done in a way that preserves the locality. We can modify
the Lagrangian density by adding local terms involving higher powers of the
field(s). The simplest example is known as the $\varphi^4$ theory, which we
consider next.

The $\varphi^4$ theory is described by the Euclidean action
\[
\mathcal A[\varphi]=\int \dd^4x\,\left[\frac12\,(\partial_\mu\varphi)^2+\frac{m^2}{2}\,\varphi^2+\frac{\lambda}{4!}\,\varphi^4\right]\,.
\]
Here $\lambda$ is a new parameter, the ``coupling constant''.

As in the KG theory, one can be interested in the correlation functions
\[
\bigl\langle\varphi(x_1)\cdots\varphi(x_N)\bigr\rangle=Z^{-1}\int\bigl[D\varphi\bigr]\,\varphi(x_1)\cdots\varphi(x_N)\,e^{-\mathcal A[\varphi]}\,,\qquad
Z=\int\bigl[D\varphi\bigr]\,e^{-\mathcal A[\varphi]}\,.
\]
Now the action is not quadratic in the integration variables $\varphi$, and
the integral cannot be evaluated in an explicit form.

\begin{exercise}
Performing the change of variables $\varphi(x)\to\varphi(x)+\varepsilon(x)$
in the functional integral defining the $N$-point correlation functions in
$\varphi^4$ theory, derive the equation
\begin{multline*}
\bigl(m^2-\partial_\mu^2\bigr)\bigl\langle\varphi(x)\,\varphi(x_1)\cdots\varphi(x_N)\bigr\rangle
+\frac{\lambda}{3!}\,\bigl\langle\varphi^3(x)\,\varphi(x_1)\cdots\varphi(x_N)\bigr\rangle\\
=\sum_{i=1}^N\delta^{(4)}(x-x_i)\,\bigl\langle\varphi(x_1)\cdots\widehat{\varphi(x_i)}\cdots\varphi(x_N)\bigr\rangle\,.
\end{multline*}
Equations of this kind are generally known as the \emph{Schwinger-Dyson
equations}.
\end{exercise}

Unlike the KG theory, the Schwinger-Dyson equation in the interacting theory is
not closed equation for the correlation function $\langle\varphi(x)\dots\rangle$,
since it involves $\langle\varphi^3(x)\dots\rangle$ which has no useful
expression in terms of the function $\langle\varphi(x)\dots\rangle$ itself.
Even the definition of the composite field $\varphi^3$ may be a non-trivial
problem (we are going to consider it later on).

Practical way of computations is the \emph{perturbation theory}, which
generates a power series expansion in the coupling parameter $\lambda$.

\subsection*{Perturbation Theory}

Perturbation theory allows one to evaluate the correlation functions as the
power series in the coupling parameter $\lambda$. One splits the action into
the free-field part (quadratic in the field variables $\varphi$), and the
interaction part, which includes the terms of higher order:
\[
\mathcal A=\mathcal A_0+\mathcal A_I\,.
\]
The exponential $\exp(-\mathcal A)$ is then expanded as
\[
e^{-\mathcal A}=e^{-\mathcal A_0}\left(1-\mathcal A_I+\frac12\,\mathcal A_I^2+\dots\right)\,.
\]
Since $\mathcal A_0$ is quadratic, each term can be evaluated using the Wick's
theorem.

To see how it works, consider the $\varphi^4$ theory. In this case the natural
way of splitting is
\[
\mathcal A_0=\mathcal A_{\mathrm{KG}}\,,\qquad \mathcal A_I=\frac{\lambda}{4!}\int\varphi^4(x)\dd^4x
\]
(later we will see that in actual calculations it is more convenient to include
certain parts of the quadratic action, the so called counterterms, into
$\mathcal A_I$, but at the moment we ignore this subtlety). With this
splitting, the power series expansion of $\exp(-\mathcal A_I)$ generates
expansion in the powers of $\lambda$.

Suppose we want to calculate such expansion of the two-point correlation
function $\langle\varphi(x_1)\varphi(x_2)\rangle$. By its definition
\[
\bigl\langle\varphi(x_1)\varphi(x_2)\bigr\rangle=Z^{-1}\int\bigl[D\varphi\bigr]\,\varphi(x_1)\varphi(x_2)\,e^{-\mathcal A_0}\,e^{-\mathcal A_I}\,.
\]
Expanding the last exponential as prescribed, we have
\[
Z^{-1}\int\bigl[D\varphi\bigr]\,\varphi(x_1)\varphi(x_2)\left(1-\frac{\lambda}{4!}\int\varphi^4(x)\dd^4x+O(\lambda^2)\right)\,e^{-\mathcal A_0}\,.
\]
The zero-order term is just the functional integral for the two-point function
of the KG theory. We can rewrite the previous expression as
\begin{equation}\label{eq:pt2}
\bigl\langle\varphi(x_1)\varphi(x_2)\bigr\rangle=\frac{Z_0}{Z}\left[\bigl\langle\varphi(x_1)\varphi(x_2)\bigr\rangle_0
-\frac{\lambda}{4!}\int\bigl\langle\varphi(x_1)\varphi(x_2)\varphi^4(x)\bigr\rangle_0\dd^4x+\dots\right]\,,
\end{equation}
where $Z_0=\int D[\varphi]\,e^{-\mathcal A_0}$ and $\langle\cdots\rangle_0$
stand for the partition function and correlation functions in the free theory.
This expression gives an example of the general situation: the perturbation
theory allows one to write the correlation functions of interacting theory as a
power series in the coupling parameter, the coefficients being certain
integrals of the correlation functions of the free theory. The role of the
factor $\frac{Z_0}{Z}$ will be clarified later.

The correlation functions of the free theory are computed using the Wick's
rules. In the above example, the first-order in $\lambda$ term involves the
$\dd^4x$ integral of the correlation function
\[
\bigl\langle\varphi(x_1)\varphi(x_2)\,\varphi^4(x)\bigr\rangle_0=\bigl\langle\varphi(x_1)\varphi(x_2)\varphi(x)\varphi(x)\varphi(x)\varphi(x)\bigr\rangle_0\,.
\]
This is a six-point correlation function of the free theory, with four points
coincident at $x$. According to the Wick's rule, it can be written as the sum
of
\[
5!!=15
\]
terms corresponding to different ways of pairwise contractions of the six field
insertions involved. Many different contractions lead to identical
contributions. There are two classes of such contractions with essentially
different contributions:
\begin{enumerate}[label=(\roman*)]
\item First, one can contract $\varphi(x_1)\varphi(x_2)\to
\contraction{}{\varphi(x_1)}{}{\varphi(x_2)}\varphi(x_1)\varphi(x_2)$, and then
contract four factors $\varphi(x)$ among themselves. There are three
possibilities for contracting. Say, one first contracts the first factor in
$\varphi(x)\varphi(x)\varphi(x)\varphi(x)$ with any one of the other three (3
ways), and once that one is chosen, there is only one way to contract the
remaining two insertions. We find that such contractions bring in the
contribution
\[
3\times
\contraction{}{\varphi(x_1)}{}{\varphi(x_2)}
\contraction{\varphi(x_1)\varphi(x_2)\,}{\varphi(x)}{}{\varphi(x)}
\contraction{\varphi(x_1)\varphi(x_2)\,\varphi(x)\varphi(x)\,}{\varphi(x)}{}{\varphi(x)}
\varphi(x_1)\varphi(x_2)\,\varphi(x)\varphi(x)\,\varphi(x)\varphi(x)
=3\,D(x_1-x_2)\,D(x-x)\,D(x-x)\,.
\]

\item Alternatively, one can contract $\varphi(x_1)$ with one of the four
identical factors $\varphi(x)$ (4 ways), and $\varphi(x_2)$ with any one of the
remaining three factors $\varphi(x)$ (3 ways); this leaves just two insertions
$\varphi(x)$, with only one way of contractions. Therefore there are
$4\times3=12$ identical contributions
\[
12\times
\contraction[2ex]{}{\varphi(x_1)}{\varphi(x_2)}{\varphi(x)}
\contraction{\varphi(x_1)}{\varphi(x_2)}{\varphi(x)}{\varphi(x)}
\contraction{\varphi(x_1)\varphi(x_2)\varphi(x)\varphi(x)}{\varphi(x)}{}{\varphi(x)}
\varphi(x_1)\varphi(x_2)\varphi(x)\varphi(x)\varphi(x)\varphi(x)
=12\,D(x_1-x)\,D(x_2-x)\,D(x-x)\,.
\]
\end{enumerate}
The resulting contribution $\sim\lambda$ in the brackets in \eqref{eq:pt2}
then reads
\begin{equation}\label{eq:pt2first}
-\frac{\lambda}{4!}\left[3\,D(x_1-x_2)\int \dd^4x\,D^2(x-x)+12\int \dd^4x\,D(x_1-x)\,D(x_2-x)\,D(x-x)\right]\,.
\end{equation}
Depicting contractions as solid lines
\[
\contraction{}{\varphi(x)}{}{\varphi(x')}\varphi(x)\varphi(x')
=\ x\;\dia{\pic{line};}\;x'\ \to\ D(x-x')
\]
the two terms are represented by two diagrams
\begin{equation}\label{eq:pt2diag}
3\times\Bigl[\,\dia{\pic{line}; \lab[left]{-1.4,0}{x_1} \lab[right]{1.4,0}{x_2} \pic at (3.4,0) {eight};}\,\Bigr]
+12\times\Bigl[\,\dia{\pic{tadline}; \lab[left]{-1.4,0}{x_1} \lab[right]{1.4,0}{x_2}}\,\Bigr]
\end{equation}
where the loop represents the contraction
\[
\contraction{}{\varphi(x)}{}{\varphi(x)}\varphi(x)\varphi(x)
=\dia{\draw[prop] (0,0.42) circle (0.42); \node[vtx] at (0,0) {};}\ \to\ D(x-x)\,.
\]
Recalling the overall factor
\[
-\frac{\lambda}{4!}=-\frac{\lambda}{24}
\]
in \eqref{eq:pt2}, we find that the two diagrams in \eqref{eq:pt2diag} come with
the factors $-\lambda/S_1$ and $-\lambda/S_2$, respectively, where the so
called ``symmetry factors'' here are
\[
\frac1{S_1}=\frac18\,,\qquad\frac1{S_2}=\frac12\,.
\]

Similar structure is found in the higher orders of the $\lambda$-expansion. The
$n$-th order term has the form
\begin{equation}\label{eq:ptn}
\frac1{n!}\left(-\frac{\lambda}{4!}\right)^n\left\langle\varphi(x_1)\varphi(x_2)\left[\int\varphi^4(y)\dd^4y\right]^n\right\rangle_0
\end{equation}
which involves the $n$-fold 4-space integration
\[
\left[\int\varphi^4(y)\dd^4y\right]^n=\int\varphi^4(y_1)\varphi^4(y_2)\cdots\varphi^4(y_n)\dd^4y_1\cdots \dd^4y_n\,.
\]
The correlation function in the integrand is that of the product of $2+4n$
factors $\varphi$, two at $x_1$ and $x_2$, and $4n$ at the integration points
$y_1,\dots,y_n$, four factors at each point.

Contractions of these field factors are represented by diagrams. The diagrams
are built from $2n+1$ contraction lines connecting two types of ``vertices'':

There are two ``external vertices'' associated with the points $x_1$ and
$x_2$, with exactly one of the contraction lines attached to each,
\[
\text{External points:}\qquad x_1\ \dia{\draw[prop] (0,0) node[vtx]{} -- (1.4,0);}
\qquad\qquad \dia{\draw[prop] (0,0) -- (1.4,0) node[vtx]{};}\ x_2
\]
and $n$ ``interaction vertices'', associated with $n$ integration points
$y_1,y_2,\dots,y_n$, with four contraction lines connected to each of them:
\[
\text{Interaction points:}\qquad
\dia{\pic{xvertex}; \lab[right]{0.3,-0.05}{y_1}}\quad
\dia{\pic{xvertex}; \lab[right]{0.3,-0.05}{y_2}}\quad\cdots\quad
\dia{\pic{xvertex}; \lab[right]{0.3,-0.05}{y_n}}
\]
Thus, in the order $\lambda^2$ the diagrams involve two external and two
interaction vertices, for example
\[
x_1\ \dia{\pic{sunline};}\ x_2
\]
Again, many distinct contractions are represented by identical diagrams. For
instance, the above diagram may represent the contraction
\[
\contraction{}{\varphi(x_1)}{}{\varphi(y_1)}
\contraction[4ex]{\varphi(x_1)\varphi(y_1)}{\varphi(y_1)}{\varphi(y_1)\varphi(y_1)\varphi(y_2)\varphi(y_2)}{\varphi(y_2)}
\contraction[3ex]{\varphi(x_1)\varphi(y_1)\varphi(y_1)}{\varphi(y_1)}{\varphi(y_1)\varphi(y_2)}{\varphi(y_2)}
\contraction[2ex]{\varphi(x_1)\varphi(y_1)\varphi(y_1)\varphi(y_1)}{\varphi(y_1)}{}{\varphi(y_2)}
\contraction{\varphi(x_1)\varphi(y_1)\varphi(y_1)\varphi(y_1)\varphi(y_1)\varphi(y_2)\varphi(y_2)\varphi(y_2)}{\varphi(y_2)}{}{\varphi(x_2)}
\varphi(x_1)\varphi(y_1)\varphi(y_1)\varphi(y_1)\varphi(y_1)\varphi(y_2)\varphi(y_2)\varphi(y_2)\varphi(y_2)\varphi(x_2)
\]
as well as any one of
\[
4\times4\times3!
\]
equivalent contractions related to this one by interchanging of the contracted
factors $\varphi$ inside $\varphi^4(y_1)$ and $\varphi^4(y_2)$. Every such
contraction corresponds to the same product of the propagators
\[
D(x_1-y_1)\,D^3(y_1-y_2)\,D(y_2-x_2)\,,
\]
and since the contributions involve the integrations over the interaction
points
\[
\int \dd^4y_1\dd^4y_2\,(\dots)
\]
the diagrams related by the interchange $y_1\leftrightarrow y_2$ yield
identical contributions. Therefore the above contraction comes with the
multiplicity factor
\[
2\cdot4\cdot4!\,.
\]
This factor almost entirely cancels the denominators in the overall factors
$(1/n!)(-\lambda/4!)^n$ in \eqref{eq:ptn}, which for $n=2$ is
\[
\frac1{2!}\left(-\frac{\lambda}{4!}\right)^2\,,
\]
leaving behind the factor $(-\lambda)^2/3!$. Finally, total contribution of
the contractions corresponding to the diagram in this example is
\[
\frac{(-\lambda)^2}{S}\int \dd^4y_1\dd^4y_2\,D(x_1-y_1)\,D^3(y_1-y_2)\,D(y_2-x_2)\,,
\]
with the ``symmetry factor'' $\frac1S=\frac1{3!}$.

This example illustrates the general rules of the diagrammatic representations
of the contributions in the perturbation theory, known as the

\subsection*{Feynman Rules (for \texorpdfstring{$\varphi^4$}{phi**4})}

To find $\lambda^n$ contribution to
\[
\bigl\langle\varphi(x_1)\cdots\varphi(x_N)\bigr\rangle
\]
\begin{enumerate}
\item Draw all diagrams with $N$ external vertices associated with the points
$x_1,\dots,x_N$, and with $n$ interaction vertices. Exclude disconnected
diagrams with ``vacuum fragments'' (see below).
\item For each diagram, associate
\begin{enumerate}[label=(\roman*)]
\item $x\ \dia{\pic{line};}\ x'=D(x-x')$;
\item $\dia{\pic{xvertex}; \lab[right]{0.3,-0.05}{y}}=-\lambda\int \dd^4y$;
\item Multiply by the ``symmetry factor'' $1/S$.
\end{enumerate}
\item Sum the contributions of all the diagrams.
\end{enumerate}

\begin{remark}
These rules are generalized in a straightforward way to interactions involving
any powers of the field $\varphi$ (possibly including terms linear and
quadratic in $\varphi$)
\begin{equation}\label{eq:genint}
\mathcal A_I=\int \dd^4y\,\left[\lambda_1(y)\,\varphi(y)+\frac{\lambda_2(y)}{2!}\,\varphi^2(y)+\frac{\lambda_3(y)}{3!}\,\varphi^3(y)+\frac{\lambda_4(y)}{4!}\,\varphi^4(y)+\cdots\right]\,.
\end{equation}
Here I allowed the coupling parameters $\lambda_k=\lambda_k(y)$ to be
arbitrary functions of $y$. (This generally breaks space-time translation and
rotation symmetries of the theory, but otherwise defines quite meaningful
quantity. Later, these generalizations will allow us to use diagrammatic
representations of the generating functional and effective action.) The
Feynman rules generalize to the case \eqref{eq:genint} by including $k$-leg
vertices
\begin{equation}\label{eq:kleg}
\begin{gathered}
\dia{\draw[prop] (-1.4,0) -- (0,0) node[vtx]{};}=-\int \dd^4y\,\lambda_1(y)\dots\,,\qquad
\dia{\draw[prop] (-1.4,0) -- (1.4,0); \node[vtx] at (0,0) {};}=-\int \dd^4y\,\lambda_2(y)\dots\,,\\[4pt]
\dia{\draw[prop] (-1,0.8) -- (0,0) -- (-1,-0.8); \draw[prop] (0,0) -- (1.3,0); \node[vtx] at (0,0) {};}=-\int \dd^4y\,\lambda_3(y)\dots\,,\qquad\cdots\,,\qquad
\dia{\foreach \a in {20,60,100,140,250,290} \draw[prop] (0,0) -- (\a:1.1);
     \draw[densely dotted] (165:0.7) arc[start angle=165,end angle=230,radius=0.7];
     \node[vtx] at (0,0) {}; \lab[below]{0.25,-0.9}{k}}=-\int \dd^4y\,\lambda_k(y)\dots\,.
\end{gathered}
\end{equation}
\end{remark}

\subsection*{Vacuum Diagrams}

Returning back to the homogeneous $\varphi^4$ theory, let us take another look
at the diagrams contributing to the two-point function. We write
\begin{equation}\label{eq:twofactors}
\bigl\langle\varphi(x_1)\varphi(x_2)\bigr\rangle
=\left[\frac1{Z_0}\int\bigl[D\varphi\bigr]\,\varphi(x_1)\varphi(x_2)\,e^{-\mathcal A_0-\mathcal A_I}\right]\times\left[\frac{Z_0}{Z}\right]\,,
\end{equation}
where $Z_0$ is the partition function of the free theory defined by
$\mathcal A_0=\mathcal A_{\mathrm{KG}}$. The diagrams we wrote down above
actually correspond to the first factor in \eqref{eq:twofactors}. For this
first factor we have
\[
\underbrace{\dia{\pic{line}; \lab[left]{-1.4,0}{x_1} \lab[right]{1.4,0}{x_2}}}_{(a)}
+\frac12\,\Bigl(\underbrace{\dia{\pic{tadline}; \lab[left]{-1.4,0}{x_1} \lab[right]{1.4,0}{x_2}}}_{(b)}\Bigr)
+\frac18\,\Bigl(\underbrace{\dia{\pic{line}; \lab[left]{-1.4,0}{x_1} \lab[right]{1.4,0}{x_2} \pic at (3.3,0) {eight};}}_{(c)}\Bigr)+\cdots\,.
\]
Observe that there are diagrams containing disconnected pieces, like the
diagram (c) above. This diagram has a fragment
\[
\dia{\pic{eight};}
\]
which is not connected to any of the ``external'' points ($x_1,x_2$ in this
case). Such fragments are called ``vacuum diagrams'' (or ``vacuum
fragments''). The value
\[
D(x_1-x_2)\times\left(-\frac{\lambda}{8}\int \dd^4x\,D(x-x)\,D(x-x)\right)
\]
associated with the diagram (c) above is the product of the values of its
parts,
\begin{equation}\label{eq:eightvalue}
\dia{\pic{line}; \lab[left]{-1.4,0}{x_1} \lab[right]{1.4,0}{x_2}}=D(x_1-x_2)\,;\qquad
\dia{\pic{eight};}=-\frac{\lambda}{8}\int \dd^4x\,D^2(x-x)\,.
\end{equation}
As follows directly from the Feynman rules above, this holds in general: the
value of any disconnected diagram is the product of the values of its
connected fragments. In fact, in computing the correlation function
\eqref{eq:twofactors} the vacuum fragments should be ignored because their
contributions are canceled exactly by the second factor in the r.h.s.\ of
\eqref{eq:twofactors}, which contains the partition function $Z$ of the full
theory. To see this, consider again the diagram (c) above. In higher orders of
the perturbation theory we obtain infinitely many diagrams in which the vacuum
fragment in (c) is repeated many times,
\[
\Bigl(\dia{\pic{line}; \lab[left]{-1.4,0}{x_1} \lab[right]{1.4,0}{x_2}}\quad
\underbrace{\dia{\pic{eight}; \pic at (1.3,0) {eight}; \node at (2.25,0) {$\cdots$}; \pic at (3.2,0) {eight};}}_{n\ \text{times}}\Bigr)
\]
The value corresponding to $n$ copies of the vacuum fragment is just
\[
\underbrace{\dia{\pic{eight}; \pic at (1.3,0) {eight}; \node at (2.25,0) {$\cdots$}; \pic at (3.2,0) {eight};}}_{n\ \text{times}}
=\frac1{n!}\,\Bigl[\dia{\pic{eight};}\Bigr]^n\,,
\]
the factor $1/n!$ being the extra symmetry factor associated with the
permutations of the fragments. The sum over $n$ therefore yields
\[
\Bigl(\dia{\pic{line}; \lab[left]{-1.4,0}{x_1} \lab[right]{1.4,0}{x_2}}\Bigr)\times\exp\Bigl(\dia{\pic{eight};}\Bigr)
\]
This simple analysis can be repeated with any other vacuum fragment. As the
result we obtain for the first factor in \eqref{eq:twofactors}
\begin{multline*}
\Bigl(\underbrace{\dia{\pic{line}; \lab[left]{-1.4,0}{x_1} \lab[right]{1.4,0}{x_2}}
+\dia{\pic{tadline}; \lab[left]{-1.4,0}{x_1} \lab[right]{1.4,0}{x_2}}
+\dia{\pic{sunline}; \lab[left]{-1.7,0}{x_1} \lab[right]{1.7,0}{x_2}}+\cdots}_{\text{sum of all connected diagrams}}\Bigr)\\
\times\exp\Bigl(\underbrace{\dia{\pic{eight};}+\dia{\pic{melon};}+\dia{\pic{chain3};}+\cdots}_{\mathclap{\text{sum of all connected vacuum diagrams}}}\Bigr)
\end{multline*}
Consider now the second factor in \eqref{eq:twofactors},
\[
\frac Z{Z_0}=\frac1{Z_0}\int\bigl[D\varphi\bigr]\,e^{-\mathcal A_0}\,e^{-\mathcal A_I}\,.
\]
Expanding in $\mathcal A_I$ we find that this quantity is given by the sum of
all vacuum diagrams, which again exponentiates in terms of the connected
vacuum diagrams. That is
\begin{equation}\label{eq:ZZ0}
Z/Z_0=\exp\Bigl(\text{sum of all connected vacuum diagrams}\Bigr)
\end{equation}
so that the second factor in \eqref{eq:twofactors} cancels the contributions of
all vacuum fragments. Thus, if we are interested in the correlation function,
the factor $Z^{-1}$ in the definition of the correlation function can indeed be
ignored if we simultaneously drop all the diagrams with disconnected vacuum
fragments.

Note that the above analysis of the ``vacuum fragments'' generalizes
straightforwardly to the case of interaction \eqref{eq:genint} with arbitrary
$x$-dependent coupling constants. In particular it is still true that the
partition function of this more general theory exponentiates in terms of the
connected vacuum diagrams, as in \eqref{eq:ZZ0}. Of course in this case the
diagrams involve all the vertices in \eqref{eq:kleg}.

Let me comment on the interpretation of the above formula \eqref{eq:ZZ0} for
$Z$. The first of the connected vacuum diagrams is given by the vacuum
fragment in \eqref{eq:eightvalue}, i.e.
\[
\dia{\pic{eight};}=-\frac{\lambda}{8}\,D^2(0)\int \dd^4x
\]
Clearly, it contains an infinite factor
\[
\int \dd^4x=V^{(4)}\,,
\]
where $V^{(4)}$ is the volume of the 4-dimensional space. This is the common
feature of the connected vacuum diagrams
\[
\dia{\pic{eight};}\,,\qquad\dia{\pic{melon};}\,,\qquad\dia{\pic{chain3};}\,,\qquad\dots
\]
As a consequence of translational invariance each of them contains this volume
factor. Therefore
\begin{equation}\label{eq:Zvac}
Z=Z_0\,\exp\bigl(-\varepsilon_I\,V^{(4)}\bigr)
\end{equation}
where $\varepsilon_I$ is just the sum of all connected vacuum diagrams with the
volume factored out (and with the sign reversed).

The functional integral
\[
Z=\int\bigl[D\varphi\bigr]\,e^{-\mathcal A}
\]
can be interpreted as a (configuration) partition function of a classical
statistical mechanics. The classical system in this interpretation is a field
$\varphi(x)$ in 4-dimensional (Euclidean) space, with the potential energy
$U[\varphi]$,
\[
\beta\,U[\varphi]:=\frac1\hbar\,\mathcal A[\varphi]\,,
\]
where $\beta$ is interpreted as the inverse temperature. ($\beta$ can be set
equal to 1, again by suitable choice of the energy units.) This is a spatially
homogeneous system with infinitely many degrees of freedom associated with all
points in the 4-dimensional space. Naturally, its free energy is an intensive
quantity proportional to $V^{(4)}$. We see that the quantity
\[
\varepsilon_0+\varepsilon_I\,,\qquad\text{where}\qquad \varepsilon_0=-\log Z_0/V^{(4)}
\]
is interpreted as the specific (per unit volume) free energy of this system,
the first term being the specific free energy of the free theory and
$\varepsilon_I$ incorporating all corrections due to the interaction.

Another way to look at the formula \eqref{eq:ZZ0} is to recall the relation
between the path integrals and the matrix elements of the (imaginary) time
evolution operator. The path integral defining $Z$ can be thought of as the
$T\to\infty$ limit of the path integral
\[
\int\bigl[D\varphi\bigr]_{\substack{\varphi(\xx,-T/2)=\varphi_i(\xx)\\\varphi(\xx,+T/2)=\varphi_f(\xx)}}\,e^{-\mathcal A}
\]
which represents the matrix element
\[
\bra{\varphi_f(\xx)}e^{-T\hat H}\ket{\varphi_i(\xx)}\,.
\]
Here $\varphi_i$ and $\varphi_f$ are arbitrary ``initial'' and ``final''
field configurations. The evolution operator inside can be written as
\[
e^{-T\hat H}=\sum_n\ket n\bra n\,e^{-TE_n}\,,
\]
where $\ket n$ is a complete set of the stationary states; its $T\to\infty$
limit then is
\[
e^{-E_0T}\,\ket0\bra0\,,\qquad\text{where}\qquad E_0=\varepsilon\,V^{(3)}
\]
($V^{(3)}$ is the volume of the 3-dimensional ``equal time'' section of our
4-dimensional space). Because $V^{(4)}=T\times V^{(3)}$, comparison with the
above formula \eqref{eq:ZZ0} shows that vacuum diagrams represent the
expansion of the ground state energy in the powers of $\lambda$ (hence the term
``vacuum diagrams''). We will see later that all the vacuum diagrams are highly
divergent at short distances. This is similar to the short-distance (high
momenta) divergence of the vacuum energy of the Klein-Gordon theory which we
observed before. Proper understanding of the physical content of these
divergent terms is important in certain problems like the Casimir effect or
interaction with gravity.

% Source pp. 73-79 (Lectures7-12.pdf).
\section{More on Feynman Diagrams. Generating Functional}

Continuing our analysis of the diagrammatic representation of the correlation
functions let us consider the diagram expansion of the four-point function
\[
\bigl\langle\varphi(x_1)\varphi(x_2)\varphi(x_3)\varphi(x_4)\bigr\rangle:=Z^{-1}\int\bigl[D\varphi\bigr]\,\varphi(x_1)\varphi(x_2)\varphi(x_3)\varphi(x_4)\,e^{-\mathcal A[\varphi]}
\]
in the $\varphi^4$ theory. First, by applying the same arguments as before it
is easy to show that in this case (and indeed for any correlation function)
the ``vacuum fragments'' again get canceled by the expansion of $Z$ in the
denominator. However, in this case the disconnected diagrams remain even after
we drop the vacuum fragments. Consider the first few diagrams for the above
four-point function, shown in figure~\ref{fig:fourpt}.
\begin{figure}[ht]
\centering\figfile{fourpt}
\caption{The first few diagrams for the four-point function in the $\varphi^4$
theory. (a)--(c) are the free-theory terms, (d) is a disconnected diagram
with a correction to one of the propagators, and (e), (f) are connected.}
\label{fig:fourpt}
\end{figure}

The first three diagrams here represent the four point function in the free
theory; they consist of the propagators which being connected to the external
points are not connected to each other. These contributions are expressed in a
simple way through the free-field two-point functions $D$. The diagram (d) and
similar diagrams are also disconnected; they represent corrections to those
two-point functions. The diagrams (e) and (f) represent the class of truly
connected diagrams. By the structure of the above diagram series it is natural
to write
\begin{multline*}
\bigl\langle\varphi(x_1)\varphi(x_2)\varphi(x_3)\varphi(x_4)\bigr\rangle
=\bigl\langle\varphi(x_1)\varphi(x_2)\bigr\rangle\bigl\langle\varphi(x_3)\varphi(x_4)\bigr\rangle
+\bigl\langle\varphi(x_1)\varphi(x_3)\bigr\rangle\bigl\langle\varphi(x_2)\varphi(x_4)\bigr\rangle\\
+\bigl\langle\varphi(x_1)\varphi(x_4)\bigr\rangle\bigl\langle\varphi(x_2)\varphi(x_3)\bigr\rangle
+\bigl\langle\varphi(x_1)\varphi(x_2)\varphi(x_3)\varphi(x_4)\bigr\rangle_{\text{conn}}
\end{multline*}
where the first three terms combine the contributions of the disconnected
diagrams; they are expressed in terms of the two-point functions. The connected
diagrams are incorporated in the last term which is called the
\emph{connected} correlation function.

Similar pattern is observed for the higher-order correlation functions. There
are disconnected diagrams which sum up into the lower-order correlation
functions, and besides there are truly connected ones, their sums define the
connected correlation functions.

In fact the disconnected diagrams can be efficiently sorted out by using the so
called Generating Functional.

\subsection*{Generating Functional (PS Sect.~9.2, 9.3)}

Perturbation theory represents any correlation function as the sum of the
Feynman diagrams. There are connected and disconnected diagrams. The vacuum
bubbles can be ignored as they are canceled out by the expansion of the
normalization factor $Z^{-1}$. The remaining disconnected diagrams are
efficiently sorted out by using the method of Generating Functional. Although
the following analysis is done for the case of $\varphi^4$ theory, it is useful
to keep in mind that all the arguments below can be repeated step by step for
the generic interaction.

Let us define the Generating Functional as follows. Introduce the functional
\begin{equation}\label{eq:zetaJ}
\zeta[J]=\sum_{n=0}^\infty\frac1{n!}\int d^4x_1\cdots d^4x_n\,\bigl\langle\varphi(x_1)\cdots\varphi(x_n)\bigr\rangle\,J(x_1)\cdots J(x_n)\,.
\end{equation}
It depends on the auxiliary function $J$, and incorporates all $n$-point
correlation functions. By the definition
\begin{equation}\label{eq:zetader}
\bigl\langle\varphi(x_1)\cdots\varphi(x_n)\bigr\rangle=\frac{\delta}{\delta J(x_1)}\cdots\frac{\delta}{\delta J(x_n)}\,\zeta[J]\Big|_{J=0}\,.
\end{equation}
The generating functional can be written as
\[
\zeta[J]=\Bigl\langle\exp\Bigl(\int J(x)\varphi(x)\,d^4x\Bigr)\Bigr\rangle
\]
and therefore it is expressed as the ratio
\[
\zeta[J]=Z[J]/Z[0]\,;
\]
where
\begin{equation}\label{eq:ZJ}
Z[J]=\int\bigl[D\varphi\bigr]\,e^{\int J(x)\varphi(x)\,d^4x}\,e^{-\mathcal A[\varphi]}
\end{equation}
is a partition function associated with the action
\[
\mathcal A_J=\mathcal A-\int J(x)\varphi(x)\,d^4x\,.
\]
For this reason the function $J(x)$ is usually referred to as an ``external
source''.

The source-dependent action above is a particular case of generic theory with
$x$-dependent coupling $\lambda_1(x)=-J(x)$,
\[
\mathcal A_I=\int d^4x\,\Bigl[-J(x)\varphi(x)+\frac{\lambda}{4!}\,\varphi^4(x)\Bigr]\,.
\]
Therefore its partition function $Z[J]$ (more precisely $Z[J]/Z_0$) is given by
a sum of all vacuum diagrams which now contain two kinds of vertices,
\[
\dia{\draw[prop] (-1.6,0) -- (0,0) node[vtx]{};}=\int d^4x\,J(x)\dots\,,\qquad\qquad
\dia{\pic{xvertex};}=-\lambda\int d^4x\dots\,.
\]
By taking the variational derivatives with respect to $J(x)$ we just convert
the one-leg vertices above into the external legs; this is the meaning of the
formula \eqref{eq:zetader} above. Now, the sum of all the vacuum diagrams
exponentiates in terms of the sum of connected vacuum
diagrams.\footnote{Formal proof of this statement can be found e.g.\ in
``Quantum Field Theory and Critical Phenomena'' by J.~Zinn-Justin, Oxford
Univ.\ Press 1996, \S6.1.} Defining
\begin{equation}\label{eq:Wdef}
\zeta[J]=\exp W[J]
\end{equation}
we conclude that the expansion of this new functional $W[J]$ in terms of $J$
will generate the \emph{connected} correlation functions, namely
\[
W[J]=\sum_{n=0}^\infty\frac1{n!}\int d^4x_1\cdots d^4x_n\,W^{(n)}(x_1,\dots,x_n)\,J(x_1)\cdots J(x_n)\,,
\]
where
\[
W^{(n)}(x_1,\dots,x_n)=\bigl\langle\varphi(x_1)\cdots\varphi(x_n)\bigr\rangle_{\text{conn}}\,.
\]
The functional $W[J]$ is called the generating functional for the
\emph{connected} correlation functions. In the KG theory one can evaluate the
Gaussian functional integral \eqref{eq:ZJ} for $Z[J]$ explicitly and find
\[
W_{\mathrm{KG}}[J]=\frac12\int d^4x\,d^4x'\,J(x)\,D(x-x')\,J(x')\,,
\]
which shows again that in the free theory all connected correlation functions
but the two-point one vanish.

\begin{exercise}
Using $\zeta[J]=\exp W[J]$, check by direct calculation that
\begin{align*}
W^{(2)}(x_1,x_2)&=\bigl\langle\varphi(x_1)\varphi(x_2)\bigr\rangle\,,\\
W^{(4)}(x_1,\dots,x_4)&=\bigl\langle\varphi(x_1)\cdots\varphi(x_4)\bigr\rangle
-\bigl\langle\varphi(x_1)\varphi(x_2)\bigr\rangle\bigl\langle\varphi(x_3)\varphi(x_4)\bigr\rangle\\
&\qquad-\text{two other disconnected terms.}
\end{align*}
\end{exercise}

In many cases it is useful to decompose the connected correlation functions
further into more elementary ``blocks''. Consider the diagrams for the connected
four-point function
\[
\bigl\langle\varphi(x_1)\cdots\varphi(x_4)\bigr\rangle_c=\figfile{fourptconn}
\]
It is clear that some of these diagrams (like (b), (c), and (e)) represent the
corrections to the ``external legs''. This shows that the connected correlation
functions can be expressed through so-called ``amputated'' correlation
functions,
\[
W^{(n)}(x_1,\dots,x_n)=\int\Bigl[\prod_{i=1}^nd^4y_i\,W(x_i-y_i)\Bigr]\,W^{(n)}_{\text{amp}}(y_1,\dots,y_n)
\]
where I used the shorthand
\[
W(x-y)\equiv W^{(2)}(x,y)=\bigl\langle\varphi(x)\varphi(y)\bigr\rangle\,.
\]
Let us give a graphical representation of this relation. Denoting the connected
$n$-point correlation function by an empty blob with $n$ legs,
\[
\bigl\langle\varphi(x_1)\cdots\varphi(x_n)\bigr\rangle_{\text{conn}}=
\diac{\pic{legblob={0.8}{0,45,90,135}{1.1}{fill=white}};
  \draw[densely dotted] (165:1.3) arc[start angle=165,end angle=330,radius=1.3];
  \node at (300:1.75) {\small$n$};}\,,
\]
we have
\[
\diac{\pic{legblob={0.8}{0,45,90,135}{1.1}{fill=white}};
  \draw[densely dotted] (165:1.3) arc[start angle=165,end angle=330,radius=1.3];
  \node at (300:1.75) {\small$n$};}
=\diac{\foreach \a in {0,45,90,135} {\draw[prop] (\a:0.9) -- (\a:2.5) node[vtx] {}; \draw[prop,fill=white] (\a:1.7) circle (0.3);}
  \draw[prop,fill=white] (0,0) circle (0.9); \node at (0,0) {\small amp};
  \draw[densely dotted] (165:1.4) arc[start angle=165,end angle=330,radius=1.4];
  \node at (300:1.85) {\small$n$};}\;;\qquad\qquad
x\ \dia{\draw[prop] (-1.6,0) node[vtx]{} -- (1.6,0) node[vtx]{}; \draw[prop,fill=white] (0,0) circle (0.4);}\ x'=W^{(2)}(x-x')\,.
\]
According to this definition the amputated correlation functions contain
neither external legs nor the corrections to them. Note that by this definition
$W^{(2)}_{\text{amp}}(x,y)=S(x-y)$ is inverse (in the operator sense) to the
two-point function, i.e.
\[
\int d^4y\,S(x-y)\,W(y-x')=\delta^{(4)}(x-x')\,.
\]
In the free theory this is the inverse of the function $D$ i.e.\ just the
(Euclidean) Klein-Gordon operator
\[
S_0(x-y)=\bigl(m^2-\Delta_x\bigr)\,\delta^{(4)}(x-y)\,.
\]

\paragraph{Proper vertices:} Finally, let us define the so called \emph{proper
vertices} (also referred to as the ``one-particle irreducible correlation
functions''). For $n>2$ the $n$-point proper vertex
$-\Gamma^{(n)}(y_1,\dots,y_n)$ is the sum of all diagrams contributing to
$W^{(n)}_{\text{amp}}(y_1,\dots,y_n)$ which cannot be made disconnected by
cutting just one line (please pay attention to the minus sign in this
definition). In particular, in the $\varphi^4$ theory
$\Gamma^{(4)}=-W^{(4)}_{\text{amp}}$ (why?), i.e.\ diagrammatically
\[
\diac{\pic{legblob={0.8}{45,135,225,315}{1.4}{fill=white}};}
=-\ \diac{\pic{dressed={0.8}{45,135,225,315}};}\,,
\]
but for the 6-point functions we have in graphic representation
\[
\diac[0.42]{\pic{legblob={0.8}{0,60,120,180,240,300}{1.4}{fill=white}};}
=-\ \diac[0.42]{\pic{dressed={0.8}{0,60,120,180,240,300}};}
+\Bigl[\ \diac[0.42]{\pic{dressed={0.8}{60,120,180,240,300}};
  \draw[prop] (0.8,0) -- (2.4,0); \draw[prop,fill=white] (1.6,0) circle (0.3);
  \begin{scope}[xshift=3.2cm]\pic{dressed={0.8}{0,60,120,240,300}};\end{scope}
  \node[vtx,minimum size=4pt] at (0.8,0) {}; \node[vtx,minimum size=4pt] at (2.4,0) {};}
\ +\ 9\ \text{similar terms}\Bigr]\,,
\]
where each shadowed blob with $n$ fat points on it stands for the proper vertex
$\Gamma^{(n)}$.

For the reasons that will become clear later the two-point proper vertex
deserves special definition, namely
\[
\Gamma^{(2)}(x,x')=S(x-x')\,.
\]
Its relation to the one-particle irreducible diagrams is as follows. Write
\[
S(x-x')=S_0(x-x')+\Sigma(x-x')\,.
\]
The function $\Sigma(x-x')$ is often referred to as the mass operator. It is
the function $-\Sigma$ (again note the sign) which is equal to the sum of all
one-particle irreducible diagrams for the two-point correlation function.
Indeed, denoting $-\Sigma$ by the two-point shadowed blob we have
\[
x\ \dia{\draw[prop,pattern=north east lines] (0,0) circle (0.6); \node[vtx,minimum size=4pt] at (-0.6,0) {}; \node[vtx,minimum size=4pt] at (0.6,0) {};}\ x'
=\underset{x,x'}{\dia{\draw[prop] (0,0.45) circle (0.45); \node[vtx] at (0,0) {};}}
+\underset{x,x'}{\dia{\draw[prop] (0,0.45) circle (0.45); \draw[prop] (0,1.35) circle (0.45); \node[vtx] at (0,0) {}; \node[vtx] at (0,0.9) {};}}
+\ x\ \dia{\pic{melon};}\ x'+\cdots\,.
\]
The full set of diagrams contributing to $W$,
\[
W(x-x')=x\,\dia{\pic{line};}\,x'
+x\,\dia{\pic{tadline};}\,x'
+x\,\dia{\draw[prop] (-2.2,0) -- (2.2,0); \draw[prop] (-0.8,0.42) circle (0.42); \draw[prop] (0.8,0.42) circle (0.42);
  \foreach \p in {-2.2,-0.8,0.8,2.2} \node[vtx] at (\p,0) {};}\,x'+\cdots\,,
\]
can be obtained by repeating the one-particle irreducible diagrams in a
chainlike manner,
\[
\dia{\draw[prop] (-1.6,0) node[vtx]{} -- (1.6,0) node[vtx]{}; \draw[prop,fill=white] (0,0) circle (0.5); \node[above] at (0,0.5) {\small$W$};}
=\dia{\pic{line};}
+\dia{\draw[prop] (-1.8,0) node[vtx]{} -- (1.8,0) node[vtx]{}; \draw[prop,pattern=north east lines] (0,0) circle (0.5);
  \node[vtx,minimum size=4pt] at (-0.5,0) {}; \node[vtx,minimum size=4pt] at (0.5,0) {};}
+\dia{\draw[prop] (-2.4,0) node[vtx]{} -- (2.4,0) node[vtx]{};
  \foreach \c in {-0.9,0.9} {\draw[prop,pattern=north east lines] (\c,0) circle (0.45);
  \node[vtx,minimum size=4pt] at (\c-0.45,0) {}; \node[vtx,minimum size=4pt] at (\c+0.45,0) {};}}
+\cdots\,.
\]
Explicitly,
\[
W=D-D*\Sigma*D+D*\Sigma*D*\Sigma*D-\cdots\,.
\]
Here $A*B$ denotes a convolution
\[
(A*B)(x-x')=\int d^4y\,A(x-y)\,B(y-x')\,.
\]
Using
\[
S_0*D=I
\]
where $I=\delta^{(4)}(x-y)$ is the unit operator, it is straightforward to check
that $S_0+\Sigma$ is indeed an inverse of $W$.

Similarly to the above generating functionals $\zeta$ and $W$ we can define a
generating functional of the proper vertices. This functional, denoted
$\Gamma[\phi]$, is called the ``Effective Action''.

% Source pp. 80-92 (Lectures7-12.pdf).
\section[Effective Action. Loop expansion]{Effective Action. Loop expansion\protect\footnote{PS Sect.~11.3--11.5.}}

The generating functional for the proper vertices
\begin{equation}\label{eq:Gammadef}
\Gamma[\phi]=\sum_{n=2}^\infty\frac1{n!}\int \dd^4x_1\cdots \dd^4x_n\,\Gamma^{(n)}(x_1,\dots,x_n)\,\phi(x_1)\cdots\phi(x_n)
\end{equation}
is called the \emph{Effective Action}. Here $-\Gamma^{(n)}(x_1,\dots,x_n)$ are
defined as the sums of one-particle irreducible diagrams for $n$-point
correlation functions, with amputated legs. It turns out that $\Gamma[\phi]$ is
related to the functional $W[J]$ by the functional version of the
\emph{Legendre transform}.

The Legendre transform must be familiar from classical mechanics and classical
thermodynamics. In the latter context it relates different thermodynamic
functions. For example, consider a magnet which in the absence of an external
field is described by a Hamiltonian $\Ham$, and let $\mu$ be its magnetic
moment. If we add an external field $B$, its internal energy is
\[
U(B)=\langle\Ham\rangle-MB
\]
where $M:=\langle\mu\rangle$ is the magnetization,
\[
M=-\frac{\partial F}{\partial B}
\]
and $F=U-TS$ is the free energy. One can solve the above equation for $B$:
$B=B(M)$, and define the thermodynamic potential (Gibbs free energy)
\[
\Phi(M)=\Bigl[F(B)+MB\Bigr]_{B=B(M)}=\langle\Ham\rangle_{B(M)}-TS\,,
\qquad\text{so that}\qquad
B(M)=\frac{\partial\Phi(M)}{\partial M}\,.
\]
In particular the equation
\[
\frac{\partial\Phi}{\partial M}=0
\]
determines spontaneous magnetization in the absence of external field. The
transformation relating $\Phi(M)$ to $F(B)$ is the \emph{Legendre transform}.

Note that the above functional $Z[J]$ is interpreted as the partition function
of a field statistical system with the Hamiltonian
\[
\mathcal A[\varphi]-\int J(x)\varphi(x)\dd^4x\,,
\]
i.e.\ the source $J(x)$ plays the role of the external field in the above
example. The generating functional $W[J]=\log\zeta[J]$ is then interpreted as
minus free energy,
\[
W[J]=-F[J]\,.
\]
In particular
\begin{equation}\label{eq:phiJ}
\phi(x,[J])\equiv\bigl\langle\varphi(x)\bigr\rangle_J=\frac{\delta W[J]}{\delta J(x)}\,,
\end{equation}
where $\langle\cdots\rangle_J$ denotes the expectation value in the presence
of the source $J$. Let us assume that the above equation can be solved for $J$,
defining the inverse functional
\[
J=J(x,[\phi]):\qquad J\bigl(x,[\phi(x,[J_0(x)])]\bigr)=J_0(x)\,,\qquad
\phi\bigl(x,[J(x,[\phi_0])]\bigr)=\phi_0(x)\,,
\]
for any functions $J_0(x)$ or $\phi_0(x)$. One can define then the analog of
the thermodynamic potential
\[
\Gamma[\phi]=\left[-W[J]+\int J(x)\phi(x)\dd^4x\right]_{J(x)=J(x,[\phi])}
\]
It is not difficult to show that for $\Gamma[\phi]$ thus defined the relation
holds
\begin{equation}\label{eq:JGamma}
J(x,[\phi])=\frac{\delta\Gamma[\phi]}{\delta\phi(x)}\,.
\end{equation}
Indeed, taking the variation of
\[
\Gamma[\phi]+W[J]-\int J(x)\phi(x)\dd^4x=0
\]
(where it is understood that $\phi$ and $J$ are related through
\eqref{eq:phiJ}) we have
\[
\int \dd^4x\,\left[\frac{\delta\Gamma}{\delta\phi(x)}\,\delta\phi(x)+\frac{\delta W}{\delta J(x)}\,\delta J(x)-\phi(x)\,\delta J(x)-J(x)\,\delta\phi(x)\right]=0
\]
The second and the third terms here cancel each other in view of
\eqref{eq:phiJ} and we obtain \eqref{eq:JGamma}.

It can be shown that the functional $\Gamma[\phi]$ coincides with the
generating functional of the proper vertices. Namely, its functional Taylor
expansion in $\phi$, Eq.~\eqref{eq:Gammadef}, has the proper vertices
$\Gamma^{(n)}$ as its coefficient functions. Let us see how it works for
lower-order correlation functions, using $\varphi^4$ theory as an example. We
have
\begin{multline*}
W[J]=\frac12\int \dd^4x_1\dd^4x_2\,W(x_1-x_2)\,J(x_1)J(x_2)\\
+\frac1{4!}\int\left[\prod_{i=1}^4\dd^4x_i\right]\,W^{(4)}(x_1,\dots,x_4)\,J(x_1)J(x_2)J(x_3)J(x_4)+O(J^6)
\end{multline*}
and
\begin{multline*}
\Gamma[\phi]=\frac12\int \dd^4x_1\dd^4x_2\,\Gamma^{(2)}(x_1-x_2)\,\phi(x_1)\phi(x_2)\\
+\frac1{4!}\int\left[\prod_{i=1}^4\dd^4x_i\right]\,\Gamma^{(4)}(x_1,\dots,x_4)\,\phi(x_1)\phi(x_2)\phi(x_3)\phi(x_4)+O(\phi^6)\,.
\end{multline*}
Using the equation $\phi(x,[J])=\delta W/\delta J(x)$ we can express $\phi$ in
terms of $J$ as
\begin{multline*}
\phi(x,[J])=\int \dd^4x_1\,W(x-x_1)\,J(x_1)\\
+\frac1{3!}\int\left[\prod_{i=1}^3\dd^4x_i\right]\,W^{(4)}(x,x_1,x_2,x_3)\,J(x_1)J(x_2)J(x_3)+O(J^5)\,.
\end{multline*}
Similarly, using $J(x,[\phi])=\delta\Gamma/\delta\phi(x)$, $J$ is expressed
through $\phi$ as
\begin{multline*}
J(x,[\phi])=\int \dd^4x_1\,\Gamma^{(2)}(x-x_1)\,\phi(x_1)\\
+\frac1{3!}\int\left[\prod_{i=1}^3\dd^4x_i\right]\,\Gamma^{(4)}(x,x_1,x_2,x_3)\,\phi(x_1)\phi(x_2)\phi(x_3)+O(\phi^5)\,.
\end{multline*}
These two relations must be inverse one to another. Substituting the first into
the second we have, schematically
\begin{multline*}
J(x)=\int_{y,y'}\Gamma^{(2)}(x-y)\,W(y-y')\,J(y')\\
+\frac1{3!}\int_{y,y_1,y_2,y_3}\Gamma^{(2)}(x-y)\,W^{(4)}(y,y_1,y_2,y_3)\,J(y_1)J(y_2)J(y_3)\\
+\frac1{3!}\int_{\substack{y'_1,y'_2,y'_3\\y_1,y_2,y_3}}\Gamma^{(4)}(x,y'_1,y'_2,y'_3)\,W(y'_1-y_1)\,W(y'_2-y_2)\\
\times W(y'_3-y_3)\,J(y_1)J(y_2)J(y_3)+O(J^5)\,.
\end{multline*}
As this must be an identity, we find again
\[
\Gamma^{(2)}(x-x')=S(x-x')\,,
\]
and
\[
W^{(4)}(x_1,\dots,x_4)=-\int_{y_1,\dots,y_4}\Gamma^{(4)}(y_1,\dots,y_4)\,W(y_1-x_1)\,W(y_2-x_2)\,W(y_3-x_3)\,W(y_4-x_4)\,,
\]
which coincide with our definitions of the proper vertices $\Gamma^{(2)}$ and
$\Gamma^{(4)}$.

\begin{exercise}
Carry out the Legendre transform to the order 6 in $J$ and $\phi$. Find
$W^{(6)}$ in terms of $\Gamma^{(6)}$, $\Gamma^{(4)}$, $\Gamma^{(2)}$.
\end{exercise}

\subsection*{Semiclassical (or Loop) Expansion}

There is a useful way to reorganize the diagrams contributing to $W[J]$ or
$\Gamma[\phi]$ known as the loop expansion. Generally, classical limit appears
as the leading term if one tries to estimate the functional integral using the
saddle point approximation. Then, the integral is dominated by the field
configurations close to the stationary point determined by solving the
equation $\delta\mathcal A[\varphi]=0$, i.e.\ the equations of motion of
classical field theory defined by the action $\mathcal A[\varphi]$. With this
in mind, let us restore the quantum parameter $\hbar$ (previously set to 1) in
the definition of the generating functional
\begin{equation}\label{eq:ZJhbar}
Z[J]=\int\bigl[D\varphi\bigr]\,e^{-(\mathcal A[\varphi]-\int_xJ(x)\varphi(x))/\hbar}
\end{equation}
The semiclassical approximation becomes exact in the limit $\hbar\to0$, while
corrections to the leading semiclassical contribution can be developed in the
form of a power series in $\hbar$. Note that as in statistical mechanics
interpretation $\hbar\sim kT$, and we can also talk here about the
low-temperature expansion. For this reason it is convenient to retain $\hbar$
as the formal parameter.\footnote{It is of course always possible to choose
units in which the actual Planck's constant is 1. Then formal parameter named
$\hbar$ in \eqref{eq:ZJhbar} is dimensionless, and it looks like an additional
parameter of the theory. In fact, such parameter can always be traded for
coupling constants determining the strength of interaction. For example, the
$\varphi^4$ action can be written as
$\mathcal A=\frac1\lambda\int\bigl(\frac12(\partial\chi)^2+\frac{m^2}2\chi^2+\frac1{4!}\chi^4\bigr)\dd^4x$
in terms of the rescaled field $\chi=\sqrt\lambda\,\varphi$.}

For the purpose of the present analysis I will consider a generic action of the
form
\begin{equation}\label{eq:genericA}
\mathcal A=\int\left[\frac12\,(\partial\varphi)^2+\frac{m^2}2\,\varphi^2+V(\varphi)\right]\,,
\end{equation}
with $V(\varphi)$ of general form
\[
V(\varphi)=\frac{\lambda_3}{3!}\,\varphi^3+\frac{\lambda_4}{4!}\,\varphi^4+\cdots
\]
with constants $\lambda_k$, although I will often give illustrations in terms
of $\varphi^4$ theory.

\subsubsection*{The leading (tree) approximation}

For small $\hbar$ we calculate the above integral \eqref{eq:ZJhbar} for $Z[J]$
using a saddle-point approximation. That is, one first finds the minimum of the
action functional
\[
\mathcal A_J[\varphi]=\mathcal A[\varphi]-\int_xJ(x)\varphi(x)\,,
\]
i.e.\ the solution of the classical equation of motion
\[
\frac{\delta\mathcal A[\varphi]}{\delta\varphi(x)}=J(x)\,.
\]
The solution depends on $J$; I will denote it as
\[
\varphi_c(x)=\varphi_c(x,[J])\,.
\]
Then in the leading approximation
\[
\log Z[J]=\frac1\hbar\Bigl(-\mathcal A[\varphi_c]+\int_xJ(x)\varphi_c(x)\Bigr)+O(\hbar^0)
\]
When $\hbar$ is explicit it is convenient to normalize the generating
functional $W[J]$ so that
\[
\frac{Z[J]}{Z[0]}=e^{W[J]/\hbar}\,.
\]
We write
\[
W[J]=W_0[J]+\hbar\,W_1[J]+\cdots\,.
\]
In this notation we have
\begin{equation}\label{eq:W0}
W_0[J]=-\mathcal A[\varphi_c]+\int J(x)\varphi_c(x)\,,
\end{equation}
where it is understood that $\varphi_c=\varphi_c[J]$.

The relation of this leading semiclassical approximation to the ordinary
perturbation theory (expansion in $\lambda_i$) is recovered by expanding the
solution $\varphi_c$ in the powers of $J$. In view of the chosen form of the
action \eqref{eq:genericA} the classical equation has the form
\[
\bigl(m^2-\Delta_x\bigr)\varphi_c(x)+V'\bigl(\varphi_c(x)\bigr)=J(x)\,.
\]
This can be transformed to the integral equation
\[
\varphi_c(x)=\int_{x'}D(x-x')\,J(x')-\int_{x'}D(x-x')\,V'\bigl(\varphi_c(x')\bigr)
\]
which can be solved iteratively, with the last term with $V'$ taken as the
perturbation. For example, if
\[
V=\frac{\lambda}{4!}\,\varphi^4\,,\qquad V'=\frac{\lambda}{3!}\,\varphi^3\,,
\]
we find
\begin{multline*}
\varphi_c(x,[J])=\int_{x'}D(x-x')\,J(x')\\
-\frac{\lambda}{3!}\int_{x',x_1,x_2,x_3}D(x-x')\,D(x'-x_1)\,D(x'-x_2)\,D(x'-x_3)\,J(x_1)J(x_2)J(x_3)+\cdots\,.
\end{multline*}
Obviously, this generates the expansion in $J$. Also, this series is
represented diagrammatically as
\[
\varphi_c(x,[J])=
x\,\dia{\node[xvtx] at (0,0) {}; \draw[prop] (0,0) -- (3,0) node[vtx] {};}
+x\,\dia{\node[xvtx] at (0,0) {}; \draw[prop] (0,0) -- (2,0); \foreach \a in {50,0,-50} \draw[prop] (2,0) -- ++(\a:1.4) node[vtx] {}; \node[vtx] at (2,0) {};}
+x\,\dia{\node[xvtx] at (0,0) {}; \draw[prop] (0,0) -- (2,0); \foreach \a in {50,0,-50} \draw[prop] (2,0) -- ++(\a:1.4) node[vtx] {}; \node[vtx] at (2,0) {};
  \path (2,0) ++(50:1.4) coordinate (b); \foreach \a in {120,60,10} \draw[prop] (b) -- ++(\a:0.9) node[vtx] {}; \node[vtx] at (b) {};}
+\cdots\,.
\]
Here I mark the ``external point'' $x$ where $\varphi_c(x)$ sits by a cross to
distinguish it from the one-leg vertex
\[
\dia{\draw[prop] (-2,0) -- (0,0) node[vtx] {};}=\int \dd^4x\,J(x)\,.
\]
It is easy to see that the diagrams appearing here have no loops (i.e.\ one
cannot travel along the diagram and come back to the starting point never going
the same line twice). Such diagrams are called the \emph{tree diagrams}.

The easiest way to obtain the functional $W_0[J]$ \eqref{eq:W0} is to observe
that
\[
\frac{\delta W_0}{\delta J}=\left[-\frac{\delta\mathcal A}{\delta\varphi_c}+J\right]\frac{\delta\varphi_c}{\delta J}+\varphi_c=\varphi_c\,,
\]
where the fact that $\varphi_c$ solves the classical equation is used.
Integrating the above expansion of $\varphi_c(x,[J])$ term by term we find
\[
W_0[J]=\underbrace{\dia{\pic{line};}}_{W_0^{(2)}}
+\underbrace{\dia{\pic{legblob={0}{45,135,225,315}{1.4}{fill=none}}; \node[vtx] at (0,0) {};}}_{W_0^{(4)}}
+\underbrace{\Bigl(\dia{\draw[prop] (-0.8,0) -- (0.8,0); \foreach \a in {60,180,300} \draw[prop] (-0.8,0) -- ++(\a:1.1) node[vtx] {};
  \foreach \a in {120,0,240} \draw[prop] (0.8,0) -- ++(\a:1.1) node[vtx] {}; \node[vtx] at (-0.8,0) {}; \node[vtx] at (0.8,0) {};}
 +\dia{\pic{legblob={0}{0,60,120,180,240,300}{1.3}{fill=none}}; \node[vtx] at (0,0) {};}\Bigr)}_{W_0^{(6)}}+\cdots,
\]
i.e.\ the leading semiclassical approximation for $W$ is just the sum of all
connected tree diagrams. For this reason the leading approximation is also
called the \emph{tree approximation}.

Let us find the effective action in this approximation. By the definition of
the Legendre transform we have to set
\[
\phi=\frac{\delta W_0}{\delta J}=\varphi_c[J]\,.
\]
Then the effective action in this approximation is
\[
\Gamma_0[\phi]=-W_0[J]+J\phi=\mathcal A[\varphi_c]-J\varphi_c+J\varphi_c=\mathcal A[\phi]\,,
\]
i.e.\ in the leading semiclassical approximation the effective action
$\Gamma_0[\phi]$ coincides with the action $\mathcal A[\phi]$. In particular
\[
\Gamma^{(2)}_0(x,x')=S_0(x-x')=\bigl(m^2-\Delta_x\bigr)\,\delta(x-x')\,;
\]
and assuming for example that
$V(\varphi)=\frac{\lambda_4}{4!}\varphi^4+\frac{\lambda_6}{6!}\varphi^6+\cdots$,
\[
\Gamma^{(4)}_0=-\ \dia{\foreach \a in {45,135,225,315} \draw[prop,densely dashed] (0,0) -- (\a:1.4); \node[vtx,rectangle,minimum size=3.5pt] at (0,0) {};}\ ;\qquad
\Gamma^{(6)}_0=-\ \dia{\foreach \a in {0,60,120,180,240,300} \draw[prop,densely dashed] (0,0) -- (\a:1.4); \node[vtx,rectangle,minimum size=3.5pt] at (0,0) {};}\ ;\qquad\dots,
\]
where I have shown where the amputated legs used to be by the dashed lines.
Clearly, these are the simplest one-particle irreducible amputated diagrams for
$\Gamma^{(n)}$.

\subsubsection*{The first order (one-loop) correction}

Let us now turn to the first subleading terms $W_1[J]$ and $\Gamma_1[\phi]$ in
the expansions
\[
\frac{W[J]}\hbar=\frac{W_0[J]}\hbar+W_1[J]+O(\hbar)\,,\qquad
\frac{\Gamma[\phi]}\hbar=\frac{\Gamma_0[\phi]}\hbar+\Gamma_1[\phi]+O(\hbar)\,.
\]
As we have previously seen, $W_0[J]$ is given by the sum of all connected tree
diagrams, while $\Gamma_0[\phi]$ coincides with the action $\mathcal A[\phi]$.

To evaluate the first correction to the effective action one must take into
account the gaussian fluctuations around the classical solution $\varphi_c$. In
the functional integral \eqref{eq:ZJhbar} defining $Z[J]$ set
\[
\varphi(x)=\varphi_c(x)+\chi(x)
\]
and expand the action $\mathcal A_J$ to the terms quadratic in $\chi$,
\[
\mathcal A_J[\varphi]=\mathcal A[\varphi_c]-J\varphi_c+\frac{\delta\mathcal A}{\delta\varphi}\Big|_{\varphi=\varphi_c}\chi-J\chi
+\int_{x,x'}\frac12\,\frac{\delta^2\mathcal A}{\delta\varphi(x)\,\delta\varphi(x')}\Big|_{\varphi=\varphi_c}\chi(x)\chi(x')+O(\chi^3)\,.
\]
As usual, the terms linear in $\chi$ cancel in view of the classical equation
$\varphi_c$ satisfies, and we obtain in this approximation
\[
Z[J]=e^{W_0[J]/\hbar}\int\bigl[D\chi\bigr]\,e^{-\mathcal A_1[\chi]/\hbar}\,,
\]
where
\begin{equation}\label{eq:A1}
\mathcal A_1[\chi]=\int \dd^4x\,\left[\frac12\,(\partial\chi)^2+\frac{m^2}2\,\chi^2+\frac12\,V''(\varphi_c)\,\chi^2\right]
\end{equation}
Since the higher-order terms in $\chi$ are neglected the integral over $\chi$
is gaussian and it can be evaluated explicitly in terms of the determinant of
the operator
\[
K(x,x')=\frac{\delta^2\mathcal A}{\delta\varphi(x)\,\delta\varphi(x')}\Big|_{\varphi=\varphi_c}
=\bigl(m^2-\Delta_x\bigr)\,\delta(x-x')+V''(\varphi_c)\,\delta(x-x')\,.
\]

\begin{exercise}
Evaluate $n$-dimensional gaussian integral
\[
\prod_{i=1}^n\left[\int_{-\infty}^\infty \dd\xi_i\right]\,\exp\left\{-\frac12\sum_{i,j=1}^ns_{ij}\,\xi_i\xi_j\right\}
\]
where $s_{ij}$ are elements of a symmetric positive matrix, and show that it
equals
\[
(2\pi)^{n/2}\,\bigl[\det(s_{ij})\bigr]^{-1/2}\,.
\]
\end{exercise}

Actually, we are interested in the ratio $Z[J]/Z[0]$, therefore we need to
evaluate the ratio of the determinants
\begin{multline*}
\left[\frac{\det\bigl(S_0(x-x')+\delta(x-x')\,V''(x')\bigr)}{\det\bigl(S_0(x-x')\bigr)}\right]^{-1/2}
\\=\exp\left[-\frac12\log\det\Bigl(\delta(x-x')+D(x-x')\,V''\bigl(\varphi_c(x')\bigr)\Bigr)\right]
\end{multline*}
This can be transformed with the help of the identity, valid for any positive
Hermitian operator $M$:
\[
\log\bigl(\det M\bigr)=\tr\bigl(\log M\bigr)\,,
\]
which is established by bringing $M$ to the diagonal form,
\[
\log\Bigl(\prod_im_i\Bigr)=\sum_i\log m_i\,.
\]
Using this property, we obtain
\begin{equation}\label{eq:W1}
W_1[J]=-\frac12\,\tr\log\bigl(I+A(\varphi_c)\bigr)\,,
\end{equation}
where $A$ is an integral operator with the kernel
\[
A(x,x')=D(x-x')\,V''\bigl(\varphi_c(x')\bigr)\,,
\]
and $I$ is the identity operator, $I(x,x')=\delta(x-x')$. The logarithm expands
as
\[
\log\bigl(I+A\bigr)=A-\frac{A^2}2+\frac{A^3}3-\dots-\frac{(-A)^n}n+\cdots\,.
\]
When taking the trace we obtain
\[
\tr A=\int_xD(x-x)\,V''(x)\,,\qquad
\tr A^2=\int_{x,x'}D(x-x')\,V''(x')\,D(x'-x)\,V''(x)\,,\qquad\text{etc.}
\]
Here $V''(x)$ is the shorthand for $V''(\varphi_c(x))$. Obviously, this
expansion can be represented in terms of diagrams
\[
W_1[J]=-\frac12\,\tr\log\bigl(I+A\bigr)
=\dia{\draw[line width=1.2pt] (0,0) circle (1); \draw[prop,fill=white] (0:1) circle (0.25);}
+\dia{\draw[line width=1.2pt] (0,0) circle (1); \foreach \a in {0,180} \draw[prop,fill=white] (\a:1) circle (0.25);}
+\dia{\draw[line width=1.2pt] (0,0) circle (1); \foreach \a in {30,150,270} \draw[prop,fill=white] (\a:1) circle (0.25);}
+\cdots\,,
\]
with the line standing for $D$, as usual, but with the special two-leg vertex
\[
\dia{\draw[prop] (-1.5,0) -- (1.5,0); \draw[prop,fill=white] (0,0) circle (0.4);}
=-\int \dd^4x\,V''\bigl(\varphi_c(x)\bigr)\,.
\]
A diagram with $n$ such vertices comes with the factor (implicit in the above
series)
\[
\frac1{2n}=\frac12\times\frac1n\,,
\]
$1/2$ already being a factor in \eqref{eq:W1}, and $1/n$ coming from the
expansion coefficient of the log function; this factor can be interpreted as
the symmetry factor of the $n$-vertex diagram. Of course the above diagram
expansion can be obtained directly from the Feynman rules for the action
\eqref{eq:A1}.

To obtain $W_1[J]$, substitute
\[
\varphi_c=\varphi_c[J]=
\dia{\node[xvtx] at (0,0) {}; \draw[prop] (0,0) -- (3,0) node[vtx] {};}
+\dia{\node[xvtx] at (0,0) {}; \draw[prop] (0,0) -- (2,0); \foreach \a in {50,0,-50} \draw[prop] (2,0) -- ++(\a:1.4) node[vtx] {}; \node[vtx] at (2,0) {};}
+\cdots\,,
\]
which, as we have seen, is a sum of the tree diagrams. For $\varphi^4$ theory
\[
V''(\varphi_c)=\frac\lambda2\,\varphi_c^2
\]
and so
\[
\dia{\draw[prop] (-1.5,0) -- (1.5,0); \draw[prop,fill=white] (0,0) circle (0.4);}
=\dia{\draw[prop] (-1.5,0) -- (1.5,0); \draw[prop] (-0.8,1.4) node[vtx] {} -- (0,0) -- (0.8,1.4) node[vtx] {}; \node[vtx] at (0,0) {};}
+\dia{\draw[prop] (-1.5,0) -- (1.5,0); \draw[prop] (-0.8,1.4) node[vtx] {} -- (0,0) -- (0.8,1.4); \node[vtx] at (0,0) {};
  \foreach \a in {120,60,10} \draw[prop] (0.8,1.4) -- ++(\a:0.8) node[vtx] {}; \node[vtx] at (0.8,1.4) {};}
+\cdots\,.
\]
Therefore $W_1$ is the sum of the diagrams of the form
\[
\dia{\draw[line width=1.2pt] (0,0) circle (1); \foreach \a in {30,150,270} \draw[prop,fill=white] (\a:1) circle (0.25);}
\ \longrightarrow\
\dia{\draw[prop] (0,0) circle (1); \foreach \a in {30,150,270} \node[vtx] at (\a:1) {};
  \draw[prop] (30:1) -- ++(60:0.9) node[vtx] {}; \draw[prop] (30:1) -- ++(0:0.9) node[vtx] {};
  \draw[prop] (270:1) -- ++(240:0.9) node[vtx] {}; \draw[prop] (270:1) -- ++(300:0.9) node[vtx] {};
  \draw[prop] (150:1) -- ++(150:0.9) coordinate (c) node[vtx] {}; \draw[prop] (150:1) -- ++(90:0.9) node[vtx] {};
  \foreach \a in {200,150,100} \draw[prop] (c) -- ++(\a:0.7) node[vtx] {};}
\]
which have exactly one loop. This is why the $\sim\hbar$ correction to $W$ is
called the ``one-loop correction''.

It is not difficult to find the effective action in this approximation. We
have
\[
\Gamma[\phi]=-W_0[J]-\hbar\,W_1[J]+\int J\phi+O(\hbar^2)\,,
\]
where we have to substitute $J=J[\phi]$. Note that by definition of the
Legendre transform, the r.h.s.\ of this equation is stationary with respect to
variations of $J$ with $\phi$ fixed. Therefore the correction $\sim\hbar$ to
the relation between $J$ and $\phi$ will produce a change of the order of
$\hbar^2$ to the r.h.s.\ of the above equation. Hence at order $\hbar$ we still
can replace $J$ by the solution of the equation
\[
\phi(x)=\varphi_c(x,[J])\,.
\]
Then, again $-W_0+\int J\phi=\mathcal A[\phi]$, and
\[
\Gamma[\phi]=\mathcal A[\phi]+\hbar\,\Gamma_1[\phi]+O(\hbar^2)\,,
\]
where
\[
\Gamma_1[\phi]=\frac12\,\tr\log\bigl(I+A(\phi)\bigr)
=-\Bigl(\dia{\draw[line width=1.2pt] (0,0) circle (1); \node[otimesv] at (0:1) {};}
+\dia{\draw[line width=1.2pt] (0,0) circle (1); \foreach \a in {0,180} \node[otimesv] at (\a:1) {};}
+\dia{\draw[line width=1.2pt] (0,0) circle (1); \foreach \a in {30,150,270} \node[otimesv] at (\a:1) {};}
+\cdots\Bigr)\,,
\]
with
\[
\dia{\node[otimesv] at (0,0) {};}
=-V''(\phi)\,.
\]
Note that all these diagrams are one-particle irreducible.

Similar analysis can be carried out in higher orders in $\hbar$. One finds that
at the order $\hbar^n$ only the diagrams with exactly $n$ loops contribute to
$W_n[J]$.

\subsection*{Effective Action vs Generating Functional}

We can now clarify the nature of the relation between $W[J]$ and
$\Gamma[\phi]$. Write again
\[
\Gamma[\phi]=\sum_{n=2}^\infty\frac1{n!}\int\Gamma^{(n)}(x_1,\dots,x_n)\,\phi(x_1)\cdots\phi(x_n)\,,
\]
and represent $\Gamma^{(n)}$ as an $n$-point blob
\[
\Gamma^{(n)}=\diac{\draw[prop,pattern=north east lines] (0,0) circle (1);
  \foreach \a/\l/\p in {180/x_1/left,125/x_2/above left,55/x_3/above right,0/x_4/right} {\node[vtx,minimum size=4pt] at (\a:1) {}; \node[\p] at (\a:1) {\small$\l$};}
  \draw[densely dotted] (200:1.35) arc[start angle=200,end angle=340,radius=1.35];}\,.
\]
The generating functional $W[J]$ is related to $\Gamma[\phi]$ by a functional
Legendre transform.
\[
W[J]=\Bigl[-\Gamma[\phi]+\int_xJ(x)\phi(x)\Bigr]_{\phi=\phi[J]}
\]
One first solves the equation
\[
J(x)=\frac{\delta\Gamma[\phi]}{\delta\phi(x)}
\]
with respect to $\phi(x)$. Exactly as was done when solving the classical
equations of motion, the solution is represented as a sum of tree diagrams
\[
\phi(x,[J])=
\dia{\node[xvtx] at (0,0) {}; \draw[prop] (0,0) -- (3.2,0) node[vtx] {}; \draw[prop,fill=white] (1.6,0) circle (0.45);}
+\dia{\node[xvtx] at (-2.3,0) {}; \draw[prop] (-2.3,0) -- (-0.7,0); \draw[prop,fill=white] (-1.5,0) circle (0.28);
  \pic{dressed={0.7}{60,0,-60}};}
+\cdots\,,
\]
this time with $W(x-x')$ (the inverse of $\Gamma^{(2)}$, i.e.\ the full
two-point function) standing instead of the propagators $D$, and the functions
$\Gamma^{(n)}$ playing the role of the vertices. In drawing the above diagrams
I have assumed that the interaction was even, $V(\varphi)=V(-\varphi)$, so that
odd order vertices are absent; of course, this limitation is not significant.
Again, integrating the equation
\[
\frac{\delta W[J]}{\delta J(x)}=\phi(x,[J])
\]
term by term in $J$, one finds that $W[J]$ is a sum of all tree diagrams, again
with $W$ for propagators and $\Gamma^{(n)}$ for vertices,
\[
W[J]=\dia{\draw[prop] (-1.6,0) node[vtx]{} -- (1.6,0) node[vtx]{}; \draw[prop,fill=white] (0,0) circle (0.45);}
+\diac{\pic{dressed={0.7}{45,135,225,315}};}
+\ \text{all tree diagrams}
\]
This is exactly what we expect to be the relation between the sum of all
connected diagrams and the sum of all one-particle irreducible diagrams: one
takes one-particle irreducible diagrams and connects them in a tree-like manner
to produce all connected diagrams.

\begin{hwproblem}{3}\label{hw:3}
In the $\varphi^3$ theory
\[
\mathcal A_I=\frac{\lambda_3}{3!}\int\varphi^3(x)\dd^4x
\]
carry out the Legendre transform of $W[J]$ to the order 4 in $J$ and $\phi$.
Define
\[
\hat\Gamma[\phi]:=\Gamma[\bar\phi+\phi]-\Gamma[\bar\phi]
\qquad\text{where}\qquad \bar\phi=W^{(1)}(x)
\]
and $\Gamma[\phi]$ is the Legendre transform of $W[J]$. Express $W^{(3)}$ and
$W^{(4)}$ in terms of $\hat\Gamma^{(4)}$, $\hat\Gamma^{(3)}$, $\hat\Gamma^{(2)}$.
Give diagrammatic representation of your result.
\end{hwproblem}

% Source pp. 93-101 (Lectures7-12.pdf).
\section{Diagram Evaluation. Divergences}

Now we turn to actual evaluation of the diagram contributions. For this purpose
in most cases it is much more convenient to use the Feynman rules in the
momentum space, evaluating the Fourier transformed correlation functions
\begin{multline}\label{eq:FTW}
\int W^{(n)}(x_1,\dots,x_n)\,e^{-ip_1\cdot x_1}\cdots e^{-ip_n\cdot x_n}\dd^4x_1\cdots \dd^4x_n\\
=(2\pi)^4\,\delta^{(4)}(p_1+\dots+p_n)\,\tilde W^{(n)}(p_1,\dots,p_n)\,.
\end{multline}
Here again $W^{(n)}$ denotes \emph{connected} $n$-point correlation function.
The delta function in the r.h.s.\ appears because of the translation
invariance: $W^{(n)}(x_1,\dots,x_n)$ depends only on the differences
$x_i-x_j$.

I will skip detailed derivation of the diagram rules in the momentum space
(look up \textbf{Sect.~4.4 of PS} for the details). Very briefly, in any
diagram one uses the Fourier transform
\[
D(x-x')=\int\frac{\dd^4p}{(2\pi)^4}\,\tilde D(p)\,e^{ip\cdot(x-x')}\,,\qquad
\tilde D(p)=\frac1{p^2+m^2}
\]
of all propagators. Then the integral over all $x$ associated with vertices, as
well as over all ``external'' points $x_i$, can be evaluated explicitly yielding
the momentum delta-functions. As the result one obtains the

\subsection*{Momentum-space Feynman rules:}
\begin{enumerate}[label=(\roman*)]
\item ``Internal'' propagators carry (directed) momentum to be integrated over,
\[
\dia{\draw[prop] (-1.5,0) -- (1.5,0); \draw[arr] (-0.5,0.35) -- (0.5,0.35) node[above,pos=0.5] {\small$p$};}
=\dia{\draw[prop] (-1.5,0) -- (1.5,0); \draw[arr] (0.5,0.35) -- (-0.5,0.35) node[above,pos=0.5] {\small$-p$};}
=\int\frac{\dd^4p}{(2\pi)^4}\,\frac1{p^2+m^2}\,.
\]
\item ``External'' legs carry fixed momenta
\[
\tilde W^{(n)}(p_1,\dots,p_n)=
\diac{\foreach \a/\l in {135/p_1,90/p_2,45/p_3,0/p_4} {\draw[prop] (0,0) -- (\a:2.2);
  \draw[arr] (\a:2.0) ++(\a+90:0.35) -- ++(\a+180:0.8); \node at ($(\a:2.2)+(\a-90:0.5)$) {\small$\l$};}
  \node[circle,draw,fill=white,inner sep=1.3pt] at (0,0) {};
  \draw[densely dotted] (165:1.4) arc[start angle=165,end angle=345,radius=1.4];}\,,
\]
\item The vertices conserve the momentum
\[
\diac{\foreach \a/\l in {135/p_1,45/p_2,315/p_3,225/p_4} {\draw[prop] (0,0) -- (\a:2);
  \draw[arr] (\a:1.9) ++(\a+90:0.3) -- ++(\a+180:0.7); \node at ($(\a:1.9)+(\a+90:0.75)$) {\small$\l$};}
  \node[vtx] at (0,0) {};}
=-\lambda\,(2\pi)^4\,\delta^{(4)}(p_1+\dots+p_4)\,,
\]
\item Integration $\int\frac{\dd^4p_i}{(2\pi)^4}$ is performed over all the
momenta associated with the ``internal'' propagators.
\end{enumerate}

\subsection*{Some properties of \texorpdfstring{$\tilde W^{(n)}$}{W(n)}}

In the momentum representation many operations simplify. For example the
``amputation'' takes very simple form
\[
\tilde W^{(n)}(p_1,\dots,p_n)=\left[\prod_{i=1}^n\tilde W(p_i)\right]\,\tilde W^{(n)}_{\text{amp}}(p_1,\dots,p_n)\,,
\]
where $\tilde W(p)$ is the Fourier transform of $W(x-x')$. Similar
simplification occurs to the relation to proper vertices.

\begin{exercise}
Express $\tilde W^{(6)}(p_1,\dots,p_6)$ in terms of
$\tilde\Gamma^{(6)}(p_1,\dots,p_6)$, $\tilde\Gamma^{(4)}$ and $\tilde W$.
\end{exercise}

The momentum space $\tilde\Gamma^{(2)}(p)$ is just the inverse of
$\tilde W(p)$ in usual sense,
\[
\tilde\Gamma^{(2)}(p)=\frac1{\tilde W(p)}=p^2+m^2+\tilde\Sigma(p)\,,
\]
where $-\tilde\Sigma(p)$ is the sum of all one-particle irreducible two-leg
diagrams in the momentum space.

\subsection*{Leading contribution to \texorpdfstring{$\tilde\Sigma$}{Sigma}. Mass renormalization}

Let us calculate some perturbative contributions to $\tilde\Sigma(p)$. In the
leading order we have
\[
\tilde\Sigma(p)=-\Bigl(\,\dia{\draw[prop,densely dashed] (-2,0) -- (2,0); \draw[prop] (0,0.6) circle (0.6);
  \draw[arr] (-0.3,1.45) -- (0.3,1.45); \node[above] at (0,1.45) {\small$k$}; \node[vtx] at (0,0) {};}\,\Bigr)
=\frac\lambda2\int\frac{\dd^4k}{(2\pi)^4}\,\frac1{k^2+m^2}
=\frac{\lambda}{16\pi^2}\int_0^\infty\frac{|k|^3\dd|k|}{|k|^2+m^2}\,.
\]
This integral is badly divergent as $k\to\infty$. What should we do about it?
We can get rid of the meaningless expression by introducing a ``cutoff'', say,
replacing
\[
\frac1{k^2+m^2}\ \to\ \frac1{k^2+m^2}\,\Phi(k^2/\Lambda^2)\,,
\]
with the ``cutoff factor'' $\Phi(\xi)$ decaying at $\xi\to\infty$ fast enough
to make the integral convergent. Then
\[
\frac\lambda2\int\frac{\dd^4k}{(2\pi)^4}\,\frac1{k^2+m^2}\,\Phi(k^2/\Lambda^2)=\frac\lambda2\,m^2\,F(\Lambda^2/m^2)\,.
\]
Just as in the case of KG vacuum energy, this move alone brings no real
progress: it replaces something meaningless by something undetermined.

In this simple diagram one can note that the contribution has no momentum
dependence, and enters the quantity $\tilde\Gamma^{(2)}(p)$ in the combination
\[
\tilde\Gamma^{(2)}(p)=p^2+m^2+\frac\lambda2\,m^2\,F(\Lambda^2/m^2)+O(\lambda^2)\,.
\]
We see that no matter what $\Phi$ and $\Lambda$ are, the parameter $m$ in our
action is not actual mass of the particle; the interaction leads to a shift
that happens to depend on the cutoff. (We will explain later that actual mass
$m_{\mathrm{ph}}$ of the particle in interacting theory is determined through
the equation $\tilde\Gamma^{(2)}(p)\big|_{-p^2=m^2_{\mathrm{ph}}}=0$.)
Therefore it is reasonable to change the notations, denoting $m_0^2$ the
coefficient in front of $\varphi^2$ in the action. Then, up to terms
$\sim\lambda^2$ and higher
\[
\tilde\Gamma^{(2)}(p)=p^2+\underbrace{m_0^2+\frac\lambda2\,m_0^2\,F(\Lambda^2/m_0^2)}_{m^2}
\]
Recall that we have already observed similar phenomenon in our path integral
representation of the free propagator $D$. We started with a discrete
approximation of the path integral (an analog of the cutoff), and we were
forced to take some parameter $m_0$ instead of $m$ in the cutoff action
\[
\mathcal A=m_0\,\text{Length}\,;
\]
we then found that $m_0$ had to be given some dependence on the cutoff
parameter $\Delta$ (which plays the same role as $\Lambda^{-1}$ here) to keep
actual mass $m$ finite in the limit $\Delta\to0$.

Similar idea can be applied here. Assume that the parameter $m_0^2$ in the
action ($m_0$ is usually called the ``bare mass'') depends on $\Lambda$ in such
a way that the actual mass
\[
m^2=m_0^2+\frac\lambda2\,m_0^2\,F(\Lambda^2/m_0^2)+O(\lambda^2)
\]
remains finite as $\Lambda\to\infty$. Then, to this order,
\[
\tilde\Gamma^{(2)}(p)=p^2+m^2+O(\lambda^2)
\]
has finite limit. The result is expressed in terms of the actual mass; all
ambiguities associated with the cutoff are ``absorbed'' by the ``bare'' mass
parameter.

This idea can be refined as follows. We start with the action
\[
\mathcal A=\int\left[\frac12\,(\partial\varphi)^2+\frac{m_0^2}2\,\varphi^2+\frac\lambda{4!}\,\varphi^4\right]\,.
\]
Let us write
\[
m_0^2=m^2+\delta m^2\,,
\]
where $m^2$ is assumed to coincide with actual mass squared; at this point the
term $\delta m^2$ is just our way to acknowledge that the mass parameter in the
action need not be equal to $m^2$. We then split the action into the free and
interaction parts as follows
\[
\mathcal A=\int\left[\underbrace{\frac12\,(\partial\varphi)^2+\frac{m^2}2\,\varphi^2}_{\mathcal A_0}
+\underbrace{\frac{\delta m^2}2\,\varphi^2+\frac\lambda{4!}\,\varphi^4}_{\mathcal A_I}\right]\,,
\]
Now the propagator is
\[
\dia{\draw[prop] (-1.5,0) -- (1.5,0); \draw[arr] (-0.5,0.35) -- (0.5,0.35) node[above,pos=0.5] {\small$p$};}
=\frac1{p^2+m^2}\,,
\]
not $\frac1{p^2+m_0^2}$, as it was in the original perturbation theory, but we
have an additional two-leg vertex
\[
\dia{\draw[prop] (-1.8,0) -- (1.8,0); \node[xvtx] at (0,0) {};
  \draw[arr] (-1.4,0.35) -- (-0.6,0.35) node[above,pos=0.5] {\small$p_1$}; \draw[arr] (1.4,0.35) -- (0.6,0.35) node[above,pos=0.5] {\small$p_2$};}
=-\delta m^2\,(2\pi)^4\,\delta^{(4)}(p_1+p_2)\,,
\]
which we will rather write in a shorter form
\[
\dia{\draw[prop] (-1.8,0) -- (1.8,0); \node[xvtx] at (0,0) {};
  \draw[arr] (-1.4,0.35) -- (-0.6,0.35) node[above,pos=0.5] {\small$p$}; \draw[arr] (0.6,0.35) -- (1.4,0.35) node[above,pos=0.5] {\small$p$};}
=-\delta m^2\,,
\]
reflecting the fact that in this case the momentum delta function is easily
integrated out.

The term $\frac{\delta m^2}2\varphi^2$ in the action is the simplest
\emph{counterterm}, and the associated two-leg vertex is called the counterterm
vertex. In this modified perturbation theory we have (up to terms
$\sim\lambda^2$ and higher)
\begin{align*}
\tilde\Sigma(p)&=-\Bigl(\,\dia{\draw[prop,densely dashed] (-2,0) -- (2,0); \draw[prop] (0,0.6) circle (0.6); \node[vtx] at (0,0) {};}
+\dia{\draw[prop,densely dashed] (-2,0) -- (2,0); \node[xvtx] at (0,0) {};}\,\Bigr)\\
&=\frac\lambda2\int\frac{\dd^4k}{(2\pi)^4}\,\frac1{k^2+m^2}\,\Phi(k^2/\Lambda^2)+\delta m^2=\frac\lambda2\,m^2\,F(\Lambda^2/m^2)+\delta m^2\,.
\end{align*}
Since we insist that $m$ is the actual mass, we have to set
\[
\delta m^2=-\frac\lambda2\,m^2\,F(\Lambda^2/m^2)+O(\lambda^2)\,.
\]
Note that this looks different from the relation
\[
m_0^2=m^2-\frac\lambda2\,m_0^2\,F(\Lambda^2/m_0^2)+O(\lambda^2)
\]
we have obtained within the original perturbation theory; however these two
relations are consistent since the difference is $\sim\lambda^2$.

As we have seen, in the modified perturbation theory, to the order $\sim\lambda$
the counterterm diagram exactly cancels the cutoff dependent bubble diagram,
\[
\dia{\draw[prop,densely dashed] (-2,0) -- (2,0); \node[xvtx] at (0,0) {}; \node[above] at (0,0.2) {\scriptsize$\delta m^2$};}
+\dia{\draw[prop,densely dashed] (-2,0) -- (2,0); \draw[prop] (0,0.6) circle (0.6); \node[vtx] at (0,0) {};}=0\,,
\]
by choice of the parameter $\delta m^2$. The result has no dependence on
$\Lambda$ and $\Phi$. Note also that this cancellation occurs inside more
complicated diagrams, like
\[
\diac{\draw[prop,densely dashed] (-3,0) -- (-1.8,0) (1.8,0) -- (3,0);
  \draw[prop] (-1.8,0) -- (1.8,0); \draw[prop] (1.8,0) arc[start angle=0,end angle=180,radius=1.8];
  \draw[prop] (0,0) arc[start angle=0,end angle=-180,radius=0.9]; \draw[prop] (1.8,0) arc[start angle=0,end angle=-180,radius=0.9];
  \draw[prop] (0,2.4) circle (0.6); \foreach \p in {(-1.8,0),(0,0),(1.8,0),(0,1.8)} \node[vtx] at \p {};}
+\diac{\draw[prop,densely dashed] (-3,0) -- (-1.8,0) (1.8,0) -- (3,0);
  \draw[prop] (-1.8,0) -- (1.8,0); \draw[prop] (1.8,0) arc[start angle=0,end angle=180,radius=1.8];
  \draw[prop] (0,0) arc[start angle=0,end angle=-180,radius=0.9]; \draw[prop] (1.8,0) arc[start angle=0,end angle=-180,radius=0.9];
  \node[xvtx] at (0,1.8) {}; \foreach \p in {(-1.8,0),(0,0),(1.8,0)} \node[vtx] at \p {};}
=0
\]
For this reason the diagrams involving the bubble can be ignored in
calculations, if we keep in mind that their contributions are canceled by the
mass counterterm.

The idea behind this modified perturbation theory can be extended to obtain the
so called \emph{renormalized perturbation theory}, in which all divergences of
original perturbation theory are eliminated. But first let's find out exactly
what diagrams diverge.

\subsection*{Divergences in \texorpdfstring{$\varphi^4$}{phi**4}}

Consider $n$-leg proper vertex
\[
\tilde\Gamma^{(n)}(p_1,\dots,p_n)=
\diac{\foreach \a/\l in {135/p_1,90/p_2,45/p_3,0/p_4} {\draw[prop] (\a:0.8) -- (\a:2.2);
  \draw[arr] (\a:2.0) ++(\a+90:0.35) -- ++(\a+180:0.8); \node at ($(\a:2.2)+(\a-90:0.5)$) {\small$\l$};}
  \draw[prop,pattern=north east lines] (0,0) circle (0.8);
  \foreach \a in {135,90,45,0} \node[vtx,minimum size=4pt] at (\a:0.8) {};
  \draw[densely dotted] (165:1.4) arc[start angle=165,end angle=345,radius=1.4];}\,.
\]
How to characterize diagrams which contribute to this vertex? Assume that the
diagram contains
\[
\begin{array}{lll}
V & \text{vertices} & \dia{\pic{xvertex};}\\[8pt]
P & \text{propagators} & \dia{\draw[prop] (0,0) -- (2,0);}\\[8pt]
I & \text{momentum integrals} & \displaystyle\int\frac{\dd^4k}{(2\pi)^4}\,.
\end{array}
\]
These numbers are not all independent. First, $V$ vertices must absorb exactly
$4V$ lines, $2P$ coming from the propagators, and $n$ from the external
(amputated) legs, i.e.\ $4V=2P+n$, or
\[
2P=4V-n\,.
\]
Next, there is a 4-momentum flowing in each propagator, but there is also a
momentum delta-function associated with each vertex. Each of these
delta-functions reduces the number of 4-momentum integrations by 1. This is
true for all delta-constraints but one, because the momentum-space connected
correlation functions must have one overall delta function left, namely
\[
(2\pi)^4\,\delta^{(4)}(p_1+\dots+p_n)\,.
\]
Therefore
\[
I=P-V+1=\bigl(P=2V-n/2\bigr)=V-n/2+1\,.
\]
Thus we have two equations
\[
\begin{cases}
2P=4V-n\\
4I=4V-2n+4
\end{cases}
\]
Now, the expression for the diagram has the following general form
\[
\int\frac{\dd^4k'_1\cdots \dd^4k'_I}{(q_1^2+m^2)\cdots(q_P^2+m^2)}\,,
\]
where $k'_1,\dots,k'_I$ are the integration momenta left over after eliminating
the delta-functions, and $q_1,\dots,q_P$ are some linear combinations of $k'_j$
and the external momenta $p_i$. For our present analysis it is irrelevant if we
have $m^2$ or $m_0^2$ in the propagators.

We are concerned with possible divergence as $k'_j\to\infty$ (which generally
implies $q_i\to\infty$). Roughly speaking, the above integral diverges if the
overall power of the momenta in the integration measure in the numerator
exceeds or equals the overall power in the denominator. Therefore let us define
the \emph{superficial degree of divergence} $D$ as
\[
D=4I-2P\,.
\]
The integral definitely diverges if $D\ge0$ (but $D<0$ does not guarantee
convergence, see below). If we introduce a cutoff momentum $\Lambda$, the
diagram of the superficial degree of divergence $D$ is $\sim\Lambda^D$.

Substituting the above expressions for $P$ and $I$, we find
\[
D=4V-2n+4-4V+n=4-n\,.
\]
We see that the number of the vertices $V$ cancels, and the superficial degree
of divergence depends on $n$ only. In addition, we see that only finite number
of the proper vertices ($n=0$, $n=2$ and $n=4$) are superficially divergent.

Incidentally, if a perturbation theory has finitely many superficially
divergent vertices (as the $\varphi^4$ perturbation theory does), it is called
\emph{renormalizable}, for the reason which will become clear later. Otherwise,
when infinitely many proper vertices are superficially divergent, we are
dealing with \emph{non-renormalizable} perturbation theory.

\begin{exercise}
Analyze the superficial degree of divergence for all proper vertices in
$\varphi^6$ theory.
\end{exercise}

In the $\varphi^4$ perturbation theory, the proper vertices with $n>4$ are
superficially convergent. This does not mean all the diagrams for these
vertices are given by convergent integrals (that is why we use the name
``superficial'' for that, essentially dimensional, analysis). The reason is
that the diagrams for $\Gamma^{(n)}$ with $n\ge6$ still can have 2- or 4-leg
subdiagrams. For example, the diagram series for $\Gamma^{(6)}$ contains the
diagrams like these
\[
\diac{\coordinate (a) at (-1.5,1.5); \coordinate (b) at (1.5,1.5); \coordinate (c) at (0,-1);
  \draw[prop] (a) -- (b) -- (c) -- (a); \draw[prop] (0,1.5) ellipse (0.55 and 0.32);
  \draw[densely dotted] (-0.9,1.0) rectangle (0.9,2.0);
  \foreach \p/\q in {a/180,a/120,b/0,b/60,c/240,c/300} \draw[prop,densely dashed] (\p) -- ++(\q:0.9);
  \foreach \p in {(a),(b),(c),(-0.55,1.5),(0.55,1.5)} \node[vtx] at \p {};}
\qquad\qquad
\diac{\coordinate (a) at (-1,1); \coordinate (b) at (1,1); \coordinate (c) at (-1,-1); \coordinate (d) at (1,-1);
  \draw[prop] (a) -- (c) -- (d) -- (b); \draw[prop] (0,1) ellipse (1 and 0.3);
  \draw[densely dotted] (-1.5,0.5) rectangle (1.5,1.6);
  \foreach \p/\q in {a/135,b/45,c/180,c/270,d/0,d/270} \draw[prop,densely dashed] (\p) -- ++(\q:0.9);
  \foreach \p in {a,b,c,d} \node[vtx] at (\p) {};}
\]
which contain divergent blocks with two or four legs. We will come back to
decide how to handle such subdiagrams later.

% Source pp. 102-110 (Lectures7-12.pdf).
\section{Renormalized Perturbation Theory}

In $\varphi^4$ theory, the vertices with positive superficial degree of
divergence $D\ge0$ (the so called \emph{primitive divergences}) are
$\Gamma^{(2)}$ and $\Gamma^{(4)}$. We have already reviewed the simplest
divergent diagram for $\Gamma^{(2)}$. The diagrams for $\Gamma^{(4)}$ have the
superficial degree of divergence $D=0$; this means these diagrams are
\emph{logarithmically divergent}. To clarify the meaning of this term, let us
consider the simplest connected diagrams for this vertex,
\[
-\tilde\Gamma^{(4)}(p_1,p_2,p_3,p_4)=
\diac{\foreach \a/\l in {135/p_1,45/p_2,315/p_3,225/p_4} {\draw[prop,densely dashed] (0,0) -- (\a:1.6);
  \node at ($(\a:1.7)+(\a+90:0.45)$) {\scriptsize$\l$};} \node[vtx] at (0,0) {};}
+\Bigl(\diac{\draw[prop] (0,0) ellipse (1.2 and 0.7);
  \foreach \a/\l in {135/p_1,225/p_4} {\draw[prop,densely dashed] (-1.2,0) -- ++(\a:1.4); \node at ($(-1.2,0)+(\a:1.6)+(\a+90:0.4)$) {\scriptsize$\l$};}
  \foreach \a/\l in {45/p_2,315/p_3} {\draw[prop,densely dashed] (1.2,0) -- ++(\a:1.4); \node at ($(1.2,0)+(\a:1.6)+(\a-90:0.4)$) {\scriptsize$\l$};}
  \draw[arr] (-0.4,1.0) -- (0.4,1.0); \node[above] at (0,1.05) {\scriptsize$p_1+p_2-k$};
  \draw[arr] (-0.4,-1.0) -- (0.4,-1.0); \node[below] at (0,-1.05) {\scriptsize$k$};
  \node[vtx] at (-1.2,0) {}; \node[vtx] at (1.2,0) {};}
+\text{permutations}\Bigr)+\cdots
\]
The explicitly shown loop diagram represents the momentum integral
\[
\frac{(-\lambda)^2}2\int\frac{d^4k}{(2\pi)^4}\,\frac1{(k^2+m^2)\bigl((p_1+p_2-k)^2+m^2\bigr)}\,.
\]
When $|k|\gg|p_1+p_2|$ and $|k|\gg m$ this reduces to
\[
\frac{\lambda^2}{2\,(2\pi)^4}\int\frac{d^4k}{k^2\,k^2}\,.
\]
Note that $d^4k=d\Omega_4\,|k|^3\,d|k|$, where $d\Omega_4$ is the surface
element of a unit 3-sphere $S^3$, $\int_{S^3}d\Omega_4=2\pi^2$. We have
\[
\frac{\lambda^2\,2\pi^2}{2\,(2\pi)^4}\int\frac{|k|^3\,d|k|}{|k|^4}=\frac{\lambda^2}{16\pi^2}\int\frac{d|k|}{|k|}\,.
\]
The integral diverges at large $|k|$ (one shouldn't worry about its divergence
at $|k|\to0$; the above simplified form is not valid at small $|k|$). If one
introduces a cutoff, say replaces
\[
\frac1{k^2+m^2}\ \to\ \frac1{k^2+m^2}\,\Phi(k^2/\Lambda^2)
\]
with some $\Phi(x)$ such that $\Phi(x)=1$ for $x\ll1$, but decays fast at
$x\gg1$, the integral evaluates as
\[
\frac{\lambda^2}{32\pi^2}\Bigl(\log\frac{\Lambda^2}{m^2}+f(p_{12})\Bigr)\,,
\]
where $f(p_{12})$ (I use $p_{12}=p_1+p_2$) has finite limit as
$\Lambda\to\infty$. Including the ``permutations'' terms, we then have
\[
\tilde\Gamma^{(4)}=\lambda-\frac{\lambda^2}{32\pi^2}\Bigl(3\log\frac{\Lambda^2}{m^2}+f(p_{12})+f(p_{13})+f(p_{14})\Bigr)\,.
\]

One can observe that the divergent term carries no dependence on the external
momenta, and in this respect comes on the same footing as the leading
contribution $\lambda$. Therefore at least at the order $\lambda^2$ the
$\Lambda$-dependent term can be absorbed into some redefinition of the coupling
constant $\lambda$. This can be formulated more precisely as follows. Just as
in the case of the mass renormalization discussed above, we should start with
the action
\[
\mathcal A=\int\Bigl[\frac12\,(\partial\varphi)^2+\frac{m_0^2}2\,\varphi^2+\frac{\lambda_0}{4!}\,\varphi^4\Bigr]\,,
\]
with some ``bare'' coupling constant $\lambda_0$ (as well as the ``bare'' mass
parameter $m_0^2$). Then it is possible to give $\lambda_0$ certain dependence
on the cutoff momentum $\Lambda$,
\[
\lambda_0=\lambda_0(\Lambda)\,,
\]
such that the limit $\Lambda\to\infty$ of $\tilde\Gamma^{(4)}$ in this order of
the perturbation theory is finite. Namely, let us set
\[
\lambda_0(\Lambda)=\lambda+\frac{3\lambda^2}{32\pi^2}\Bigl(\log\frac{\Lambda^2}{m^2}+C\Bigr)+O(\lambda^3)\,,
\]
where new parameter $\lambda$ is called the \emph{renormalized} coupling
constant, and $C$ is arbitrary number. We then have
\begin{align*}
\tilde\Gamma^{(4)}&=\lambda_0-\frac{\lambda_0^2}{32\pi^2}\Bigl(3\log\frac{\Lambda^2}{m^2}+\tilde f(p)\Bigr)+O(\lambda_0^3)\\
&=\lambda+\frac{3\lambda^2}{32\pi^2}\Bigl(\log\frac{\Lambda^2}{m^2}+C\Bigr)-\frac{\lambda^2}{32\pi^2}\Bigl(3\log\frac{\Lambda^2}{m^2}+\tilde f(p)\Bigr)+O(\lambda^3)\,.
\end{align*}
(Here $\tilde f(p)=f(p_{12})+f(p_{13})+f(p_{14})$.) The $\Lambda$-dependent term
cancels out. Assuming that $\lambda$ is fixed in the limit $\Lambda\to\infty$,
the above proper vertex has finite limit when the cutoff is removed.

Later we shall formulate this renormalization prescription with full
precision, and in particular explain the meaning of this constant $C$. But for
now we still have one more divergence to deal with.

Consider the diagram
\[
\diac{\draw[prop,densely dashed] (-2.4,0) -- (-1.2,0) (1.2,0) -- (2.4,0);
  \draw[prop] (-1.2,0) -- (1.2,0); \draw[prop] (0,0) ellipse (1.2 and 0.75);
  \node[above] at (-2.0,0) {\small$p$};
  \draw[arr] (-0.4,1.05) -- (0.4,1.05); \node[above] at (0,1.1) {\small$p-k_1-k_2$};
  \draw[arr] (-0.4,0.25) -- (0.4,0.25); \node[below] at (0,0.05) {\scriptsize$k_2$};
  \draw[arr] (-0.4,-1.05) -- (0.4,-1.05); \node[below] at (0,-1.1) {\small$k_1$};
  \node[vtx] at (-1.2,0) {}; \node[vtx] at (1.2,0) {};}\,,
\]
contributing to $-\tilde\Sigma(p)$ in the order $\lambda^2$. I will denote this
contribution as $-\tilde\Sigma_2(p)$. According to the general counting above,
this diagram has superficial degree of divergence $D=2$, i.e.\ with the cutoff
introduced it behaves as $\Lambda^2$ as $\Lambda\to\infty$. We can see this
writing the contribution explicitly,
\[
-\tilde\Sigma_2(p)=\frac{(-\lambda)^2}{3!}\int\frac{d^4k_1}{(2\pi)^4}\,\frac{d^4k_2}{(2\pi)^4}\,
\frac1{(k_1^2+m^2)(k_2^2+m^2)\bigl((p-k_1-k_2)^2+m^2\bigr)}
\]
(note the symmetry factor $1/3!$). For $k_1,k_2\gg p,m$ the integrand
simplifies as
\[
\sim\int\frac{d^4k_1\,d^4k_2}{k_1^2\,k_2^2\,(k_1+k_2)^2}\sim\Lambda^2\,.
\]
Let us write
\[
\tilde\Sigma_2(p)=\tilde\Sigma_2(0)+\bigl(\tilde\Sigma_2(p)-\tilde\Sigma_2(0)\bigr)\,.
\]
Here
\[
\tilde\Sigma_2(0)=-\frac{\lambda^2}{3!}\int\frac{d^4k_1\,d^4k_2}{(2\pi)^8}\,\frac1{(k_1^2+m^2)(k_2^2+m^2)\bigl((k_1+k_2)^2+m^2\bigr)}\,.
\]
This divergent term does not depend on $p$; in this respect it is similar to
the contribution of the diagram
\[
\dia{\draw[prop,densely dashed] (-2,0) -- (2,0); \draw[prop] (0,0.6) circle (0.6); \node[vtx] at (0,0) {};}
\]
Clearly, it also can be absorbed into the mass renormalization, by suitable
modification of the counterterm
\[
\frac{\delta m^2(\Lambda)}2\,\varphi^2\,.
\]
The difference
\begin{multline*}
\tilde\Sigma_2(p)-\tilde\Sigma_2(0)\\
=-\frac{\lambda^2}{3!}\int\frac{d^4k_1\,d^4k_2}{(2\pi)^8}\,\frac1{(k_1^2+m^2)(k_2^2+m^2)}\,
\frac{2(pk_1)+2(pk_2)-p^2}{\bigl((k_1+k_2)^2+m^2\bigr)\bigl((k_1+k_2-p)^2+m^2\bigr)}
\end{multline*}
still diverges (logarithmically). If the cutoff is introduced, this
contribution is written as
\[
\sim p^2\log\frac{\Lambda^2}{m^2}+\text{finite}\,.
\]
This divergent term cannot be absorbed by the mass renormalization, since it
depends on the momentum $p$. Note however that this dependence has simple form
$\sim p^2$.

Recall that the proper vertex $\Gamma^{(2)}$ has the form
\[
\tilde\Gamma^{(2)}(p)=p^2+m^2+\text{loop corrections}\,,
\]
where the $p^2$ term originates from the term
\[
\frac12\,(\partial\varphi)^2
\]
term in the original action. This suggests that the above divergent term
$\sim p^2$ can be understood in terms of the ``field renormalization''.
Originally, the field $\varphi$ appeared in the theory as the integration
variable in the functional integral. Let us rename this integration variable as
$\varphi_0$, and call it the ``bare field''. The new divergent term shows that
the correlation functions of this bare field retains the dependence on the
cutoff parameter $\Lambda$, even after the mass and coupling renormalizations
are carried out. The way this dependence appears (through the $p^2$ term)
suggests that one can define renormalized field,
\[
\varphi(x)=Z^{-1/2}(\Lambda)\,\varphi_0(x)\,,
\]
with certain cutoff-dependent renormalization constant $Z(\Lambda)$, such that
the correlation functions
\[
\bigl\langle\varphi(x_1)\cdots\varphi(x_n)\bigr\rangle
\]
admit finite $\Lambda\to\infty$ limit.

We are ready to formulate the \emph{renormalization program}.

\subsection*{Renormalization program (\texorpdfstring{$\varphi^4$}{phi**4} theory)}

1. We start with the action
\[
\mathcal A=\int d^4x\,\Bigl[\frac12\,(\partial\varphi_0)^2+\frac{m_0^2}2\,\varphi_0^2+\frac{\lambda_0}{4!}\,\varphi_0^4\Bigr]\,,
\]
which involves the bare field $\varphi_0$, the bare mass $m_0$, and the bare
coupling constant $\lambda_0$.

2. We introduce some cutoff, with the cutoff momentum $\Lambda$, to make all
diagram contributions finite but cutoff-dependent (this can be done in many
ways, as we shall discuss later).

3. The bare parameters $m_0^2$, $\lambda_0$, and the field renormalization
constant $Z$ should be given certain dependence on the cutoff momentum
$\Lambda$,
\[
m_0^2(\Lambda)\,,\qquad\lambda_0(\Lambda)\,,\qquad Z(\Lambda)\,,
\]
such that the correlation functions of the renormalized field
$\varphi(x)=Z^{-1/2}\varphi_0(x)$ will have finite $\Lambda\to\infty$ limit.

In practice, one uses the perturbation theory to evaluate the correlation
functions. The above program can be further reformulated as the

\subsection*{Renormalized perturbation theory}

Let us rewrite the action as
\[
\mathcal A=\int d^4x\,\Bigl[\underbrace{\frac12\,(\partial\varphi)^2+\frac{m^2}2\,\varphi^2}_{\text{free part}}
+\underbrace{\frac\lambda{4!}\,\varphi^4}_{\text{interaction part}}
+\underbrace{\frac{\delta Z}2\,(\partial\varphi)^2+\frac{\delta m^2}2\,\varphi^2+\frac{\delta\lambda}{4!}\,\varphi^4}_{\text{counterterms}}\Bigr]\,,
\]
where $m$ is the actual mass and $\lambda$ is suitably defined renormalized
coupling constant (to be discussed later), and $\varphi$ is the renormalized
field. The last three terms are the \emph{counterterms}. The identity with the
original action implies
\[
1+\delta Z=Z\,,\qquad m^2+\delta m^2=Z\,m_0^2\,,\qquad\lambda+\delta\lambda=Z^2\,\lambda_0\,.
\]
In perspective, the renormalized perturbation theory is an expansion in the
renormalized coupling constant $\lambda$. We should assume that the
counterterm coefficients $\delta Z$, $\delta m^2$, and $\delta\lambda$
themselves depend perturbatively (i.e.\ as the power series) on $\lambda$,
\[
\delta Z=z_2\,\lambda^2+z_3\,\lambda^3+\cdots\,,\qquad
\delta m^2=b_1\,\lambda+b_2\,\lambda^2+\cdots\,,\qquad
\delta\lambda=a_2\,\lambda^2+a_3\,\lambda^3+\cdots\,,
\]
with $\Lambda$-dependent coefficients. We must show that these coefficients can
be adjusted in such a way that the correlation functions of the renormalized
field $\varphi$ are $\Lambda$-independent order by order in $\lambda$.

Let me stress that this renormalization program is based on our observation
that the $\varphi^4$ perturbation theory contains only finitely many basic
types of divergent diagrams (the ``primitive divergences''). Namely, the
primitively divergent diagrams are in $\Gamma^{(2)}$ and $\Gamma^{(4)}$. The
diagrams for higher vertices $\Gamma^{(n)}$ can be divergent only through
subdiagrams of the $\Gamma^{(2)}$ and $\Gamma^{(4)}$ types. The idea of the
renormalization program is that the three primitive divergences discussed above
can be absorbed in three renormalizations
\[
m_0^2(\Lambda)=Z^{-1}(\Lambda)\,\bigl(m^2+\delta m^2(\Lambda)\bigr)\,;\qquad
\lambda_0(\Lambda)=Z^{-2}(\Lambda)\,\bigl(\lambda+\delta\lambda(\Lambda)\bigr)\,;\qquad
\varphi_0=Z^{\frac12}(\Lambda)\,\varphi\,,
\]
order by order in the renormalized coupling constant $\lambda$.

Before we turn to execution of this program in $\varphi^4$ theory let us briefly
discuss divergences and renormalizations in a scalar field theory whose
interaction term is more general polynomial in $\varphi$.

\subsection*{Divergences in Scalar Theories}

It is instructive to consider such field theory in the space of $d$
dimensions, with generic $d$. We will consider general scalar theories with the
actions of the form
\begin{equation}\label{eq:genscalar}
\mathcal A=\int d^dx\,\Bigl[\frac12\,(\partial\varphi_0)^2+\frac{m_0^2}2\,\varphi_0^2+\sum_{k=3}^N\frac{\lambda_{0\,k}}{k!}\,\varphi_0^k\Bigr]\,.
\end{equation}
The Feynman rules remain the same as in 4-dimensional theory, except for the
momentum integrations go over $d$-dimensional momentum space,
\[
\frac{d^4k}{(2\pi)^4}\ \to\ \frac{d^dk}{(2\pi)^d}\,,
\]
and of course the diagrams contain $k$-leg vertices associated with the
couplings $\lambda_{0\,k}$.

As before, we can consider a generic diagram contributing to
$\tilde\Gamma^{(n)}$, containing $P$ propagators and $I$ $d$-momentum
integrations. Because now the momenta are $d$-dimensional, the superficial
degree of divergence of this diagram is
\[
D=d\cdot I-2\cdot P\,.
\]
Assume that the diagram contains $V_k$ $k$-leg vertices. Analysis similar to
that we have made for $\varphi^4$ theory reveals two identities
\[
2P+n=\sum_kk\cdot V_k\,,\qquad I=P-\sum_kV_k+1\,.
\]
Combining these equations one finds
\[
D=\sum_k\Bigl(\frac{d-2}2\,k-d\Bigr)V_k-\frac{d-2}2\,n+d\,.
\]
This formula admits simple interpretation in terms of dimensional counting.
Note that in our units $c=\hbar=1$ there is only one independent unit, which we
take to be mass unit. I will denote by $[X]$ the mass dimension of a quantity
$X$, for example
\[
[\textit{mass}]=1\,,\qquad[\textit{length}]=-1\,.
\]
The action is dimensionless and therefore it follows from its form that
\begin{equation}\label{eq:engdim}
[\varphi_0]=\frac{d-2}2\,,\qquad[m_0^2]=2\,,\qquad[\lambda_{0\,k}]=d-\frac{d-2}2\,k\,.
\end{equation}
I want to note in passing that this simple dimensional analysis is applied to
the \emph{bare} quantities. We will see later that because of renormalization
constant $Z(\Lambda)$ the renormalized field $\varphi$ can have different
dimension. This is why the above dimensions in \eqref{eq:engdim} are often
referred to as canonical (or ``engineering'') dimensions.

It is easy to check that
\[
\bigl[\tilde\Gamma^{(n)}\bigr]=d-\frac{d-2}2\,n\,.
\]
Therefore the equation for the superficial degree can be written as
\begin{equation}\label{eq:Ddim}
D=\bigl[\tilde\Gamma^{(n)}\bigr]-\sum_k\bigl[\lambda_{0\,k}\bigr]\cdot V_k\,.
\end{equation}
Now the meaning of this equation must be obvious. For $\Lambda\gg|p_i|$ the
dominating contribution of the diagram with $V_k$ vertices $\lambda_k$ is
\begin{equation}\label{eq:GammaLambda}
\tilde\Gamma^{(n)}\sim\Bigl[\prod_k\lambda_{0\,k}^{V_k}\Bigr]\,\Lambda^D\,,
\end{equation}
and our equation simply describes the balance of dimensions.

It is clear from the Eq.~\eqref{eq:Ddim} that the mass dimensions of the
coupling constants $\lambda_{0\,k}$ must play the key role in the analysis of
the perturbative divergences. Suppose some coupling $\lambda_{0\,k}$ in the
action have strictly negative mass dimension, $[\lambda_{0\,k}]<0$. Then there
are divergent contributions to $\tilde\Gamma^{(n)}$ with any $n$ which come
from the diagrams with sufficiently large $V_k$. In other words such theory has
infinitely many primitive divergences which of course cannot be absorbed by any
finite number of counterterms. Quantum field theories of this type are called
(perturbatively) \emph{non-renormalizable}. Overall consistency of
non-renormalizable theories is questionable (but not impossible). From pure
pragmatic point of view, the necessity to introduce infinitely many
counterterms brings in also infinitely many free parameters, and predictive
power of such theories is limited.

If mass dimensions of all coupling constants in the action are non-negative,
the equation \eqref{eq:Ddim} shows that there is only finite number of
primitively divergent proper vertices (this is true if $d>2$, the case $d=2$ is
exceptional and must be studied separately); these divergences then can be
absorbed by finitely many renormalization constants. For this reason such
theories are called (perturbatively) \emph{renormalizable}. Moreover, if all
$\lambda_{0\,k}$ have strictly positive mass dimensions there is only finite
number of divergent diagrams. The theories of this last kind are often referred
to as \emph{super-renormalizable}, in which case the term ``renormalizable'' is
reserved for the theories with infinitely many divergent diagrams but finitely
many primitive divergences; in view of \eqref{eq:Ddim} this requires that at
least one of the coupling constants is dimensionless. We will see later when
studying renormalization group that super-renormalizable scalar theories (with
the polynomial $V(\varphi_0)$ bounded from below) are consistent field
theories. Overall consistency of renormalizable theories requires more subtle
analysis, but at least they make sense perturbatively.

One observes from \eqref{eq:engdim} that for $d>6$ there are no (interacting)
renormalizable scalar theories. For $6\ge d>2$ renormalizable field theories
are
\[
\begin{array}{ll}
\varphi^3 & \text{in}\quad d=6\quad\text{dimensions,}\\
\varphi^3,\ \varphi^4 & \text{in}\quad d=4\quad\text{dimensions,}\\
\varphi^3,\ \varphi^4,\ \varphi^5,\ \varphi^6 & \text{in}\quad d=3\quad\text{dimensions.}
\end{array}
\]
It is not difficult to show that any polynomial in $\varphi$ is
super-renormalizable in $d=2$ dimensions.

I want to stress here that nonexistence of perturbatively renormalizable field
theories in high dimensionalities does not imply that consistent field
theories are limited to low space-time dimensions; there may exist perfectly
consistent field theories which are just ``too far'' from free field theory to
admit meaningful perturbative interpretation. I hope to discuss this
possibility in greater detail later.

% Source pp. 111-134 (Lectures13-18.pdf).
\section{Systematics of Renormalization}

% local diagram pieces for this section
\newcommand{\legs}[1]{\foreach \p/\a/\n in {#1} {\draw[prop,densely dashed] (\p) -- ++(\a:1.1) node[pos=1.18,font=\tiny] {\n};}}
\newcommand{\bubS}{\diac[0.4]{\coordinate (L) at (-0.9,0); \coordinate (R) at (0.9,0);
  \draw[prop] (0,0) ellipse (0.9 and 0.55); \legs{L/135/1,L/225/2,R/45/3,R/315/4}
  \node[vtx] at (L) {}; \node[vtx] at (R) {};}}
\newcommand{\bubT}{\diac[0.4]{\coordinate (T) at (0,0.8); \coordinate (B) at (0,-0.8);
  \draw[prop] (0,0) ellipse (0.55 and 0.8); \legs{T/150/1,T/30/3,B/210/2,B/330/4}
  \node[vtx] at (T) {}; \node[vtx] at (B) {};}}
\newcommand{\bubU}{\diac[0.4]{\coordinate (T) at (0,0.8); \coordinate (B) at (0,-0.8);
  \draw[prop] (0,0) ellipse (0.55 and 0.8);
  \draw[prop,densely dashed] (T) -- ++(150:1.1) node[pos=1.18,font=\tiny] {1};
  \draw[prop,densely dashed] (B) -- ++(210:1.1) node[pos=1.18,font=\tiny] {2};
  \draw[prop,densely dashed] (T) .. controls (0.9,0.6) and (0.7,-1.4) .. (1.2,-1.6) node[right,font=\tiny] {4};
  \draw[prop,densely dashed] (B) .. controls (0.9,-0.6) and (0.7,1.4) .. (1.2,1.6) node[right,font=\tiny] {3};
  \node[vtx] at (T) {}; \node[vtx] at (B) {};}}
\newcommand{\ctX}[1]{\diac[0.4]{\coordinate (O) at (0,0); \legs{O/135/1,O/225/2,O/45/3,O/315/4}
  \node[otimesv,minimum size=5pt] at (O) {}; \node[font=\tiny] at (0,0.55) {(#1)};}}
\newcommand{\ctline}[1]{\dia[0.4]{\draw[prop,densely dashed] (-2,0) -- (2,0); \node[otimesv,minimum size=5pt] at (0,0) {}; \node[font=\tiny] at (0,0.55) {(#1)};}}

If the theory under consideration is renormalizable, like $\varphi^4$ theory in
4 dimensions, one can develop renormalized perturbation theory. According to the
program outlined above one has to start with the action in its ``bare'' form,
implicitly equipped with some ``regularization'', or ``cutoff'', i.e.\ with some
modification of the theory at the momenta $\ge\Lambda$ which makes all integrals
in perturbation theory convergent. There are many possible implementations of
the cutoff which can be more or less convenient depending on circumstances. Let
us discuss some more frequently used ones.

\subsection*{Regularization methods}

\paragraph{Lattice regularization:} We already discussed this regularization in
relation with the definition of the functional integral. In this approach
continuous space $\R^d$ is replaced by $d$-dimensional, say hypercubic, lattice
with some lattice spacing $\Delta$ which plays the role of inverse cutoff
parameter $\Lambda^{-1}$,
\[
x\to x_n=\Delta\sum_{\mu=1,\dots,d}e_\mu\,n_\mu\,,
\]
where $e_\mu$ is the unit vector in the direction $\mu$ and $n_1,\dots,n_d$ are
integers. The lattice action is obtained by replacing derivatives by finite
differences,
\begin{equation}\label{eq:Alat}
\mathcal A_{\text{lat}}=\Delta^d\sum_{x\in\Delta\mathbb Z^d}\left[\frac12\sum_{\mu=1,\dots,d}\left(\frac{\varphi_0(x+\Delta e_\mu)-\varphi_0(x)}\Delta\right)^2
+\frac{m_0^2}2\,\varphi_0^2(x)+V\bigl(\varphi_0(x)\bigr)\right]
\end{equation}
and the integration measure in the functional integral is taken as
\[
[D\varphi]\to\prod_{x\in\Delta\mathbb Z^d}d\varphi_0(x)\,.
\]
Advantage of this method is that it makes sense non-perturbatively. The lattice
theory \eqref{eq:Alat} is perfectly meaningful in classical statistical
mechanics. Although it is hardly possible to calculate the integral
analytically, one can try to apply numerical approach, say using the
Monte-Carlo method. This approach is known as the lattice Field Theory. At the
same time perturbative calculations on the lattice are relatively complicated.
In addition, continuous space symmetries are broken, and one only expects them
to be restored after the cutoff is removed.

\paragraph{Proper-time regularization:} We have seen before that the
momentum-space propagator $\tilde D(p)$ admits Schwinger's proper-time
representation
\[
\tilde D(p)=\int_0^\infty ds\,e^{-s(p^2+m^2)}\,,
\]
where $s$ is interpreted as the rescaled length of the path of relativistic
particle in the Euclidean space-time. One can exclude the paths which are
``too short'' by replacing $\tilde D(p)\to\tilde D_\Lambda(p)$,
\[
\tilde D_\Lambda(p)=\int_{1/\Lambda^2}^\infty ds\,e^{-s(p^2+m^2)}\,.
\]
Note that this is particular case of the regularized propagator we considered
before,
\begin{equation}\label{eq:regprop}
\frac1{p^2+m^2}\,\Phi(p^2/\Lambda^2)\,,
\end{equation}
with $\Phi(x)=e^{-x}$ (up to the constant factor $e^{-m^2/\Lambda^2}$, which
tends to $1$ as $\Lambda\to\infty$).

\paragraph{Pauli-Villars regularization} is another version of
\eqref{eq:regprop}, with
\begin{equation}\label{eq:PV}
\Phi_{PV}(p^2/\Lambda^2)=\frac{\Lambda^2}{p^2+\Lambda^2}\,.
\end{equation}
In this case the cutoff propagator takes the form
\[
\tilde D(p)\to\frac1{p^2+m^2}\,\Phi_{PV}(p^2/\Lambda^2)\to\frac1{p^2+m^2}-\frac1{p^2+\Lambda^2}\,,
\]
equivalent to \eqref{eq:PV} when $\Lambda\gg m$.

\paragraph{Dimensional regularization:} This is the most technically advanced
(although perhaps the least physically transparent) regularization method. Its
advantage is that it usually preserves important symmetries (in particular the
gauge symmetry in the gauge theories) and significantly simplifies
calculations, and for these reasons this method is used today in perturbative
calculations more frequently than others. The idea is this. We have already
observed that lowering the space dimension generally improves the
large-momentum convergence of Feynman diagrams. Suppose that for $d$
sufficiently low the momentum integral associated with given diagram is
convergent without any additional cutoff. Suppose that in addition we have
managed to calculate this integral as an analytic function of $d$. Then we can
analytically continue the result to the ``physical value'' $d=4$; the
4-dimensional divergences manifest themselves as singularities (poles) in the
variable $d$ at $d=4$. On more formal level one can define the dimensional
continuation of momentum integrals by three conditions, obviously valid for
convergent integrals, and postulated for any $d$:
\[
\begin{array}{lll}
\text{(i)} & \displaystyle\int d^dk\,F(k+p)=\int d^dk\,F(k) & \text{translation}\\[10pt]
\text{(ii)} & \displaystyle\int d^dk\,F(Ck)=|C|^{-d}\int d^dk\,F(k) & \text{dilation}\\[10pt]
\text{(iii)} & \displaystyle\int d^{d_1}k_1\,d^{d_2}k_2\,F(k_1)G(k_2)=\int d^{d_1}k\,F(k_1)\int d^{d_2}k_2\,G(k_2)\,,\quad d=d_1+d_2 & \text{factorization}
\end{array}
\]
Let me illustrate this approach by simple example. Consider the integral
\begin{equation}\label{eq:dimreg1}
\int\frac{d^dk}{(2\pi)^d}\,\frac1{(k^2+m^2)^2}\,.
\end{equation}
This integral converges for $d<4$. For any $d<4$ it can be transformed as
follows
\begin{equation}\label{eq:dimreg2}
\int_0^\infty s\,ds\int\frac{d^dk}{(2\pi)^d}\,e^{-s(k^2+m^2)}=\int_0^\infty s\,ds\,(4\pi s)^{-\frac d2}\,e^{-m^2s}\,,
\end{equation}
where the properties (iii) and (ii) were used in evaluating the $k$-integral.
Using the definition of the Euler's Gamma function
\[
\Gamma(z)=\int_0^\infty dt\,t^{z-1}\,e^{-t}
\]
we obtain for $d<4$
\begin{equation}\label{eq:dimreg3}
\frac{(m^2)^{\frac d2-2}}{(4\pi)^{\frac d2}}\,\Gamma\left(2-\frac d2\right)\,.
\end{equation}
As is known the Gamma function $\Gamma(z)$ is analytic (actually, meromorphic)
function of $z$ with poles at $z=0,-1,-2,\dots$. In particular, the above
expression has pole at $d=4$, i.e.\ exactly where the integral diverges
logarithmically. For $d=4-\epsilon$, $\epsilon\to0$
\[
\Gamma\left(2-\frac d2\right)=\frac2\epsilon-\gamma+O(\epsilon)
\]
where $\gamma=0.5772\ldots$ is the Euler's constant.

It is interesting to explain on this example the interplay between dimensional
continuation and cutoff regularization. If we regularize the propagators in
\eqref{eq:dimreg1}, say by proper-time regularization (which amounts to
replacing $\int_0^\infty ds\to\int_{1/\Lambda^2}^\infty ds$ in
\eqref{eq:dimreg2}) the integral \eqref{eq:dimreg1} becomes regular function of
$d$ for all $d$ including $d=4$. In the neighborhood of $d=4$ one would have
instead of \eqref{eq:dimreg2}
\[
\frac{(m^2)^{\frac d2-2}}{(4\pi)^{\frac d2}}\int_{m^2/\Lambda^2}^\infty dt\,t^{1-\frac d2}\,e^{-t}\,,\qquad t=m^2s\,.
\]
The integral here can be transformed using integration by parts,
\[
\int_{m^2/\Lambda^2}^\infty dt\,t^{1-\frac d2}\,e^{-t}=\frac2{4-d}\Bigl(t^{2-\frac d2}\,e^{-t}\Big|_{m^2/\Lambda^2}^\infty+\int_{m^2/\Lambda^2}^\infty dt\,t^{2-\frac d2}\,e^{-t}\Bigr)\,.
\]
For $\epsilon=4-d\ll1$ and $m^2\ll\Lambda^2$ the cutoff integral becomes
\[
\frac1{(4\pi)^{\frac d2}}\,\frac2{4-d}\,\Bigl((m^2)^{\frac d2-2}-(\Lambda^2)^{\frac d2-2}\Bigr)\,.
\]
If at fixed $d<4$ we send the cutoff momentum $\Lambda$ to infinity we obtain
the ``dimensionally continued'' expression with the pole at $d=4$. On the other
hand, if at fixed $\Lambda$ we take the limit $d\to4$ we get a finite result
with $\log\Lambda^2$ replacing the pole in $4-d$.

With this understanding one can often use dimensional continuation to do the
calculations efficiently and then interpret the results in terms of cutoff
regularization. Instead of giving the bare parameters $m_0^2$, $\lambda_0$ and
$Z$ dependence on the cutoff momentum $\Lambda$, within the framework of
dimensional regularization one adjusts their dependence on $\epsilon=4-d$, so
that the correlation functions of renormalized fields $\varphi$ have finite
limit $d\to4$.

Some generalizations of the integral \eqref{eq:dimreg1} will be useful in
future calculations. For instance,
\begin{equation}\label{eq:dimregalpha}
\int\frac{d^dk}{(2\pi)^d}\,\frac1{(k^2+m^2)^\alpha}=\frac1{\Gamma(\alpha)}\int_0^\infty ds\,s^{\alpha-1}\int\frac{d^dk}{(2\pi)^d}\,e^{-s(k^2+m^2)}
=\frac{(m^2)^{\frac d2-\alpha}}{(4\pi)^{\frac d2}}\,\frac{\Gamma(\alpha-\frac d2)}{\Gamma(\alpha)}\,,
\end{equation}
\begin{equation}\label{eq:dimregk2}
\begin{split}
\int\frac{d^dk}{(2\pi)^d}\,\frac{k^2}{(k^2+m^2)^\alpha}&=\int\frac{d^dk}{(2\pi)^d}\left[\frac1{(k^2+m^2)^{\alpha-1}}-\frac{m^2}{(k^2+m^2)^\alpha}\right]\\
&=\frac{(m^2)^{\frac d2-\alpha+1}}{(4\pi)^{\frac d2}}\,\frac d2\,\frac{\Gamma(\alpha-\frac d2-1)}{\Gamma(\alpha)}\,.
\end{split}
\end{equation}
In deriving \eqref{eq:dimregk2} the defining property of the $\Gamma$ function
$\Gamma(z+1)=z\,\Gamma(z)$ was used.

\subsection*{Renormalizations in \texorpdfstring{$\varphi^4$}{phi**4} theory, revisited}

Let us recall the basic idea of the (perturbative) renormalization program. To
develop renormalized perturbation theory, one starts with the action
\[
\mathcal A=\int d^dx\,\left[\frac12\,(\partial\varphi)^2+\frac{m^2}2\,\varphi^2+\frac\lambda{4!}\,\varphi^4+\text{counterterms}\right]\,,
\]
expressed in terms of the renormalized quantities. I have written it in the
$d$-dimensional form having in mind dimensional regularization. This action
contains counterterms
\[
\text{counterterms}=\frac{\delta Z}2\,(\partial\varphi)^2+\frac{\delta m^2}2\,\varphi^2+\frac{\delta\lambda}{4!}\,\varphi^4\,.
\]
We are planning to treat the theory perturbatively, with
\[
\mathcal A_I=\int_x\left[\frac\lambda{4!}\,\varphi^4+\text{counterterms}\right]
\]
taken as the perturbation.

The counterterms bring in additional vertices, so the diagrams will contain the
following elements (momentum $\delta$-functions are implicit)
\[
\dia{\draw[prop] (-1.5,0) -- (1.5,0); \draw[arr] (-0.5,0.35) -- (0.5,0.35) node[above,pos=0.5] {\small$p$};}=\frac1{p^2+m^2}\,,\qquad
\dia{\pic{xvertex};}=-\lambda\,,\qquad
\dia{\draw[prop,densely dashed] (-1.6,0) -- (1.6,0); \node[otimesv] at (0,0) {};}=-\bigl(p^2\,\delta Z+\delta m^2\bigr)\,,\qquad
\dia{\draw[prop] (-0.6,-0.6) -- (0.6,0.6); \draw[prop] (-0.6,0.6) -- (0.6,-0.6); \node[otimesv] at (0,0) {};}=-\delta\lambda\,.
\]
The counterterms $\delta Z$, $\delta m^2$, $\delta\lambda$ are assumed to be
power series in $\lambda$, i.e.
\begin{align*}
\delta Z&=z^{(1)}\,\lambda+z^{(2)}\,\lambda^2+\cdots\,,\\
\delta m^2&=b^{(1)}\,\lambda+b^{(2)}\,\lambda^2+\cdots\,,\\
\delta\lambda&=a^{(1)}\,\lambda+a^{(2)}\,\lambda^2+\cdots\,,
\end{align*}
with the coefficients dependent on $\Lambda$, or, in dimensional
regularization, on
\[
\epsilon=4-d\,,
\]
in such a way that all diagram contributions to the correlation functions of
renormalized fields have finite limit as $\Lambda\to\infty$, or, respectively,
$\epsilon\to0$ in dimensional regularization, order by order in $\lambda$. This
is general strategy of renormalized perturbation theory. Specific calculations
depend on the choice of the so called ``Renormalization Scheme''.

\subsection*{Renormalization Schemes}

Obviously, just satisfying this requirement (that the counterterms cancel all
singular contributions of ordinary diagrams) is not sufficient to fix the
counterterms uniquely. For instance, the renormalized field $\varphi$ always can
be renormalized further by multiplying it by a finite (i.e.\ regular at
$\epsilon=0$) constant
\[
\varphi=Z_{\text{finite}}^{\frac12}\,\tilde\varphi\,;
\]
similar finite renormalizations of $m^2$ and $\lambda$ are possible. Obviously,
such renormalizations do not spoil finiteness of renormalized perturbation
theory. We see that there is a whole class of renormalizable perturbation
theories, related one to another by finite renormalizations. These different
renormalized perturbation theories are called the \emph{Renormalization
Schemes}. Of course all these renormalized perturbation theories describe the
same field theory, and yield the correlation functions with identical physical
content. What is different, besides finite renormalizations of the correlation
functions due to the field renormalizations, is \emph{parametrization} of the
correlation functions. If $\varphi$, $m^2$, $\lambda$ and $\tilde\varphi$,
$\tilde m^2$, $\tilde\lambda$ are renormalized quantities in two different
renormalization schemes, the corresponding correlation functions are related as
\[
\Gamma^{(n)}(p_i|m^2,\lambda)=\Bigl[Z_{\text{finite}}(m^2,\lambda)\Bigr]^{-\frac n2}\,
\tilde\Gamma^{(n)}\bigl(p_i|\tilde m^2(m^2,\lambda),\tilde\lambda(m^2,\lambda)\bigr)\,.
\]
One can think of the parameters $(m^2,\lambda)$ as coordinates in the ``space
of $\varphi^4$ theories''. Different renormalization schemes correspond to
different coordinate systems in that space.

A renormalization scheme is usually defined by certain \emph{Normalization
conditions}.

\subsection*{Normalization conditions}

These are conditions that specify the choice of the renormalized parameters
$m^2$, $\lambda$, as well as the overall normalization of $\varphi$. This is
usually done by relating the renormalized parameters to the physical quantities
(this is not so in the so called ``minimal subtraction scheme'' which has more
technical definition; we will discuss it in some details later). In practice,
the schemes where these relations are simple are more convenient.

One such normalization condition we have discussed already when we defined the
renormalized mass parameter $m^2$ as the physical mass squared. Namely,
assuming that $\tilde\Gamma^{(2)}(p^2)$ has zero at some negative value of
$p^2$ we have defined $-m^2$ as the position of this zero. Note that this
condition does not depend on the overall normalization of $\varphi$. The latter
must be fixed by a separate condition. One can fix for instance the slope of
$\tilde\Gamma^{(2)}(p^2)$ at this zero $p^2=-m^2$; a convenient normalization
condition is to set this slope equal 1, i.e.
\begin{equation}\label{eq:SIa}
\tilde\Gamma^{(2)}(p^2)=p^2+m^2+O\bigl((p^2+m^2)^2\bigr)\qquad\text{as}\qquad p^2+m^2\to0\,.
\end{equation}
Equivalently, this condition means that the two-point correlation function
\begin{equation}\label{eq:SIW}
\tilde W(p)=\frac1{p^2+m^2}+O(1)\qquad\text{as}\qquad p^2+m^2\to0
\end{equation}
has a pole at $p^2=-m^2$ with the residue one. We also need to fix
normalization of $\lambda$. This is usually done by imposing certain condition
on $\tilde\Gamma^{(4)}$. The condition often taken in high-energy theory
calculations is
\begin{equation}\label{eq:SIb}
\tilde\Gamma^{(4)}(p,p,-p,-p)\Big|_{p^2=-m^2}=\lambda
\end{equation}
(the quantity in the left-hand side is related to $2\to2$ forward scattering
amplitude at zero energy). The equations \eqref{eq:SIW} and \eqref{eq:SIb} above
specify one possible renormalization scheme (I will refer to it as the scheme
SI).

Of course there are other schemes. Another scheme which is more frequently used
in statistical mechanics context is specified by the normalization conditions
\begin{align}
\tilde\Gamma^{(2)}(p^2)&=\tilde m^2+p^2+O(p^4)\,,\label{eq:SIIa0}\\
\tilde\Gamma^{(4)}(0,0,0,0)&=\tilde\lambda\,.\label{eq:SIIb0}
\end{align}
Here I have used the notations $\tilde m^2$ and $\tilde\lambda$ for the
associated renormalized parameters to stress that these parameters are
different from those defined in the scheme (SI) (in particular $\tilde m^2$
does not coincide with the physical mass), i.e.\ it defines different
``coordinates''. Henceforth this scheme is referred to as SII.

Let us see how this program works in few leading orders in the coupling
constant. The following calculation is equivalent to the one in PS, \S10.2,
10.5. There it is done using the normalization condition (SI). For the sake of
diversity I will instead adopt (SII). To facilitate the comparison I will use
notations $\tilde m^2$ and $\tilde\lambda$ for the renormalized parameters.
Also, to simplify notations I will write $\Gamma^{(n)}$ for the \emph{momentum
space} proper vertices (omitting the tilde).

\subsection*{Renormalization of \texorpdfstring{$\varphi^4$}{phi**4} at the leading (one loop) order}

Let me describe in some details the renormalization of the $\varphi^4$ theory
using dimensional regularization. For this illustration I will adopt the
renormalization scheme SII, as was described above. For reference, that scheme
is defined by the normalization conditions
\begin{align}
\tilde\Gamma^{(2)}(p^2)&=\tilde m^2+p^2+O(p^4)\,,\label{eq:SIIa}\\
\tilde\Gamma^{(4)}(0,0,0,0)&=\tilde\lambda\,.\label{eq:SIIb}
\end{align}
The normalization conditions specify the meaning of the renormalized parameters
$\tilde m^2$, $\tilde\lambda$, as well as the normalization of the field
$\varphi(x)$.

Consider first the two-point function,
\[
\Gamma^{(2)}(p^2)=p^2+\tilde m^2+\Sigma(p^2)\,.
\]
The normalization condition \eqref{eq:SIIa} states
\begin{equation}\label{eq:SIIsigma}
\Sigma(0)=0\,;\qquad\Sigma'(0)=0\,.
\end{equation}
In the leading order in $\tilde\lambda$ there are two diagrams contributing to
$\Sigma$,
\[
-\Bigl(\underbrace{\dia[0.4]{\draw[prop,densely dashed] (-2,0) -- (2,0); \draw[prop] (0,0.6) circle (0.6); \node[vtx] at (0,0) {};}}_{(a)}
+\underbrace{\ctline{1}}_{(b)}\Bigr)\,,
\]
where the diagram (b) comes from the $\sim\tilde\lambda$ terms in $\delta Z$ and
$\delta\tilde m^2$,
\[
\ctline{1}=-\tilde\lambda\,\bigl(z^{(1)}\,p^2+b^{(1)}\bigr)\,;
\]
The coefficients $z^{(1)}$ and $b^{(1)}$ are to be adjusted to satisfy
\eqref{eq:SIIsigma}.

We already have considered the diagram (a). It is a constant independent of
$p^2$; its value is
$(a)=-\frac{\tilde\lambda}2\int\frac{d^dk}{(2\pi)^d}\frac1{k^2+\tilde m^2}$.
Therefore to satisfy \eqref{eq:SIIsigma} we have to choose
\[
z^{(1)}=0\,;\qquad b^{(1)}=\frac{(a)}{\tilde\lambda}=-\frac12\int\frac{d^dk}{(2\pi)^d}\,\frac1{k^2+\tilde m^2}
=-\frac12\,\frac{(\tilde m^2)^{\frac d2-1}}{(4\pi)^{\frac d2}}\,\Gamma(1-d/2)\,,
\]
where the dimensional regularization was used to evaluate the integral. Note
that at this order there is a trivial cancellation between the diagrams (a) and
(b). Because of this, in considering higher-order diagrams we can in fact set
$b^{(1)}=0$ if we at the same time decide to ignore the ``bubble'' contributions
(a) in all propagators.

Consider now the leading perturbative corrections to the four-point function
$\Gamma^{(4)}$. At the tree level $\Gamma^{(4)}=\tilde\lambda$ and the
normalization condition \eqref{eq:SIIb} is satisfied at this order if we set
$a^{(1)}=0$, i.e.\ $\delta\tilde\lambda=O(\tilde\lambda^2)$. Nontrivial
contributions appear at the order $\tilde\lambda^2$; the corresponding diagrams
are
\[
\Gamma^{(4)}(p_1,\dots,p_4)=\tilde\lambda-\Bigl(\underset{(a)}{\bubS}+\underset{(b)}{\bubT}+\underset{(c)}{\bubU}+\underset{(d)}{\ctX{2}}\Bigr)
\]
where again the diagram (d) represents the contribution of the counterterm
$\delta\lambda$ in the order $\tilde\lambda^2$,
\begin{equation}\label{eq:a2vertex}
\ctX{2}=-\tilde\lambda^2\,a^{(2)}\,.
\end{equation}
Each of the diagrams (a-c) can be expressed through the integral
\begin{equation}\label{eq:Ip2}
I(p^2)=\int\frac{d^dk}{(2\pi)^d}\,\frac1{(k^2+\tilde m^2)\bigl((p+k)^2+\tilde m^2\bigr)}\,,
\end{equation}
namely
\[
(a)+(b)+(c)=\frac{\tilde\lambda^2}2\,\Bigl(I(p_{12}^2)+I(p_{13}^2)+I(p_{14}^2)\Bigr)\,,
\]
where $p_{12}=p_1+p_2$, $p_{13}=p_1+p_3$, $p_{14}=p_1+p_4$.

The normalization condition \eqref{eq:SIIb} requires that
\[
\Bigl[(a)+(b)+(c)+(d)\Bigr]_{p_i=0}=0\,,
\]
which fixes the value of the counterterm constant in \eqref{eq:a2vertex},
\[
a^{(2)}=\frac32\,I(0)\,.
\]
With this we obtain
\begin{equation}\label{eq:G4oneloop}
\Gamma^{(4)}(p_1,\dots,p_4)=\tilde\lambda-\frac{\tilde\lambda^2}2\,\Bigl(I_r(p_{12}^2)+I_r(p_{13}^2)+I_r(p_{14}^2)\Bigr)+O(\tilde\lambda^3)\,,
\end{equation}
where
\[
I_r(p^2)=I(p^2)-I(0)\,.
\]
Our goal now is to show that while the integral \eqref{eq:Ip2} is singular at
$d=4$ the combination appearing in \eqref{eq:G4oneloop} has finite limit at
$d=4$.

The integrand in \eqref{eq:Ip2} can be transformed using the so called
\emph{Feynman parametrization}. The identity\footnote{More general identity
\begin{equation}\label{eq:FeynParGen}
\frac1{A^\alpha}\,\frac1{B^\beta}=\frac{\Gamma(\alpha+\beta)}{\Gamma(\alpha)\Gamma(\beta)}\int_0^1\frac{dw\,w^{\alpha-1}(1-w)^{\beta-1}}{\bigl[wA+(1-w)B\bigr]^{\alpha+\beta}}
\end{equation}
is often useful. To verify it, change the integration variable to
$w=Bz/(Bz+A(1-z))$ to reduce the integral to the Euler's ``beta integral''
\[
B(\alpha,\beta):=\int_0^1dz\,z^{\alpha-1}(1-z)^{\beta-1}=\frac{\Gamma(\alpha)\Gamma(\beta)}{\Gamma(\alpha+\beta)}\,.
\]}
\[
\frac1{AB}=\int_0^1\frac{du}{\bigl(uA+(1-u)B\bigr)^2}
\]
which can be checked by explicit evaluation of the integral in the r.h.s.,
allows one to write the integral \eqref{eq:Ip2} as
\begin{equation}\label{eq:IFeyn}
I(p^2)=\int_0^1du\int\frac{d^dk}{(2\pi)^d}\,\frac1{\bigl(\tilde m^2+k^2+2u\,(kp)+u\,p^2\bigr)^2}\,.
\end{equation}
The combination in the denominator can be put in the form
\[
\tilde m^2+k^2+2u\,(kp)+u\,p^2=\tilde m^2+u(1-u)\,p^2+(k+up)^2\,,
\]
and by applying the translation property (i) of the dimensionally continued
integrals we have
\[
I(p^2)=\int_0^1du\int\frac{d^dk}{(2\pi)^d}\,\frac1{\bigl(k^2+u(1-u)\,p^2+\tilde m^2\bigr)^2}\,.
\]
The momentum integral has the form \eqref{eq:dimregalpha} with $\alpha=2$, hence
\[
I(p^2)=\frac{\Gamma(2-\frac d2)}{(4\pi)^{\frac d2}}\int_0^1du\,\bigl[\tilde m^2+u(1-u)\,p^2\bigr]^{\frac d2-2}\,.
\]
We see that due to the factor $\Gamma(2-d/2)=2/\epsilon-\gamma+O(\epsilon)$,
$\epsilon:=4-d$, this expression has pole at $d=4$. However the combination
\[
I_r(p^2)=I(p^2)-I(0)=\frac{\Gamma(2-\frac d2)}{(4\pi)^{\frac d2}}\int_0^1du\,\Bigl(\bigl[\tilde m^2+u(1-u)\,p^2\bigr]^{\frac d2-2}-\bigl[\tilde m^2\bigr]^{\frac d2-2}\Bigr)
\]
has finite value at $d=4$ as the result of compensation between the pole of the
gamma function and zero of the integrand. One finds
\begin{align*}
I_r(p^2)\Big|_{d=4}&=-\frac1{(4\pi)^2}\int_0^1du\,\log\left(1+u(1-u)\,\frac{p^2}{\tilde m^2}\right)\\
&=-\frac1{16\pi^2}\left[\sqrt{\frac{p^2+4\tilde m^2}{p^2}}\,\log\frac{\sqrt{p^2+4\tilde m^2}+\sqrt{p^2}}{\sqrt{p^2+4\tilde m^2}-\sqrt{p^2}}-2\right]\,.
\end{align*}
I leave evaluating the last explicit expression as an \textbf{Exercise}. We see
that by giving the counterterm coefficient $a^{(2)}$ certain dependence on $d$
which is singular at $d=4$,
\begin{equation}\label{eq:a2}
a^{(2)}=\frac32\,I(0)=\frac32\,\frac{\Gamma(2-\frac d2)}{(4\pi)^{\frac d2}}\,(\tilde m^2)^{\frac d2-2}\,,
\end{equation}
we compensated the singularities of the diagrams (a), (b), (c) above.

\begin{exercise}
Compare the above calculation with the one in \S10.2 of PS. Find the relation
(to the order $\lambda^2$) between our renormalized parameters $\tilde m^2$,
$\tilde\lambda$ and parameters $m^2$, $\lambda$ defined there, which correspond
to the scheme (SI).
\end{exercise}

It is instructive to study renormalization at the next, two loop order. The two
loop contributions correspond to the terms $\sim\lambda^2$ in $\Gamma^{(2)}$
and $\lambda^3$ in $\Gamma^{(4)}$. Analysis of the two loop terms in
$\Gamma^{(4)}$ can be found in PS Sect.~10.5. I repeat it in the Appendix
below, along with the brief discussion of the two-loop divergences in
$\Gamma^{(2)}$.

To complete our analysis, we have to check that no new counterterms are
required to compensate for divergences in higher proper vertices
$\Gamma^{(6)}$, $\Gamma^{(8)}$, etc. At one-loop level the higher vertices are
finite without any counterterms since the diagrams
\[
-\Gamma^{(2n)}_{\text{1-loop}}=\diac[0.62]{\draw[prop] (0,0) circle (1.1);
  \foreach \a in {20,110,200,290} {\node[vtx] at (\a:1.1) {}; \draw[prop,densely dashed] (\a:1.1) -- ++(\a+30:0.9); \draw[prop,densely dashed] (\a:1.1) -- ++(\a-30:0.9);}
  \node[font=\tiny] at (245:1.6) {$\cdot\,\cdot$};}
\]
are convergent at $d=4$. The two-loop contributions to these vertices involve
divergent diagrams, with the divergences coming from the 2- or 4-leg
subdiagrams. For example, $\Gamma^{(6)}_{\text{2-loop}}$ receives contribution
from the diagram
\[
-\Gamma^{(6)}=\cdots+\diac[0.62]{\draw[prop] (0,0) circle (1.1);
  \draw[prop] (125:1.1) .. controls (0,0.2) .. (55:1.1);
  \foreach \a/\b in {125/150,55/30} {\node[vtx] at (\a:1.1) {}; \draw[prop,densely dashed] (\a:1.1) -- ++(\b:0.9);}
  \foreach \a/\b/\c in {200/180/240,320/290/350} {\node[vtx] at (\a:1.1) {}; \draw[prop,densely dashed] (\a:1.1) -- ++(\b:0.9); \draw[prop,densely dashed] (\a:1.1) -- ++(\c:0.9);}}
\]
which is singular as $1/\epsilon$. However, it is possible to check that the
one-loop counterterms
\[
\frac{\lambda\,b^{(1)}(\epsilon)}2\,\varphi^2+\frac{\lambda^2\,a^{(2)}(\epsilon)}{4!}\,\varphi^4
\]
are sufficient to cancel all such divergences. For instance, the above divergent
diagram comes along with two other divergent diagrams
\[
\diac[0.62]{\draw[prop] (0,0) circle (1.1); \draw[prop] (125:1.1) .. controls (0,0.2) .. (55:1.1);
  \foreach \a/\b/\l in {125/150/1,55/30/2} {\node[vtx] at (\a:1.1) {}; \draw[prop,densely dashed] (\a:1.1) -- ++(\b:0.9) node[pos=1.3,font=\scriptsize] {\l};}
  \foreach \a/\b/\c/\l/\m in {200/180/240/6/5,320/290/350/4/3} {\node[vtx] at (\a:1.1) {}; \draw[prop,densely dashed] (\a:1.1) -- ++(\b:0.9) node[pos=1.3,font=\scriptsize] {\l}; \draw[prop,densely dashed] (\a:1.1) -- ++(\c:0.9) node[pos=1.3,font=\scriptsize] {\m};}}
\qquad
\diac[0.62]{\draw[prop] (0,0) circle (1.1); \draw[prop] (125:1.1) .. controls (0,0.2) .. (55:1.1);
  \foreach \a/\b/\l in {125/150/2,55/30/1} {\node[vtx] at (\a:1.1) {}; \draw[prop,densely dashed] (\a:1.1) -- ++(\b:0.9) node[pos=1.3,font=\scriptsize] {\l};}
  \foreach \a/\b/\c/\l/\m in {200/180/240/6/5,320/290/350/4/3} {\node[vtx] at (\a:1.1) {}; \draw[prop,densely dashed] (\a:1.1) -- ++(\b:0.9) node[pos=1.3,font=\scriptsize] {\l}; \draw[prop,densely dashed] (\a:1.1) -- ++(\c:0.9) node[pos=1.3,font=\scriptsize] {\m};}}
\qquad
\diac[0.62]{\draw[prop] (0,0) circle (1.1); \draw[prop] (0,1.55) circle (0.45);
  \node[vtx] at (0,1.1) {}; \node[vtx] at (0,2.0) {};
  \draw[prop,densely dashed] (0,2.0) -- ++(135:0.9) node[pos=1.3,font=\scriptsize] {1}; \draw[prop,densely dashed] (0,2.0) -- ++(45:0.9) node[pos=1.3,font=\scriptsize] {2};
  \foreach \a/\b/\c/\l/\m in {200/180/240/6/5,320/290/350/4/3} {\node[vtx] at (\a:1.1) {}; \draw[prop,densely dashed] (\a:1.1) -- ++(\b:0.9) node[pos=1.3,font=\scriptsize] {\l}; \draw[prop,densely dashed] (\a:1.1) -- ++(\c:0.9) node[pos=1.3,font=\scriptsize] {\m};}}
\]
The sum of the singular part of these three diagrams is canceled by the diagram
\[
\diac[0.62]{\draw[prop] (0,0) circle (1.1);
  \node[otimesv] at (0,1.1) {};
  \draw[prop,densely dashed] (0,1.1) -- ++(135:0.9) node[pos=1.3,font=\scriptsize] {1}; \draw[prop,densely dashed] (0,1.1) -- ++(45:0.9) node[pos=1.3,font=\scriptsize] {2};
  \foreach \a/\b/\c/\l/\m in {200/180/240/6/5,320/290/350/4/3} {\node[vtx] at (\a:1.1) {}; \draw[prop,densely dashed] (\a:1.1) -- ++(\b:0.9) node[pos=1.3,font=\scriptsize] {\l}; \draw[prop,densely dashed] (\a:1.1) -- ++(\c:0.9) node[pos=1.3,font=\scriptsize] {\m};}}
\]
involving the one-loop counterterm vertex.

Similar situation persists to higher loops. Suppose we have renormalized the
theory (i.e.\ determined all the counterterms as the functions of $\epsilon$) up
to $L-1$ loops. Then contributions of all $L$-loop diagrams, with all
$L-1$-loop counterterms, to the effective action $\Gamma[\phi]$ can be written
as
\[
\Gamma^{\text{div}}_L[\phi]+\Gamma^{\text{finite}}_L[\phi]\,,
\]
where $\Gamma^{\text{div}}_L$ has the form
\[
-\tilde\lambda^L\int d^dx\,\left[\frac{Z^{(L)}(\epsilon)}2\,(\partial\phi)^2+\frac{B^{(L)}(\epsilon)}2\,\phi^2+\frac{\tilde\lambda\,A^{(L)}(\epsilon)}{4!}\,\phi^4\right]\,,
\]
with coefficients having $L$th order pole in $\epsilon$,
\begin{equation}\label{eq:divpoles}
Z^{(L)}=\sum_{k=1}^{L-1}\frac{Z^{(L)}_{-k}}{\epsilon^k}\,;\qquad
B^{(L)}=\sum_{k=1}^L\frac{B^{(L)}_{-k}}{\epsilon^k}\,;\qquad
A^{(L)}=\sum_{k=1}^L\frac{A^{(L)}_{-k}}{\epsilon^k}\,,
\end{equation}
and $\Gamma^{\text{finite}}$ is finite at $\epsilon=0$. It is clear then that the
divergent part can be compensated for if one adds to the action the $L$-loop
counterterm
\[
\mathcal A^{(L)}=\tilde\lambda^L\int d^dx\,\left[\frac{z^{(L)}(\epsilon)}2\,(\partial\varphi)^2+\frac{b^{(L)}(\epsilon)}2\,\varphi^2+\frac{\tilde\lambda\,a^{(L)}(\epsilon)}{4!}\,\varphi^4\right]\,,
\]
with
\[
z^{(L)}(\epsilon)=Z^{(L)}(\epsilon)+Z_0^{(L)}\,;\qquad
b^{(L)}(\epsilon)=B^{(L)}(\epsilon)+B_0^{(L)}\,;\qquad
a^{(L)}(\epsilon)=A^{(L)}(\epsilon)+A_0^{(L)}\,,
\]
where $Z_0^{(L)},B_0^{(L)},A_0^{(L)}$ are finite numbers.

The freedom in the choice of these finite numbers represents the freedom in the
choice of the renormalization scheme. In given scheme they are determined by the
normalization conditions, like \eqref{eq:SIW}, \eqref{eq:SIb} or
\eqref{eq:SIIa0}, \eqref{eq:SIIb0}, which usually directly relate the parameters
$\lambda$ and $m^2$ to physical quantities. The form \eqref{eq:divpoles} of the
divergent part of the effective action suggests another class of renormalization
schemes. In so called \emph{Minimal Subtraction Scheme} one simply sets
$Z_0^{(L)}=0$, $B_0^{(L)}=0$, $A_0^{(L)}=0$. In other words, in this scheme one
chooses the $L$-loop counterterm to be exactly equal to the divergent part (i.e.\
the singular part of the Laurent expansion around $\epsilon=0$) of the effective
action renormalized to $L-1$ loops, as defined above, i.e.
\[
\mathcal A^{(L)}[\varphi]=\Gamma^{\text{div}}_L[\varphi]\,.
\]
In this scheme the parameters $m^2_{MS}$ and $\lambda_{MS}$ entering the action
\[
\mathcal A=\int d^dx\,\left[\frac12\,(\partial\varphi)^2+\frac{m^2_{MS}}2\,\varphi^2+\frac{\lambda_{MS}}{4!}\,\varphi^4+\text{Counterterms in MS Scheme}\right]\,,
\]
are not a priori related to physical quantities such as the particle mass and
scattering amplitudes. This relation has to be calculated order by order in
perturbation theory.

Full proof that such procedure works to all orders is rather involved (it can be
found e.g.\ in \emph{J.~Zinn-Justin, Quantum Field Theory and Critical
Phenomena}, or in original papers quoted on p.~338 of PS). According to these
general theorems similar renormalization procedure works to all orders in any
field theory whose perturbation theory has only finitely many primitive
divergences. We have seen that this is true if and only if the mass dimensions
of all couplings are non-negative.

\subsection*{Appendix A: Two loop analysis of \texorpdfstring{$\varphi^4$}{phi**4} (PS Sect.~10.5)}

Here we carry out two-loop renormalization analysis in the $\varphi^4$ theory.
Let us now start with the $\tilde\lambda^3$ contributions to $\Gamma^{(4)}$. I
will still assume the renormalization scheme SII. Also, I will disregard the
$b^{(1)}$ term in $\delta m^2$ counterterm and the ``trivial bubble diagram''; we
saw that these cancel each other exactly. Then we have 16 diagrams contributing
to $\Gamma^{(4)}$ in this order:
\begin{multline*}
-\Gamma^{(4)}(p_1,\dots,p_4)=\cdots+\Bigl(\underset{(a)}{\diac[0.4]{\coordinate (L) at (-1.6,0); \coordinate (M) at (0,0); \coordinate (R) at (1.6,0);
  \draw[prop] (-0.8,0) ellipse (0.8 and 0.5); \draw[prop] (0.8,0) ellipse (0.8 and 0.5);
  \legs{L/135/1,L/225/2,R/45/3,R/315/4} \foreach \p in {L,M,R} \node[vtx] at (\p) {};}}
+\underset{(b)}{\diac[0.4]{\coordinate (L) at (-0.8,0); \coordinate (T) at (0.55,0.6); \coordinate (B) at (0.55,-0.6);
  \draw[prop] (0,0) circle (0.8); \draw[prop] (T) .. controls (0.25,0) .. (B);
  \legs{L/135/1,L/225/2,T/45/3,B/315/4} \foreach \p in {L,T,B} \node[vtx] at (\p) {};}}
+\underset{(c)}{\diac[0.4]{\coordinate (R) at (0.8,0); \coordinate (T) at (-0.55,0.6); \coordinate (B) at (-0.55,-0.6);
  \draw[prop] (0,0) circle (0.8); \draw[prop] (T) .. controls (-0.25,0) .. (B);
  \legs{T/135/1,B/225/2,R/45/3,R/315/4} \foreach \p in {R,T,B} \node[vtx] at (\p) {};}}\\
+\underset{(d)}{\diac[0.4]{\coordinate (L) at (-0.9,0); \coordinate (R) at (0.9,0);
  \draw[prop] (0,0) ellipse (0.9 and 0.55); \legs{L/135/1,L/225/2,R/45/3,R/315/4}
  \node[vtx] at (L) {}; \node[otimesv,minimum size=5pt] at (R) {}; \node[font=\tiny] at (1.25,0) {(2)};}}
+\underset{(e)}{\diac[0.4]{\coordinate (L) at (-0.9,0); \coordinate (R) at (0.9,0);
  \draw[prop] (0,0) ellipse (0.9 and 0.55); \legs{L/135/1,L/225/2,R/45/3,R/315/4}
  \node[otimesv,minimum size=5pt] at (L) {}; \node[vtx] at (R) {}; \node[font=\tiny] at (-1.25,0) {(2)};}}\Bigr)\\
+\Bigl(\text{10 diagrams with permutations of the external legs}\Bigr)+\Bigl(\ctX{3}\Bigr)+\cdots
\end{multline*}
The diagrams (d) and (e) (and 4 similar diagrams with permutations of the
external legs) involve the counterterm vertex $-\tilde\lambda^2a^{(2)}$ with the
coefficient $a^{(2)}$ which was already determined in the one-loop analysis, see
Eq.~\eqref{eq:a2} above.

The last diagram shown above represents the contribution of the term
$\tilde\lambda^3$ in $\delta\lambda$ counterterm. We must see that it is indeed
possible to adjust this single coefficient $a^{(3)}$ in such a way that all
divergences (i.e.\ all singularities at $d=4$) cancel. This last contribution is
momentum independent constant. For our renormalization program to be successful
to this order we must prove that all momentum-dependent divergences cancel out
among the first 15 diagrams above.

Consider 5 diagrams (a-e) above. It is useful to divide them into three groups
as follows (the counterterm diagrams (d) and (e) are shared among the groups,
with the weights shown)
\begin{align*}
\text{I)}&\quad\diac[0.4]{\coordinate (L) at (-1.6,0); \coordinate (M) at (0,0); \coordinate (R) at (1.6,0);
  \draw[prop] (-0.8,0) ellipse (0.8 and 0.5); \draw[prop] (0.8,0) ellipse (0.8 and 0.5);
  \legs{L/135/1,L/225/2,R/45/3,R/315/4} \foreach \p in {L,M,R} \node[vtx] at (\p) {};}
+\frac13\,\diac[0.4]{\coordinate (L) at (-0.9,0); \coordinate (R) at (0.9,0);
  \draw[prop] (0,0) ellipse (0.9 and 0.55); \legs{L/135/1,L/225/2,R/45/3,R/315/4}
  \node[otimesv,minimum size=5pt] at (L) {}; \node[vtx] at (R) {}; \node[font=\tiny] at (-1.25,0) {(2)};}
+\frac13\,\diac[0.4]{\coordinate (L) at (-0.9,0); \coordinate (R) at (0.9,0);
  \draw[prop] (0,0) ellipse (0.9 and 0.55); \legs{L/135/1,L/225/2,R/45/3,R/315/4}
  \node[vtx] at (L) {}; \node[otimesv,minimum size=5pt] at (R) {}; \node[font=\tiny] at (1.25,0) {(2)};}\,,\\
\text{II)}&\quad\diac[0.4]{\coordinate (R) at (0.8,0); \coordinate (T) at (-0.55,0.6); \coordinate (B) at (-0.55,-0.6);
  \draw[prop] (0,0) circle (0.8); \draw[prop] (T) .. controls (-0.25,0) .. (B);
  \legs{T/135/1,B/225/2,R/45/3,R/315/4} \foreach \p in {R,T,B} \node[vtx] at (\p) {};}
+\frac23\,\diac[0.4]{\coordinate (L) at (-0.9,0); \coordinate (R) at (0.9,0);
  \draw[prop] (0,0) ellipse (0.9 and 0.55); \legs{L/135/1,L/225/2,R/45/3,R/315/4}
  \node[otimesv,minimum size=5pt] at (L) {}; \node[vtx] at (R) {}; \node[font=\tiny] at (-1.25,0) {(2)};}\,,\\
\text{III)}&\quad\diac[0.4]{\coordinate (L) at (-0.8,0); \coordinate (T) at (0.55,0.6); \coordinate (B) at (0.55,-0.6);
  \draw[prop] (0,0) circle (0.8); \draw[prop] (T) .. controls (0.25,0) .. (B);
  \legs{L/135/1,L/225/2,T/45/3,B/315/4} \foreach \p in {L,T,B} \node[vtx] at (\p) {};}
+\frac23\,\diac[0.4]{\coordinate (L) at (-0.9,0); \coordinate (R) at (0.9,0);
  \draw[prop] (0,0) ellipse (0.9 and 0.55); \legs{L/135/1,L/225/2,R/45/3,R/315/4}
  \node[vtx] at (L) {}; \node[otimesv,minimum size=5pt] at (R) {}; \node[font=\tiny] at (1.25,0) {(2)};}\,.
\end{align*}
Let us show that all momentum dependent divergences cancel separately within
each group.

Consider first the group I. The first diagram there is easily calculated to be
\[
-\frac{\tilde\lambda^3}4\,\bigl[I(p_{12}^2)\bigr]^2\,,
\]
the 4 in the denominator being the symmetry factor, $p_{12}=p_1+p_2$, and again
\[
I(p^2)=\int\frac{d^dk}{(2\pi)^d}\,\frac1{(k^2+\tilde m^2)\bigl((p+k)^2+\tilde m^2\bigr)}\,.
\]
Now, each of the other two diagrams in this group contribute
\[
\frac13\,\frac{\tilde\lambda^3}2\,a^{(2)}\,I(p_{12}^2)\,,
\]
so that the sum of all three diagrams in this group is
\[
\text{Group I}=-\frac{\tilde\lambda^3}4\,\bigl[I(p_{12}^2)\bigr]^2+\frac23\,\frac{\tilde\lambda^3}2\,\frac32\,I(0)\,I(p_{12}^2)
=-\frac{\tilde\lambda^3}4\,\bigl[I(p_{12}^2)-I(0)\bigr]^2+\frac{\tilde\lambda^3}4\,\bigl[I(0)\bigr]^2
\]
where the expression $a^{(2)}=\frac32I(0)$, Eq.~\eqref{eq:a2}, is used. We know
from the one-loop analysis that the difference $I_r(p^2)=I(p^2)-I(0)$ is finite
at $d=4$ and therefore the divergent parts of the diagrams in the group I indeed
sum up to a momentum independent constant (the last term in the expression
above). Note that this divergent constant has second order pole at $d=4$, i.e.\
it behaves as $1/\epsilon^2$ as $\epsilon\to0$.

Consider now the group II. It contains the diagram
\[
\diac[0.45]{\coordinate (R) at (0.8,0); \coordinate (T) at (-0.55,0.6); \coordinate (B) at (-0.55,-0.6);
  \draw[prop] (0,0) circle (0.8); \draw[prop] (T) .. controls (-0.25,0) .. (B);
  \foreach \p/\a/\n in {T/135/p_2,B/225/p_1,R/45/p_3,R/315/p_4} {\draw[prop,densely dashed] (\p) -- ++(\a:1.1) node[pos=1.25,font=\tiny] {$\n$};}
  \draw[arr] (60:1.15) arc[start angle=60,end angle=100,radius=1.15]; \node[font=\tiny] at (80:1.45) {$k$};
  \draw[arr] (260:1.15) arc[start angle=260,end angle=300,radius=1.15]; \node[font=\tiny] at (280:1.5) {$p_{12}+k$};
  \foreach \p in {R,T,B} \node[vtx] at (\p) {};}
=-\frac{\tilde\lambda^3}2\int\frac{d^dk}{(2\pi)^d}\,\frac1{k^2+\tilde m^2}\,\frac1{(p_{12}+k)^2+\tilde m^2}\,I\bigl((k+p_1)^2\bigr)\,,
\]
where $I(p^2)$ is the integral \eqref{eq:Ip2}, which we already found equal to
\[
I(p^2)=\frac{\Gamma(2-d/2)}{(4\pi)^{\frac d2}}\int_0^1du\,\Bigl[\tilde m^2+u(1-u)\,p^2\Bigr]^{\frac d2-2}\,.
\]
Using this formula and combining the propagators with the help of Feynman
parametrization we obtain for the above diagram
\begin{multline*}
-\frac{\tilde\lambda^3}2\,\frac{\Gamma(2-d/2)}{(4\pi)^{\frac d2}}
\int_0^1du\int_0^1dv\int\frac{d^dk}{(2\pi)^d}\,\frac1{\bigl(k^2+2v\,(kp_{12})+v\,p_{12}^2+\tilde m^2\bigr)^2}\\
\times\frac1{\bigl(\tilde m^2+u(1-u)\,(k+p_1)^2\bigr)^{2-\frac d2}}\,.
\end{multline*}
The $k$ integral here can be evaluated by using the transformation generalizing
the Feynman parametrization. Namely, consider the identity
\eqref{eq:FeynParGen},
\begin{equation}\label{eq:FeynParGen2}
\frac1{A^\alpha\,B^\beta}=\frac{\Gamma(\alpha+\beta)}{\Gamma(\alpha)\Gamma(\beta)}\int_0^1\frac{dw\,w^{\alpha-1}(1-w)^{\beta-1}}{\bigl[wA+(1-w)B\bigr]^{\alpha+\beta}}\,.
\end{equation}
It can be derived by explicit evaluation of the integral in the right-hand side:
projective transformation of the integration variable $w=Bz/(Bz+A(1-z))$ brings
it to the form of the Beta-integral
\[
B(\alpha,\beta)=\int_0^1dz\,z^{\alpha-1}(1-z)^{\beta-1}=\frac{\Gamma(\alpha)\Gamma(\beta)}{\Gamma(\alpha+\beta)}
\]
(details can be found in \S\,10.5 of PS). Using this identity (with
$\alpha=2-\frac d2$, $\beta=2$) one can write the above integral as
\begin{multline*}
-\frac{\tilde\lambda^3}2\,\frac{\Gamma(4-d/2)}{(4\pi)^{\frac d2}}\int_0^1du\,dv\,dw\\
\times\int\frac{d^dk}{(2\pi)^d}\,\frac{w^{1-\frac d2}\,(1-w)}{\Bigl[w\bigl(\tilde m^2+u(1-u)\,(k+p_1)^2\bigr)+(1-w)\bigl(\tilde m^2+k^2+2v\,(kp_{12})+v\,p_{12}^2\bigr)\Bigr]^{4-\frac d2}}\,.
\end{multline*}
By appropriate shift of the $d$-momentum integration variable the denominator
can be brought to the form
\[
A\,q^2+P^2+\tilde m^2
\]
where
\[
A=(1-w)+w\,u(1-u)\,,\qquad A\big|_{w=0}=1\,,
\]
and $P$ is complicated function of $p_{12}$, $p_1$ and the parameters $u,v,w$.
We only need to know that
\[
P^2\Big|_{w=0}=v(1-v)\,p_{12}^2\,.
\]
In this representation the integral over $q$ is readily evaluated,
\[
-\frac{\tilde\lambda^3}2\,\frac{\Gamma(4-d)}{(4\pi)^d}\int_0^1du\,dv\,dw\,\frac{w^{1-\frac d2}\,(1-w)}{A^{\frac d2}}\,\bigl(\tilde m^2+P^2\bigr)^{d-4}\,.
\]
Evidently this expression is singular at $d=4$. First, there is a pole at $d=4$
in the factor
\[
\Gamma(4-d)=\Gamma(\epsilon)=\frac1\epsilon+O(1)
\]
in front. In addition, at $d=4$ the integral over $w$ diverges at $w\to0$
because of the factor $w^{1-\frac d2}$ in the integrand; as the result the
integral over $w$ has additional pole at $d=4$, i.e.\ the whole diagram again
has a second order pole at this point.

To isolate the singular contributions let us write the integral over $w$ as
\[
\int_0^1dw\,w^{1-\frac d2}\,f(w)=f(0)\int_0^1dw\,w^{1-\frac d2}+\int_0^1dw\,w^{1-\frac d2}\,\bigl[f(w)-f(0)\bigr]\,.
\]
The second piece is
\[
-\frac{\tilde\lambda^3}2\,\frac{\Gamma(4-d)}{(4\pi)^d}\int_0^1du\,dv\,dw\,w^{1-\frac d2}\,
\left[\frac{1-w}{A^{\frac d2}}\,\bigl(\tilde m^2+P^2(w)\bigr)^{d-4}-\bigl(\tilde m^2+P^2(w=0)\bigr)^{d-4}\right]\,.
\]
Here the integral over $u,v,w$ is convergent and brings in no additional
singularity at $d=4$. Therefore the above expression has only simple pole at
$\epsilon=0$. To evaluate the residue at this simple pole one can set $d=4$
everywhere except in $\Gamma(4-d)$. It is clear then that this residue does not
depend on external momenta, i.e.\ this part of the above diagram has the form
\[
\frac{C_0}\epsilon+C_1(p_1,\dots,p_4)\,,
\]
where $C_0$ does not depend on $p_i$, and the $C_1$ is finite at $\epsilon=0$.

Let us take a look at the other piece, proportional to $f(0)$. It is
\[
f(0)\int_0^1dw\,w^{1-\frac d2}=\frac2\epsilon\,f(0)=-\frac2\epsilon\,\frac{\tilde\lambda^3}2\,\frac{\Gamma(\epsilon)}{(4\pi)^d}\int_0^1du\,dv\,\bigl(\tilde m^2+v(1-v)\,p_{12}^2\bigr)^{-\epsilon}
\]
The $u$ integral here is trivial. Now, for $\epsilon\to0$ the integral can be
written as
\[
\int_0^1dv\,\Bigl(1-\epsilon\,\log\bigl(\tilde m^2+v(1-v)\,p_{12}^2\bigr)+O(\epsilon^2)\Bigr)\,,
\]
and therefore this piece of the diagram is
\[
-\frac{\tilde\lambda^3}{(4\pi)^4}\,\frac1\epsilon\left(\frac1\epsilon-\gamma+\log4\pi-\int_0^1dv\,\log\bigl[\tilde m^2+v(1-v)\,p_{12}^2\bigr]\right)+\text{finite terms.}
\]
We have to add here the second diagram in the group II, which is
\[
\frac23\,\frac{\tilde\lambda^3}2\,a^{(2)}\,I(p_{12}^2)=\frac{\tilde\lambda^3}2\,I(0)\,I(p_{12}^2)\,,
\]
where the expression \eqref{eq:a2} for the coefficient $a^{(2)}$ is used. Recall
that
\begin{align*}
I(p^2)&=\frac{\Gamma(2-\frac d2)}{(4\pi)^{\frac d2}}\int_0^1du\,\bigl[\tilde m^2+u(1-u)\,p^2\bigr]^{\frac d2-2}\\
&=\frac1{(4\pi)^{\frac d2}}\left(\frac2\epsilon-\gamma-\int_0^1du\,\log\bigl(\tilde m^2+u(1-u)\,p^2\bigr)+O(\epsilon)\right)\,.
\end{align*}
Therefore the contribution of this diagram is
\[
\frac{\tilde\lambda^3}{(4\pi)^4}\,\frac1\epsilon\left(\frac2\epsilon-2\gamma+2\log4\pi-\log\tilde m^2-\int_0^1du\,\log\bigl[\tilde m^2+u(1-u)\,p_{12}^2\bigr]\right)+\text{finite terms.}
\]
Adding the two diagrams of the group II we see that the $p_{12}$-dependent
singular terms cancel out,
\[
\text{Group II}=\frac{\tilde\lambda^3}{(4\pi)^4}\,\frac1\epsilon\left(\frac1\epsilon-\gamma+\log4\pi-\log\tilde m^2\right)+\text{finite terms.}
\]
The group III can be analyzed similarly. We find that all singular parts of the
first 15 diagrams are pole terms with the coefficients independent of external
momenta,
\begin{equation}\label{eq:15diagrams}
\text{15 diagrams}=\tilde\lambda^3\left(\frac{c_{-2}}{\epsilon^2}+\frac{c_{-1}}\epsilon+F(p_1,\dots,p_4)\right)\,,
\end{equation}
where $c_{-2},c_{-1}$ do not depend on $p_i$, and $F$ is finite at $d=4$.
Therefore the counterterm coefficient $a^{(3)}$ appearing in the 16th diagram
always can be appropriately adjusted to compensate the singular terms. In fact,
according to our normalization condition \eqref{eq:SIIb}
($\Gamma^{(4)}(0,0,0,0)=\tilde\lambda$), this coefficient is determined as
\[
a^{(3)}=\frac{c_{-2}}{\epsilon^2}+\frac{c_{-1}}\epsilon+F(0,\dots,0)
\]
and hence
\[
\text{All 16 diagrams}=\tilde\lambda^3\Bigl(F(p_1,\dots,p_4)-F(0,\dots,0)\Bigr)\,.
\]

\begin{exercise}
Verify these two-loop calculations. Evaluate $c_{-2}$ and $c_{-1}$.
\end{exercise}

Let us briefly consider the $\tilde\lambda^2$ contributions to $\Gamma^{(2)}$.
These are given by three diagrams
\[
\Gamma^{(2)}(p^2)=p^2+\tilde m^2-\Bigl(\underset{(a)}{\diac[0.45]{\draw[prop,densely dashed] (-2.3,0) -- (-1.1,0) (1.1,0) -- (2.3,0);
  \draw[prop] (-1.1,0) -- (1.1,0); \draw[prop] (0,0) ellipse (1.1 and 0.7);
  \node[above,font=\tiny] at (0,0.75) {$p+k_1+k_2$}; \node[font=\tiny] at (0,0.25) {$k_1$}; \node[below,font=\tiny] at (0,-0.75) {$k_2$};
  \node[vtx] at (-1.1,0) {}; \node[vtx] at (1.1,0) {};}}
+\underset{(b)}{\diac[0.45]{\draw[prop,densely dashed] (-1.8,0) -- (1.8,0); \draw[prop] (0,0.6) circle (0.6);
  \node[otimesv,minimum size=5pt] at (0,0) {}; \node[font=\tiny] at (0,0.6) {(2)};}}
+\underset{(c)}{\ctline{2}}\Bigr)
\]
The diagram (b) involves the counterterm vertex $-\tilde\lambda^2a^{(2)}$, and
the diagram (c) contains the $\tilde\lambda^2$ part of the counterterms $\delta Z$
and $\delta m^2$,
\[
\ctline{2}=-\tilde\lambda^2\,\bigl(z^{(2)}\,p^2+b^{(2)}\bigr)\,.
\]
Again, the coefficients here must be determined from the normalization condition
\eqref{eq:SIIa}.

Let us denote $\Sigma_a(p^2)$ the contribution of the diagram (a) to the mass
operator. One can write it as
\[
\Sigma_a(p^2)=\Sigma_a(0)+p^2\,\Sigma_a'(0)+\Sigma_r(p^2)\,,
\]
where by the definition
\[
\Sigma_r(p^2)=O(p^4)\qquad\text{as}\qquad p^2\to0\,.
\]
Let us analyse the diagram (a) and show that $\Sigma_r(p^2)$ has finite value at
$d=4$. The contribution of the diagram (a) is given by the integral
\[
\Sigma_a(p^2)=-\frac{\tilde\lambda^2}{3!}\int\frac{d^dk_1\,d^dk_2}{(2\pi)^{2d}}\,\frac1{(k_1^2+\tilde m^2)(k_2^2+\tilde m^2)\bigl((p+k_1+k_2)^2+\tilde m^2\bigr)}\,.
\]
One of the $k$-integrations, say over $k_2$, can be done exactly the way we did
one-loop integrals for $\Gamma^{(4)}$:
\begin{align*}
&\int\frac{d^dk_2}{(2\pi)^d}\,\frac1{(k_2^2+\tilde m^2)\bigl((p+k_1+k_2)^2+\tilde m^2\bigr)}=I\bigl((p+k_1)^2\bigr)\\
&\qquad=\frac{\Gamma(2-d/2)}{(4\pi)^{\frac d2}}\int_0^1du\,\Bigl[\tilde m^2+u(1-u)\,(p+k_1)^2\Bigr]^{\frac d2-2}\,.
\end{align*}
We find
\[
\Sigma_a(p^2)=-\frac{\tilde\lambda^2}{3!}\,\frac{\Gamma(2-d/2)}{(4\pi)^{\frac d2}}\int_0^1du\int\frac{d^dk_1}{(2\pi)^d}\,
\frac1{\bigl(u(1-u)\,(p+k_1)^2+\tilde m^2\bigr)^{2-d/2}\,(k_1^2+\tilde m^2)}\,.
\]
The remaining integral over $k_1$ can be done using \eqref{eq:FeynParGen2} which
allows one to transform the integral as
\[
\Sigma_a(p^2)=-\frac{\tilde\lambda^2}{3!}\,\frac{\Gamma(2-d/2)}{(4\pi)^{\frac d2}}\,\frac{\Gamma(3-d/2)}{\Gamma(2-d/2)\,\Gamma(1)}\int_0^1du\int_0^1dw\,J(u,w)\,,
\]
\[
J(u,w)=\int\frac{d^dk_1}{(2\pi)^d}\,\frac{w^{1-\frac d2}}{\bigl[\tilde m^2+(1-w)\,k_1^2+w\,u(1-u)\,(k_1+p)^2\bigr]^{3-\frac d2}}\,.
\]
The integral over $k_1$ can be evaluated after the shift of this momentum
variable
\[
q=k_1+\frac{w\,u(1-u)}A\,p\,,
\]
where again
\[
A=(1-w)+w\,u(1-u)\,.
\]
One obtains
\[
J(u,w)=\int\frac{d^dq}{(2\pi)^d}\,\frac{w^{1-d/2}}{\Bigl(A\,q^2+w(1-w)\,u(1-u)\,p^2/A+\tilde m^2\Bigr)^{3-\frac d2}}\,.
\]
With appropriate rescaling of $q$ one can bring this integral to the form
\eqref{eq:dimregalpha}, and obtain
\[
J(u,w)=\frac{A^{3-\frac{3d}2}}{(4\pi)^{\frac d2}}\,\frac{\Gamma(3-d)}{\Gamma(3-d/2)}\,w^{1-d/2}\,\Bigl[w(1-w)\,u(1-u)\,p^2+A\,\tilde m^2\Bigr]^{d-3}\,.
\]
Combining these results we get
\[
\Sigma_a(p^2)=-\frac{\tilde\lambda^2}{3!}\,\frac{\Gamma(3-d)}{(4\pi)^d}\int_0^1dw\,du\,w^{1-\frac d2}\,A^{3-\frac{3d}2}\,\Bigl[u(1-u)\,w(1-w)\,p^2+A\,\tilde m^2\Bigr]^{d-3}\,.
\]
This quantity has second order pole in $\epsilon$, one singular factor coming
from the Gamma function
\[
\Gamma(3-d)=-\frac1\epsilon+O(1)\,,
\]
and another appearing due to the divergence of the integral over $w$ at
$w\to0$. However, subtracting $\Sigma_a(0)$ and $p^2\Sigma_a'(0)$ one obtains
\[
\Sigma_r(p^2)=-\frac{\tilde\lambda^2}{3!}\,\frac{\Gamma(3-d)}{(4\pi)^d}\,(\tilde m^2)^{d-3}\int_0^1du\,dw\,w^{1-\frac d2}\,A^{-\frac d2}\,H(u,w)\,,
\]
where
\[
H(u,w)=(1+p^2B)^{d-3}-1-(d-3)\,p^2B\,,\qquad\text{and}\qquad B=\frac{u(1-u)\,w(1-w)}{\tilde m^2\,A}\,.
\]
Note that $H(u,w=0)=0$ and therefore the $w$ integration does not produce
singularity in this expression. Moreover, one can see that $H=O(\epsilon)$,
\[
H(u,w)=\epsilon\,\Bigl[p^2B-(1+p^2B)\,\log(1+p^2B)\Bigr]+O(\epsilon^2)\,,
\]
and therefore $\Sigma_r$ has finite value at $d=4$,
\[
\Sigma_r(p^2)\Big|_{d=4}=\frac{\tilde\lambda^2}{3!}\,\frac{\tilde m^2}{(4\pi)^4}\int_0^1\frac{du\,dw}{w\,A^2}\,\Bigl[p^2B-(1+p^2B)\,\log(1+p^2B)\Bigr]
\]
At the same time the subtracted terms $\Sigma_a(0)$ and $\Sigma_a'(0)$ are
singular at $d=4$. The contributions of two diagrams (a) and (b) to $\Sigma$ then
have the form
\[
\Sigma_a(p^2)+\Sigma_b(p^2)=\Sigma_r(p^2)+\tilde\lambda^2\,\bigl(M(\epsilon)+p^2N(\epsilon)\bigr)\,,
\]
where the constants $M(\epsilon)$ and $N(\epsilon)$ have poles in $\epsilon$,
\begin{equation}\label{eq:MN}
M(\epsilon)=\frac{M_{-2}}{\epsilon^2}+\frac{M_{-1}}\epsilon+O(1)\,;\qquad
N(\epsilon)=\frac{N_{-1}}\epsilon+O(1)\,.
\end{equation}
The diagram (c) brings the contribution
\[
\Sigma_c(p^2)=\tilde\lambda^2\,\bigl(p^2z^{(2)}+b^{(2)}\bigr)\,,
\]
and our normalization condition
\[
\Sigma(0)=\Sigma'(0)=0
\]
requires the choice
\[
b^{(2)}=-M(\epsilon)\,;\qquad z^{(2)}=-N(\epsilon)
\]
of the counterterm coefficients at this order. For the proper vertex we then
have simply
\[
\Gamma^{(2)}(p^2)=p^2+\tilde m^2+\Sigma_r(p^2)+O(\tilde\lambda^3)\,.
\]

\begin{hwproblem}{4}\label{hw:4}
Within the original perturbation theory (i.e.\ with the bare mass in the
propagator $1/(p^2+m_0^2)$, and no counterterm vertices) consider the sum of all
``cactus diagrams''
\[
-\Sigma_{\text{cactus}}=
\dia[0.35]{\draw[prop,densely dashed] (-2,0) -- (2,0); \draw[prop] (0,0.6) circle (0.6); \node[vtx] at (0,0) {};}
+\dia[0.35]{\draw[prop,densely dashed] (-2,0) -- (2,0); \draw[prop] (0,0.6) circle (0.6); \draw[prop] (0,1.8) circle (0.6); \node[vtx] at (0,0) {}; \node[vtx] at (0,1.2) {};}
+\dia[0.35]{\draw[prop,densely dashed] (-2,0) -- (2,0); \draw[prop] (0,0.6) circle (0.6); \draw[prop] (-0.75,1.6) circle (0.5); \draw[prop] (0.75,1.6) circle (0.5);
  \node[vtx] at (0,0) {}; \node[vtx] at (-0.42,1.02) {}; \node[vtx] at (0.42,1.02) {};}
+\dia[0.35]{\draw[prop,densely dashed] (-2,0) -- (2.6,0); \draw[prop] (0,0.6) circle (0.6); \draw[prop] (-0.75,1.6) circle (0.5); \draw[prop] (0.75,1.6) circle (0.5);
  \draw[prop] (1.75,1.6) circle (0.5); \node[vtx] at (0,0) {}; \node[vtx] at (-0.42,1.02) {}; \node[vtx] at (0.42,1.02) {}; \node[vtx] at (1.25,1.6) {};}
+\dots
\]
and ignore all other diagrams for the mass operator. In this ``cactus
approximation'', find the relation between $m_0^2$ and $m^2$ to all orders in
$\lambda$. Use Schwinger's regularization $\Phi(p^2/\Lambda^2)=\exp(-p^2/\Lambda^2)$.
\end{hwproblem}

% Source pp. 135-143 (Lectures13-18.pdf).
\section[Renormalization of Composite Fields (\texorpdfstring{$\varphi^4$}{phi**4} theory)]{Renormalization of Composite Fields (\texorpdfstring{$\varphi^4$}{phi**4} theory)}

% insertion vertices
\newcommand{\insv}[2][]{\dia[0.4]{\draw[prop,wavy] (-2.2,0) -- (0,0);
  \draw[prop] (0,0) -- (1.4,1.4); \draw[prop] (0,0) -- (1.4,-1.4); #2 \node[font=\tiny] at (-2.6,0.45) {#1};}}

In many cases one may be interested in correlation functions involving composite
fields, such as powers of $\varphi^n(x)$, or energy-momentum tensor. Such
correlation functions can be computed perturbatively, as the power series in the
coupling constant $\lambda$. However new renormalizations, in addition to the
basic renormalizations of the parameters $m^2$, $\lambda$ and the field
$\varphi$, have to be made to make such correlation functions finite. Let us
demonstrate this for the simplest case of the composite field
\[
\varphi^2(x)
\]
in the $\varphi^4$ theory.

Consider the 3-point connected correlation function
\[
\Bigl\langle\frac12\,\varphi^2(x)\,\varphi(x_1)\,\varphi(x_2)\Bigr\rangle_c
\]
where $\varphi(x)$ stands for the renormalized field, $\varphi=Z^{-1/2}\varphi_0$,
$Z=1+\delta Z$, with $Z$ chosen according to some scheme, say SII. Diagrams can
be used to describe perturbative contributions to this correlation function if
one introduces additional ``external'' vertex
\[
\insv[$\varphi^2$]{\node[vtx,minimum size=4.5pt] at (0,0) {};}=1
\]
associated with the insertion $\frac12\varphi^2$. As before, it is useful to
switch to the momentum representation, and introduce the associated amputated
vertex
\begin{equation}\label{eq:Gphi2}
-\Gamma^{(\varphi^2,2)}(p_1,p_2)=\frac{W^{(\varphi^2,2)}(p_1,p_2)}{W(p_1)\,W(p_2)}\,,
\end{equation}
where
\[
(2\pi)^4\delta^{(4)}(q-p_1-p_2)\,W^{(\varphi^2,2)}(p_1,p_2)=\text{f.t.}\,\Bigl\langle\frac12\,\varphi^2(x)\,\varphi(x_1)\varphi(x_2)\Bigr\rangle_c\,,
\]
with $q,p_1,p_2$ being the momenta corresponding to $x,x_1,x_2$, respectively.
It is given by the sum of all one-particle irreducible diagrams
\begin{multline*}
-\Gamma^{(\varphi^2,2)}(p_1,p_2)=
\underset{(a)}{\diac[0.4]{\draw[prop,wavy] (-2,0) -- (0,0); \draw[prop,densely dashed] (0,0) -- (1.2,1.2); \draw[prop,densely dashed] (0,0) -- (1.2,-1.2); \node[vtx] at (0,0) {};}}
+\underset{(b)}{\diac[0.4]{\draw[prop,wavy] (-2.8,0) -- (-1.6,0); \draw[prop] (-0.8,0) ellipse (0.8 and 0.5);
  \draw[prop,densely dashed] (0,0) -- (1.2,1.2); \draw[prop,densely dashed] (0,0) -- (1.2,-1.2); \node[vtx] at (-1.6,0) {}; \node[vtx] at (0,0) {};}}
+\underset{(c)}{\diac[0.4]{\draw[prop,wavy] (-4.4,0) -- (-3.2,0); \draw[prop] (-2.4,0) ellipse (0.8 and 0.5); \draw[prop] (-0.8,0) ellipse (0.8 and 0.5);
  \draw[prop,densely dashed] (0,0) -- (1.2,1.2); \draw[prop,densely dashed] (0,0) -- (1.2,-1.2); \foreach \x in {-3.2,-1.6,0} \node[vtx] at (\x,0) {};}}\\
+\underset{(d)}{\diac[0.4]{\draw[prop,wavy] (-2.8,0) -- (-1.6,0); \coordinate (T) at (-0.25,0.6); \coordinate (B) at (-0.25,-0.6);
  \draw[prop] (-0.8,0) circle (0.8); \draw[prop] (T) .. controls (-0.55,0) .. (B);
  \draw[prop,densely dashed] (T) -- ++(45:1.2); \draw[prop,densely dashed] (B) -- ++(-45:1.2); \foreach \p in {(-1.6,0),(T),(B)} \node[vtx] at \p {};}}
+\underset{(e)}{\diac[0.4]{\draw[prop,wavy] (-2.8,0) -- (-1.6,0); \draw[prop] (-0.8,0) ellipse (0.8 and 0.5);
  \draw[prop,densely dashed] (0,0) -- (1.2,1.2); \draw[prop,densely dashed] (0,0) -- (1.2,-1.2); \node[vtx] at (-1.6,0) {}; \node[otimesv,minimum size=5pt] at (0,0) {};}}
+\cdots
\end{multline*}
Most of these diagrams are singular at $d=4$, and a quick look reveals that these
singularities can not be eliminated by renormalizations of the parameters $m^2$,
$\lambda$, and the field $\varphi$. This is already seen in the order
$\sim\lambda$. All the counterterm vertices that we have introduced already show
up at the order $\lambda^2$ (diagram (e)) and higher. But the diagram (b), which
is $\sim\lambda$, is singular. Indeed its contribution is
\[
-\Gamma^{(\varphi^2,2)}(p_1,p_2)=1-\frac\lambda2\,I(p_{12}^2)+O(\lambda^2)
\]
where $I(p^2)$ is the same integral \eqref{eq:Ip2} which we encountered
previously. We know that
\[
I(p_{12}^2)=I(0)+I_r(p_{12}^2)\,,
\]
where $I_r(p^2)$ has finite limit at $d=4$, while $I(0)$ diverges as $1/\epsilon$.
In fact this shows what we need to do. We write
\[
-\Gamma^{(\varphi^2,2)}(p_1,p_2)=\Bigl(1-\frac\lambda2\,I(0)+O(\lambda^2)\Bigr)\Bigl(1-\frac\lambda2\,I_r(p_{12}^2)+O(\lambda^2)\Bigr)+O(\lambda^2)\,.
\]
Therefore, at least to the leading order, we can handle the problem by defining
the renormalized field $[\varphi^2]_R$,
\begin{equation}\label{eq:phi2R}
\varphi^2=Z_{\varphi^2}^{-1}\,[\varphi^2]_R\,,
\end{equation}
where
\[
Z_{\varphi^2}^{-1}=1-\frac\lambda2\,I(0)+O(\lambda^2)\,,\qquad
I(0)=\frac{\Gamma(2-d/2)}{(4\pi)^{d/2}}\,(m^2)^{d/2-2}\,.
\]
Then
\[
\Gamma^{(\varphi^2,2)}=Z_{\varphi^2}^{-1}\,\Gamma^{([\varphi^2]_R,2)}\,,
\]
where the vertex function of the renormalized field $[\varphi^2]_R$ is finite,
to the order $\lambda$.

It is convenient to invert the above relation,
\begin{equation}\label{eq:phi2Rinv}
[\varphi^2]_R(x)=Z_{\varphi^2}\,\varphi^2=Z_{\varphi^2}\,Z^{-1}\,\varphi_0^2\,.
\end{equation}
(Note that in the ratio \eqref{eq:Gphi2} defining $\Gamma^{(\varphi^2,2)}(p_1,p_2)$
the field renormalization factor $Z$ cancels out.) As with other renormalization
parameters, we assume that $Z_{\varphi^2}$ expands in power series in $\lambda$,
\[
Z_{\varphi^2}=z^{(0)}_{\varphi^2}+\lambda\,z^{(1)}_{\varphi^2}+\lambda^2\,z^{(2)}_{\varphi^2}+\dots
\]
In the above calculation we have set $z^{(0)}_{\varphi^2}=1$; with this we have
\[
z^{(1)}_{\varphi^2}=\frac12\,I(0)+\text{finite}\,.
\]
By inserting the renormalized field $[\varphi^2]_R(x)$ instead of $\varphi^2(x)$
we replace the $\varphi^2$ vertex with
\[
\insv[$\varphi^2$]{\node[vtx,minimum size=4.5pt] at (0,0) {};}=1\ \longrightarrow\
\insv{\node[circle,draw,fill=white,inner sep=0pt,minimum size=6pt] at (0,0) {}; \fill (0,0) circle (0.7pt);}=Z_{\varphi^2}\,,
\]
In the perturbation theory (i.e.\ expansion in $\lambda$), this new vertex is
represented as the series in $\lambda$,
\[
\insv{\node[circle,draw,fill=white,inner sep=0pt,minimum size=6pt] at (0,0) {}; \fill (0,0) circle (0.7pt);}
=\insv{\node[circle,draw,fill=white,inner sep=0pt,minimum size=6pt] at (0,0) {}; \fill (0,0) circle (0.7pt); \node[font=\tiny] at (0,0.6) {(0)};}
+\insv{\node[circle,draw,fill=white,inner sep=0pt,minimum size=6pt] at (0,0) {}; \fill (0,0) circle (0.7pt); \node[font=\tiny] at (0,0.6) {(1)};}
+\insv{\node[circle,draw,fill=white,inner sep=0pt,minimum size=6pt] at (0,0) {}; \fill (0,0) circle (0.7pt); \node[font=\tiny] at (0,0.6) {(2)};}
+\cdots
\]
where
\[
\insv{\node[circle,draw,fill=white,inner sep=0pt,minimum size=6pt] at (0,0) {}; \fill (0,0) circle (0.7pt); \node[font=\tiny] at (0,0.6) {$(n)$};}
=\lambda^n\,z^{(n)}_{\varphi^2}
\]
In computing the correlation functions with the insertion of $[\varphi^2]_R$ one
has to take into account these new field renormalization vertices along with
the counterterm vertices associated with the counterterm part of the action. The
field renormalization constants $z^{(n)}_{\varphi^2}$ are to be adjusted, order
by order in the perturbation theory, to cancel the divergent parts of the
diagrams.

As in the case of the renormalizations of the parameters of the action, there
is freedom of finite renormalizations
\[
[\varphi^2]_R\to Z_{\text{finite}}\,[\varphi^2]_R\,.
\]
This ambiguity can be fixed by imposing some normalization condition. For
instance, one can fix the normalization of the field $[\varphi^2]_R$ by
\[
-\Gamma^{([\varphi^2]_R,2)}(0,0)=1\,,
\]
where $-\Gamma^{([\varphi^2]_R,2)}(p_1,p_2)$ is the momentum-space amputated
correlation function
\[
-\Gamma^{([\varphi^2]_R,2)}(p_1,p_2)\,W(p_1)\,W(p_2)=\text{f.t.}\,\Bigl\langle\frac12\,[\varphi^2]_R(x)\,\varphi(x_1)\varphi(x_2)\Bigr\rangle_c
\]
Under this condition we have
\[
z^{(0)}_{\varphi^2}=1\,,\qquad z^{(1)}_{\varphi^2}=\frac12\,I(0)\,.
\]
Renormalizations of higher composite fields follow the same idea but are more
complicated. In particular, it generally involves the phenomenon of ``operator
mixing'': renormalizations of higher composite fields may require forming
linear combinations of them, with renormalization factors becoming matrices.
Brief discussion of the simplest case of this phenomenon, mixing of the fields
$\varphi^3$ and $\Delta\varphi$, is given below. More systematically this
phenomenon is discussed in the context of \emph{Renormalization Group}.

\subsection*{Higher composite fields. Operator mixing}

Simple multiplicative renormalization of the field $\varphi^3$, i.e.\ the form
$[\varphi^3]_R=Z_3(\epsilon)\,\varphi^3$, is not sufficient for removing
divergences in all correlation functions. Instead, one has to look for the
renormalized field in the form of the linear combination
\begin{equation}\label{eq:phi3R}
\frac1{3!}\,[\varphi^3]_R=\frac{Z_3}{3!}\,\varphi^3+Z_{3,1}\,\varphi-Z_{3,1''}\,\partial_\mu^2\varphi\,.
\end{equation}
The reason for adding the term $\sim\partial^2\varphi$ will become clear in a
moment. As usual, the coefficients $Z$ are assumed to be certain power series in
the renormalized coupling constant,
\begin{align*}
Z_3&=1+\lambda\,z^{(1)}_3+\cdots\,,\\
Z_{3,1}&=z^{(0)}_{3,1}+\lambda\,z^{(1)}_{3,1}+\cdots\,,\\
Z_{3,1''}&=\lambda\,z^{(1)}_{3,1''}+\cdots\,.
\end{align*}
Therefore, insertion of the renormalized field $\frac1{3!}[\varphi^3]_R$ gives
rise to series of vertices,
\[
\dia[0.4]{\draw[prop,wavy] (-2.2,0) -- (0,0); \draw[prop] (0,0) -- (1.4,1.4); \draw[prop] (0,0) -- (1.8,0); \draw[prop] (0,0) -- (1.4,-1.4);
  \node[vtx,minimum size=4.5pt] at (0,0) {}; \node[font=\tiny] at (-2.6,0.45) {$\varphi^3$};}=1\,,\qquad
\dia[0.4]{\draw[prop,wavy] (-2.2,0) -- (0,0); \draw[prop] (0,0) -- (1.4,1.4); \draw[prop] (0,0) -- (1.8,0); \draw[prop] (0,0) -- (1.4,-1.4);
  \node[otimesv,minimum size=5pt] at (0,0) {}; \node[font=\tiny] at (0,0.6) {(1)}; \node[font=\tiny] at (-2.6,0.45) {$\varphi^3$};}=\lambda\,z^{(1)}_3\,,\qquad\text{etc}
\]
and
\[
\dia[0.4]{\draw[prop,wavy] (-2.2,0) -- (0,0); \draw[prop] (0,0) -- (2,0); \node[otimesv,minimum size=5pt] at (0,0) {}; \node[font=\tiny] at (0,0.6) {(0)}; \node[font=\tiny] at (-2.6,0.45) {$\varphi^3$};}=z^{(0)}_{3,1}\,,\qquad
\dia[0.4]{\draw[prop,wavy] (-2.2,0) -- (0,0); \draw[prop] (0,0) -- (2,0); \node[otimesv,minimum size=5pt] at (0,0) {}; \node[font=\tiny] at (0,0.6) {(1)}; \node[font=\tiny] at (-2.6,0.45) {$\varphi^3$};}
=\lambda\,\bigl(z^{(1)}_{3,1}+p^2\,z^{(1)}_{3,1''}\bigr)\,,\qquad\text{etc.}
\]
Then, for the vertex $\Gamma^{(\varphi^3_R,1)}$ we have the following series of
diagrams
\begin{multline*}
\underbrace{\Bigl(\dia[0.4]{\draw[prop,wavy] (-2.2,0) -- (0,0); \draw[prop,densely dashed] (0,0) -- (1.8,0); \draw[prop] (0,0.6) circle (0.6);
  \node[vtx,minimum size=4.5pt] at (0,0) {}; \node[font=\tiny] at (-2.6,0.45) {$\varphi^3$};}
+\dia[0.4]{\draw[prop,wavy] (-2.2,0) -- (0,0); \draw[prop,densely dashed] (0,0) -- (1.8,0); \node[otimesv,minimum size=5pt] at (0,0) {}; \node[font=\tiny] at (0,0.6) {(0)}; \node[font=\tiny] at (-2.6,0.45) {$\varphi^3$};}\Bigr)}_{\sim\lambda^0}\\
+\underbrace{\Bigl(\dia[0.4]{\draw[prop,wavy] (-2.2,0) -- (0,0); \draw[prop] (0,0) -- (2,0); \draw[prop] (1,0) ellipse (1 and 0.55); \draw[prop,densely dashed] (2,0) -- (3.2,0);
  \node[vtx,minimum size=4.5pt] at (0,0) {}; \node[vtx] at (2,0) {}; \node[font=\tiny] at (-2.6,0.45) {$\varphi^3$};}
+\dia[0.4]{\draw[prop,wavy] (-2.2,0) -- (0,0); \draw[prop,densely dashed] (0,0) -- (1.8,0); \draw[prop] (0,0.6) circle (0.6);
  \node[otimesv,minimum size=5pt] at (0,0) {}; \node[font=\tiny] at (0,0.6) {(1)}; \node[font=\tiny] at (-2.6,0.45) {$\varphi^3$};}
+\dia[0.4]{\draw[prop,wavy] (-2.2,0) -- (0,0); \draw[prop,densely dashed] (0,0) -- (1.8,0); \node[otimesv,minimum size=5pt] at (0,0) {}; \node[font=\tiny] at (0,0.6) {(1)}; \node[font=\tiny] at (-2.6,0.45) {$\varphi^3$};}\Bigr)}_{\sim\lambda^1}+\cdots
\end{multline*}
We already know from the free theory that
\begin{equation}\label{eq:z031}
z^{(0)}_{3,1}=-\frac12\,D(0)+\text{finite}=-\frac12\int\frac{d^dk/(2\pi)^d}{k^2+m^2}+\text{finite}\,.
\end{equation}
The finite parts here, and in higher order counterterms are arbitrary in that
they are not constrained by general requirement of locality. After all, finite
renormalization of $[\varphi^3]_R$, or adding finite portions of (already
renormalized) fields $\varphi$ and $\partial^2\varphi$ produces equally good
local field that one can take for renormalized $\varphi^3$. One can fix this
ambiguity by adopting minimal subtraction scheme, or by imposing some
normalization conditions, say
\begin{equation}\label{eq:phi3norm}
-\Gamma^{\varphi^3_R,1}(p^2)=O(p^4)\,,
\end{equation}
in addition to the one we have taken already,
\[
-\Gamma^{\varphi^3_R,3}(0,0,0)=1\,.
\]
With this normalization condition, the finite term in \eqref{eq:z031} must be
zero, so that the two diagrams $\sim\lambda^0$ cancel exactly. At the order
$\lambda^1$, the first of the above diagrams is
\[
\dia[0.4]{\draw[prop,wavy] (-2.2,0) -- (0,0); \draw[prop] (0,0) -- (2,0); \draw[prop] (1,0) ellipse (1 and 0.55); \draw[prop,densely dashed] (2,0) -- (3.2,0);
  \node[vtx,minimum size=4.5pt] at (0,0) {}; \node[vtx] at (2,0) {}; \node[font=\tiny] at (-2.6,0.45) {$\varphi^3$};}
=-\frac\lambda{3!}\,\Pi(p^2)\,,
\]
where $\Pi(p^2)$ is already familiar integral
\[
\Pi(p^2)=\int\frac{d^dk_1\,d^dk_2}{(2\pi)^{2d}}\,\frac1{(k_1^2+m^2)(k_2^2+m^2)\bigl((p+k_1+k_2)^2+m^2\bigr)}\,.
\]
We have found that
\[
\Pi_r(p^2)=\Pi(p^2)-\Pi(0)-p^2\,\Pi'(0)
\]
is finite, and the above normalization condition is satisfied by the choice
\[
z^{(1)}_{3,1}=\frac1{3!}\,\Pi(0)-\frac{z^{(1)}_3}2\,D(0)\,,\qquad
z^{(1)}_{3,1''}=\frac1{3!}\,\Pi'(0)\,,
\]
where $z^{(1)}_3=\frac32I(0)$ is the first-order counterterm associated with
$Z_3$. Under this choice all singular parts cancel, and the above three diagrams
$\sim\lambda^1$ sum to finite quantity $-\frac\lambda{3!}\Pi_r(p^2)$.

We have seen few lectures ago that the fields in the correlation functions
satisfy the equation of motion in the ``weak sense'',
\begin{equation}\label{eq:weakEOM}
\bigl(m_0^2-\partial_\mu^2\bigr)\varphi_0(x)+\frac{\lambda_0}{3!}\,\varphi_0^3(x)\simeq0\,,
\end{equation}
the ``weak'' equation $F(x)\simeq0$ means that
\[
\bigl\langle F(x)\,\mathcal O_1(x_1)\cdots\mathcal O_n(x_n)\bigr\rangle=0\qquad\text{for}\quad x\neq x_1,\dots,x_n\,,
\]
with any local fields $\mathcal O_i$. The above ``weak'' equation was obtained by
formal manipulations with the functional integral, namely by performing the
change of variables
\[
\varphi_0(x)\to\varphi_0(x)+\epsilon(x)
\]
in the functional integral written in its ``bare'' form. It involves therefore
the bare parameters and the bare field $\varphi_0$. As it stands it makes sense
only if some regularization is present. One can repeat the same manipulation with
the functional integral expressed through the renormalized field and
renormalized parameters. This way one obtains
\[
\bigl(m^2+\delta m^2\bigr)\varphi(x)-(1+\delta Z)\,\partial_\mu^2\varphi(x)+\frac{\lambda+\delta\lambda}{3!}\,\varphi^3(x)\simeq0\,,
\]
which of course is just the same ``bare'' equation \eqref{eq:weakEOM} expressed
through renormalized parameters and renormalized $\varphi$. It contains the
singular counterterm coefficients as well as the field $\varphi^3$ which, before
renormalization, is also singular. Again, it makes sense only in regularized
theory. However now we can use \eqref{eq:phi3R} to exclude the field $\varphi^3$
in favor of the renormalized field $[\varphi^3]_R$. One then obtains the
equation
\begin{equation}\label{eq:renEOM}
m^2B_{m^2}\,\varphi(x)-\partial_\mu^2\varphi(x)+\frac\lambda{3!}\,B_\lambda\,[\varphi^3]_R(x)\simeq0\,,
\end{equation}
where
\[
B_{m^2}=\frac1{m^2}\,\frac{(m^2+\delta m^2)\,Z_3-(\lambda+\delta\lambda)\,Z_{3,1}}{(1+\delta Z)\,Z_3-(\lambda+\delta\lambda)\,Z_{3,1''}}\,,\qquad
B_\lambda=\frac{1+\delta\lambda/\lambda}{(1+\delta Z)\,Z_3-(\lambda+\delta\lambda)\,Z_{3,1''}}\,.
\]
The equation \eqref{eq:renEOM} states certain relation between the correlation
functions of fully renormalized fields which remain finite when the
regularization is removed (i.e.\ when $\epsilon$ is set equal 0). This suggests
that the coefficients $B_{m^2}$ and $B_\lambda$, which appear as complicated
combination of singular counterterms, are in fact finite. This finiteness is
direct consequence of the existence of renormalized perturbation theory. Let us
check it in the leading order in $\lambda$. As $\delta\lambda/\lambda=a^{(2)}\lambda+O(\lambda^2)$,
$Z_3=1+z_3^{(1)}\lambda+O(\lambda^2)$ and $\delta Z=O(\lambda^2)$ we find for
example
\[
B_\lambda=1+\bigl(a^{(2)}-z^{(1)}_3\bigr)\,\lambda+O(\lambda^2)\,.
\]
It is not difficult to check that the difference $a^{(2)}-z^{(1)}_3$ is finite
in any renormalization scheme.

Note that canonical mass dimension of $\varphi^3$ at $d=4$ is
\[
[\varphi_0^3]=3\,,\qquad\text{and}\qquad[\varphi]=1\,,\qquad[\partial^2\varphi]=3\,.
\]
The renormalized field $[\varphi^3]_R$ is a linear combination of the bare fields
of canonical dimension 3 or lower, which have the same symmetries as
$\varphi^3$. It is not difficult to check that there are no other composite
fields, besides $\varphi_0$, $\partial^2\varphi_0$, and $\varphi_0^3$ itself,
which i) have canonical dimension $\le3$, ii) transform as scalars under
coordinate rotations, and iii) change sign under obvious symmetry of the
$\varphi^4$ theory: $\varphi\to-\varphi$. This is a general feature of
renormalization of composite fields. In order to eliminate divergences, it is
sufficient to look for renormalized fields in the forms of linear combinations of
the bare fields of the same or lower canonical dimensions, and of the same
symmetries. In doing so one usually assumes that the fields of lower dimensions
are already renormalized, and so renormalizations of composite fields goes step
by step, from lower dimensions up.

To illustrate this further, let us briefly discuss renormalization of the fields
of canonical dimension 4. There are two new fields of this dimension,
\begin{equation}\label{eq:dim4ops}
\mathcal O_I=\frac12\,(\partial\varphi)^2\,,\qquad\mathcal O_{II}=\frac1{4!}\,\varphi^4\,.
\end{equation}
Along with these we must consider two other fields
\[
\mathcal O_{III}=\frac12\,\partial^2\varphi^2\,,\qquad\mathcal O_{IV}=\frac12\,\varphi^2\,,
\]
because the first of them also has canonical dimension 4, and canonical
dimension of the second is 2 which is lower than 4, and both have the same
$\varphi\to-\varphi$ symmetry. Of course renormalization of these last two fields
is completely determined by the renormalization of the field $\varphi^2$ which we
already considered. At the same time renormalized fields corresponding to
\eqref{eq:dim4ops} must be obtained as linear combinations of the bare fields
$\mathcal O_{A,0}$, $A=I,II,III,IV$, i.e.
\[
\mathcal O_{A,R}=\sum_B\bigl[Z^{-1}\bigr]_{AB}\,\mathcal O_{B,0}\,,
\]
where $Z$ is $4\times4$ matrix
\[
Z=\begin{pmatrix}\mathbf A&\mathbf B\\0&\mathbf C\end{pmatrix}\,.
\]
Here $\mathbf C=\operatorname{diag}(Z_{\varphi^2},Z_{\varphi^2})$, and $2\times2$
matrices $\mathbf A$ and $\mathbf B$ must be found from suitable normalization
conditions order by order in $\lambda$. Note that this matrix has triangular (in
terms of dimensions) form: fields of higher dimensions do not mix into the fields
of lower dimensions. Renormalizations of composite fields of yet higher canonical
dimensions, and renormalizations of tensor fields, follow the same pattern.

% Source pp. 144-160 (Lectures13-18.pdf).
\section[Perturbative Renormalization Group]{Perturbative Renormalization Group\protect\footnote{PS Sect's 12.2, 12.3.}}

\subsection*{Finite renormalizations}

Renormalized perturbation theory allows one to obtain the correlation functions
of $\varphi^4$ theory as power series in the renormalized coupling parameter
$\lambda$ with well defined finite coefficients. This concerns the correlation
functions of the field $\varphi$ as well as the ones involving appropriately
renormalized composite fields. This is achieved by adjusting the counterterm
coefficients $\delta m^2$, $\delta\lambda$, $\delta Z$ in the action (and,
additionally, the renormalization factors in the composite fields) as the
functions of regularization parameter ($\epsilon=4-d$ in dimensional
regularization, or $\Lambda$ if some cutoff regularization is assumed).
Renormalized correlation functions appearing in this way depend on two
parameters, $\lambda$ and $m^2$, named the ``renormalized coupling'' and the
``renormalized mass parameter''. Precise relation of these parameters to
physical quantities, as well as the field normalizations, are specified
(explicitly or implicitly) by the \emph{renormalization scheme}.

By doing calculations in two different renormalization schemes we are not
solving two different field theories; rather, we are describing the same theory
in terms of different parameters. One can think of two-dimensional (because two
parameters, $m^2$ and $\lambda$, are involved) ``manifold of $\varphi^4$
theories'', with the parameters $m^2$ and $\lambda$ playing the role of (local)
coordinates on this manifold (the latter is often called the ``theory space'').

Each point of this manifold represents a field theory, i.e.\ a collection of all
correlation functions, defined up to finite field renormalizations
\[
\bigl\langle\varphi(x_1)\cdots\varphi(x_n)\bigr\rangle\ \sim\ Z_{\text{finite}}^{n/2}\,\bigl\langle\varphi(x_1)\cdots\varphi(x_n)\bigr\rangle\,.
\]
To simplify things one can consider the ratios like
\[
R(x_1,x_2;x_3,\dots,x_n)=\bigl\langle\varphi(x_1)\varphi(x_2)\cdots\varphi(x_{2n})\bigr\rangle\big/\bigl\langle\varphi(x_1)\varphi(x_2)\bigr\rangle^n
\]
which do not depend on the normalizations of the fields. Different
renormalization schemes correspond to different choices of the coordinate system
on the above manifold. Once the coordinate system (i.e.\ renormalization scheme)
is fixed, that is the points are somehow ``labeled'' by two parameters $m^2$ and
$\lambda$, the above ratios become functions of these parameters,
\[
R(x_i|m^2,\lambda)\,.
\]
Relation between the parameters $(m^2,\lambda)$ and $(\tilde m^2,\tilde\lambda)$,
associated with two different renormalization schemes, is just a coordinate
transformation,
\[
m^2\to\tilde m^2(m^2,\lambda)\,,\qquad\lambda\to\tilde\lambda(m^2,\lambda)\,.
\]
At $d=4$ this relation has the form (by dimensional analysis)
\begin{equation}\label{eq:schemechange}
\tilde m^2=m^2\,f(\lambda)\,,\qquad\tilde\lambda=g(\lambda)\,,
\end{equation}
where the functions $f,g$ can in principle be computed order by order in
$\lambda$,
\begin{align*}
f(\lambda)&=f_0+\lambda\,f_1+\lambda^2\,f_2+\cdots\,,\\
g(\lambda)&=\lambda+\lambda^2\,g_2+\lambda^3\,g_3+\cdots\,,
\end{align*}
the coefficients $f_k$, $g_k$ being finite numbers. The correlation functions
(or rather the above ratios $R$) computed in the two schemes are related as
\[
\tilde R\bigl(x_i|m^2f(\lambda),g(\lambda)\bigr)=R(x_i|m^2,\lambda)\,,
\]
so that these are the same correlation functions expressed in terms of different
parameters, that is, coordinates in the space of $\varphi^4$ theories.

\subsection*{Short distance problem}

Renormalized perturbation theory, no matter what scheme is assumed, yields the
correlation functions as formal power series in the coupling parameter
$\lambda$. Although in each given order these series give finite and well
behaved correlation functions, the problem of overall \emph{non-perturbative}
consistency of the theory is a separate question. There is an issue of
convergence of the perturbative expansion. It is possible to show that the
series in $\varphi^4$ have in fact zero radius of convergence, and at the best
they can be understood as asymptotic series. If so, there is an important
question if these series can be understood as asymptotic expansions of some
non-perturbative correlation functions which are compatible with requirements of
\emph{locality} and with \emph{quantum mechanical interpretation} (the property
usually called \emph{positivity}). Although one can show that these properties
(locality and positivity) are preserved at the perturbative level, i.e.\ in each
given order in $\lambda$, it is conceivable that they can break down after
summing the series up.

We have seen that in Euclidean version of quantum field theory locality implies
in particular that the correlation function
\[
\bigl\langle\varphi(x_1)\varphi(x_2)\cdots\varphi(x_{2n})\bigr\rangle\,,
\]
viewed as the function of, say, $x_1$, has only point-like singularities at
$x_1=x_2,x_3,\dots,x_{2n}$. (We have observed this in the free Klein-Gordon
theory (see Section~4) for the two-point function, but the arguments we have
applied in that analysis depend only on general properties of the
energy-momentum spectrum, and are generalized straightforwardly to general
interacting theories; I hope to come back to this point later.) This is because
upon continuation to complex coordinates and ultimately to the Minkowski
space-time these singularities translate into the singularities at the light
cones associated with the points $x_2,\dots,x_{2n}$; we have seen that these
singularities are directly related to commutators of the field operators (see
again the analysis in Section~4). Locality states that the commutators of
Heisenberg field operators must vanish outside the light cone,
\[
[\hat\varphi(\xx,t),\hat\varphi(\mathbf0,0)]=0\qquad\text{for}\qquad t^2<\xx^2\,.
\]
This property of locality is preserved in the renormalized perturbation theory at
each order in $\lambda$ but there is no guarantee it will not break down in full
non-perturbative theory. It is very well conceivable that upon summing the
perturbative series up singularities at finite space-like separations may be
generated, leading to violations of the local commutativity.

The following example demonstrates this possibility. Consider the power series
\[
\sum_{n=0}^\infty(-\lambda)^n\,\Bigl(\log m|x_1-x_2|\Bigr)^n\,,
\]
which models a perturbative series for some correlation function. At each given
order the above expression has point-like singularity at $x_1=x_2$. However its
sum is
\[
\frac1{1+\lambda\log m|x_1-x_2|}\,.
\]
This function is singular at
\[
|x_1-x_2|=x_0\equiv m^{-1}\,e^{-1/\lambda}\,.
\]
As the result, if one continues to the Minkowski space-time the light cone
``fattens'' and causality is broken at short scales $\simeq x_0$.\footnote{One
can take different attitudes toward such possibility. One pragmatic attitude is
that we do not know what describes short-distance physics anyway. After all, at
short distances quantum gravity must play important role, perhaps along with
extra dimensions, superstrings, D-branes, and whatnot. And therefore we do not
care about possible inconsistencies at these short distances. With this attitude
quantum field theory is viewed as a way of producing (renormalized) perturbative
expansions in the coupling parameter. Such expansions of course are useful only
if the coupling parameter is small. In quantum electrodynamics (QED) where the
analog of $\lambda$ is the fine structure constant which is small indeed, this is
valid point of view. And if the interaction is not small in any sense, according
to this philosophy one would conclude that quantum field theory is not useful at
all.\par
In the past, at certain stage of development, such was prevailing attitude.
Discovery of \emph{Asymptotic Freedom} in nonabelian gauge theories, along with
better understanding of \emph{Renormalization Group} has changed that. On the
face of it, the discovery of the asymptotic freedom was pure perturbative result,
obtained by computing few diagrams, but through understanding the Renormalization
Group one can see that it actually guaranteed full nonperturbative consistency of
quantum nonabelian gauge theories. More generally, renormalization allows one to
understand quantum field theory in terms of criticality, or vice versa, to the
extent that from theoretical point of view quantum field theory and theory of
critical phenomena are much one and the same thing. Moreover, main idea of the
Wilson's renormalization group has proved very useful in theoretical physics well
beyond quantum field theory. This is why I hope to devote significant part of this
course to it.}

The problem of the short-distance consistency of a quantum field theory requires
analysis of its behavior under changes of the length scale. The latter is
generally understood in terms of \emph{Renormalization Group}. Here I will
briefly discuss the basics of this approach, in its perturbative version, and
applied to $\varphi^4$ theory. The analysis is substantially simpler in the case
of \emph{massless theory}.

\subsection*{Massless \texorpdfstring{$\varphi^4$}{phi**4}}

As a preparation, let us discuss what happens with our renormalized perturbation
if one sets the renormalized mass parameter $m^2$ to zero. Although precise
meaning of $m^2$ depends on the renormalization scheme, the condition
\[
m^2=0
\]
is \emph{scheme independent}, i.e.\ once it is satisfied in one scheme, it is
satisfied in any other scheme, as the Eq.~\eqref{eq:schemechange} suggests.
There are at least two reasons why this case is interesting.
\begin{itemize}
\item Potential problem with locality can occur at short scales, much shorter
than $m^{-1}$. The mass term is not important at such scales, and by analyzing
this case one can understand the general situation.
\item We will see that in applications to critical phenomena the case $m^2=0$
describes physics at exactly critical point, and therefore it is of central
interest.
\end{itemize}

Taking a look at our previous perturbative results we discover that it is not
possible to just set $m^2=0$ there, since nearly all quantities we have computed
are singular at this point. Consider for instance the 4-point vertex,
\[
\Gamma^{(4)}(p_1,\dots,p_4)=\lambda-\frac{\lambda^2}2\Bigl(I_r(p_{12}^2)+I_r(p_{13}^2)+I_r(p_{14}^2)\Bigr)+O(\lambda^3)\,,
\]
where again
\[
I_r(p^2)=-\frac1{(4\pi)^2}\int_0^1\dd u\,\log\left(1+u(1-u)\,\frac{p^2}{m^2}\right)\,.
\]
Obviously, there is a singularity at $m^2=0$. This can signal one of two things.
Either the theory itself becomes sick at this point, or our coordinate system
does not extend smoothly to this case. I want to argue that $m^2=0$
singularities of the perturbative correlation functions originate from a
coordinate singularity.\footnote{Think of some a function $f(x,y)$ written in new
coordinates $(x',y')$, $x=x'$, $y=y'+y'^2\log x'^2$.}

Indeed, a particular coordinate system is determined by the normalization
conditions. For example, in our scheme SII we have assumed
\begin{equation}\label{eq:SIIrecall}
\Gamma^{(2)}(p^2)\Big|_{p^2=0}=m^2\,,\qquad\text{and}\qquad\frac{\dd\Gamma^{(2)}(p^2)}{\dd p^2}\Big|_{p^2=0}=1\,.
\end{equation}
At $m^2=0$ the first of these conditions remains useful, just stating that
$\Gamma^{(2)}(p^2)$ turns to zero at $p^2=0$, which can be considered as a
scheme-independent definition of the ``massless theory'', but at $m^2=0$ this is
no longer a simple zero, being replaced by more complicated singular point. We
will see in a moment that at $m^2=0$ $\Gamma^{(2)}$ contains terms like
$p^2\log p^2$ and therefore the second of the conditions \eqref{eq:SIIrecall}
cannot be imposed. The coordinate system associated with our scheme is singular
at $m^2=0$.

It is useful to have a scheme which admits smooth limit $m^2\to0$. Such scheme
must involve dimensional parameter different from $m^2$. Let us denote it
$\mu^2$, $[\mu]=1$. Since $\mu$ can be chosen arbitrarily, we will have families
of schemes, parameterized by auxiliary scale $\mu$. Let us take one such
$\mu$-dependent scheme, defined by the normalization conditions
\begin{align}
\frac{\dd\Gamma^{(2)}(p^2)}{\dd p^2}\Big|_{p^2=\mu^2}&=1\,,\label{eq:mu1}\\
\Gamma^{(4)}(p_1,\dots,p_4)\Big|_{p_i=\mu e_i}&=\lambda_{(\mu)}\,,\label{eq:mu2}
\end{align}
where $e_i$ is some set of four fixed 4-vectors, which satisfy
$e_1+e_2+e_3+e_4=0$, but have no other linear dependencies. To be specific, we
can choose symmetric arrangement,
\[
e_i^2=3\qquad\text{and}\qquad e_i\cdot e_j=-1\quad\text{for}\quad i\neq j\,.
\]
I write $\lambda_{(\mu)}$ for the renormalized coupling constant as the reminder
about the scale $\mu$ associated with this scheme. It is conventional to refer to
this auxiliary scale as the ``normalization scale''. It is perfectly legitimate
to use such scheme in general case, but we will apply it to massless theory,
which is singled out by the condition
\begin{equation}\label{eq:mu3}
\Gamma^{(2)}(0)=0\,.
\end{equation}

We can now develop renormalized perturbation theory for massless $\varphi^4$
theory taking exact same steps we did before. We start with the action in its
bare form,
\[
\mathcal A=\int \dd^dx\,\left[\frac12\,(\partial\varphi_0)^2+\frac{m_0^2}2\,\varphi_0^2+\frac{\lambda_0}{4!}\,\varphi_0^4\right]\,,
\]
with some bare mass parameter (yet to be adjusted), and write it as
\[
\mathcal A=\int \dd^dx\,\left[\frac12\,(\partial\varphi)^2+\frac{\lambda_{(\mu)}}{4!}\,\varphi^4+\text{counterterms}\right]\,,
\]
where
\[
\varphi=Z_{(\mu)}^{-1/2}\,\varphi_0\,,
\]
and the counterterms are
\[
\text{counterterms}=\frac{\delta Z_{(\mu)}}2\,(\partial\varphi)^2+\frac{\delta m^2_{(\mu)}}2\,\varphi^2+\frac{\delta\lambda_{(\mu)}}{4!}\,\varphi^4\,,
\]
with
\[
\delta Z_{(\mu)}=Z_{(\mu)}-1\,,\qquad
\delta m^2_{(\mu)}=m_0^2\,Z_{(\mu)}\,,\qquad
\delta\lambda_{(\mu)}=\lambda_0\,Z_{(\mu)}^2-\lambda_{(\mu)}\,.
\]
Again, we assume that all the counterterm coefficients are power series in the
renormalized coupling $\lambda_{(\mu)}$,
\begin{align*}
\delta Z_{(\mu)}&=z^{(2)}_{(\mu)}\,\lambda_{(\mu)}^2+\cdots\,,\\
\delta\lambda_{(\mu)}&=a^{(2)}_{(\mu)}\,\lambda_{(\mu)}^2+\cdots\,,\\
\delta m^2_{(\mu)}&=b^{(1)}_{(\mu)}\,\lambda_{(\mu)}+\cdots\,,
\end{align*}
with the coefficients determined from the normalization conditions
\eqref{eq:mu1}, \eqref{eq:mu2}, \eqref{eq:mu3}.

A remark is in order. We have put two parameters, $m_0^2$ and $\lambda_0$, in the
bare action. Renormalized massless theory depends on a single parameter, the
renormalized coupling $\lambda_{(\mu)}$, only. The normalization scale $\mu$ is
not a parameter of the theory, in that it is an attribute of the renormalization
scheme, rather than the theory itself. We have imposed three normalization
conditions Eqs.~\eqref{eq:mu1}--\eqref{eq:mu3}. While Eq.~\eqref{eq:mu2} defines
the renormalized coupling $\lambda_{(\mu)}$ and Eq.~\eqref{eq:mu1} fixes
normalization of $\varphi$, the Eq.~\eqref{eq:mu3} states that the theory is
massless, and thus it imposes restriction on the theory itself. In principle,
this last condition can be translated into certain relation between the bare
parameters $m_0^2$ and $\lambda_0$. The bare theory is implicitly equipped with
some regularization, and precise form of such relation depends on the
regularization method. For instance, if some cutoff regularization is implied,
with the cutoff momentum $\Lambda$, at $d=4$ the relation has general form
\begin{equation}\label{eq:critbare}
m_0^2=\Lambda^2\,C(\lambda_0)\,.
\end{equation}
Precise shape of $C(\lambda_0)$ heavily depends on specific implementation of the
cutoff. Subset of $\varphi^4$ theories selected by this condition corresponds to
critical domain in the theory of second order phase transitions. In dimensional
regularization there is no cutoff parameter $\Lambda$, and simple dimensional
analysis shows that this relation must take simple form
\[
m_0^2=0\qquad\text{(dimensional regularization)}.
\]
This means in particular that $\delta m^2=0$, and we do not have to bother with
the mass renormalization at all. This is one of the technical advantages of the
dimensional regularization technique.

Again, thinking in terms of the ``two-dimensional manifold of $\varphi^4$
theories'', the massless $\varphi^4$ theories correspond to certain
one-dimensional ``critical'' submanifold (critical curve, in our case),
figure~\ref{fig:critcurve}.
\begin{figure}[ht]
\centering\figfile{critcurve}
\caption{The two-dimensional space of $\varphi^4$ theories, with coordinates
$(m^2,\lambda)$. The massless theories form the critical curve $m^2=0$.}
\label{fig:critcurve}
\end{figure}
In statistical mechanics, the ``bare action'' models microscopic interaction, and
the parameters $m_0^2$ and $\lambda_0$, as well as the cutoff energy $\Lambda$,
are related to the temperature and the microscopic physical characteristics of
the matter under consideration (e.g.\ molecular interactions in a gas/liquid).
When these parameters cross the ``critical submanifold'' the macroscopic system
undergoes a second order phase transition.

\subsection*{Renormalization as the scale transformation. Callan-Symanzik equation}

Although in general $\varphi^4$ theory the coordinates $(m^2,\lambda)$ depend on
the renormalization scheme, the ``critical submanifold'' $m^2=0$ is scheme
independent. It is because the condition \eqref{eq:mu3}, being a physical
statement, is scheme independent. But even within the critical curve there still
is a freedom in the choice of the coordinate $\lambda$ along the curve. Within
the above class of schemes \eqref{eq:mu1}--\eqref{eq:mu2} this choice is
controlled by the choice of the normalization scale $\mu$. If $\mu$ and
$\tilde\mu$ are two different values of the normalization scale, the associated
renormalized couplings constants
\[
\lambda=\lambda_{(\mu)}\qquad\text{and}\qquad\tilde\lambda=\lambda_{(\tilde\mu)}
\]
are related as
\[
\tilde\lambda=g(\lambda)=g(\lambda|\tilde\mu/\mu)\,,
\]
where at $d=4$, by dimensional counting, the function $g$ can depend on
$\lambda$ and the dimensionless ratio $\tilde\mu/\mu$. And if $\tilde\mu$ is
infinitesimally close to $\mu$, i.e.\ $\tilde\mu=\mu+\delta\mu$, then
\[
\tilde\lambda=\lambda+\frac{\delta\mu}\mu\,\beta(\lambda)\,,\qquad\text{where}\qquad
\beta(\lambda)=\frac{\dd g(\lambda|\xi)}{\dd\xi}\Big|_{\xi=1}
\]
is a function of $\lambda$ alone.

Of course the correlation functions are essentially independent of the
renormalization scheme, in the present context of the choice of the scale $\mu$.
More precisely, the ratios of the correlation functions, e.g.
\begin{equation}\label{eq:Rratio}
R(x_i|\lambda_{(\mu)},\mu)=\frac{\bigl\langle\varphi(x_1)\cdots\varphi(x_{2n})\bigr\rangle_{(\mu)}}{\langle\varphi\varphi\rangle^n_{(\mu)}}
\end{equation}
which are insensitive to the normalization of $\varphi$, do not change under
substitutions
\[
\mu\to\tilde\mu\,,\qquad\lambda_{(\mu)}\to\lambda_{(\tilde\mu)}\,.
\]
Here the index $\mu$ in $\langle\cdots\rangle_{(\mu)}$ signifies that the
correlation function is calculated in the scheme with the normalization scale
$\mu$. This independence can be expressed in the form of partial differential
equation
\[
\left(\mu\frac\partial{\partial\mu}+\beta(\lambda)\,\frac\partial{\partial\lambda}\right)R(x_i|\lambda,\mu)=0\,,
\]
where now $\lambda$ and $\mu$ are viewed as two independent variables.

These equations can be extended to the correlation functions themselves if one
recalls that a change in renormalization scheme in general induces a finite
change in normalization of $\varphi$,
\[
\varphi_{(\tilde\mu)}=Z^{-1/2}_{\text{finite}}(\lambda|\tilde\mu/\mu)\,\varphi_{(\mu)}\,,
\]
where again at $d=4$ the constant $Z_{\text{finite}}$ can depend on $\lambda$ and
dimensionless ratio $\tilde\mu/\mu$. For an infinitesimal change
$\tilde\mu=\mu+\delta\mu$
\[
\varphi_{(\mu+\delta\mu)}=\left(1-\frac{\delta\mu}\mu\,\gamma(\lambda)\right)\,\varphi_{(\mu)}\,,
\qquad\text{with}\qquad
\gamma(\lambda)=\frac12\,\frac{\dd Z_{\text{finite}}(\lambda,\xi)}{\dd\xi}\Big|_{\xi=1}\,.
\]
Taking into account this transformation of $\varphi$ we obtain the so-called
\emph{Callan-Symanzik equation}:
\begin{equation}\label{eq:CS}
\left(\mu\frac\partial{\partial\mu}+\beta(\lambda)\,\frac\partial{\partial\lambda}+n\,\gamma(\lambda)\right)\bigl\langle\varphi(x_1)\cdots\varphi(x_n)\bigr\rangle_{(\mu)}=0
\end{equation}
where the correlation function is viewed as depending on two parameters,
$\lambda$ and $\mu$.

At the first glance the Callan-Symanzik equation contains strictly zero physical
information. Indeed, in order to make a meaningful renormalized perturbation
theory we have introduced an auxiliary scale $\mu$ having in mind that nothing
physical can depend on this arbitrary scale, and then the equation \eqref{eq:CS}
just states that yes indeed nothing physical depends on $\mu$. This is
absolutely true if no information about the functions $\beta(\lambda)$ and
$\gamma(\lambda)$ is added. However, if something is known about those functions
the Callan-Symanzik equation becomes useful.

To see this let us first notice that by dimensional counting any correlation
function of the $d=4$ massless theory can be written as
\[
\bigl\langle\varphi(x_1)\cdots\varphi(x_n)\bigr\rangle_{(\mu)}=\mu^n\,C^{(n)}(\mu x_1,\dots,\mu x_n|\lambda)\,,
\]
where $C^{(n)}$ is dimensionless function of dimensionless variables $\mu x_i$
and $\lambda$. Therefore the Callan-Symanzik equation \eqref{eq:CS} can be written
as
\begin{equation}\label{eq:CSdil}
\left(\sum_{i=1}^nx_i\frac\partial{\partial x_i}+n\bigl(1+\gamma(\lambda)\bigr)+\beta(\lambda)\,\frac\partial{\partial\lambda}\right)\bigl\langle\varphi(x_1)\cdots\varphi(x_n)\bigr\rangle=0\,,
\end{equation}
where the coordinate derivative operators have the following exact meaning
\[
x_i\frac\partial{\partial x_i}=\sum_{\mu=1}^4x_i^\mu\,\frac\partial{\partial x_i^\mu}\,.
\]
Note that the differential operator $\sum_ix_i\partial/\partial x_i$ describes
infinitesimal dilation of the coordinates $x_i$ and therefore the equation
\eqref{eq:CSdil} is the statement that such dilation can be essentially
compensated (up to the field renormalization) by appropriate change of the
coupling parameter $\lambda$. More precisely, an infinitesimal dilation of the
coordinates is equivalent to certain infinitesimal shift of $\lambda$ plus an
infinitesimal renormalization of $\varphi$,
\[
x^\mu\to(1+\delta L)\,x^\mu\quad\Longleftrightarrow\quad
\lambda\to\lambda-\beta(\lambda)\,\delta L\,,\quad\varphi\to\varphi-\bigl(1+\gamma(\lambda)\bigr)\,\varphi\,\delta L\,.
\]
Let $\lambda(L)$ be a solution of the following differential equation and initial
condition
\begin{equation}\label{eq:flow1}
L\frac{\dd}{\dd L}\,\lambda(L)=-\beta\bigl(\lambda(L)\bigr)\,,\qquad\lambda(1)=\lambda\,.
\end{equation}
Then the Callan-Symanzik equation states that
\begin{equation}\label{eq:CSsol}
Z^{n/2}(L)\,\bigl\langle\varphi(x_1/L)\cdots\varphi(x_n/L)\bigr\rangle_{\lambda(L)}=\bigl\langle\varphi(x_1)\cdots\varphi(x_n)\bigr\rangle_\lambda\,,
\end{equation}
where the index $\lambda$ now shows the value of renormalized coupling constant,
and $Z(L)$ solves the linear differential equation
\begin{equation}\label{eq:ZL}
L\frac{\dd}{\dd L}\,Z^{1/2}(L)=-\Bigl(1+\gamma\bigl(\lambda(L)\bigr)\Bigr)\,Z^{1/2}(L)\,,\qquad Z^{1/2}(1)=1\,.
\end{equation}
It is assumed that the correlation functions in \eqref{eq:CSsol} are computed
within the same scheme, with the same normalization scale $\mu$, but with
different values, $\lambda(L)$ and $\lambda$, of the renormalized coupling
parameter. Note that $\mu$ does not appear in \eqref{eq:CSsol}, i.e.\ in this form
the Callan-Symanzik equation is not a statement about scheme dependence but
rather determines how the theory behaves when one goes from one scale to
another.

\subsection*{Renormalization group flow}

We can give this equation the following geometric interpretation. The coupling
constants $\lambda$ and $\lambda(L)$ are coordinates representing two points, $P$
and $P(L)$, of the critical submanifold (figure~\ref{fig:flowpoints}).
\begin{figure}[ht]
\centering\figfile{flowpoints}
\caption{Two points $P$ and $P(L)$ on the critical curve, with coordinates
$\lambda$ and $\lambda(L)$.}
\label{fig:flowpoints}
\end{figure}
The equation \eqref{eq:CSsol} then states that the two field theories $P$ and
$P(L)$ are related by the scale transformation
\begin{equation}\label{eq:scaletr}
x\to L\,x
\end{equation}
(the field renormalizations obviously do not affect the physical content of the
theory). The equation
\begin{equation}\label{eq:RGflow}
L\frac{\dd}{\dd L}\,\lambda(L)=-\beta\bigl(\lambda(L)\bigr)
\end{equation}
describes, in given coordinate system, the ``flow'' in the space of field
theories generated by the scale transformations. It is called the
\emph{renormalization group flow}, and the equation \eqref{eq:RGflow} is usually
referred to as the \emph{renormalization group flow equation}. The quantity
$\lambda(L)$ is often referred to as ``running'', or ``scale-dependent'' coupling
constant. The fact that the scale transformations \eqref{eq:scaletr} generate a
``flow'' is not limited to the massless theory; we will see shortly that the
Callan-Symanzik equation generalizes to the case $m^2\neq0$.

We have seen that the functions $\beta(\lambda)$ and $\gamma(\lambda)$ contain
most important information about the scale dependence of the massless theory.
Unfortunately, these functions are not known exactly. However we can compute them
perturbatively, order by order in the coupling parameter $\lambda$.

There are many ways to do this. The simplest is just to compute the correlation
functions perturbatively and then determine the functions $\beta$ and $\gamma$
from the Callan-Symanzik equation \eqref{eq:CS} itself. Let us first note that the
Callan-Symanzik equation in the form analogous to \eqref{eq:CS} applies to the
momentum-space correlation functions as well as to coordinate-space ones (the
Fourier transform relating the two depends neither on $\lambda$ nor on $\mu$).
Moreover, simple calculation shows that this equation can be written as the
equation
\begin{equation}\label{eq:CSGamma}
\left(\mu\frac\partial{\partial\mu}+\beta(\lambda)\,\frac\partial{\partial\lambda}-n\,\gamma(\lambda)\right)\Gamma^{(n)}(p_1,\dots,p_n)=0
\end{equation}
involving the proper vertices.

To the order $\lambda^2$ the proper vertices $\Gamma^{(2)}$ and $\Gamma^{(4)}$,
when computed in the scheme \eqref{eq:mu1}--\eqref{eq:mu2}, are (I ask you to
carry out the calculations as the homework Problem~\ref{hw:5})
\begin{align}
\Gamma^{(2)}(p^2)&=p^2+\frac1{12}\,\frac{\lambda^2}{(4\pi)^4}\,p^2\,\Bigl(\log(\mu^2/p^2)+1\Bigr)+O(\lambda^3)\,,\label{eq:G2mu}\\
\Gamma^{(4)}(p_1,\dots,p_4)&=\lambda+\frac{\lambda^2}{2\,(4\pi)^2}\left[\log\left(\frac{p_{12}^2}{4\mu^2}\right)+\log\left(\frac{p_{13}^2}{4\mu^2}\right)+\log\left(\frac{p_{14}^2}{4\mu^2}\right)\right]+O(\lambda^3)\,,\label{eq:G4mu}
\end{align}
where $p_{12}^2=(p_1+p_2)^2$, etc. Let us first plug \eqref{eq:G2mu} into
\eqref{eq:CSGamma}. The result is
\[
\frac2{12}\,\frac{\lambda^2}{(4\pi)^4}\,p^2+\beta(\lambda)\,\frac2{12}\,\frac\lambda{(4\pi)^4}\,p^2\,\bigl[\log(\mu^2/p^2)+1\bigr]-2\gamma(\lambda)\,p^2+O(\lambda^3)=0\,,
\]
where the first term came from the $\mu\partial/\partial\mu$ term in
\eqref{eq:CSGamma}. The second term contains $p^2\log\mu^2/p^2$; this dependence
cannot be canceled by two other terms exposed in the above equation. The only way
to satisfy this equation in the order $\lambda^2$ is to assume that
$\beta(\lambda)\sim\lambda^2$ so that the second term has to be attributed to the
$O(\lambda^3)$ corrections. We then find
\begin{equation}\label{eq:gamma1}
\gamma(\lambda)=\frac1{12}\,\frac{\lambda^2}{(4\pi)^4}+O(\lambda^3)\,.
\end{equation}
Similarly, substitute the above perturbative result for $\Gamma^{(4)}$ into
\eqref{eq:CSGamma}. We have
\[
-3\,\frac{\lambda^2}{(4\pi)^2}+\beta(\lambda)-4\lambda\,\gamma(\lambda)+O(\lambda^3)=0\,.
\]
We already know that $\gamma(\lambda)\sim\lambda^2$ and so the term involving this
function is $O(\lambda^3)$. Therefore
\begin{equation}\label{eq:beta1}
\beta(\lambda)=3\,\frac{\lambda^2}{(4\pi)^2}+O(\lambda^3)\,.
\end{equation}

Note that the leading term \eqref{eq:beta1} of the $\beta$-function is
\emph{positive}. This is the feature of $\varphi^4$ theory. The flow equation
\eqref{eq:flow1} then shows that for small $\lambda$ the scale-dependent coupling
constant $\lambda(L)$ \emph{decreases} with $L$. The equation \eqref{eq:CSsol}
shows that larger $L$ correspond to \emph{longer} length scales (this is clear if
one rewrites \eqref{eq:CSsol} in equivalent form
\[
\bigl\langle\varphi(x_1)\cdots\varphi(x_n)\bigr\rangle_{\lambda(L)}\sim\bigl\langle\varphi(Lx_1)\cdots\varphi(Lx_n)\bigr\rangle_\lambda
\]
showing that the correlation functions computed with $\lambda(L)$ are equivalent
to the correlation functions computed with $\lambda$, with the separations
``stretched'' by the scale factor $L$).

If one ignores the higher order corrections in \eqref{eq:beta1} the flow equation
\eqref{eq:flow1} is easily integrated,
\begin{equation}\label{eq:running}
\lambda(L)=\frac\lambda{1+3\frac\lambda{(4\pi)^2}\log L}\,.
\end{equation}
At sufficiently small $L$ (i.e.\ sufficiently short distances) the function
$\lambda(L)$ diverges. This signals potential trouble with locality which we
discussed earlier. Of course when $\lambda(L)$ becomes large one cannot ignore the
higher-order corrections in \eqref{eq:beta1}; in fact in order to understand
short-distance behavior of $\varphi^4$ theory we need to know the shape of the
$\beta$-function at large $\lambda$. We see that the question of consistency of
$\varphi^4$ theory cannot be settled on the perturbative level.

In nonabelian gauge theories the leading term of the $\beta$-function turns out
to be \emph{negative}, meaning that the running coupling constant
\emph{decreases} at short distances (instead it increases at large distances); as
a result, although nonabelian gauge theory usually cannot be solved by means of
perturbation theory, the question of its \emph{consistency} as a local field
theory can be positively answered on the basis of perturbative calculations. We
will discuss this phenomenon known as \emph{asymptotic freedom} in greater details
later.

\subsection*{Mass perturbation}

Let us refine our techniques further. So far we have considered the
renormalization group flow in massless $\varphi^4$ theory. We would like to
generalize this picture to the massive case. The most natural way to do this is
to consider the massive $\varphi^4$ theory as the massless theory perturbed by
adding the term $\frac{m^2}2\varphi^2$ to the action density. We already know that
the field $\varphi^2$ requires renormalization. Let us discuss how this field
renormalizes in our $\mu$-dependent scheme.

In the massless theory the normalization condition
$\Gamma^{(\varphi^2_R,2)}(0,0)=1$ we have used before has to be replaced by some
nonsingular $\mu$ dependent condition. One can take for example
\begin{equation}\label{eq:phi2mu}
\Gamma^{(\varphi^2_R,2)}(p_1,p_2)\Big|_{\substack{p_i^2=2\mu^2\\p_1\cdot p_2=-\mu^2}}=1\,.
\end{equation}
The renormalized field $[\varphi^2]_{(\mu)}$ defined by this condition differs
from the bare field $\varphi_0^2$ by the renormalization
constant\footnote{Note that this definition of the $\varphi^2$ renormalization
constant differs from that introduced in Section~14, Eq.~\eqref{eq:phi2Rinv}.
There, we have defined the constant $Z_{\varphi^2}$ as the coefficient between
$[\varphi^2]_R$ and $\varphi^2$. The new $Z^{\varphi^2}$ differs from
$Z_{\varphi^2}$ in Section~14 by the factor of $Z$.},
\[
[\varphi^2]_{(\mu)}=\bigl(Z^{\varphi^2}_{(\mu)}\bigr)^{-1}\,\varphi_0^2\,.
\]
Simple calculation similar to that discussed in the previous Section shows that
in the leading order
\begin{equation}\label{eq:Zphi2mu}
Z^{\varphi^2}_{(\mu)}=1-\frac12\,\frac\lambda{(4\pi)^{d/2}}\,(2\mu^2)^{d/2-2}\,\frac{\Gamma(2-d/2)\,\Gamma^2(d/2-1)}{\Gamma(d-2)}+O(\lambda^2)
\end{equation}
Although $Z^{\varphi^2}_{(\mu)}$ is divergent at $d=4$, the ratio
$Z^{\varphi^2}_{(\tilde\mu)}/Z^{\varphi^2}_{(\mu)}$ is expected to be finite. An
infinitesimal change of the renormalization scale $\mu\to\mu+\delta\mu$ then leads
to infinitesimal renormalization
\[
[\varphi^2]_{(\mu+\delta\mu)}=\left(1-\frac{\delta\mu}\mu\,\gamma_{\varphi^2}(\lambda)\right)\,[\varphi^2]_{(\mu)}\,,
\]
where
\begin{equation}\label{eq:gammaphi2}
\gamma_{\varphi^2}(\lambda)=\bigl(Z^{\varphi^2}_{(\mu)}\bigr)^{-1}\,\mu\frac{\dd}{\dd\mu}\,Z^{\varphi^2}_{(\mu)}
\end{equation}
is finite function of $\lambda$. From \eqref{eq:Zphi2mu} we obtain
\begin{equation}\label{eq:gammaphi2val}
\gamma_{\varphi^2}(\lambda)=\frac\lambda{(4\pi)^2}+O(\lambda^2)\,.
\end{equation}
Repeating our previous calculations we find that a correlation function with $N$
insertions of $[\varphi^2]_{(\mu)}$ and $n$ insertions of $\varphi$, satisfies the
modified Callan-Symanzik equation
\[
\left(\mu\frac\partial{\partial\mu}+\beta(\lambda)\,\frac\partial{\partial\lambda}+n\,\gamma(\lambda)+N\,\gamma_{\varphi^2}(\lambda)\right)
\bigl\langle\varphi(x_1)\cdots\varphi(x_n)\,[\varphi^2]_{(\mu)}(y_1)\cdots[\varphi^2]_{(\mu)}(y_N)\bigr\rangle_{(\mu)}=0\,.
\]

Let us now introduce the mass term as
\[
\mathcal A_{M^2}=\mathcal A_0+\int \dd^dx\,\frac{M^2}2\,[\varphi^2]_{(\mu)}\,,
\]
where $\mathcal A_0$ is the massless action
\[
\mathcal A_0=\int \dd^dx\,\left[\frac12\,(\partial\varphi)^2+\frac\lambda{4!}\,\varphi^4+\frac{\delta Z_{(\mu)}}2\,(\partial\varphi)^2+\frac{\delta\lambda_{(\mu)}}{4!}\,\varphi^4\right]\,.
\]
Here $\delta Z_{(\mu)}$ and $\delta\lambda_{(\mu)}$ are counterterms computed in the
massless theory, within the above $\mu$-dependent renormalization scheme (and I
assumed dimensional regularization and therefore the $\delta m^2$ counterterm is
not included). Recalling that
\[
[\varphi^2]_{(\mu)}=\frac{Z_{(\mu)}}{Z^{\varphi^2}_{(\mu)}}\,\varphi^2\,,\qquad Z_{(\mu)}=1+\delta Z_{(\mu)}\,,
\]
the massive action can be written as
\[
\mathcal A_{M^2}=\int \dd^dx\,\left[\frac12\,(\partial\varphi)^2+\frac{M^2}2\,\varphi^2+\frac\lambda{4!}\,\varphi^4+\frac{\delta Z_{(\mu)}}2\,(\partial\varphi)^2+\frac{\delta M^2_{(\mu)}}2\,\varphi^2+\frac{\delta\lambda_{(\mu)}}{4!}\,\varphi^4\right]\,,
\]
where
\begin{equation}\label{eq:deltaM2}
\delta M^2_{(\mu)}=M^2\left(\frac{Z_{(\mu)}}{Z^{\varphi^2}_{(\mu)}}-1\right)\,.
\end{equation}
This action can be used now to develop a renormalized perturbation theory in the
massive case. Note that the above definition of the mass parameter corresponds to
specific renormalization scheme in which
\begin{itemize}
\item The counterterms $\delta Z_{(\mu)}$ and $\delta\lambda_{(\mu)}$ are borrowed
from the massless theory; in particular these coefficients do not depend on $M^2$.
\item The counterterm $\delta M^2_{(\mu)}$ is expressed by \eqref{eq:deltaM2}
through the $\varphi^2$-renormalization constant $Z^{\varphi^2}_{(\mu)}$; again,
this constant has to be computed upfront in the massless theory, and
$\delta M^2_{(\mu)}$ depends on $M^2$ linearly.
\end{itemize}
Although it is possible to show that the counterterms thus defined indeed
compensate for the divergences in the renormalized perturbation theory, no
explicit normalization conditions are imposed in the massive theory. Relation
between $M^2$, $\lambda$ and physical quantities in this scheme is not specified
in advance but must be computed order by order in the perturbation theory. In
this sense this renormalization scheme is similar to the ``minimal subtraction''
scheme.

Note that in this scheme the correlation functions depend on two mass
parameters, $M$ and our auxiliary scale $\mu$, although we know only one is
needed. This is the price we pay for defining specific coordinate system
$(M^2,\lambda)$ which is nonsingular on the critical submanifold $M^2=0$. It is
not difficult to check that correlation functions computed within this scheme
admit smooth limit $M^2=0$; in fact it is clear in advance: at $M^2=0$ this
perturbation theory simply reduces to the massless renormalized perturbation
theory.

Like in the massless case, nothing except for normalizations depends on the
auxiliary scale $\mu$. Under the infinitesimal change
\[
\mu\to\mu+\delta\mu
\]
the renormalized field $[\varphi^2]_{(\mu)}$ transforms as
\[
[\varphi^2]_{(\mu+\delta\mu)}=\left(1-\frac{\delta\mu}\mu\,\gamma_{\varphi^2}(\lambda)\right)\,[\varphi^2]_{(\mu)}\,.
\]
This brings in certain change in the action, which however is compensated if we
simultaneously change the mass parameter,
\[
M^2\to\left(1+\frac{\delta\mu}\mu\,\gamma_{\varphi^2}(\lambda)\right)\,M^2\,.
\]
Invariance of the action w.r.t.\ the above simultaneous change in $\mu$ and $M^2$
leads to the equation
\[
\left(\mu\frac\partial{\partial\mu}+\beta(\lambda)\,\frac\partial{\partial\lambda}+n\,\gamma(\lambda)+\gamma_{\varphi^2}(\lambda)\,M^2\frac\partial{\partial M^2}\right)\bigl\langle\varphi(x_1)\cdots\varphi(x_n)\bigr\rangle=0\,,
\]
or, equivalently (using $\langle\varphi(x_1)\cdots\varphi(x_n)\rangle=\mu^n\,C^{(n)}(\mu x_i,M^2/\mu^2|\lambda)$),
\[
\left(\sum_ix_i\frac\partial{\partial x_i}+n\bigl(1+\gamma(\lambda)\bigr)+\beta(\lambda)\,\frac\partial{\partial\lambda}+\bigl(-2+\gamma_{\varphi^2}\bigr)\,M^2\frac\partial{\partial M^2}\right)\bigl\langle\varphi(x_1)\cdots\varphi(x_n)\bigr\rangle=0\,.
\]
This equation shows that in the massive case the scale transformations induce the
``flow'' in the two-dimensional space of $\varphi^4$ theories,
\begin{equation}\label{eq:massflow}
L\frac{\dd}{\dd L}\,\lambda=-\beta(\lambda)\,,\qquad L\frac{\dd}{\dd L}\,M^2=\bigl(2-\gamma_{\varphi^2}(\lambda)\bigr)\,M^2\,.
\end{equation}
Note that $\beta(\lambda)$ and $\gamma(\lambda)$ here are exactly the same as in the
massless theory, while $\gamma_{\varphi^2}(\lambda)$ is also determined from the
massless theory via Eq.~\eqref{eq:gammaphi2}.

\begin{hwproblem}{5}\label{hw:5}
In the massless $\varphi^4$ theory calculate $\Gamma^{(2)}$ and $\Gamma^{(4)}$
explicitly, to the order $\lambda_{(\mu)}^2$. Use dimensional regularization, and
the $\mu$-dependent scheme defined by Eqs.~\eqref{eq:mu1}, \eqref{eq:mu2}. In
particular, calculate the counterterm coefficients $z^{(2)}_{(\mu)}$ and
$a^{(2)}_{(\mu)}$, and confirm that $b^{(1)}_{(\mu)}=b^{(2)}_{(\mu)}=0$.

In solving this problem the following dimensionally continued integral is useful
\[
\int\frac{\dd^dk}{k^2}=0\,.
\]
\end{hwproblem}

\begin{hwproblem}{6}\label{hw:6}
Consider a scalar field theory with $N$-component field
$\bm{\varphi}=(\varphi_1,\dots,\varphi_N)$, and $O(N)$-symmetric action
\[
\mathcal A=\int \dd^dx\,\left[\frac12\,\bigl(\partial\bm{\varphi}\bigr)^2+\frac{M^2}2\,\bm{\varphi}^2+\frac\lambda8\,\bigl(\bm{\varphi}^2\bigr)^2\right]\,.
\]
\begin{hwparts}
\item Use the $O(N)$ symmetry to show that the two- and four-point correlation
functions $\tilde W^{(2)}_{ij}(p)$ and $\tilde W^{(4)}_{ijkl}(p_1,p_2,p_3,p_4)$
defined as
\begin{align*}
\tilde W^{(2)}_{ij}(p)&:=\text{f.t.}\,\bigl\langle\varphi_i(x)\varphi_j(x')\bigr\rangle\,,\\
\tilde W^{(4)}_{ijkl}(p_1,p_2,p_3,p_4)&:=\text{f.t.}\,\bigl\langle\varphi_i(x_1)\varphi_j(x_2)\varphi_k(x_3)\varphi_l(x_4)\bigr\rangle
\end{align*}
enjoy the following structures
\begin{align*}
\tilde W^{(2)}_{ij}(p)&=\delta_{ij}\,\tilde W^{(2)}(p)\,,\\
\tilde W^{(4)}_{ijkl}(p_1,p_2,p_3,p_4)&=\delta_{ij}\delta_{kl}\,\tilde W^{(4)}(p_1,p_2;p_3,p_4)+\delta_{ik}\delta_{jl}\,\tilde W^{(4)}(p_1,p_3;p_2,p_4)\\
&\qquad+\delta_{il}\delta_{jk}\,\tilde W^{(4)}(p_1,p_4;p_2,p_3)
\end{align*}
with scalar functions $\tilde W^{(2)}(p)$ and $\tilde W^{(4)}(p_1,p_2;p_3,p_4)$.
Define the scalar proper vertices $\Gamma^{(2)}(p)$ and
$\Gamma^{(4)}(p_1,p_2,p_3,p_4)$.
\item Introduce all elements of Feynman diagram technique (in the momentum space)
for this model, i.e.\ propagators and vertices.
\item For the massless theory write down the Callan-Symanzik equations for the
proper vertices $\Gamma^{(2)}(p)$ and $\Gamma^{(4)}(p_1,p_2,p_3,p_4)$.
\item At $d=4$ compute the functions $\beta(\lambda)$ and $\gamma(\lambda)$ to the
order $\lambda^2$.
\end{hwparts}
\end{hwproblem}

% Source pp. 161-172 (Lectures13-18.pdf).
\section{Wilson's Renormalization Group}

\subsection*{On Second Order Phase Transitions}

Let us briefly discuss the so called \emph{continuous} or \emph{second order}
phase transitions, because it is the theory of continuous phase transitions in
statistical mechanics which is related to (Euclidean) quantum field theory in
most close way.

Consider a magnet, i.e.\ a piece of a magnetic material. It can be thought of
as consisting of ``micromagnets'', i.e.\ small elements having nonzero magnetic
moments. Although in reality the structure of ``micromagnets'' is usually rather
complex, for the sake of this discussion we can think of, say, a cubic lattice
with each cell representing a ``micromagnet'' having a magnetic moment
(proportional to the volume of the cell),
\[
\diac[0.55]{\draw[prop] (0,0) rectangle (2,2); \draw[prop] (0,2) -- (0.8,2.8) -- (2.8,2.8) -- (2,2); \draw[prop] (2.8,2.8) -- (2.8,0.8) -- (2,0);
  \draw[densely dashed] (0,0) -- (0.8,0.8) -- (2.8,0.8) (0.8,0.8) -- (0.8,2.8);
  \draw[{Stealth[length=1.4mm]}-{Stealth[length=1.4mm]}] (3.2,0.8) -- node[right] {\small$a$} (3.2,2.8);}
\qquad\longrightarrow\qquad a^3\cdot\mathbf m_i
\]
The interaction between the ``micromagnets'' is such that they tend to align,
thus producing an overall magnetic moment
\[
V\,\mathbf M=V\,\langle\mathbf m_i\rangle\,,
\]
where $\mathbf M$ denotes specific (per unit volume) magnetization, and the
expectation value in the r.h.s.\ is taken over thermal equilibrium state. The
thermal fluctuations tend to put the directions of $\mathbf m_i$ at random, and
as the result of this competition the ferromagnet has two phases, the low
temperature, or \emph{ordered} phase with nonvanishing local magnetization
$\mathbf M$, and high-temperature \emph{disordered} phase with $\mathbf M=0$.

The most important qualitative features of this phase transition, which occurs
at certain critical temperature
\[
T=T_c\,,
\]
can be understood through the \emph{Landau Theory}.

\paragraph{Landau Theory:} Consider the thermodynamic potential
\[
V\,\Omega(\mathbf M,T)
\]
as the function of specific magnetization $\mathbf M$ and temperature $T$;
$\Omega$ denotes the \emph{specific} (per unit volume) thermodynamic potential.
The equilibrium value of $\mathbf M$ is determined by the equation
\[
\frac{\partial\Omega(\mathbf M,T)}{\partial\mathbf M}=0\,.
\]
Consider the high-temperature phase $T>T_c$, where the solution of this equation
is $\mathbf M=0$, and assume that one can expand $\Omega(\mathbf M,T)$ in powers
of $\mathbf M$ near this minimum,
\[
\Omega(\mathbf M,T)=A(T)+B(T)\,\mathbf M^2+C(T)\,\mathbf M^4+\cdots
\]
(I assume here isotropy of the problem, hence the expansion is in powers of
$\mathbf M^2$). It is clear that $\mathbf M=0$ is the right solution (i.e.\ the
minimum of $\Omega$) only as long as
\[
B(T)>0\ \to\ \mathbf M=0\,.
\]
When $B(T)$ becomes negative (as it happens at $T<T_c$, according to this
theory), the minimum occurs at nonzero $\mathbf M$ such that
\[
|\mathbf M|=\sqrt{\frac{-B}{2C}}\,;
\]
in fact there is a continuum of degenerate minima corresponding to different
directions of $\mathbf M$. This situation represents the low-temperature phase
where spontaneous magnetization appears. Thus, in the Landau theory the critical
temperature $T_c$ is the point where $B(T)$ changes its sign,
\[
\begin{array}{lll}
B(T)>0 & \text{at} & T>T_c\\
B(T)<0 & \text{at} & T<T_c
\end{array}
\]
Landau has assumed that near $T_c$ one can write $B(T)$ as
\[
B(T)=b\,(T-T_c)+\cdots\,,
\]
with some positive constant $b$. Then it follows that
\[
\begin{array}{lll}
|\mathbf M|=0 & \text{at} & T>T_c\\
|\mathbf M|\sim(T_c-T)^{1/2} & \text{at} & T<T_c
\end{array}
\]
It is not difficult to calculate the entropy at $T\approx T_c$,
\[
S/V=-\frac{\partial\Omega(\mathbf M,T)}{\partial T}\approx S_0/V+
\begin{cases}0 & \text{for}\quad T>T_c\\ \frac{b^2}{2C}\,(T-T_c) & \text{for}\quad T<T_c\end{cases}
\]
where $S_0$ refers to the critical point. The entropy (and hence the internal
energy) is continuous at $T_c$, so there is no latent heat. This is the
characteristic property of a second order phase transition. The specific heat
\[
C=T\,\frac{\partial S}{\partial T}
\]
has finite discontinuity at $T=T_c$.

\paragraph{Critical behavior:} Although the Landau theory explains existence and
the ``continuous'' nature of the critical point, in real systems the singularity
of the thermodynamic functions at $T_c$ turns out to be more complicated than the
above arguments predict. Experiment shows that
\[
|\mathbf M|\sim(T_c-T)^\beta\qquad\text{for}\quad T<T_c\,,
\]
with $\beta\approx0.35\pm0.03$, definitely different from $0.5$. The specific heat
is not only discontinuous, but in fact diverges near $T_c$,
\[
C\sim|T_c-T|^{-\alpha}\,,
\]
with $\alpha\approx0.12$, not 0.

The Landau theory does not take into account prominent role of the
\emph{long-range fluctuations} of the magnetization near critical point.
According to the thermodynamic theory of fluctuations, the probability to have
some non-equilibrium value $\mathbf M$ of the magnetization is (for $T>T_c$)
\[
P(\mathbf M)\sim e^{-\frac{V\Omega(\mathbf M,T)}{kT}}\sim e^{-\frac{VB(T)}{kT}\,\mathbf M^2}
\]
If $B>0$, in the limit $V\to\infty$ the fluctuations are totally suppressed.
Vanishing of $B(T)$ at $T=T_c$ indicates that in the critical domain the
fluctuations must play significant role.

\subsection*{Fluctuation Theory of Phase Transitions}

Previously, we have given a brief discussion of critical behavior near a second
order phase transition. We have pointed out singular character of the
thermodynamic functions near the critical point, and suggested that the origin of
the singularities is in the enhanced role of fluctuations.

The qualitative picture of the fluctuations is as follows. In the vicinity of
$T_c$ the competition between the energetic tendency of the micromagnets to
align, and the thermal fluctuations leads to formation of large domains of
approximately collinear micromagnets (figure~\ref{fig:domains}),
\begin{figure}[ht]
\centering\figfile{domains}
\caption{Two distant domains near $T_c$. Inside each domain the micromagnets are
nearly aligned; the average directions in different domains differ.}
\label{fig:domains}
\end{figure}
so that if one looks inside one domain the magnets are almost aligned, but their
average directions can be largely different in two distant domains. Near $T_c$
the average sizes of the domains become much larger than any microscopic scales,
and the domains involve macroscopic numbers of the micromagnets. The orientations
of the domains experience fluctuations which vary slowly in space (in the
microscopic scale). These long-range fluctuations lead to formation of the
critical singularities at $T_c$.

Having this picture in mind, one can consider local magnetization
$\mathbf M(x)$, a slow varying function of the spatial coordinates (meaning that
the distance at which $\mathbf M(x)$ changes significantly must involve
macroscopic number of the micromagnets). One then replaces the homogeneous
magnetization $\mathbf M$ by the local magnetization
\[
V\,\Omega(\mathbf M,T)\ \to\ \int d^3x\,\Omega\bigl(\mathbf M(x),T\bigr)\,.
\]
Once $\mathbf M$ is not a constant, one may include the spatial derivatives of
$\mathbf M(x)$, and assume that the long-range fluctuations have the probability
distribution
\[
P[\mathbf M(x)]\sim Z^{-1}\exp\left(-\frac{\Ham[\mathbf M(x)]}{kT}\right)
\]
with the ``effective energy'' of the form
\[
\Ham[\mathbf M(x)]=\int d^3x\,\Omega\bigl(\mathbf M(x),\nabla\mathbf M(x),\dots;T\bigr)
\]
where $\Omega$ now is a local density. In the spirit of the Landau theory we can
write
\[
\Omega=A(T)+B(T)\,\mathbf M^2(x)+C(T)\,\mathbf M^4(x)+D(T)\,\bigl(\nabla\mathbf M(x)\bigr)^2+\dots
\]
where the dots may involve higher powers of $\mathbf M^2$ as well as terms with
higher derivatives. Neglecting all but the terms explicitly shown above
corresponds to the so-called Landau-Ginzburg model.

With this probability distribution, the free energy of the magnet is to be
computed as some sort of functional integral
\[
Z=e^{-\frac F{kT}}=\int\bigl[D\mathbf M(x)\bigr]\,e^{-\frac1{kT}\Ham[\mathbf M(x)]}
\]
where the integration is over the functions $\mathbf M(x)$ which vary slowly in
terms of the microscopic scale.

Note that this expression differs mostly in notations from the functional
integral in Euclidean scalar field theory. After the rescaling
\[
\sqrt{\frac{2D(T)}{kT}}\,\mathbf M(x)=\varphi_0(x)
\]
it becomes
\[
Z=\int\bigl[D\varphi_0\bigr]\,e^{-\mathcal A[\varphi_0]}\,,
\]
where in the Landau-Ginzburg approximation
\[
\mathcal A=\int\left[\frac12\,(\partial\varphi_0)^2+\frac{m_0^2}2\,\varphi_0^2+\frac{\lambda_0}{4!}\,\varphi_0^4\right]
\]
has the familiar form, with
\[
m_0^2=m_0^2(T)\,,\qquad\lambda_0=\lambda_0(T)
\]
some functions of $T$, whose form depends on microscopic details of the magnetic
system. On the superficial level the main difference is in the dimensionality of
the space (3 instead of 4), and the number of components of the scalar field (3
instead of 1). As for the difference in $d$, we have already seen that the idea
of studying QFT in spaces of different (even non-integer) dimensions $d$ is
useful. The problematic of the phase transition theory, in its application to
condensed matter physics, usually corresponds to the case $d=3$. But it is
interesting and instructive to study criticality at any $d$, in particular to see
how the characteristics of a phase transition change with $d$. Also, in the above
simple model of magnet, the direction of $\mathbf M$ is completely decoupled from
any space direction. Therefore formally the dimensionality of the vector
$\mathbf M$, or equivalently the number of components of
\[
\varphi_0=\{\varphi_{0,i}\}\,,\qquad i=1,2,\dots,N
\]
can be any number, not necessarily $d$. In fact, there are magnets with
anisotropy of the type of ``easy plane'' or ``easy axis'', which correspond to
$N=2$ and $N=1$, respectively. In most of the following discussion I will assume
$N=1$ for simplicity of notations.

On a deeper level, one can raise two objections against this picture. First, the
above formulation assumes integration only over the ``slow varying'' fields
$\varphi_0(x)$. This is similar to the idea of cutoff in field theory, but here
it seems to be dictated by physics of the problem. Indeed, in our approach the
``thermodynamic potential''
\[
\Ham([\mathbf M(x)],T)
\]
already takes into account ``microscopic'' thermal fluctuations in our system,
while the ``macroscopic'' fluctuations are left to be accounted for in the
functional integral over $D[\mathbf M(x)]$. But how these two types of
fluctuations are separated? We have said that $\mathbf M(x)$ can vary on the
scale ``much larger than the microscopic scales''. But how much larger? Should we
take the cutoff parameter $\Lambda=(100\,a)^{-1}$, or $(10^4\,a)^{-1}$? The second
objection is that it is not clear what makes it possible to neglect the higher
order (both in $\mathbf M$ and derivatives) terms in $\Ham$.

These questions are at the heart of the Wilson's approach to the Renormalization
Group.

\subsection*{Wilson's Approach\protect\footnote{Discussion in this and following
Sections has overlap with \S12.1 in PS. More systematic treatment is in K.~Wilson,
J.~Kogut, ``The RG and epsilon expansion'', Phys.\ Rep.\ \textbf{12C} (1974)
(highly recommended!). See also L.~Kadanoff, ``Statistics, Dynamics, and
Renormalization'', WS.}}

To come to the idea of this approach let us ask instead \emph{what difference
does it make} if we take
\[
\Lambda=\Lambda_0=(100\,a)^{-1}\qquad\text{or}\qquad\Lambda=\Lambda_1=(10^4\,a)^{-1}\,.
\]
In other words, if we consider say $\varphi^4$ theory, and write down the
functional integral
\[
\int D[\varphi_0]\,e^{-\mathcal A[\varphi_0]}\,(\dots)
\]
in its bare form, with some cutoff regularization implied, how the theory depends
on the value of the cutoff parameter $\Lambda$?

Because basic ideas of the Wilson's RG do not depend very much on precise way how
the cutoff regularization is implemented, let me assume specific one. Let me
assume that the field $\varphi_0$ can be written as the Fourier-like integral
\[
\varphi_0(x)=\int_{|k|<\Lambda}\frac{d^dk}{(2\pi)^d}\,\phi(k)\,e^{ikx}\,,
\]
where only the wave-vectors with $|k|<\Lambda$ contribute to the integral. This
can be named (as I will) the ``Sharp momentum cutoff'', because in perturbation
theory it corresponds to the cutoff propagator
\[
D_\Lambda(k)=\frac{\Phi(k^2/\Lambda^2)}{k^2+m_0^2}
\]
with $\Phi(x)=\Theta(1-x)$ being the step function. (We did not consider this
regularization before because it is not very convenient in evaluating usual
Feynman diagrams).

So our question now is how the functional integral depends on $\Lambda$, or what
is the difference between the integral with $\Lambda=\Lambda_0$ and
$\Lambda=\Lambda_1$. Let us temporarily use the notations
\[
\varphi_0(x)=\int_{|k|<\Lambda_0}\phi(k)\,e^{ikx}\,,\qquad
\varphi_1(x)=\int_{|k|<\Lambda_1}\phi(k)\,e^{ikx}\,.
\]
If $\Lambda_0>\Lambda_1$ the difference of course is that the integration over
$\varphi_0$ includes the integration over all modes with the momenta
\[
\Lambda_0>|k|>\Lambda_1\,,
\]
while the integral over $\varphi_1$ does not. We can write
\[
\varphi_0(x)=\varphi_1(x)+\tilde\varphi(x)\,,\qquad\text{where}\qquad
\tilde\varphi(x)=\int_{\Lambda_0>|k|>\Lambda_1}\phi(k)\,e^{ikx}\,.
\]
In this context, the parts $\varphi_1(x)$ and $\tilde\varphi$ are often called
the ``slow part'' and the ``fast part'', respectively. We have
\[
D[\varphi_0]=D[\varphi_1]\,D[\tilde\varphi]\,.
\]
Suppose we want to calculate some correlation function in the theory with the
cutoff $\Lambda_0$,
\[
\bigl\langle\varphi_0(x_1)\cdots\varphi_0(x_n)\bigr\rangle=\frac1{Z_0}\int D[\varphi_0]\,\varphi_0(x_1)\cdots\varphi_0(x_n)\,e^{-\mathcal A_0[\varphi_0]}\,,
\]
where $\mathcal A_0$ is some action, say $\varphi^4$ action above. We are
interested in low-momentum behaviour of the associated proper vertices, with the
momenta conjugated to $x_1,\dots,x_n$ much smaller than both $\Lambda_0$ and
$\Lambda_1$. Therefore in the insertions we can replace
\[
\varphi_0(x_1)\cdots\varphi_0(x_n)\to\varphi_1(x_1)\cdots\varphi_1(x_n)\,.
\]
Then the functional integral for this correlation function becomes
\[
\frac1{Z_0}\int D[\varphi_1]\,\varphi_1(x_1)\cdots\varphi_1(x_n)\int D[\tilde\varphi]\,e^{-\mathcal A_0[\varphi_1+\tilde\varphi]}\,.
\]
Now, the idea is to integrate over the fast part $\tilde\varphi$ first. This
operation defines the new ``effective action'' $\mathcal A_1[\varphi_1]$ for the
slow part,
\[
\int D[\tilde\varphi]\,e^{-\mathcal A_0[\varphi_1+\tilde\varphi]}=\Const\,e^{-\mathcal A_1[\varphi_1]}\,.
\]
The constant factor here is not important, because after the same operation it
appears in $Z_0$ as well. Thus, we have
\[
\text{f.t.}\,\bigl\langle\varphi_0(x_1)\cdots\varphi_0(x_n)\bigr\rangle\Big|_{\mathcal A_0,\Lambda_0}
=\text{f.t.}\,\bigl\langle\varphi_1(x_1)\cdots\varphi_1(x_n)\bigr\rangle\Big|_{\mathcal A_1,\Lambda_1}\,,\qquad p_i\ll\Lambda\,.
\]
In other words, if only the low-momentum correlation functions are of concern,
the theory with cutoff $\Lambda_0$ and action $\mathcal A_0$ is equivalent to the
theory with lower cutoff $\Lambda_1$ and new action $\mathcal A_1$,
\[
\mathcal A_0,\Lambda_0\ \to\ \mathcal A_1,\Lambda_1\,.
\]
Note that I am not writing $\leftrightarrow$ here, because this operation can not
be undone. Information about the short-distance interactions, of the scales
between $\Lambda_0^{-1}$ and $\Lambda_1^{-1}$ is lost.

The new action $\mathcal A_1$ is defined through the functional integral over the
fast part $\tilde\varphi$. Although at this level of discussion precise form of
$\mathcal A_0$ is not important, to get some idea what this operation amounts to,
let us assume for the moment that $\mathcal A_0$ is the $\varphi^4$ action.
Substituting
\[
\varphi_0=\varphi_1+\tilde\varphi
\]
we find
\[
\mathcal A_0[\varphi_0]=\mathcal A_0[\varphi_1]+\tilde{\mathcal A}[\tilde\varphi,\varphi_1]\,,
\]
and after the integration over $\tilde\varphi$
\[
\mathcal A_1[\varphi_1]=\mathcal A_0[\varphi_1]+\mathcal A'[\varphi_1]\,,\qquad\text{where}\qquad
\Const\,e^{-\mathcal A'[\varphi_1]}=\int[D\tilde\varphi]\,e^{-\tilde{\mathcal A}[\tilde\varphi,\varphi_1]}\,.
\]
For $\mathcal A_0$ of the $\varphi^4$ form
\[
\tilde{\mathcal A}=\int d^dx\,\left[\frac12\,(\partial\tilde\varphi)^2+\frac{m_0^2}2\,\tilde\varphi^2
+\frac{\lambda_0}{4!}\,\bigl(4\varphi_1^3\tilde\varphi+6\varphi_1^2\tilde\varphi^2+4\varphi_1\tilde\varphi^3+\tilde\varphi^4\bigr)\right]
\]
The terms $\partial\varphi_1\partial\tilde\varphi$ and $\varphi_1\tilde\varphi$ do
not appear because by definition $\varphi_1$ and $\tilde\varphi$ carry different
momenta.

The integral over $\tilde\varphi$ is complicated non-gaussian integral. The
functional $\tilde{\mathcal A}[\tilde\varphi,\varphi_1]$ in this context should be
regarded as the action for $\tilde\varphi$, with $\varphi_1$ determining the
``source terms'' (linear, quadratic and cubic in $\tilde\varphi$).

As usual, the functional integral over $\tilde\varphi$ can be represented as a sum
of the Feynman diagrams. Let us denote the $\tilde\varphi$-propagator by a solid
line,
\[
\langle\tilde\varphi\tilde\varphi\rangle\ \to\
\dia{\draw[line width=0.9pt] (-1.6,0) -- (1.6,0); \node[above] at (0,0.05) {\small$k$};}
=\frac{\Theta(\Lambda_0>|k|>\Lambda_1)}{k^2+m_0^2}\,,
\]
with obvious meaning of the step-function $\Theta$. The remaining terms in
$\tilde{\mathcal A}$ give rise to the vertices
\begin{align*}
\dia[0.45]{\foreach \a in {150,180,210} \draw[prop,densely dashed] (0,0) -- (\a:1.4); \draw[line width=0.9pt] (0,0) -- (1.6,0); \node[vtx] at (0,0) {};}
&=-\frac{\lambda_0}{3!}\int_x\varphi_1^3(x)\,,&
\dia[0.45]{\foreach \a in {135,225} \draw[prop,densely dashed] (0,0) -- (\a:1.4); \foreach \a in {45,315} \draw[line width=0.9pt] (0,0) -- (\a:1.4); \node[vtx] at (0,0) {};}
&=-\frac{\lambda_0}2\int_x\varphi_1^2(x)\,,\\
\dia[0.45]{\draw[prop,densely dashed] (-1.4,0) -- (0,0); \foreach \a in {45,0,-45} \draw[line width=0.9pt] (0,0) -- (\a:1.4); \node[vtx] at (0,0) {};}
&=-\lambda_0\int_x\varphi_1(x)\,,&
\dia[0.45]{\foreach \a in {45,135,225,315} \draw[line width=0.9pt] (0,0) -- (\a:1.4); \node[vtx] at (0,0) {};}
&=-\lambda_0\int_x\,.
\end{align*}
Here the dashed line represents the field $\varphi_1$, the number of the dashed
lines showing the power of $\varphi_1$. Each dashed line carries momentum with
$|k|<\Lambda_1$.

We have to integrate over $\tilde\varphi$, with $\varphi_1$ fixed. Therefore all
the solid lines in the diagrams are contracted, while all the dashed lines are
left as external legs.

As usual, the sum of all the diagrams exponentiate in terms of the sum of
connected diagrams, i.e.\ the functional $\mathcal A'[\varphi_1]$ is given by the
sum of all connected diagrams,
\[
\mathcal A'[\varphi_1]=-\bigl(\text{sum of all connected diagrams}\bigr)
\]
Also, the diagrams with no external legs can be ignored, because they only
contribute to the constant factor in $\exp(-\mathcal A')$.

It is obvious that the new action $\mathcal A_1=\mathcal A_0+\mathcal A'$ does not
have simple form of the $\varphi^4$ action. The contributions to $\mathcal A'$
contain all powers of $\varphi_1$, from diagrams like this
\[
\diac[0.45]{\draw[line width=0.9pt] (0,0) circle (1); \coordinate (a) at (180:1); \coordinate (b) at (50:1); \coordinate (c) at (-60:1); \coordinate (d) at (-115:1);
  \draw[line width=0.9pt] (d) .. controls (0,-0.2) .. (c);
  \foreach \p/\q in {a/135,a/225,b/90,b/20,c/-20,d/-110} \draw[prop,densely dashed] (\p) -- ++(\q:1.3);
  \foreach \p in {a,b,c,d} \node[vtx] at (\p) {};}
\ \sim\varphi^6
\]
Moreover, the action $\mathcal A_1$ is not even formally local. For example, the
diagram
\[
\diac[0.45]{\draw[line width=0.9pt] (0,0) ellipse (1 and 0.7); \coordinate (L) at (-1,0); \coordinate (R) at (1,0);
  \foreach \p/\q in {L/135,L/225,R/45,R/315} \draw[prop,densely dashed] (\p) -- ++(\q:1.4);
  \node[vtx] at (L) {}; \node[vtx] at (R) {};}
\]
brings in the contribution
\[
-\frac{\lambda_0^2}{16}\int d^dx\,d^dx'\,\varphi_1^2(x)\,\tilde D^2(x-x')\,\varphi_1^2(x')\,,
\]
where $\tilde D$ is the $\tilde\varphi$ propagator,
\[
\tilde D(x)=\int_{\Lambda_0>|k|>\Lambda_1}\frac{d^dk}{(2\pi)^d}\,\frac{e^{ikx}}{k^2+m_0^2}\,.
\]
This non-locality, albeit unpleasant, is not disastrous. Recall that we have
already violated the locality by introducing the cutoff. It is important however
that the additional nonlocality in $\mathcal A_1$ is of short-range nature. More
precise statement is that the new non-local terms generated in this way admit
convergent derivative expansion.

This means the following. Suppose we have some diagram evaluated in the momentum
space,
\[
\diac[0.45]{\draw[line width=0.9pt] (0,0) circle (0.8); \node[font=\tiny] at (0,0) {diagram};
  \foreach \a/\l in {135/p_n,90/p_1,45/p_2,0/p_3} {\draw[prop] (\a:0.8) -- (\a:2.2); \node[vtx] at (\a:0.8) {}; \node at (\a:2.6) {\small$\l$};}
  \draw[densely dotted] (165:1.5) arc[start angle=165,end angle=330,radius=1.5];}
=\mathcal D(p_1,\dots,p_n)\,.
\]
One can expand $\mathcal D(p_1,\dots,p_n)$ in a simultaneous power series in all
the external momenta $p_i$, and the series will converge all the way up to the
momenta $|p_i|\le\Lambda_1$. This convergence is the consequence of the fact that
all the integration momenta in the diagram lie above $\Lambda_1$. Note that
should the integration involve all momenta, as in usual Feynman diagrams, similar
expansion would converge only up to $|p_i|<m_0$, at the best.

This convergence of the expansions of $\mathcal D(p_1,\dots,p_n)$ implies that the
new action $\mathcal A_1$ admits meaningful expansion in terms of local
composites,
\[
\mathcal A_1[\varphi_1]=\int d^dx\,\left[\sum_{k=1}^\infty\frac{u_{2k}}{2k!}\,\varphi_1^{2k}(x)+\sum_{k=0}^\infty\frac{v_{2k}}{2k!}\,\varphi_1^{2k}(x)\,\bigl(\partial\varphi_1(x)\bigr)^2+\text{higher terms}\right]
\]
where the ``higher terms'' involve both the higher powers of $\partial\varphi_1$,
and the terms with the higher derivatives of $\varphi_1$. The most significant
difference from truly local action is that these higher terms include derivatives
of arbitrarily high order. However it is assumed that the full expansion
converges for the fields $\varphi_1$, in which only the Fourier modes with
$|k|<\Lambda_1$ are present. I will call this form ``quasi-local''.

Our conclusion at this point is that by integrating out the fast part
$\tilde\varphi$, the theory described originally by the action $\mathcal A_0$,
with the cutoff $\Lambda_0$, can be equivalently described by new action
$\mathcal A_1$, with the lower cutoff $\Lambda_1$,
\[
\mathcal A_0,\Lambda_0\ \to\ \mathcal A_1,\Lambda_1\,.
\]
The new action however has complicated quasi-local form.

At first glance it may look as a step in wrong direction. But it becomes very
powerful move when combined with another step, which is to make a variable
transformation
\[
\varphi_1(x)=Z^{-1/2}(L)\,\varphi'_0(x/L)
\]
in the functional integral
\[
\frac1{Z_1}\int[D\varphi_1]\,\varphi_1(x_1)\cdots\varphi_1(x_n)\,e^{-\mathcal A_1[\varphi_1]}\,,
\]
where $L$ is the ratio
\[
L=\Lambda_0/\Lambda_1>1\,,
\]
and $Z^{1/2}(L)$ is a factor to be discussed later. This is the dilation with the
factor $1/L$, which returns the cutoff from $\Lambda_1$ back to the original value
$\Lambda_0$. We have
\[
\varphi_1(x)=\int_{|k|<\Lambda_1}\phi(k)\,e^{ikx}=\int_{|k|<\Lambda_1}\phi(k)\,e^{i(Lk)(x/L)}=L^{-d}\int_{|q|<\Lambda_0}\phi(q/L)\,e^{iq(x/L)}\,,
\]
so that
\[
\varphi'_0(x)=L^{-d}\,Z^{1/2}(L)\int_{|k|<\Lambda_0}\phi(k/L)\,e^{ikx}
\]
indeed contains all the wave-vectors with $|k|<\Lambda_0$.

For the action this amounts to transformation
\[
\mathcal A_1[\varphi_1]=\mathcal A'_0[\varphi'_0]\,,
\]
where $\mathcal A'_0$ is obtained from $\mathcal A_1$ simply by re-scaling of the
coordinates, plus the field renormalization. We then have for correlation
functions
\[
\bigl\langle\varphi_0(x_1)\cdots\varphi_0(x_n)\bigr\rangle\Big|_{\mathcal A_0,\Lambda_0}
=Z^{-n/2}(L)\,\bigl\langle\varphi_0(x_1/L)\cdots\varphi_0(x_n/L)\bigr\rangle\Big|_{\mathcal A'_0,\Lambda_0}\,,
\]
i.e.\ the theories described by $\mathcal A_0$ and $\mathcal A'_0$, \emph{with the
same cutoff} $\Lambda_0$ are different only by the scale transformation, with the
scale factor $L$. The transformation
\[
RG(L):\quad\mathcal A_0\to\mathcal A'_0
\]
is the Wilson's RG transformation. It represents the scale transformation in the
space of effective actions.

% Source pp. 173-181 (Lectures13-18.pdf).
\section{More on Wilson's RG}

The last time we have defined the \emph{Wilson's Renormalization Group
Transformation}. This transformation acts on the action functionals
$\mathcal A[\varphi_0]$,
\[
\mathcal A\to\mathcal A'\,,
\]
which represent \emph{cutoff field theories}, i.e.\ field theories equipped with
some cutoff regularization, with cutoff parameter $\Lambda$. The transformation
depends on a parameter $L\ge1$ and a function $Z(L)$, and consists of two steps.

\paragraph{STEP 1:} Integrating out the ``fast'' degrees of freedom.

The field $\varphi_0(x)$ is split into its ``slow'' and ``fast'' parts,
\[
\varphi_0(x)=\varphi_1(x)+\tilde\varphi(x)\,,
\]
where
\begin{equation}\label{eq:slowfast}
\varphi_1(x)=\int_{|k|<\Lambda/L}\phi(k)\,e^{ikx}\,,\qquad
\tilde\varphi(x)=\int_{\Lambda/L<|k|<\Lambda}\phi(k)\,e^{ikx}\,.
\end{equation}
Integrating the part $\tilde\varphi$ out one defines new action
\begin{equation}\label{eq:A1def}
\Const\,e^{-\mathcal A_1[\varphi_1]}=\int[D\tilde\varphi]\,e^{-\mathcal A[\varphi_1+\tilde\varphi]}
\end{equation}
which depends only on the ``slow'' degrees of freedom, with the cutoff parameter
$\Lambda/L$.

\paragraph{STEP 2:} Rescaling of the coordinates $x\to x/L$.

One makes the field variable transformation
\begin{equation}\label{eq:rescale}
\varphi_1(x)=Z^{-\frac12}(L)\,\varphi_0(x/L)
\end{equation}
in the functional integral over $\varphi_1$, defining new integration variable
$\varphi_0$. This transformation returns the cutoff parameter to its original
value $\Lambda$, and yields new action
\begin{equation}\label{eq:Aprime}
\mathcal A'[\varphi_0(x)]=\mathcal A_1\bigl[Z^{-\frac12}(L)\,\varphi_0(x/L)\bigr]\,.
\end{equation}

In principle, the factor $Z^{\frac12}(L)$ in \eqref{eq:rescale}, \eqref{eq:Aprime}
is arbitrary. Field renormalizations $\varphi_0\to Z^{\frac12}\varphi_0$ do not
change physical content of field theory, it can be thought of as some sort of
gauge transformation.\footnote{Much more general class of ``gauge
transformations'' can be applied to field theories, the so called ``field
redefinitions''. One considers generic local change of variables of the form
$\varphi_0(x)\to f\bigl(\varphi_0(x),\partial_\mu\varphi_0(x),\cdots\bigr)$, which
leaves the quasilocal nature of the action \eqref{eq:quasilocal} intact. Here we
will not consider implications of this huge ambiguity. Of course, it is not to be
mixed with the usual gauge transformations in gauge field theories.} When we
characterize field theories by their action functionals, two actions different
only in such renormalization are equivalent,
\[
\mathcal A\bigl[Z^{\frac12}\varphi_0\bigr]\sim\mathcal A[\varphi_0]\,.
\]
In practice this redundancy is usually eliminated by some field normalization
condition which plays role of gauge fixing condition. In this case the factor
$Z^{\frac12}(L)$ must be adjusted so that the transformation
$\mathcal A\to\mathcal A'$ does preserve that condition.

It follows from the above definition that if $\mathcal A'$ is obtained from
$\mathcal A$ by the Wilson's RG transformation, the corresponding correlation
functions are related as,
\[
\text{f.t.}\,\bigl\langle\varphi_0(x_1)\cdots\varphi_0(x_n)\bigr\rangle_{\mathcal A'}
=\text{f.t.}\,Z^{\frac n2}(L)\,\bigl\langle\varphi_0(x_1L)\cdots\varphi_0(x_nL)\bigr\rangle_{\mathcal A}
\]
where it is understood that the Fourier transforms of the correlation functions
are taken and the comparison is made at low values of the momenta
$|p_i|\ll\Lambda$. In other words, the theory described by the action
$\mathcal A'$ differs from that described by $\mathcal A$ by a scale
transformation $x\to x/L$ only. The Wilson's RG shows what it takes to make such
scale transformation in field theory. Just rescaling the coordinates in the
action, as is done in \eqref{eq:rescale}, is not sufficient. The ``naive''
scaling transformation, represented by the step 2 above, must be supplemented by
the integration over the fast modes in the step 1. Roughly speaking, the dilation
$x\to x/L$ increases the spatial density of the field degrees of freedom, and one
has to ``dilute'' them back to the original density by performing the step 1. In
this sense the step 1 can be thought of as representing the transformation of the
measure in the functional integral.

One remark is in order here. In general, implementation of the step 1 of the
Wilson's RG (elimination of the fast modes) depends on particular form of the
cutoff regularization implied. The definition \eqref{eq:slowfast} corresponds to
the ``sharp momentum cutoff''. We know that other cutoff regularizations are
possible. For example, in the lattice regularization, where the field degrees of
freedom are associated with the lattice cells, the procedure described above is
hardly natural. Instead, one can consider a sublattice whose cells consist of
$L^d$ elementary cells ($L$ must be integer in this case). The ``slow'' field
$\varphi_1$ is then defined as the original field averaged over these greater
cells, and the step 1 amounts to integrating over $\varphi_0$ with $\varphi_1$
fixed. I am mentioning this here to stress that the two step procedure of the
Wilson's RG can be considered as a general idea rather than an exact
prescription, the idea whose implementation depends on the problem in hand.
However, in this discussion we will follow the particular implementation
described in \eqref{eq:slowfast}.

We have seen that even if the original action $\mathcal A$ has some simple form,
say $\varphi^4$, the result of the Wilson's RG transformation is usually a
complicated ``quasilocal'' action which contains all possible terms with
arbitrarily high powers of the field and with arbitrarily high powers and orders
of the derivatives. On the other hand, if $\mathcal A$ is generic quasilocal
action, one can expect that the action $\mathcal A'$ will again be of the
quasilocal form. Under this assumption, the Wilson's RG transformation acts in
the space of all quasilocal action functionals. In what follows I will use the
notation $\Sigma$ for this infinite-dimensional space, so that
\[
RG:\ \Sigma\to\Sigma\,.
\]
Also, I will denote the RG transformation as
\[
\mathcal A'=RG_l\{\mathcal A\}\,,
\]
using more convenient parameter
\[
l=\log L\ge0
\]
instead of $L$. It is not difficult to understand that these transformations
obey the composition law
\begin{equation}\label{eq:semigroup}
RG_{l+l'}\{\mathcal A\}=RG_l\{RG_{l'}\{\mathcal A\}\}\,,\qquad RG_0\{\mathcal A\}=\mathcal A\,.
\end{equation}
These properties define one-parameter \emph{semigroup}. It is not a group
because a group contains inverse elements, and inverse to $RG_l$ in general
cannot be defined. The reason is simple but important: the step 1 of the RG
transformation leads to a loss of information about the interactions at the
cutoff scales; this step cannot be ``undone''.

In general exact evaluation of the RG transformation for arbitrary $l$ is at
least as difficult as exact evaluation of the full functional integral, i.e.\
usually hopeless. If the interaction is not small no reliable approximations are
even available. One can hope however to find good approximations for $RG_l$ with
$l$ small enough. Then according to the composition law \eqref{eq:semigroup} the
transformation $RG_{Nl}$ can be obtained as $N$th iteration of $RG_l$. While it
is not clear how to evaluate the functional integral at once, maybe it is
possible to do it step by step, integrating out the shortest scale degrees of
freedom, and then repeating this again and again.

It is usually most convenient to consider the infinitesimal RG transformation
\[
RG_{\delta l}\{\mathcal A\}=\mathcal A+B\{\mathcal A\}\,\delta l+O(\delta l^2)\,,
\qquad\text{where}\qquad
B\{\mathcal A\}=\frac{\dd}{\dd l}\,RG_l\{\mathcal A\}\Big|_{l=0}\,.
\]
It follows from \eqref{eq:semigroup} that the one-parameter family of actions
\[
\mathcal A_l=RG_l\{\mathcal A_0\}\,,
\]
where $\mathcal A_0$ is an arbitrary ``initial'' action, satisfies the differential
equation
\begin{equation}\label{eq:WilsonFlow}
\frac{\dd}{\dd l}\,\mathcal A_l=B\{\mathcal A_l\}\,,
\end{equation}
which is usually referred to as \emph{Wilson's RG flow equation}. It generates the
scale transformations in the space of quasilocal actions $\Sigma$, namely, two
actions $\mathcal A_{l_1}$ and $\mathcal A_{l_2}$ belonging to the same integral
curve $\mathcal A_l$ (sometimes called the \emph{RG trajectory}) describe the
theories which differ by the scale transform $x\to x\exp(l_1-l_2)$. Thus the
problem splits into two parts: (i) Evaluation of the RG operator
$B\{\mathcal A\}$ and (ii) Integration of the RG flow equation
\eqref{eq:WilsonFlow}.

The RG flow equation \eqref{eq:WilsonFlow} is written above in very abstract
form. To make it substantial one has to introduce some coordinates on $\Sigma$.
Generic quasilocal action can be written as (I assume $\varphi_0\to-\varphi_0$
symmetry)
\begin{multline}\label{eq:quasilocal}
\mathcal A[\varphi_0]=\int \dd^dx\,\Bigl[\sum_{k=1}^\infty\frac{u_{2k}}{2k!}\,\varphi_0^{2k}(x)+\sum_{k=0}^\infty\frac{v_{2k}}{2k!}\,\varphi_0^{2k}\,(\partial\varphi_0)^2(x)\\
+\sum_{k=0}^\infty\frac{w_{2k+1}}{2k!}\,\varphi_0^{2k+1}\,(\partial\varphi_0)^2\,\partial^2\varphi_0(x)+\cdots\Bigr]\,.
\end{multline}
The coefficients (``coupling constants'') $u_{2k},v_{2k},w_{2k+1},\dots$ can be
taken as such coordinates. In general there are infinitely many couplings and
therefore the space $\Sigma$ is infinite-dimensional. One of these parameters
can be eliminated by appropriate renormalization of $\varphi_0$, for instance one
can use this freedom to set
\[
v_0=1
\]
(the ``gauge fixing'' discussed above).

In terms of these coordinates the RG flow equation \eqref{eq:WilsonFlow} can be
written as
\begin{equation}\label{eq:Bfunctions}
\begin{aligned}
\frac{\dd}{\dd l}\,u_{2k}(l)&=B_{u_{2k}}\bigl(\{u_{2k}(l)\},\{v_{2k}(l)\},\{w_{2k+1}(l)\},\dots\bigr)\,;\\
\frac{\dd}{\dd l}\,v_{2k}(l)&=B_{v_{2k}}\bigl(\{u_{2k}(l)\},\{v_{2k}(l)\},\{w_{2k+1}(l)\},\dots\bigr)\,;\\
&\ \ \text{etc}
\end{aligned}
\end{equation}
where $u_{2k}(l),v_{2k}(l),\dots$ are the coupling constants corresponding to the
action $\mathcal A_l$, and $B$'s are some functions of all these infinitely many
coordinates. Of course it is impossible to calculate exactly all these functions.
Neither can one hope to integrate exactly all the RG flow equations. In practical
calculations some approximations are necessary. Most importantly one must find
some finite-dimensional approximation to the infinite-dimensional space $\Sigma$.
But before doing any calculations it is useful to think about what are the
properties of the RG flow we must be interested in.

\subsection*{Topological properties of RG flow\protect\footnote{The discussion in
this section follows closely that in K.~Wilson, J.~Kogut, ``The RG and epsilon
expansion'', Phys.\ Rep.\ \textbf{12C} (1974).}}

Some general topological properties of RG flow can be discussed in an abstract
form without reference to specific form of the RG transformation. Some general
assumptions will be made concerning the actions generated by the RG
transformations. Most importantly, we will make an assumption already mentioned
above, namely, that the RG transformation $RG_l$ applied to a quasilocal action
$\mathcal A$ produces an action $\mathcal A_l$ which is quasilocal. In other
words, our main assumption is that $RG:\Sigma\to\Sigma$. This is not a trivial
assumption. I remind you that our term ``quasilocal action'' means not only an
action which admits formal derivative expansion \eqref{eq:quasilocal}, but also
that this expansion converges for all fields $\varphi_0$ which contain only
Fourier components with $|k|\lesssim\Lambda$. Roughly speaking quasilocal action
can contain nonlocal interactions of the range $\sim\Lambda^{-1}$ but not much
larger. There is no proof to this statement, but it holds for many known
implementations of the RG procedure, so we accept it.

Under this assumption, if $\mathcal A_0\in\Sigma$, the associated RG trajectory
$\mathcal A_l$, i.e.\ the solution of the flow equation \eqref{eq:WilsonFlow}
with the initial condition $\mathcal A_{l=0}=\mathcal A_0$, lies inside $\Sigma$
for all $l>0$ (figure~\ref{fig:trajectory}).
\begin{figure}[ht]
\centering\figfile{trajectory}
\caption{An RG trajectory $\mathcal A_l$ starting at $\mathcal A_0$ stays inside
the space $\Sigma$ of quasilocal actions.}
\label{fig:trajectory}
\end{figure}

The trajectory $\mathcal A_l$, as a function of the parameter $l$, gives a set of
actions which all describe essentially the same physical system, the only
difference being in the coordinate rescalings. Namely one must make an $l$
dependent change of scale when relating $\mathcal A_l$ to physics.

It is useful to discuss in terms of physical mass $m$, or, as it is more common in
statistical mechanics, in terms of the ``correlation length'' $R_c=m^{-1}$. This
is defined as follows. Consider the two-point correlation function. Its typical
large-distance behavior in Euclidean space is (recall for instance your solution
of the behavior of the Klein-Gordon propagator $D(x)$ at $|x|\to\infty$)
\[
\bigl\langle\varphi_0(x)\varphi_0(0)\bigr\rangle\simeq e^{-\frac{|x|}{R_c}}\qquad\text{as}\qquad|x|\to\infty\,.
\]
Let $R_c$ be the correlation length associated with our system described by the
action $\mathcal A_0$, with the cutoff $\Lambda$. Of course $\mathcal A_l$
describes the same system, but in the units in which the cutoff associated with
$\mathcal A_l$ is still $\Lambda$, the correlation length will be
\[
R_c(l)=e^{-l}\,R_c\,.
\]
One can define various subspaces, or ``surfaces'', in $\Sigma$. The elements of
$\Sigma$ can be separated according to their correlation length.
Figure~\ref{fig:Rcsurfaces} shows schematically the surfaces of constant $R_c$.
\begin{figure}[ht]
\centering\figfile{Rcsurfaces}
\caption{Surfaces of constant correlation length in $\Sigma$. The thick curve is
the critical surface $R_c=\infty$.}
\label{fig:Rcsurfaces}
\end{figure}

Of special interest is the ``critical surface'' with $R_c=\infty$, or $m=0$. We
denote it $\Sigma_c$.

Important subspace is the subspace $\Sigma(l)$ generated by applying $RG_l$ to
$\Sigma$ itself,
\[
\Sigma(l)=RG_l\{\Sigma\}\,,
\]
see figure~\ref{fig:nested}.
\begin{figure}[ht]
\centering\figfile{nested}
\caption{The nested images $\Sigma\supset\Sigma(l)\supset\Sigma(l+l')\supset\dots$
of $\Sigma$ under the RG transformations, shrinking to $\Sigma(\infty)$ (thick
curve).}
\label{fig:nested}
\end{figure}
As we have assumed that $RG_l$ acts in $\Sigma$,
\[
\Sigma(l)\subset\Sigma\,.
\]
Obviously the subspaces $\Sigma(l)$ are nested, i.e.
\[
\Sigma(l+l')\subset\Sigma(l)\qquad\text{for any}\qquad l,l'>0
\]
(this follows from $\Sigma(l+l')=RG_l\{\Sigma(l')\}$).

Strictly speaking the space $\Sigma$ is infinite-dimensional. However we assume
that in many important respects it can be treated as finite dimensional one. This
is equivalent to the assumption that it is possible to build a system of
approximations to the RG transformation in which $\Sigma$ is replaced by spaces of
large but finite dimensionalities, and that such approximations converge to exact
RG transformation when this dimensionality goes to infinity.\footnote{Of course
these statements are formulated in a way which is less than satisfactory from
mathematical point of view. Mathematical rigor can be achieved only in some
simplified models of the RG transformation.} If dimensionality of $\Sigma$ is
understood in this sense, the spaces $\Sigma(l)$ for finite $l$ have the same
dimensionality as $\Sigma$ itself. However, the space $\Sigma(\infty)$, obtained
as the limit $l\to\infty$ of $\Sigma(l)$, can be much smaller. The basic idea of
the Wilson's RG approach is in the hope that $\Sigma(\infty)$ has far fewer
dimensions than $\Sigma$. The reason why $\Sigma(\infty)$ is expected to be much
smaller than $\Sigma$ is the following: For large $l$ the action $\mathcal A_l$
describes the physical system with actual interactions of much shorter range than
the range of interactions which appear in the action $\mathcal A_l$ itself.
Remember, $\mathcal A_l$ is physically equivalent to $\mathcal A_0$, with the
scales changed by the factor $e^{-l}$. It is unlikely that arbitrary action in
$\Sigma$ has this property. Thus $\Sigma(\infty)$ is expected to be only a small
subset in $\Sigma$.

The differential equation \eqref{eq:WilsonFlow}, i.e.
\[
\frac{\dd}{\dd l}\,\mathcal A=B\{\mathcal A\}
\]
can be integrated backward in $l$ as well as forward. Suppose $\mathcal A$ is in
$\Sigma$ but not in $\Sigma(l)$. Then $\mathcal A_{-l}$, obtained by integrating
backwards, can not lie in $\Sigma$\footnote{If $\mathcal A_{-l}\in\Sigma$ then
$\mathcal A=RG_l\{\mathcal A_{-l}\}\in\Sigma(l)$.} (figure~\ref{fig:backward}).
\begin{figure}[ht]
\centering\figfile{backward}
\caption{An action $\mathcal A\in\Sigma$ outside $\Sigma(l)$: integrated backwards
by $l$, its trajectory leaves $\Sigma$.}
\label{fig:backward}
\end{figure}

The reason $\mathcal A_{-l}$ is not in $\Sigma$ is likely to be one of the
following: $\mathcal A_{-l}$ can fail to be quasilocal, i.e.\ $\mathcal A_{-l}$
may contain long-range interactions. Or $\mathcal A_{-l}$ can have other
pathological properties, like some coefficients which control convergence of the
functional integral can become such that the functional integral with
$\mathcal A_{-l}$ no longer exists (like $\lambda$ in $\varphi^4$ theory can
become negative), so that the action $\mathcal A_{-l}$ does not define meaningful
theory. Anyhow, for large $l$ the actions $\mathcal A_{-l}$ are typically outside
$\Sigma$.

Important exceptions from this general rule are the actions belonging to the
subspace $\Sigma(\infty)$. If one takes
\[
\mathcal A\in\Sigma(\infty)\,,
\]
and obtains $\mathcal A_{-l}$ by integrating backwards,
\[
\mathcal A_{-l}\in\Sigma(\infty)\,.
\]
This is true, for if $\mathcal A_{-l}$ is in $\Sigma$ but not in
$\Sigma(\infty)$ we have
\[
RG_l\{\mathcal A_{-l}\}=\mathcal A\,,
\]
so that $\mathcal A$ can be reached in finite ``time'' $l$ from some point in
$\Sigma$ outside $\Sigma(\infty)$, but this contradicts the definition of
$\Sigma(\infty)$. In short, starting in $\Sigma(\infty)$ one can integrate back
and forth, still remaining within $\Sigma(\infty)$.

For this reason the subspace $\Sigma(\infty)$ is of central importance in local
Quantum Field Theory. The problem of QFT is to obtain finite local theory in the
limit $\Lambda\to\infty$, with appropriately adjusted ``bare action''
$\mathcal A(\Lambda)$. ``Integrating backward'' is essentially considering the
actions which lead to the same physics as the action $\mathcal A_0$ one starts
with, but equipped with larger and larger cutoff $\Lambda$, measured in the units
of $m$. It is then clear that finding renormalized local theories is equivalent
to finding actions $\mathcal A$ from which one can ``integrate backward'' without
limits. For this reason it is natural to identify $\Sigma(\infty)$ with the space
of local QFT. We will come back to this point many times over.

% Source pp. 182-200 (Lectures13-18.pdf).
\section{Fixed Points of RG}

At generic point in $\Sigma$ the pattern of the RG flow is rather boring: it just
flows. Something interesting is expected to happen near singular points of the
flow. Of these the simplest and probably the most important ones are so called
\emph{Fixed Points}.

A fixed point is an action $\mathcal A_*$ which satisfies the equation
\[
B\{\mathcal A_*\}=0\,.
\]
Associated trajectory is trivial,
\[
\mathcal A_l=\mathcal A_*\qquad\text{for all}\quad l\,.
\]
As the correlation length measured in the units $\Lambda$ is
\[
R_c(l)=e^{-l}\,R_c
\]
for $\mathcal A_l$, this means that for $\mathcal A_*$ we must have
\[
\text{either}\quad R_c=0\qquad\text{or}\quad R_c=\infty\,.
\]
By the definition, $\mathcal A_*$ with $R_c=\infty$ belongs to the ``critical
surface'',
\[
\mathcal A_*\in\Sigma_c\qquad\text{if}\qquad R_c=\infty\,.
\]
We call the fixed points $\mathcal A_*$ with $R_c=\infty$ the ``critical'' fixed
points.

To illustrate the topological ideas related to the Wilson's RG, I will describe
the simplest possible form of topology. There are many ways in which topology of
the RG flow can be more complicated but complications can be understood once the
simplest picture is grasped. In this analysis I will assume continuity of the RG
flow, i.e.\ the functions $B$ in \eqref{eq:Bfunctions} are continuous functions of
their arguments. Again, general proof of this statement is not available, but
various examples support this assumption.

Assume there are only two fixed points, one with $R_c=0$ and one with
$R_c=\infty$. Denote them $P_0$ and $P_\infty$ (figure~\ref{fig:twofixed}, left).
\begin{figure}[ht]
\centering\figfile{twofixed}
\caption{Left: surfaces of constant $R_c$ with the two fixed points $P_0$
($R_c=0$) and $P_\infty$ on the critical surface ($R_c=\infty$). Right: the
trajectory $U$ leaving $\Sigma_c$ at $P_\infty$, and a generic trajectory $G$
starting close to $\Sigma_c$.}
\label{fig:twofixed}
\end{figure}

As $R_c(l)=e^{-l}R_c$, any trajectory which starts at finite $R_c$ will approach
$R_c=0$ and then $P_0$. On the other hand trajectories which start on $\Sigma_c$
will approach $P_\infty$. The only point where some trajectory can leave
$\Sigma_c$ is $P_\infty$.\footnote{Suppose a trajectory ``leaves'' $\Sigma_c$ at
some $P_c\in\Sigma_c$. Obviously $P_c=\lim_{l\to\infty}RG_{-l}\{P\}$ for any $P$ on
such trajectory.} So there must be a trajectory (let's call it $U$) going from
$P_\infty$ to $P_0$ (figure~\ref{fig:twofixed}, right).

Consider also a trajectory $G$ which starts close to $\Sigma_c$. This one first
approaches $P_\infty$ while staying close to $\Sigma_c$, then spends very long
``time'' $l$ in the vicinity of $P_\infty$ (by continuity $B\{\mathcal A\}$ is
small when $\mathcal A$ is close to $P_\infty$) and then departs along some way
close to $U$ (recall that $U$ is the only trajectory that leaves $\Sigma_c$),
finally going towards $P_0$.

Now let us look for the space $\Sigma(\infty)$ in this picture. It is clear that
$\Sigma(\infty)$ is $U$, with the points $P_\infty$ and $P_0$ included. Starting
with any action $\mathcal A$ on the curve $U$ one can go back along the trajectory
arbitrary amount of ``time'' $l$ without leaving the space $\Sigma$ (actually,
without leaving $U$). The reason is that $R_c=\infty$ on $U$ only at the point
$P_\infty$. To reach $P_\infty$ from any point on $U$ with finite $R_c$ takes
infinite ``time'' $l$ (when moving backwards). Starting from other point in
$\Sigma$ and integrating backwards, one generally leaves $\Sigma$ in sufficient
amount of ``time''.

We see from this example the role played by critical fixed points in the
formation of the space $\Sigma(\infty)$. Therefore finding critical fixed points
is of central importance in the RG approach.

\subsection*{Critical Fixed Points}

Suppose $\mathcal A_*$ is a critical fixed point, i.e.
\[
B\{\mathcal A_*\}=0\,.
\]
and associated $R_c=\infty$. Consider vicinity of $\mathcal A_*$. Let
\[
\mathcal A=\mathcal A_*+\delta\mathcal A\,,
\]
where $\delta\mathcal A$ is infinitesimal. The variation $\delta\mathcal A$ must be
thought of as a vector in the tangent space $T\Sigma$ taken at the point
$\mathcal A_*$ (figure~\ref{fig:tangent}),
\[
\delta\mathcal A\in T\Sigma|_{\mathcal A_*}\,.
\]
\begin{figure}[ht]
\centering\figfile{tangent}
\caption{A variation $\delta\mathcal A$ of the fixed-point action is a vector in
the tangent space $T\Sigma|_{\mathcal A_*}$.}
\label{fig:tangent}
\end{figure}
One can linearize the RG transformation near $\mathcal A_*$,
\[
B\{\mathcal A_*+\delta\mathcal A\}=K\{\delta\mathcal A\}+O(\delta\mathcal A^2)\,,
\]
where $K$ is a linear operator acting in this tangent space,
\[
K:\ T\Sigma|_{\mathcal A_*}\to T\Sigma|_{\mathcal A_*}\,.
\]
As it is conventional in the theory of linear operators, I will write
\begin{equation}\label{eq:Klin}
K\{\delta\mathcal A\}\equiv K\,\delta\mathcal A\,.
\end{equation}
The RG flow in this close vicinity of $\mathcal A_*$ is then described by linear
differential equation
\[
\frac d{dl}\,(\delta\mathcal A)=K\,\delta\mathcal A\,.
\]
This equation is easily solved provided one can find eigenvectors $\Psi_\alpha$
and associated eigenvalues $\kappa_\alpha$ of the operator $K$,
\[
K\,\Psi_\alpha=\kappa_\alpha\,\Psi_\alpha\,.
\]
The solution takes the form
\begin{equation}\label{eq:linsol}
\delta\mathcal A(l)=\sum_\alpha C_\alpha\,e^{\kappa_\alpha l}\,\Psi_\alpha\,,
\end{equation}
where the constants $C_\alpha$ are determined by the initial vector
$\delta\mathcal A(0)=\delta\mathcal A_0$.

It is easy to understand that the pattern of the RG flow in the vicinity of
$\mathcal A_*$ is directly related to the eigenvalues $\kappa_\alpha$. To
illustrate this let us assume that there are only two eigenvalues (this
corresponds to a two-dimensional approximation to $\Sigma$), both real, one
positive and one negative,
\[
\kappa_+>0\,,\qquad\kappa_-<0\,.
\]
Figure~\ref{fig:saddle} shows the pattern of corresponding RG flow in the vicinity
of the fixed point. There is a natural coordinate system on
$T\Sigma|_{\mathcal A_*}$ associated with the eigenvectors $\Psi_+$ and $\Psi_-$.
\begin{figure}[ht]
\centering\figfile{saddle}
\caption{RG flow near a critical fixed point with one positive and one negative
eigenvalue. The $\Psi_-$ axis lies in the critical surface $\Sigma_c$ and flows
into the fixed point; the $\Psi_+$ axis is the trajectory $U$ leaving it. A
generic trajectory $G$ passes near the fixed point and leaves along $U$.}
\label{fig:saddle}
\end{figure}

All trajectories starting at the $\Psi_-$ axis never leave it; they flow towards
the fixed point $\mathcal A_*$. As the correlation length can only decrease with
$l$, the entire axis $\Psi_-$ must belong to $\Sigma_c$. The trajectories
starting at the $\Psi_+$ also remain at this axis, but they flow away from the
fixed point. For sufficiently large $l$ these trajectories leave the domain of
validity of the linear approximation \eqref{eq:Klin}, and \eqref{eq:linsol} says
nothing about their ultimate destination. The same is true for any generic
trajectory $G$ starting at neither axis: at sufficiently large $l$ it leaves the
vicinity of the fixed point and its fate cannot be predicted within the linear
approximation. However there is an important difference between the generic
trajectories like $G$ and any trajectory $U$ that lies on the axis $\Psi_+$.
Starting at any point on the axis $\Psi_+$ one can ``integrate backward'' for
indefinite negative $l$ without leaving the vicinity of the fixed point
$\mathcal A_*$; in fact corresponding $\mathcal A_{-l}$ approaches $\mathcal A_*$
as $l\to\infty$. This is in contrast with the negative $l$ behavior of generic
trajectories; associated $\mathcal A_{-l}$ leave the vicinity of $\mathcal A_*$
for sufficiently large $l$, and the question if they remain in $\Sigma$ for
$l\to\infty$ cannot be answered within the linear approximation. On the other
hand the linear approximation is good enough to recognize in this example at
least one-dimensional component of the subspace $\Sigma(\infty)$ associated with
the fixed point $\mathcal A_*$; the axis $\Psi_+$ is the linear approximation to
this component. Precise statement is
\[
\Psi_+\subset T\Sigma(\infty)|_{\mathcal A_*}\,.
\]

It is not difficult to generalize the conclusions obtained in the example above.
It is useful to define two special subspaces of the tangent space
$T\Sigma|_{\mathcal A_*}$. Let $\mathcal F_+$ be the linear envelope of the
eigenvectors of $K$ corresponding to positive eigenvalues $\kappa$; with some
abuse of notations we can write
\[
\mathcal F_+=\oplus_{\alpha:\kappa_\alpha>0}\,\Psi_\alpha\,.
\]
Similarly, let $\mathcal F_-$ be the sum of eigenspaces corresponding to negative
$\kappa_\alpha$, i.e.
\[
\mathcal F_-=\oplus_{\alpha:\kappa_\alpha<0}\,\Psi_\alpha\,.
\]
This separation leaves unaccounted very interesting possibility of having zero
eigenvalues, later we will discuss it separately. Also, it is assumed that all
eigenvalues of $K$ are real. At this point this assumption looks arbitrary. We
will see later that having complex eigenvalues contradicts quantum mechanical
interpretation of field theory.

Simple analysis of \eqref{eq:linsol} shows that
\begin{itemize}
\item The fixed point $\mathcal A_*$ belongs to a critical surface $\Sigma_c$ whose
dimensionality coincides with the dimensionality of $\mathcal F_-$, and
\[
\mathcal F_-=T\Sigma_c|_{\mathcal A_*}\,.
\]
\item The fixed point $\mathcal A_*$ also belongs to some component of
$\Sigma(\infty)$; the dimensionality of this component is equal or greater than
the dimensionality of $\mathcal F_+$, and
\[
\mathcal F_+\subset T\Sigma(\infty)|_{\mathcal A_*}\,.
\]
\end{itemize}

Let us define the subspace $U(\mathcal A_*)\subset\Sigma(\infty)$ by the condition
that all points of $U(\mathcal A_*)$ lie on RG trajectories which, when
``integrated backward'', converge to the fixed point $\mathcal A_*$, i.e.
\[
\mathcal A_0\in U(\mathcal A_*)\ \Rightarrow\ \mathcal A_l\to\mathcal A_*\quad\text{as}\quad l\to-\infty
\]
The subspace $U(\mathcal A_*)$ is called the \emph{unstable manifold of the fixed
point $\mathcal A_*$}. Obviously $\mathcal A_*\in U(\mathcal A_*)$. If we ignore
the possibility of zero eigenvalues then the tangent space to $U(\mathcal A_*)$
taken at the fixed point $\mathcal A_*$ itself coincides with the space
$\mathcal F_+$, i.e.
\[
\mathcal F_+=TU(\mathcal A_*)|_{\mathcal A_*}\,.
\]

The reason why $\Sigma(\infty)$ can have dimensionality greater than that of the
unstable manifold $U(\mathcal A_*)$ is illustrated by the following simple
example. Suppose we have two fixed points, the $\mathcal A_*$ and another one,
$\mathcal A'_*$, such that $\mathcal A_*\in U(\mathcal A'_*)$
(figure~\ref{fig:twofp}).
\begin{figure}[ht]
\centering\figfile{twofp}
\caption{Two fixed points with $\mathcal A_*\in U(\mathcal A'_*)$. One
trajectory of $U(\mathcal A'_*)$ flows into $\mathcal A_*$.}
\label{fig:twofp}
\end{figure}
Then there must be at least one RG trajectory in $U(\mathcal A'_*)$ which flows to
$\mathcal A_*$ (this is by the definition of unstable manifold $U$). This
trajectory has $R_c=\infty$ and therefore it simultaneously belongs to
$\Sigma(\infty)$ and $\Sigma_c$. It follows that the unstable manifold
$U(\mathcal A_*)$ has at least one dimension less than $U(\mathcal A'_*)$.

Understanding the global topological properties of the RG flow, like the one
illustrated in the above example, is beyond the capabilities of linear analysis.
However, as we have seen, it is within the power of linear analysis to prove that
each fixed point $\mathcal A_*$ generates a manifold $U(\mathcal A_*)$ which
belongs to $\Sigma(\infty)$. The dimensionality $n$ of this manifold equals the
number of positive eigenvalues of the operator $K$. In other words, the fixed
point $\mathcal A_*$ generates $n$-parameter family of finite local QFT.

The above discussion was somewhat abstract. It is useful to understand the
meaning of the abstract notions of the above linear analysis, like tangent
spaces, in terms of local field degrees of freedom.

Recall that generic quasilocal action $\mathcal A$ can be written as integral
\[
\mathcal A=\int d^dx\,\Lag(x)\,,
\]
where
\[
\Lag(x)=\Lag\bigl(\varphi_0(x),\partial_\mu\varphi_0(x),\partial_\mu\partial_\nu\varphi_0(x),\dots\bigr)
=\sum_{k=1}^\infty\frac{u_{2k}}{2k!}\,\varphi_0^{2k}(x)+\sum_{k=0}^\infty\frac{v_{2k}}{2k!}\,\varphi_0^{2k}(x)\,(\partial\varphi_0)^2(x)+\cdots
\]
is a scalar function of the $\varphi_0(x)$ and all its derivatives. The variation
$\delta\mathcal A$ has similar form,
\[
\delta\mathcal A=\int d^dx\,\delta\Lag(x)\,,\qquad
\delta\Lag(x)=\sum_{k=1}^\infty\frac{\delta u_{2k}}{2k!}\,\varphi_0^{2k}(x)+\sum_{k=0}^\infty\frac{\delta v_{2k}}{2k!}\,\varphi_0^{2k}(x)\,(\partial\varphi_0)^2(x)+\cdots\,.
\]
In other words $\delta\mathcal A$ is an integral of (infinite) linear combination
of local composite fields; we can write it as
\[
\delta\mathcal A=\int d^dx\,\sum_\alpha\delta\lambda_\alpha\,\mathcal O_\alpha(x)\,,
\]
where
$\{\mathcal O_\alpha\}=\{\varphi_0^{2n},\varphi_0^{2n}(\partial\varphi_0)^2,\dots\}$,
and $\{\delta\lambda_\alpha\}$ are variations of corresponding coupling
constants. Thus the tangent space $T\Sigma|_{\mathcal A_*}$ can be identified
with the space of scalar composite fields spanned by $\{\mathcal O_\alpha\}$. I
will use notation $\mathcal F$ for this space of local fields; we have
\[
T\Sigma|_{\mathcal A_*}=\mathcal F\,.
\]

How the RG transformation acts on these composite fields? Consider a correlation
function involving the field $\mathcal O_\alpha(x)$, say
\[
\bigl\langle\mathcal O_\alpha(x)\cdots\bigr\rangle=\frac1Z\int[D\varphi_0]\,\bigl[\mathcal O_\alpha(x)\,\varphi_0(x_1)\cdots\varphi_0(x_n)\bigr]\,e^{-\mathcal A[\varphi_0]}\,,
\]
To perform the step 1 of the Wilson's RG transformation we must split the field
$\varphi_0$ into its ``slow'' and ``fast'' parts,
\[
\varphi_0(x)=\varphi_1(x)+\tilde\varphi(x)\,,
\]
and then integrate the ``fast'' part $\tilde\varphi$ out.

Recall that $\mathcal O_\alpha(x)$ in general is a polynomial function of
$\varphi_0$ and its derivatives at the point $x$,
\[
\mathcal O_\alpha(x)=\mathcal O_\alpha\bigl(\varphi_0(x),\partial_\mu\varphi_0(x),\dots\bigr)
\]
and therefore
\[
\mathcal O_\alpha\bigl(\varphi_1(x)+\tilde\varphi(x),\ \text{derivatives}\bigr)=\sum_\beta Y_\alpha{}^\beta\bigl(\tilde\varphi(x),\ \text{derivatives}\bigr)\,\mathcal O_\beta\bigl(\varphi_1(x),\ \text{derivatives}\bigr)\,,
\]
where $Y_\alpha{}^\beta$ are polynomials of the ``fast'' field $\tilde\varphi(x)$ and
its derivatives. Example is the field $\varphi_0^n$,
\[
\bigl(\varphi_1(x)+\tilde\varphi(x)\bigr)^n=\sum_{k=0}^nY_n{}^k\bigl(\tilde\varphi(x)\bigr)\,\varphi_1^{n-k}(x)\,,\qquad Y_n{}^k(\tilde\varphi)=C_n^k\,\tilde\varphi^k\,.
\]
Assuming again that we are interested in the case when all external momenta
$p_1,\dots,p_n\ll\Lambda$, one can write the result of $\tilde\varphi$ integration
as
\begin{multline}\label{eq:Ostep1}
\frac1Z\int[D\tilde\varphi]\,\bigl[\mathcal O_\alpha\bigl(\varphi_1(x)+\tilde\varphi(x),\ \text{derivatives}\bigr)\,\varphi_0(x_1)\cdots\varphi_0(x_n)\bigr]\,e^{-\mathcal A[\varphi_1+\tilde\varphi]}\\
=\frac1{Z_1}\Bigl[\sum_\beta\bar Y_\alpha{}^\beta\bigl(L|[\varphi_1]\bigr)\,\mathcal O_\beta\bigl(\varphi_1(x),\ \text{derivatives}\bigr)\,\varphi_1(x_1)\cdots\varphi_1(x_n)\Bigr]\,e^{-\mathcal A_1[\varphi_1]}\,,
\end{multline}
where $\mathcal A_1$ is the action appearing in the step 1 of the Wilson's RG
transformation with the parameter $L$, and
\[
\bar Y_\alpha{}^\beta\bigl(L|[\varphi_1]\bigr)=\bigl\langle Y_\alpha{}^\beta(\tilde\varphi,\ \text{derivatives})\bigr\rangle_{\tilde\varphi,c}
\]
are the expectation values of $Y_\alpha{}^\beta$ over the ensemble of the fast mode
fluctuations taken at fixed $\varphi_1$; the subscript $c$ signifies that all
diagrams disconnected from the insertion point $x$ must be disregarded (they are
already accounted for by the factor $\exp(-\mathcal A_1)$). Although these
expectation values depend on the slow modes $\varphi_1$, this dependence is
``quasilocal'', i.e.\ they can be written as series in powers of $\varphi_1(x)$
and its derivatives taken at the point $x$. Therefore the result of
$\tilde\varphi$ integration has the form
\[
\frac1{Z_1}\sum_\beta y_\alpha^\beta(L)\,\mathcal O_\beta\bigl(\varphi_1(x),\ \text{derivatives}\bigr)\,\varphi_1(x_1)\cdots\varphi_1(x_n)\,e^{-\mathcal A_1[\varphi_1]}\,,
\]
where in general the sum involves infinitely many terms, and $y_\alpha^\beta(L)$
are complicated sums of Feynman diagrams. Thus, in the step 1 of RG procedure the
insertion of the composite field $\mathcal O_\alpha(x)$ gets replaced by certain
linear combinations of slow-mode composite fields, with $L$-dependent
coefficients.

This structure holds true after we perform the step 2, which amounts to the change
of variables
\[
\varphi_1(x)=Z^{-\frac12}(L)\,\varphi_0(x/L)\,.
\]
As $\mathcal O_\beta$ in \eqref{eq:Ostep1} are polynomials of $\varphi_1$ and its
derivatives taken at $x$, we have
\begin{multline*}
\sum_\beta y_\alpha^\beta(L)\,\mathcal O_\beta\left(Z^{-\frac12}(L)\,\varphi_0(x/L),\ Z^{-\frac12}(L)\,\frac\partial{\partial x^\mu}\varphi_0(x/L),\dots\right)\\
=\sum_\beta z_\alpha{}^\beta(L)\,\mathcal O_\beta\left(\varphi_0(x/L),\ \frac\partial{\partial(x^\mu/L)}\varphi_0(x/L),\dots\right)\,,
\end{multline*}
i.e.\ under the full RG transformation
\begin{equation}\label{eq:RGfields}
\bigl\langle\mathcal O_\alpha(x)\cdots\bigr\rangle\big|_{\mathcal A}=\Bigl\langle\sum_\beta z_\alpha{}^\beta(L)\,\mathcal O_\beta(x/L)\cdots\Bigr\rangle\Big|_{\mathcal A'}\,,
\end{equation}
where $\mathcal A'$ is RG transformed action,
\[
\mathcal A'=RG_l\{\mathcal A\}\,,\qquad l=\log L\,.
\]
We will write
\begin{equation}\label{eq:RGop}
RG_{l,\mathcal A}\,\mathcal O_\alpha=\sum_\beta z_\alpha{}^\beta(L)\,\mathcal O_\beta\,,\qquad L=e^l\,,
\end{equation}
stressing the fact that $RG$ transformation acts on the fields linearly;
additional index $\mathcal A$ is written down to remind that this transformation
depends on the action $\mathcal A$.

Let us assume now that $\mathcal A$ is the fixed-point action,
\[
\mathcal A=\mathcal A_*\,.
\]
Then $\mathcal A'$ coincides with $\mathcal A$, and \eqref{eq:RGop} yields linear
relation for the correlation functions of the composite fields. Recall that by the
composition law of the RG transformations,
\[
RG_{l+l'}\{*\}=RG_l\{RG_{l'}\{*\}\}\,,
\]
in this case the coefficients $z$ must obey the equation
\[
z_\alpha{}^\beta(l+l')=\sum_\gamma z_\alpha{}^\gamma(l')\,z_\gamma{}^\beta(l)\,,
\]
where I changed the notations $z(l)=z(L=e^l)$. In other words these coefficients,
as the functions of $l$, must satisfy linear differential equation with constant
coefficients,
\[
\frac d{dl}\,z_\alpha{}^\beta(l)=-\sum_\gamma D_\alpha{}^\gamma\,z_\gamma{}^\beta(l)\,,
\qquad\text{where}\qquad
D_\alpha{}^\beta=-\frac d{dl}\,z_\alpha{}^\beta(l)\Big|_{l=0}\,.
\]
The coefficients $D_\alpha{}^\beta$ are matrix elements of the linear operator $D$
corresponding to the infinitesimal transformation \eqref{eq:RGop}
\[
RG_{\delta l,\mathcal A_*}=I-\delta l\,D+O(\delta l^2)\,.
\]
Let us denote $\Phi_\alpha(x)$ the eigenvectors of the operator $D$, i.e.\ special
linear combinations of the composite fields $\mathcal O_\alpha(x)$ which satisfy
the equations
\[
D\,\Phi_\alpha\equiv D_\alpha{}^\beta\,\Phi_\beta=D_\alpha\,\Phi_\alpha\,,
\]
where I use the notation $D_\alpha$ for the corresponding eigenvalues of $D$. Then
\[
RG_{l,\mathcal A_*}\,\Phi_\alpha=e^{-lD_\alpha}\,\Phi_\alpha\,,
\]
and for the fixed-point theory the equation \eqref{eq:RGop} reads
\begin{equation}\label{eq:scaling}
\bigl\langle\Phi_{\alpha_1}(x_1)\cdots\Phi_{\alpha_n}(x_n)\bigr\rangle\big|_{\mathcal A_*}
=L^{-D_{\alpha_1}}\cdots L^{-D_{\alpha_n}}\,\bigl\langle\Phi_{\alpha_1}(x_1/L)\cdots\Phi_{\alpha_n}(x_n/L)\bigr\rangle\big|_{\mathcal A_*}\,,
\end{equation}
i.e.\ the above correlation functions are homogeneous functions of the coordinates
$x_1,\dots,x_n$, of the degree $-(D_{\alpha_1}+\dots+D_{\alpha_n})$. For the
two-point correlation functions, which depend only on the distance between the two
points, it follows from \eqref{eq:scaling}
\[
\bigl\langle\Phi_\alpha(x)\Phi_\beta(x')\bigr\rangle\big|_{\mathcal A_*}=\frac{G_{\alpha,\beta}}{|x-x'|^{D_\alpha+D_\beta}}\,,
\]
where $G_{\alpha,\beta}$ are constants. In general, the equation expresses the
\emph{scale invariance} of the theory described by the fixed-point action
$\mathcal A_*$. The eigenvalues $D_\alpha$ are called the \emph{scale dimensions}
(or \emph{anomalous scale dimensions}) of the fields $\Phi_\alpha$. If
$\mathcal A_*$ is not a free-field theory, there is no reason to expect them to
coincide with ``canonical dimensions''. I'd like to stress that these anomalous
dimensions as well as the basic ``scaling'' fields $\Phi_\alpha$ are defined
relative to given fixed point $\mathcal A_*$. We will see that in the theory of
continuous phase transitions these dimensions determine the critical exponents.

We have seen that the space $\mathcal F$ is identical to the tangent space
$T\Sigma_{\mathcal A_*}$ and therefore the above operator $D$ must be closely
related to the operator $K$ which is the linearization of the RG generator
$B\{\mathcal A\}$ near the fixed point,
\begin{equation}\label{eq:BK}
B\{\mathcal A_*+\delta\mathcal A\}=K\,\delta\mathcal A\,.
\end{equation}
The relation is very simple. Any eigenvector $\Psi_\alpha$ of $K$,
\[
K\,\Psi_\alpha=\kappa_\alpha\,\Psi_\alpha\,,
\]
is the integral
\[
\Psi_\alpha=\int d^dx\,\Phi_\alpha(x)\,.
\]
Therefore the eigenvalues are related as
\[
\kappa_\alpha=d-D_\alpha\,,
\]
where $d$ is the number of the space dimensions. The term $d$ appears because in
the step 2 of the RG transformation of the integral $\int d^dx\,\Phi(x)$ we have to
make additional transformation of the integration measure $d^dx$,
\[
\int d^dx\,\Phi(x/L)=L^d\int d^dx\,\Phi(x)\,.
\]

As we have seen, the roles played by positive and negative eigenvalues of $K$ in
the formation of the RG flow near $\mathcal A_*$ are very different. The
negative-eigenvalue subspace $\mathcal F_-$ corresponds to the critical surface,
whereas the positive-eigenvalue subspace $\mathcal F_+$ is the linear
approximation to the unstable manifold $U(\mathcal A_*)\subset\Sigma(\infty)$. The
common terminology is as follows. One refers to
\[
\begin{array}{llll}
\Phi_\alpha & \text{with}\quad D_\alpha>d & \text{(i.e.\ with }\kappa_\alpha<0) & \text{as the ``irrelevant'' field,}\\
\Phi_\alpha & \text{with}\quad D_\alpha=d & \text{(i.e.\ with }\kappa_\alpha=0) & \text{as the ``marginal'' field,}\\
\Phi_\alpha & \text{with}\quad D_\alpha<d & \text{(i.e.\ with }\kappa_\alpha>0) & \text{as the ``relevant'' field.}
\end{array}
\]
If marginal fields are absent, the dimensionality of $U(\mathcal A_*)$ equals the
number of relevant fields.

The situation is more complex if the marginal fields are present. Let me assume
that there is only one marginal field $\Phi_M$. By definition, the relevant
(irrelevant) fields correspond to deformations $\delta\mathcal A$ of the
fixed-point action which increase (decrease) with the RG flow in the linear
approximation in $\delta\mathcal A$. The deformation
\[
\delta\mathcal A_M=\lambda_M\int d^dx\,\Phi_M(x)
\]
corresponding to the marginal field does not change in the linear approximation,
but it may well increase or decrease with the RG flow if we go beyond the linear
approximation, say take into account the terms $\sim\delta\mathcal A_M^2$. If
$\delta\mathcal A_M$ increases with the RG flow for $\lambda_M$ small enough, the
deformation $\delta\mathcal A_M$ (and the field $\Phi_M$) is called
\emph{marginally relevant}; similarly, if it decreases, one refers to $\Phi_M$ as
\emph{marginally irrelevant} field.

In many important respects the marginally relevant (marginally irrelevant) field
can be treated on the same footing with the relevant (irrelevant) ones. Thus, if a
marginally relevant field is present, the unstable manifold $U(\mathcal A_*)$ is
generated by the subspace $\mathcal F'_+$ which includes this marginally relevant
field along with all the relevant fields,
\[
\mathcal F'_+=TU(\mathcal A_*)|_{\mathcal A_*}\,.
\]
We are going to have more precise statements after studying some examples.

\subsection*{Gaussian fixed point}

Let us consider the simplest fixed point for which all calculations can be
performed explicitly, the so called \emph{Gaussian fixed point}. It is described
by the massless free action
\[
\mathcal A_{*,G}=\int d^dx\,\frac12\,(\partial\varphi_0)^2
\]
It is straightforward to check that it is indeed a fixed point. If we split
$\varphi_0$ into the slow and the fast components, these modes completely
decouple,
\[
\mathcal A_{*,G}=\int d^dx\,\left(\frac12\,(\partial\varphi_1)^2+\frac12\,(\partial\tilde\varphi)^2\right)\,,
\]
and therefore the step 1 of the RG procedure does not change the action
$\mathcal A_{*,G}$ at all,
\[
\mathcal A_1=\int d^dx\,\frac12\,(\partial\varphi_1)^2
\]
The step 2 is also easy to perform. Setting
\[
\varphi_1(x)=Z^{-\frac12}(L)\,\varphi_0(x/L)
\]
we obtain
\[
\mathcal A'[\varphi_0(x)]=\mathcal A_1\bigl[Z^{-\frac12}(L)\,\varphi_0(x/L)\bigr]\,,
\]
\[
\mathcal A'=Z^{-1}(L)\int d^dx\,\frac12\left(\frac\partial{\partial x^\mu}\varphi_0(x/L)\right)^2=\frac{L^{d-2}}{Z(L)}\int d^dx\,\frac12\left(\frac\partial{\partial x^\mu}\varphi_0(x)\right)^2\,.
\]
where the change of the coordinates $x\to Lx$ is made in the last step. Under the
choice
\[
Z^{\frac12}(L)=L^{\frac{d-2}2}
\]
$\mathcal A'$ coincides with the original action $\mathcal A_{*,G}$, so it indeed
is a fixed point.

It is not too difficult to find the RG action on the composite fields in this
fixed-point theory. It will be sufficient for us to consider the fields
\[
\mathcal O_n=\frac1{n!}\,\varphi_0^n\,.
\]
The step 1 of the RG transformation yields for $\mathcal O_n$
\[
\frac1{n!}\,\varphi_1^n+\sum_{k=1}^n\frac1{k!}\,\langle\tilde\varphi^k\rangle_{\tilde\varphi}\,\frac1{(n-k)!}\,\varphi_1^{n-k}\,.
\]
In the case of infinitesimal transformation $L=1+\delta l$ we have
\[
\langle\tilde\varphi^2\rangle_{\tilde\varphi}=\int_{\Lambda(1-\delta l)<|k|<\Lambda}\frac{d^dk}{(2\pi)^d}\,\frac1{k^2}=N_d\,\Lambda^{d-2}\,\delta l\,,
\]
where
\[
N_d=\frac{\Omega_d}{(2\pi)^d}\,,\qquad\text{and}\qquad\Omega_d=\frac{2\pi^{\frac d2}}{\Gamma(\frac d2)}
\]
is the area of $S^{d-1}$. Therefore we have in this case
\[
\frac1{n!}\,\varphi_1^n+\frac1{(n-2)!}\,\varphi_1^{n-2}\,\frac{N_d\Lambda^{d-2}}2\,\delta l\,.
\]
As this field does not involve derivatives the step 2 amounts to multiplying
$\varphi_1$ by $Z^{-\frac12}(L)\sim1-\frac{d-2}2\,\delta l$. Therefore the operator
$D$ acts on $\mathcal O_n$ as
\begin{equation}\label{eq:DOn}
D\,\mathcal O_n=n\,\frac{d-2}2\,\mathcal O_n-\frac{N_d\Lambda^{d-2}}2\,\mathcal O_{n-2}\,.
\end{equation}
It is not difficult to diagonalize this operator.

\begin{exercise}
Find the linear combinations
\[
\Phi_n=\mathcal O_n+\sum_{k=1}^na_k\,\mathcal O_{n-k}
\]
such that
\[
D\,\Phi_n=D_n\,\Phi_n\,.
\]
Show that
\begin{equation}\label{eq:Phi4}
\Phi_4=\frac1{4!}\,\varphi_0^4-\frac14\,\langle\varphi_0^2\rangle\,\varphi_0^2+\frac18\,\langle\varphi_0^2\rangle^2\,,
\qquad\langle\varphi_0^2\rangle=\int_{|k|<\Lambda}\frac{d^dk}{(2\pi)^d}\,\frac1{k^2}=\frac{N_d\Lambda^{d-2}}{d-2}\,,
\end{equation}
and generally
\[
\Phi_n=\frac1{n!}\,{:}\varphi_0^n{:}\ :=\ \frac1{n!}\,\varphi_0^n-\frac1{2!\,(n-2)!}\,\langle\varphi_0^2\rangle\,\varphi_0^{n-2}-\text{further self-contractions.}
\]
Calculate $\langle\Phi_n(x)\Phi_{n'}(0)\rangle$.
\end{exercise}

The eigenvalues $D_n$ are found immediately, without computing exact form of the
eigenvectors $\Phi_n$. Indeed, the eq.~\eqref{eq:DOn} shows that the operator $D$
acts on the fields $\mathcal O_n$ as triangular matrix (i.e.\ it mixes
$\varphi_0^n$ with itself and lower powers of $\varphi_0$), and so its eigenvalues
coincide with its diagonal elements. We find
\[
D_n=n\,\frac{d-2}2\,.
\]
Note that $D_1=\frac{d-2}2$ is the canonical dimension of $\varphi_0$ in
$d$-dimensional space. The dimensions $D_n$ coincide with the canonical dimensions
of the composite fields $\varphi_0^n$, which is not surprising because we are
considering the free massless field theory $\mathcal A_{*,G}$.

Dimensionality of the space $\mathcal F_+$ (i.e.\ the space of relevant fields)
associated with the fixed point $\mathcal A_{*,G}$ depends on $d$. The field
$\Phi_2=\frac12{:}\varphi_0^2{:}$ is always relevant; its dimension
\[
D_2=d-2<d\,.
\]
It remains the only relevant field for all $d>6$. For $d<6$ the field
$\Phi_3=\frac1{3!}{:}\varphi_0^3{:}$ with the dimension
\[
D_3=\frac32\,(d-2)
\]
becomes relevant, too. For $d<4$ it is joined by the field
$\Phi_4=\frac1{4!}{:}\varphi_0^4{:}$. The situation is summarized in the table
\begin{center}
\begin{tabular}{cc}
$d$ & relevant fields\\[4pt]
$d>6$ & $\varphi_0^2$\\
$6>d>4$ & $\varphi_0^2,\ \varphi_0^3$\\
$4>d>10/3$ & $\varphi_0^2,\ \varphi_0^3,\ \varphi_0^4$\\
$10/3>d>3$ & $\varphi_0^2,\ \varphi_0^3,\ \varphi_0^4,\ \varphi_0^5$\\
$\cdots$ & \\
$\frac{2n}{n-2}>d>\frac{2(n+1)}{n-1}$ & $\varphi_0^2,\ \varphi_0^3,\dots,\varphi_0^n$
\end{tabular}
\end{center}
At $d=2$ all polynomials of $\varphi_0$ are relevant.

Note that the above classification coincides with our earlier classification of
the perturbatively renormalizable scalar field theories. We see that in generic
dimensionality $d$ the perturbatively renormalizable field theories correspond to
the unstable manifold $U(\mathcal A_{*,G})$ of the gaussian fixed point. Also, we
realize now that this classification of renormalizable field theories is relevant
only to the vicinity of the gaussian fixed point (after all the perturbative
expansion is developed around this point); if we manage to find more complex fixed
point, its own unstable manifold will give rise to a family of finite local field
theories which may be not too similar to the perturbatively defined renormalizable
field theories we considered before.

Thus we have a correspondence between the unstable manifold $U(\mathcal A_{*,G})$
associated with the Gaussian fixed point and perturbatively renormalizable scalar
field theories in different dimensions $d$. This correspondence is more subtle if
$d$ takes special values such that one of the fields $\Phi_n$ acquires the
dimension $D_n=d$, i.e.\ becomes marginal. This happens for instance at $d=6$ with
$\Phi_3\sim\varphi_0^3$ and at $d=4$ with $\Phi_4\sim\varphi_0^4$. The last case is
of most interest to us, and we take it as main example.

\subsection*{\texorpdfstring{$\varphi^4$}{phi**4} flow at \texorpdfstring{$d=4$}{d=4}}

As at $d=4$ the field $\Phi_4$ is marginal, we have
\begin{equation}\label{eq:marginal4}
RG_l\Bigl\{\mathcal A_{*,G}+\lambda\int d^4x\,\Phi_4(x)\Bigr\}=\mathcal A_{*,G}+\lambda\int d^4x\,\Phi_4(x)+O(\lambda^2)\,,
\end{equation}
i.e.\ in the order $\lambda$ the perturbed action does not flow at all. However,
nontrivial flow can be generated beyond the linear approximation, in the higher
orders in $\lambda$. Let us estimate this flow in the order $\lambda^2$.

To perform the step 1 of the RG transformation we have to integrate over the fast
component $\tilde\varphi$ in
\[
\varphi_0=\varphi_1+\tilde\varphi\,.
\]
Consider the action $\mathcal A_0$ of the form
\begin{equation}\label{eq:A0phi4}
\mathcal A_0=\mathcal A_{*,G}+\lambda\int d^4x\,\Phi_4(x)=\int d^4x\,\left(\frac12\,(\partial\varphi_0)^2+\frac\lambda{4!}\,\varphi_0^4-\frac{\lambda\,b}4\,\varphi_0^2+\Const\right)
\end{equation}
where
\[
b=\langle\varphi_0^2\rangle_{\mathcal A_{*,G}}=\int_{|k|<\Lambda}\frac{d^4k}{(2\pi)^4\,k^2}=\frac{\Lambda^2}{16\pi^2}
\]
(recall that $\Phi_4=\frac1{4!}\varphi_0^4-\frac14\langle\varphi_0^2\rangle_{*,G}\,\varphi_0^2+\Const$,
as we should have found in the Exercise above). We have
\begin{multline}\label{eq:A1phi4}
e^{-\mathcal A_1[\varphi_1]}\sim e^{-\mathcal A_0[\varphi_1]}\times\int[D\tilde\varphi]\\
\times\exp\left\{-\int d^4x\,\left[\frac12\,(\partial\tilde\varphi)^2+\lambda\left(\frac1{4!}\,\tilde\varphi^4-\frac b4\,\tilde\varphi^2+\frac16\,\varphi_1\tilde\varphi^3+\frac14\,\varphi_1^2\tilde\varphi^2+\frac16\,\varphi_1^3\tilde\varphi\right)\right]\right\}\,.
\end{multline}
As we discussed already in Section~16, integration over $\tilde\varphi$ yields
\[
\mathcal A_1=\mathcal A_0-\text{sum of all connected diagrams;}
\]
the diagrams being built from the vertices
\[
\dia[0.45]{\draw[line width=0.9pt] (-1.4,0) -- (1.4,0); \node[vtx] at (0,0) {};}\qquad
\dia[0.45]{\foreach \a in {45,135,225,315} \draw[line width=0.9pt] (0,0) -- (\a:1.4); \node[vtx] at (0,0) {};}\qquad
\dia[0.45]{\draw[prop,densely dashed] (-1.4,0) -- (0,0); \foreach \a in {45,0,-45} \draw[line width=0.9pt] (0,0) -- (\a:1.4); \node[vtx] at (0,0) {};}\qquad
\dia[0.45]{\foreach \a in {135,225} \draw[prop,densely dashed] (0,0) -- (\a:1.4); \foreach \a in {45,315} \draw[line width=0.9pt] (0,0) -- (\a:1.4); \node[vtx] at (0,0) {};}\qquad
\dia[0.45]{\foreach \a in {150,180,210} \draw[prop,densely dashed] (0,0) -- (\a:1.4); \draw[line width=0.9pt] (0,0) -- (1.6,0); \node[vtx] at (0,0) {};}\,,
\]
where the solid lines correspond to the fast field $\tilde\varphi$ (and thus all
solid lines must be contracted), and the dashed lines represent the field
$\varphi_1$ which plays the role of an external field (i.e.\ the diagrams involve
no contractions of the dashed lines).

It is not difficult to check that the equation \eqref{eq:marginal4} indeed holds
to the order $\lambda$. The constant $b$ in the $\varphi^2$ term in
\eqref{eq:A0phi4} is such that the contribution of the diagram
\[
\dia[0.45]{\draw[prop,densely dashed] (-2,0) -- (2,0); \draw[line width=0.9pt] (0,0.6) circle (0.6); \node[vtx] at (0,0) {};}\ \sim\ \varphi_1^2
\]
is compensated exactly by the transformation of the term
$\frac{\lambda b}4\varphi_1^2$ in the step 2 of the RG procedure.

In the order $\lambda^2$ many new terms in $\mathcal A_1$ are generated. For
example the diagram
\begin{equation}\label{eq:tree6}
\diac[0.45]{\coordinate (L) at (-1.2,0); \coordinate (R) at (1.2,0); \draw[line width=0.9pt] (L) -- (R);
  \foreach \a in {120,180,240} \draw[prop,densely dashed] (L) -- ++(\a:1.1);
  \foreach \a in {60,0,-60} \draw[prop,densely dashed] (R) -- ++(\a:1.1);
  \node[vtx] at (L) {}; \node[vtx] at (R) {};}
\end{equation}
gives rise to the terms $\sim\varphi_1^6$. All such fields are irrelevant. The
diagram contributing to the terms $\sim\varphi_1^4$ is
\begin{equation}\label{eq:bubble4}
\diac[0.45]{\coordinate (L) at (-1,0); \coordinate (R) at (1,0); \draw[line width=0.9pt] (0,0) ellipse (1 and 0.7);
  \foreach \a in {135,225} \draw[prop,densely dashed] (L) -- ++(\a:1.1);
  \foreach \a in {45,-45} \draw[prop,densely dashed] (R) -- ++(\a:1.1);
  \node[vtx] at (L) {}; \node[vtx] at (R) {};}
\end{equation}
This diagram admits Taylor expansion in the external momenta. The zeroth term of
this expansion brings in the term $\varphi_1^4$ (without derivatives) into
$\mathcal A_1$, while the higher terms give rise to the irrelevant contributions
like $\varphi_1^2(\partial\varphi_1)^2$. I am going to use approximation in which
all irrelevant terms in $\mathcal A$ are neglected. Intuitive reason for this
approximation is that irrelevant terms die out very fast with the RG flow. Within
this approximation the diagram \eqref{eq:tree6} can be neglected, and the diagram
\eqref{eq:bubble4} can be evaluated at zero external momenta only, yielding the
following contribution to $\mathcal A_1$
\[
\frac{\Delta\lambda}{4!}\int d^4x\,\varphi_1^4(x)\,,
\]
where
\begin{equation}\label{eq:Dlambda}
\frac{\Delta\lambda}{4!}=-\frac{\lambda^2}{16}\int_{\Lambda>|k|>\Lambda/L}\frac{d^4k}{(2\pi)^4}\,\frac1{k^2}\,\frac1{k^2}\,.
\end{equation}
The factor $\lambda^2/16$ in this formula appears as follows. The term in
\eqref{eq:A1phi4} responsible for the vertices in the diagram \eqref{eq:bubble4}
is $\frac\lambda4\varphi_1^2\tilde\varphi^2$, so each vertex comes with the factor
$-\lambda/4$. One has to expand the exponential to the second order to get the
diagram \eqref{eq:bubble4}, this brings in additional factor $1/2!$. Finally, a
factor of 2 comes from the combinatorics of contractions in
$\langle\tilde\varphi^2\tilde\varphi^2\rangle$. The momentum integral in
\eqref{eq:Dlambda} is easy to evaluate, and we obtain
\[
\Delta\lambda=-3\,\frac{\lambda^2}{(4\pi)^2}\,\log L\,.
\]
In this order in $\lambda$ the step 2 is trivial. For the marginal term
$\int d^4x\,\varphi_1^4(x)$ the factor $Z^{-2}(L)$ which comes from the
transformation
\[
\varphi_1(x)=Z^{-\frac12}(L)\,\varphi_0(x/L)
\]
is exactly compensated by the transformation of the integration measure $d^4x$
(recall that $Z^{\frac12}=L$ at $d=4$; it is easy to see that possible corrections
to this formula are $\sim\lambda^2$).

Thus, in our approximation the result of the RG transformation of the action
\eqref{eq:A0phi4} again has the form \eqref{eq:A0phi4} with $\lambda$ replaced by
\[
\lambda'=\lambda-3\,\frac{\lambda^2}{(4\pi)^2}\,\log L\,.
\]
Taking $L=1+\delta l$ we see that $\lambda$ \emph{decreases} with the RG flow,
according to the equation
\begin{equation}\label{eq:WilsonBeta}
-\frac d{dl}\,\lambda(l)=3\,\frac{\lambda(l)^2}{(4\pi)^2}\,,
\end{equation}
i.e.\ at $d=4$ the field $\Phi_4$ is \emph{marginally irrelevant}. Note that the
r.h.s.\ of \eqref{eq:WilsonBeta} coincides with the $\lambda^2$ term
\eqref{eq:beta1} of the Callan-Symanzik $\beta$ function of $\varphi^4$ theory
which we obtained in Section~15.

% Source pp. 201-207 (Lectures19-25.pdf).
\section{Some Properties of RG Flow}

In the previous Section we have discussed the basic properties of the Gaussian
fixed point $\mathcal A_{*G}$ at different dimensions $d$. At special values of
$d$ the fixed point theory has marginal operators. Thus, at $d=4$ the field
$\Phi_4\sim\varphi^4$ is marginal, and as more careful calculation shows, it is
\emph{marginally irrelevant}. This is expressed by the RG flow equation
\eqref{eq:WilsonBeta} (note that it is equivalent to the leading term in the beta
function \eqref{eq:beta1}).

\subsection*{RG flow and Poincar\'e-Dulac normal form}

It is interesting to see, in general, what is special about the marginal
deformations from the mathematical point of view. Discussing on the abstract
level, let us assume that we have generic critical fixed point
$\mathcal A_*\in\Sigma$. Let $y^\alpha$ be some local coordinates on $\Sigma$,
covering some neighborhood of $\mathcal A_*$ and chosen in such a way that the
fixed point corresponds to $y^\alpha=0$. One can think of $y^\alpha$ as the
collection of coupling constants $u_{2k}-u_{2k,*}$, $v_{2k}-v_{2k,*}$, $\dots$,
where $u$'s, $v$'s, etc, are the coefficients appearing in writing $\mathcal A$ in
terms of various monomials in $\varphi_0$ and its derivatives, and $u_*$'s, etc,
are their fixed-point values. Anyhow, the RG flow equation
\[
\frac d{dl}\,\mathcal A=B\{\mathcal A\}
\]
can be written in such coordinates as
\begin{equation}\label{eq:flowy}
\frac d{dl}\,y^\alpha=B^\alpha(\{y\})\,,
\end{equation}
where $B^\alpha$ are some functions of all $y$'s. As $y=0$ is assumed to be a fixed
point, these functions satisfy the condition
\[
B^\alpha(\{0\})=0\,.
\]
Let us assume that these functions can be expanded into Taylor series around the
fixed point,
\[
B^\alpha(\{y\})=K^\alpha_\beta\,y^\beta+C^\alpha_{\beta\gamma}\,y^\beta y^\gamma+\cdots\,,
\]
where $K$, $C$, etc, are some constants, and summation over repeated indices is
assumed. The coefficients $K^\alpha_\beta$ are the matrix elements of the operator
$K$, the linearization of the RG generator $B$ near the fixed point,
\[
B\{\mathcal A_*+\delta\mathcal A\}=K\,\delta\mathcal A+O(\delta\mathcal A^2)\,.
\]
By appropriate linear transformation of the coordinates $y$:
$y^\alpha\to M^\alpha_\beta y^\beta$ the matrix $K^\alpha_\beta$ can be brought to
diagonal form.\footnote{Non-diagonalizable forms contradict unitarity (positivity
of norms).} In these coordinates the flow equation \eqref{eq:flowy} takes the form
\begin{equation}\label{eq:flowdiag}
\frac d{dl}\,y^\alpha=\kappa_\alpha\,y^\alpha+c^\alpha_{\beta\gamma}\,y^\beta y^\gamma+\cdots\,,
\end{equation}
where $\kappa_\alpha$ are the eigenvalues of $K$, and no summation over $\alpha$ is
assumed in the first term in the r.h.s.; the coefficients $c$ are related to $C$
by a linear transformation, in an obvious way.

No further simplification of the linear term in the r.h.s.\ of \eqref{eq:flowdiag}
is possible. However one may ask if the quadratic and higher order terms in the
r.h.s.\ can be simplified by nonlinear coordinate transformation
\[
y^\alpha=y^\alpha(\{\tilde y\})=\tilde y^\alpha+a^\alpha_{\beta\gamma}\,\tilde y^\beta\tilde y^\gamma+\cdots\,.
\]
Under coordinate transformations the functions $B^\alpha$ transform as components
of a vector field on $\Sigma$. Simple calculation shows that in terms of $\tilde y$
the RG flow equations take the form
\[
\frac d{dl}\,\tilde y^\alpha=\kappa_\alpha\,\tilde y^\alpha+\tilde c^\alpha_{\beta\gamma}\,\tilde y^\beta\tilde y^\gamma+\cdots\,,
\qquad\text{where}\qquad
\tilde c^\alpha_{\beta\gamma}=c^\alpha_{\beta\gamma}+(\kappa_\alpha-\kappa_\beta-\kappa_\gamma)\,a^\alpha_{\beta\gamma}\,.
\]
One can completely eliminate the quadratic term by taking
\[
a^\alpha_{\beta\gamma}=-\frac{c^\alpha_{\beta\gamma}}{\kappa_\alpha-\kappa_\beta-\kappa_\gamma}\,;
\]
with this choice
\[
\frac d{dl}\,\tilde y^\alpha=\kappa_\alpha\,\tilde y^\alpha+O(\tilde y^3)\,.
\]
This is true provided the combinations of the eigenvalues appearing in the
denominators do not vanish. However if some triplet of the eigenvalues satisfies
the \emph{resonance condition}
\[
\kappa_\alpha-\kappa_\beta-\kappa_\gamma=0
\]
the ``resonance'' term
\[
c^\alpha_{\beta\gamma}\,y^\beta y^\gamma
\]
(no summation) cannot be eliminated.

Similar situation takes place in the higher orders. If components $B^\alpha$ of a
vector field $B$ can be written as power series
\begin{equation}\label{eq:Bseries}
B^\alpha(\{y\})=\kappa_\alpha\,y^\alpha+c^\alpha_{\beta_1\beta_2}\,y^{\beta_1}y^{\beta_2}+c^\alpha_{\beta_1\beta_2\beta_3}\,y^{\beta_1}y^{\beta_2}y^{\beta_3}+\cdots\,.
\end{equation}
in some local coordinates $y^\alpha$, then by means of appropriate analytic
coordinate transformation
\begin{equation}\label{eq:normaltr}
y^\alpha\to y^\alpha+a^\alpha_{\beta_1\beta_2}\,y^{\beta_1}y^{\beta_2}+a^\alpha_{\beta_1\beta_2\beta_3}\,y^{\beta_1}y^{\beta_2}y^{\beta_3}+\cdots
\end{equation}
one can eliminate all nonlinear terms in \eqref{eq:Bseries}, except for the
so-called resonance terms. A term
\begin{equation}\label{eq:resterm}
c^\alpha_{\beta_1\beta_2\cdots\beta_n}\,y^{\beta_1}y^{\beta_2}\cdots y^{\beta_n}
\end{equation}
in the function $B^\alpha$ is called the \emph{resonance term} if the following
\emph{resonance condition}
\begin{equation}\label{eq:rescond}
\kappa_\alpha=\kappa_{\beta_1}+\kappa_{\beta_2}+\dots+\kappa_{\beta_n}
\end{equation}
is satisfied. One can always adjust the coefficients $a$ in \eqref{eq:normaltr},
order by order in $y$, in such a way that in new coordinates
\begin{equation}\label{eq:normalform}
\frac d{dl}\,y^\alpha=\kappa_\alpha\,y^\alpha+\text{Resonance Terms}
\end{equation}
If there are no resonances, all nonlinear terms in \eqref{eq:Bseries} can be
eliminated. This statement is known as \emph{Poincar\'e-Dulac
theorem}.\footnote{See e.g.\ V.I.~Arnold, Geometrical Methods in the Theory of
Ordinary Differential Equations, Springer, 1988.} I will call ``canonical'' the
coordinates $y$ in which the flow equation takes the form \eqref{eq:normalform}.

Two remarks are in order. First, the above theorem \emph{does not} suggest that
the global structure of the RG flow is in any sense ``trivial'', even in
non-resonance cases. The series \eqref{eq:normaltr} which defines the
transformation converges at best in some finite neighborhood of the fixed point
$y=0$. Other interesting irregularities the global RG flow may possess (like
other fixed points) are located no closer than at the boundary of this domain of
convergence; in the new coordinates where $B^\alpha$'s are ``trivial'' they are
``pushed out to the infinity''. To get some idea consider the following

\begin{exercise}
Study the flow equation with just one coordinate,
\[
\frac d{dl}\,y=y-y^2
\]
(which has two fixed points $y=0$ and $y=1$). Find transformation
$y=f(\tilde y)=\tilde y+\cdots$ such that
\[
\frac d{dl}\,\tilde y=\tilde y\,.
\]
Study the properties of the map $y\to\tilde y$.
\end{exercise}

Second remark concerns very simple interpretation which the above resonance
conditions have in terms of ordinary dimensional analysis. In the neighborhood of
$\mathcal A_*$ where the coordinates $y$ are defined, the action can be written as
\[
\mathcal A=\mathcal A_*+y^\alpha\int d^dx\,\Phi_\alpha(x)+O(y^2)\,,
\]
where $\Phi_\alpha(x)$ are the fields which diagonalize the operator $D$
associated with the fixed point $\mathcal A_*$. The eigenvalues $D_\alpha$ (i.e.\
the dimensions of the fields $\Phi_\alpha$) are related to $\kappa_\alpha$ as
\[
D_\alpha+\kappa_\alpha=d\,.
\]
Therefore the numbers $\kappa_\alpha$ can be interpreted as the mass dimensions of
the ``coupling constants'' $y^\alpha$. The resonance condition \eqref{eq:rescond}
then states that the resonance term \eqref{eq:resterm} has ``right dimension''. In
this sense, the coordinate system in which the RG flow equations have the form
\eqref{eq:normalform} respects the straightforward dimensional analysis.

From this point of view, what distinguishes the marginal case is that the presence
of a marginal field (i.e.\ the field $\Phi_\alpha$ with $D_\alpha=d$, or
$\kappa_\alpha=0$) always leads to infinitely many resonances. Indeed, if, say,
$\kappa_0=0$, the resonance condition \eqref{eq:rescond} with any $\alpha$ can be
satisfied by taking $\beta_1=\alpha$ and $\beta_2=\beta_3=\dots=\beta_n=0$.

To illustrate this point let us consider an example with only two coupling
constants $y^0$, $y^1$, such that $\kappa_0=0$ and $\kappa_1=2$. I will use
suggestive notations $y^0=\lambda$ and $y^1=M^2$. Then in the flow equation for
$\lambda$ all terms $\lambda^n$ are resonant, so the coordinates exist in which
this equation takes the form
\begin{equation}\label{eq:normallambda}
\frac d{dl}\,\lambda=\sum_{n=2}^\infty b_n\,\lambda^n=-\beta(\lambda)\,.
\end{equation}
Similarly, in the equation for $M^2$ all terms of the form $M^2\lambda^n$ are
resonant, and in canonical coordinates one can write
\begin{equation}\label{eq:normalM}
\frac d{dl}\,M^2=2M^2-\sum_{n=1}^\infty\gamma_n\,M^2\lambda^n=\bigl(2-\gamma_{\varphi^2}(\lambda)\bigr)\,M^2\,.
\end{equation}
Note that with these notations these equations have exactly the form of the
coupling and the mass ``flow'' equations associated with the Callan-Symanzik
equation of the $d=4$ $\varphi^4$ theory.

\begin{remark}
In fact, in this important example a freedom still remains for further changes of
coordinates $\lambda\to\tilde\lambda(\lambda)$, $M^2\to\tilde M^2=M^2f(\lambda)$,
which doesn't change the above ``canonical'' form of the flow equations.
\end{remark}

\begin{exercise}
Check (to the order $n=$ when you get bored, or see the pattern) that all terms
$\sim\lambda^n$ with $n\ge4$ can be eliminated order by order, by suitable choice
of coefficients $a_n$ in
\[
\tilde\lambda=\lambda+a_3\,\lambda^3+a_4\,\lambda^4+\cdots\,.
\]
Further, check that suitable redefinition of $M^2\to\tilde M^2$,
\[
\tilde M^2=M^2\,\bigl(1+c_1\,\lambda+c_2\,\lambda^2+\cdots\bigr)
\]
allows one to eliminate the terms with $n\ge2$ in the r.h.s.\ of
\eqref{eq:normalM}. In reference to renormalization of the $\varphi^4$ theory at
$d=4$, these results imply that a scheme exists in which all $n$th order
corrections to the beta-function vanish starting from $n=4$, and all corrections
to $\gamma_{\varphi^2}$ vanish starting from $n=2$. At the same time the
coefficients $b_2$ and $b_3$, and $\gamma_1$, are scheme-independent.
\end{exercise}

Note that the statement of the Poincar\'e-Dulac theorem concerns formal power
series. Existence of ``canonical coordinates'', even locally, depends on
convergence of the series which in general is not guaranteed. For general RG flow
in the full space $\Sigma$ the situation is indeed difficult. One of the problems
lies again in the ``resonance denominators''
$\kappa_\alpha-\kappa_{\beta_1}-\dots-\kappa_{\beta_n}$ which the coefficients $a$
in the series \eqref{eq:normaltr} have. When both positive and negative
$\kappa_\alpha$ are present, even if no resonance conditions \eqref{eq:rescond} are
satisfied exactly, they often are ``almost satisfied'', in which case some
denominators come very close to zero giving rise to anomalously large coefficients
in \eqref{eq:normaltr}.

The situation becomes much more favorable if, instead of the flow in the full
space $\Sigma$, we concentrate attention on the RG flow in an unstable manifold
$U(\mathcal A_*)$ associated with some fixed point $\mathcal A_*$. First,
$U(\mathcal A_*)$ typically has finite dimensionality (as our study of
$\mathcal A_{*,G}$ illustrates). Also, by the definition of unstable manifold all
corresponding $\kappa$ are either positive (relevant) or zero (marginal). In this
case the resonance denominators are under much better control and the series
\eqref{eq:normaltr} defining the canonical coordinates usually converges in finite
neighborhood of $\mathcal A_*$. Recall that $U(\mathcal A_*)$ is a part of
$\Sigma(\infty)$, i.e.\ it describes finite local QFT. For local QFT generated by
the fixed point $\mathcal A_*$ the canonical coordinates do exist.

\subsection*{Callan-Symanzik equation revisited}

The preceding analysis allows one to re-derive the Callan-Symanzik equation which
we have obtained before on the basis of renormalized perturbation theory.

Let $\mathcal A_*$ be a fixed point. Assuming that its unstable manifold is
$n$-dimensional, let $\{\lambda^i\}$, $i=1,\dots,n$ be coordinates on
$U(\mathcal A_*)$, such that $\mathcal A_*$ corresponds to $\lambda=0$. This
unstable manifold describes $n$-parameter family of finite local field theories. I
will denote the corresponding correlation functions as
\[
\bigl\langle\varphi(x_1)\cdots\varphi(x_N)\bigr\rangle_\lambda\,,
\]
indicating their dependence on $n$ parameters $\lambda^i$.

In the coordinates $\lambda$ the RG flow equation can be written as
\begin{equation}\label{eq:flowlambdai}
\frac d{dl}\,\lambda^i=-\beta^i(\lambda)\,,
\end{equation}
where $-\beta^i$ are components of the vector field $B^\alpha$ on
$U(\mathcal A_*)$ (by definition of $U$, $B(U)$ is tangent to $U$), the minus sign
is taken to match our previous notation for the beta-function. According to the
definition of the Wilson's RG, if $\lambda(l)$ is a solution of
\eqref{eq:flowlambdai} with some initial condition $\lambda(0)=\lambda$, the field
theories corresponding to $\lambda$ and $\lambda(l)$ are related by the scale
transformation $x\to x/L$. In particular, one can write for the correlation
functions
\[
\bigl\langle\varphi(x_1)\cdots\varphi(x_N)\bigr\rangle_{\lambda(l)}=Z^{\frac N2}(L)\,\bigl\langle\varphi(Lx_1)\cdots\varphi(Lx_N)\bigr\rangle_{\lambda(0)}\,,
\]
where $L=e^l$. Taking infinitesimal transformation $l=\delta l$, $L=1+\delta l$,
$\lambda(l)=\lambda-\beta(\lambda)\,\delta l$, we have
\begin{equation}\label{eq:CSgeneral}
\Bigl[x\frac\partial{\partial x}+N\,D(\lambda)+\sum_{i=1}^n\beta^i(\lambda)\,\frac\partial{\partial\lambda^i}\Bigr]\bigl\langle\varphi(x_1)\cdots\varphi(x_N)\bigr\rangle_\lambda=0\,,
\end{equation}
where $x\frac\partial{\partial x}$ is the shorthand notation for
\[
x\frac\partial{\partial x}=\sum_{k=1}^Nx_k^\mu\,\frac\partial{\partial x_k^\mu}\,,
\]
and
\[
D(\lambda)=L\frac d{dL}\,Z^{\frac12}(L)\Big|_{L=1}\,.
\]
Note that $D(0)=D_\varphi$, the dimension of $\varphi$ at the fixed point
$\mathcal A_*$. $D(\lambda)$ can be written as
\[
D(\lambda)=D_\varphi+\gamma(\lambda)\,,
\]
so that $\gamma(0)=0$.

The equation \eqref{eq:CSgeneral} generalizes straightforwardly to correlation
functions involving composite fields $\mathcal O_\alpha$. Repeating the same
analysis we obtain
\begin{multline}\label{eq:CScomposite}
\sum_{k=1}^N\bigl\langle\mathcal O_{\alpha_1}(x_1)\cdots(\mathcal D\mathcal O)_{\alpha_k}(x_k)\cdots\mathcal O_{\alpha_N}(x_N)\bigr\rangle_\lambda\\
+\sum_{i=1}^n\beta^i(\lambda)\,\frac\partial{\partial\lambda^i}\bigl\langle\mathcal O_{\alpha_1}(x_1)\cdots\mathcal O_{\alpha_N}(x_N)\bigr\rangle_\lambda=0\,,
\end{multline}
where
\[
(\mathcal D\mathcal O)_\alpha(x)=x^\mu\frac\partial{\partial x^\mu}\,\mathcal O_\alpha(x)+D_\alpha^\beta(\lambda)\,\mathcal O_\beta(x)\,,
\qquad\text{and}\qquad
D_\alpha^\beta(\lambda)=-L\frac d{dL}\,z_\alpha^\beta(L)\Big|_{L=1}\,;
\]
$z_\alpha^\beta(L)$ describes the RG transformation of the composite fields
$\mathcal O_\alpha$ (see Sect.~18). Again,
\[
D_\alpha^\beta=D_\alpha^\beta(\lambda)\big|_{\lambda=0}
\]
coincides with the matrix of dimensions at $\mathcal A_*$, whose eigenvalues are
the dimensions of the fields $\Phi_\alpha$.

We see that the CS equation is very general consequence of Wilson's RG theory and
not just the property of the perturbation theory around the Gaussian fixed point
(i.e.\ free field theory).

% Source pp. 208-213 (Lectures19-25.pdf).
\section{Fixed Points and Criticality}

It is possible to establish direct relation between the singular behavior of
thermodynamic quantities near critical points and the fixed points of the
Wilson's RG. This relation explains experimentally observed \emph{universality}
of critical singularities.

Let us consider again a magnet. As we have seen few lectures ago, the long-range
fluctuations of local magnetization $M(x)$ can be taken into account by
considering the statistical integral
\[
Z=\int[D\varphi_0]\,e^{-\mathcal A_0[\varphi_0]}\,,
\]
where the field $\varphi_0(x)$ is proportional to $M(x)$,
\[
\varphi_0(x)=C(T)\,M(x)\,,
\]
with some temperature-dependent coefficient. It is assumed that this functional
integral has cutoff, with the cutoff parameter $\Lambda$ such that
\[
\Lambda^{-1}\gg a\,,
\]
the atomic scale. Exact form of the action $\mathcal A_0$ is not known; it
combines statistical contributions of all microscopic fluctuations and depends
significantly on all details of the microscopic interactions. In Landau-Ginzburg
approximation one takes
\[
\mathcal A_0=\int d^dx\,\left[\frac12\,(\partial\varphi_0)^2+\frac{m_0^2(T)}2\,\varphi_0^2+\frac{\lambda_0(T)}{4!}\,\varphi_0^4\right]\,,
\]
where $m_0^2(T)$, $\lambda_0(T)$ are unknown functions of $T$, but it can be any
other combination of powers of $\varphi_0$ and its derivatives. What is important,
all the coefficients are some (generally unknown) functions of the temperature
$T$. This means that the actions describing real physical system, with all its
microscopic details, being functions of the temperature, form one-dimensional
submanifold in $\Sigma$; I will call it the ``physical curve'' (Wilson used the
term ``canonical curve'').

Now let us assume that there is a fixed point $\mathcal A_*$ out there in
$\Sigma$. At this point we actually will consider a subspace
$\Sigma_{\text{symm}}\subset\Sigma$ which contains only actions invariant under
$\varphi_0\to-\varphi_0$. We assume that there is a fixed point
$\mathcal A_*\in\Sigma_{\text{symm}}$ such that it has only one (symmetric)
relevant (or marginally relevant) field $\Phi_0$, with dimension $D_0<d$. In
other words, the unstable manifold in $\Sigma_{\text{symm}}$ is one-dimensional,
and correspondingly the critical surface $\Sigma_c$ has codimension one. This
means that the physical curve generally intersects this critical surface at some
point. This is illustrated in figure~\ref{fig:physcurve}.
\begin{figure}[ht]
\centering\figfile{physcurve}
\caption{The critical surface $\Sigma_c$ with the fixed point $\mathcal A_*$ and
its one-dimensional unstable manifold $U(\mathcal A_*)$. The physical curve,
parametrized by the temperature, crosses $\Sigma_c$ at $T=T_c$.}
\label{fig:physcurve}
\end{figure}

Let us denote $T_c$ the temperature corresponding to the intersection point. The
large-distance properties of the system can be understood if we can follow the RG
flow which starts at any point of the physical curve. Let us consider a typical
pattern of the RG flow in this situation. Besides the critical fixed point
$\mathcal A_*$, there are two non-critical ($R_c=0$) fixed points,
$P_{\text{high}}$ and $P_{\text{low}}$\footnote{For a magnetic transition
$P_{\text{high}}$ corresponds to the probability distribution
$P_{\text{high}}[\varphi_0]\sim\prod_x\delta(\varphi_0(x))$, and
$P_{\text{low}}[\varphi_0]\sim\prod_x\delta(\varphi_0(x)-\bar\varphi_0)+\prod_x\delta(\varphi_0(x)+\bar\varphi_0)$.}
(figure~\ref{fig:critflow}).
\begin{figure}[ht]
\centering\figfile{critflow}
\caption{RG flow near a critical point. Trajectories starting on the physical
curve above $T_c$ go to $P_{\text{high}}$, those below $T_c$ go to
$P_{\text{low}}$. The trajectory $C$ starting at $T_c$ stays in $\Sigma_c$ and
ends at $\mathcal A_*$; a trajectory $C_\tau$ starting close to $T_c$ follows $C$,
lingers near $\mathcal A_*$, and leaves along the unstable manifold.}
\label{fig:critflow}
\end{figure}

A trajectory starting at any point on the physical curve corresponding to $T>T_c$,
and far enough from $T_c$, converges quickly to the noncritical fixed point
$P_{\text{high}}$. Similarly, starting at any point with $T<T_c$ one quickly
reaches $P_{\text{low}}$. The only trajectory which behaves qualitatively
differently is the one that starts at the point $T_c$ itself. This trajectory $C$
lies entirely at the critical surface, and converges to the critical fixed point
$\mathcal A_*$ as $l\to\infty$. This means that the large-distance physics exactly
at the critical temperature is completely determined by the fixed-point action
$\mathcal A_*$. In particular, the two-point correlation function will have the
scale-invariant form
\[
\bigl\langle M(x)M(x')\bigr\rangle_{T=T_c}=\frac\Const{|x-x'|^{2D_\varphi}}
\]
for $|x-x'|\gg a\,e^{l_0}$, where $l_0$ is some finite amount of the RG ``time''
$l$ needed to arrive at close proximity of $\mathcal A_*$. Here $D_\varphi$ is the
dimension of $\varphi$ at this fixed point.

Consider now a trajectory $C_\tau$ which starts at some point $T$ very close to
$T_c$,
\begin{equation}\label{eq:tau}
\tau=\frac{T_c-T}{T_c}\ll1\,.
\end{equation}
To be definite I will assume that $T<T_c$. There are three distinct stages of the
RG ``evolution'' for such trajectory. At the first stage it goes very close to $C$
approaching $\mathcal A_*$, until it comes to close proximity of this fixed point.
It then stays near $\mathcal A_*$ for very long RG ``time'' $l$. This is the
second stage. Finally, in the third stage it departs toward $P_{\text{low}}$
staying close to the unstable manifold $U$.

The second stage is the most important one; that is where the critical
singularities are formed. Because at this stage the trajectory stays close to the
fixed point, we can use linear analysis. Let us assume that in some finite amount
$l_0$ of RG ``time'' the trajectory $C_\tau$ comes close to $\mathcal A_*$, where
the corresponding action can be written as
\begin{equation}\label{eq:nearfp}
\mathcal A=\mathcal A_*+y_0\int d^dx\,\Phi_0(x)+\sum_{\alpha\neq0}y_\alpha\int d^dx\,\Phi_\alpha(x)\,,
\end{equation}
where the constants $y$ depend on $\tau$, the deviation \eqref{eq:tau} from the
critical point. It is important to note that this dependence is analytic because
it comes from integrating the RG flow equations for finite ``time'' $l_0$. As for
$\tau=0$ one must stay at the critical surface, which in the linear approximation
corresponds to $y_0=0$, for small $\tau$ we can assume
\begin{equation}\label{eq:y0T}
y_0=\Const\times(T_c-T)\,,
\end{equation}
and neglect the $\tau$-dependence of all $y_\alpha$ with $\alpha\neq0$.

In the second stage the RG flow is described by the equations
\[
\frac d{dl}\,y_0=\kappa_0\,y_0\,,\qquad\frac d{dl}\,y_\alpha=\kappa_\alpha\,y_\alpha\,,
\]
where the coefficients $\kappa$ are related to the dimensions $D_\alpha$ of the
fields $\Phi_\alpha$,
\[
\kappa_0=d-D_0>0\,,\qquad\kappa_\alpha=d-D_\alpha<0\quad\text{for}\quad\alpha\neq0\,.
\]
During this second stage all irrelevant terms in \eqref{eq:nearfp} die out, while
the coefficient $y_0$ increases as
\begin{equation}\label{eq:y0l}
y_0(l)=y_0\,e^{\kappa_0(l-l_0)}\,.
\end{equation}
Recall that if $R_c$ is the correlation length corresponding to the initial action
on the physical curve, the correlation length corresponding to the time $l$ along
the trajectory is
\[
R_c(l)=R_c\,e^{-l}\,.
\]
From this and \eqref{eq:y0l} one can find\footnote{Observe that within the linear
approximation \eqref{eq:nearfp} the ratio $R_c(l)/y_0^{-\nu}(l)$ is a constant
independent of $l$.}
\begin{equation}\label{eq:Rcnu}
R_c\sim y_0^{-\nu}\,,\qquad\text{where}\qquad\nu=\frac1{\kappa_0}=\frac1{d-D_0}\,,
\end{equation}
and from the linear dependence \eqref{eq:y0T}
\[
R_c\sim|T_c-T|^{-\nu}\,.
\]
Thus, the correlation length diverges near $T_c$, with the exponent determined by
the dimension $D_0$ of the relevant field $\Phi_0$. This divergence is in full
agreement with the idea that it is the large-scale fluctuations which are
responsible for the critical behavior.

Similar way one can relate other critical exponents to the dimensions of various
fields at the fixed point $\mathcal A_*$. For example, the spontaneous
magnetization
\[
\bar M\sim\langle\varphi_0\rangle
\]
(nonzero in the low-$T$ phase) must satisfy the CS equation, which in this case
reads
\[
\left[D_\varphi-\kappa_0\,y_0\,\frac\partial{\partial y_0}\right]\bar M=0\,,
\]
i.e.
\begin{equation}\label{eq:Mbar}
\bar M\sim y_0^{\frac{D_\varphi}{\kappa_0}}\,.
\end{equation}
Again, taking into account \eqref{eq:y0T} we find (for $T<T_c$)
\[
\bar M\sim(T_c-T)^\beta\,,\qquad\text{where}\qquad\beta=\frac{D_\varphi}{d-D_0}\,.
\]
It is useful to note that the relations \eqref{eq:Rcnu}, \eqref{eq:Mbar} can be
obtained by simple dimensional counting, if one assumes that the coupling constant
$y_0$ has the dimension $\kappa_0$ and the field $\varphi_0$ has the dimension
$D_\varphi$. This follows from the CS equation with the linear approximation for
the RG functions $\beta$; in fact it is valid beyond the linear approximation, as
it follows from the existence of ``canonical coordinates''.

As another example of this dimensional counting let us find some other critical
exponents in terms of dimensions. If, at $T=T_c$, one adds an external field,
i.e.\ the term
\[
H\int d^dx\,\varphi_0(x)\,,
\]
the dimensional counting implies that $H$ has dimension $d-D_\varphi$, and so the
magnetization as the function of $H$ must have the following singular behavior
\[
\bar M\sim H^{\frac1\delta}\,,\qquad\text{where}\qquad\delta=\frac{d-D_\varphi}{D_\varphi}\,.
\]
Similar way one finds for the specific heat at $H=0$
\[
C\sim|T_c-T|^{-\alpha}
\]
with
\begin{equation}\label{eq:alpha}
\alpha=2-d\,\nu\,.
\end{equation}
Note that all these exponents are expressed through just two numbers, the
dimensions of the fields $\varphi$ and $\Phi_0$. Therefore no matter what the
values of these dimensions are, the critical exponents must obey certain
\emph{scaling relations}, the relation \eqref{eq:alpha} being one example. These
relations are confirmed in experiment.

Note that the above analysis relies on the properties of the fixed point
$\mathcal A_*$, as well as on general topological properties of RG. Very little
indeed depends on the details of the original actions on the ``physical curve'',
and hence on the details of the microscopic structure of the magnet. One expects
therefore that the most important characteristics of the critical singularities
must not be different for phase transitions in different magnetic materials (with
the same symmetry). One fixed point may represent the criticality in the whole
class of physical systems. It is said it describes a \emph{universality class} of
critical behavior. In fact, systems of very different physical nature can belong
to the same universality class. For example, the liquid-gas critical point is
likely to represent the same universality class as the phase transition in axial
ferromagnet.

This all shows that understanding the fixed points is of crucial importance both
in statistical mechanics and in quantum field theory. Unfortunately, analysis and
classification of all fixed points are not yet achieved. In fact, at present very
few nontrivial fixed points can be analyzed on a systematic level.

% Source pp. 214-226 (Lectures19-25.pdf).
\section{Wilson-Fisher Fixed Point}

Previously, we have studied the simplest example of fixed point, the Gaussian
fixed point corresponding to free massless scalar field. It is easy to argue,
however, that the Gaussian fixed point cannot be responsible for the ordinary
transition in a magnet in three dimensions (the Gaussian fixed point describes the
so called ``tricritical behavior'', as we will briefly discuss later). Indeed, it
was very important for the arguments given in the previous lecture that the
critical surface of $\mathcal A_*$ has codimension one in $\Sigma_{\text{symm}}$,
i.e.\ that there is only one relevant field which respects the symmetry of the
problem ($\varphi\to-\varphi$ in our case). However, as we have seen, the gaussian
fixed point has more than one symmetric relevant field for all $d<4$. In
particular, for $d=3$ two fields, $\Phi_2\sim\varphi^2$ and $\Phi_4\sim\varphi^4$,
are relevant. Therefore there must be another fixed point describing the
universality class of the magnetic phase transition.

The best theoretical evidence for the existence of another nontrivial fixed point
at $d=3$ (on a formal level, at all $2<d<4$) was given by Wilson and Fisher, and
it is usually referred to as the Wilson-Fisher fixed point $\mathcal A_{*WF}$.

At $d<4$ the topology of the RG flow relating the fixed points $\mathcal A_{*G}$
and $\mathcal A_{*WF}$ is likely to be the following. The unstable manifold
$U(\mathcal A_{*G})$ at $d<4$ has dimensionality two; it is generated by the two
relevant fields $\Phi_2$ and $\Phi_4$ (again, we restrict attention to the
symmetric space $\Sigma_{\text{symm}}$). The fixed point $\mathcal A_{*WF}$ is
likely to lie in $U(\mathcal A_{*G})$, so that the RG flow on this unstable
manifold would look qualitatively as shown in figure~\ref{fig:wfflow}.
\begin{figure}[ht]
\centering\figfile{wfflow}
\caption{RG flow on the unstable manifold $U(\mathcal A_{*G})$ at $d<4$. The
trajectory $C$ along the critical surface $\Sigma_c$ ($M^2=0$) connects the
Gaussian fixed point to the Wilson-Fisher fixed point.}
\label{fig:wfflow}
\end{figure}

According to this picture there is a trajectory $C$ connecting the fixed point
$\mathcal A_{*G}$ to $\mathcal A_{*WF}$. Therefore $\mathcal A_{*G}$ lies on the
critical surface $\Sigma_c$ which also passes through $\mathcal A_{*WF}$. The
unstable manifold $U(\mathcal A_{*WF})$ is then one-dimensional, as expected.

Note that the whole unstable manifold $U(\mathcal A_{*G})$ in this picture, being
an element of $\Sigma(\infty)$, represents local field theories. Namely, it
represents two-parameter family of $\varphi^4$ field theories. These field
theories can be characterized by two parameters, the renormalized mass parameter
$M^2$ and the renormalized coupling constant $\lambda$. Therefore these parameters
can be taken for coordinates in the above picture. Although there is an ambiguity
in the choice of these coordinates (the scheme dependence), the critical surface
corresponding to $M^2=0$ is defined by the scheme-independent condition
\begin{equation}\label{eq:masslesscond}
\Gamma^{(2)}(p^2)\big|_{p^2=0}=0\,.
\end{equation}
In these coordinates the fixed point $\mathcal A_{*WF}$ lies at $M^2=0$ and
$\lambda=\lambda_*$, where the value of $\lambda_*$ depends on the scheme.

We can formally consider the $\varphi^4$ theory in $d$ dimensional space, treating
$d$ as continuous parameter. What happens when $d$ approaches 4? We know that at
$d=4$ the field $\varphi^4$ becomes marginally irrelevant, and for $d>4$ it is
strictly irrelevant, and so for all $d\ge4$ the gaussian fixed point has only one
relevant field $\varphi^2$. This behavior emerges naturally if one assumes that
when $d$ approaches the value 4 from below the fixed point $\mathcal A_{*WF}$ gets
very close to the gaussian fixed point, and at $d=4$ these two fixed points merge.
At $d>4$ $\lambda_*$ becomes negative, and it does not correspond to field theory
in usual sense.\footnote{The ``magnetic criticality'' at $d\ge4$ then is described
by the gaussian fixed point.} This idea leads to the patterns of the RG flow at
$d=4$ and at $d>4$ shown in figure~\ref{fig:wfmerge}.
\begin{figure}[ht]
\centering\figfile{wfmerge}
\caption{The RG flow near the Gaussian fixed point at $d=4$ (left), where
$\varphi^4$ is marginally irrelevant, and at $d>4$ (right), where the
Wilson-Fisher fixed point has moved to negative $\lambda$.}
\label{fig:wfmerge}
\end{figure}

Note that the pattern at $d=4$ is exactly what one expects to see when $\varphi^4$
is marginally irrelevant, in which case we know that the flow in the vicinity of
the gaussian fixed point is described by the equations
\begin{align*}
\frac{\dd}{\dd l}\,\lambda&=-\beta(\lambda)=-3\,\frac{\lambda^2}{(4\pi)^2}+O(\lambda^3)\,,\\
\frac{\dd}{\dd l}\,M^2&=\bigl(2-\gamma_{\varphi^2}(\lambda)\bigr)\,M^2=\left(2-\frac\lambda{(4\pi)^2}+O(\lambda^2)\right)\,M^2\,.
\end{align*}
On the basis of these arguments Wilson and Fisher have assumed that if one
formally takes $d=4-\epsilon$ with $\epsilon$ small, say $d=3.99$, the fixed point
$\mathcal A_{*WF}$ can be observed within the framework of renormalized
perturbation theory in $\varphi^4$ field theory.

Let us briefly discuss what changes in the renormalized perturbation theory of
$\varphi^4$ if one takes $d=4-\epsilon$ instead of $d=4$. As before, we start with
the action in its ``bare'' form
\[
\mathcal A=\int \dd^dx\,\left[\frac12\,(\partial\varphi_0)^2+\frac{m_0^2}2\,\varphi_0^2+\frac{\lambda_0}{4!}\,\varphi_0^4\right]
\]
Note that in $d$ dimensions the canonical dimension of $\varphi_0$ is $(d-2)/2$
and hence coupling constant $\lambda_0$ has positive mass dimension,
\[
[\varphi_0]=\frac{d-2}2=1-\frac\epsilon2\,;\qquad[\lambda_0]=4-d=\epsilon\,.
\]
To develop renormalized perturbation theory we rewrite this action in terms of the
renormalized field and renormalized parameters,
\begin{equation}\label{eq:Aren21}
\mathcal A=\int \dd^dx\,\left[\frac12\,(\partial\varphi)^2+\frac{M^2}2\,\varphi^2+\frac\lambda{4!}\,\varphi^4
+\frac{\delta Z}2\,(\partial\varphi)^2+\frac{\delta M^2}2\,\varphi^2+\frac{\delta\lambda}{4!}\,\varphi^4\right]\,,
\end{equation}
where the counterterms are determined from suitably chosen normalization
conditions.

It is worth stressing here that for $d<4$ the diagrams describing the field and
the coupling constant renormalizations do not diverge; the only divergent
renormalization is the one of the mass parameter. Therefore the counterterms
$\delta Z$ and $\delta\lambda$ have finite limit when the cutoff parameter
$\Lambda$ goes to infinity. Therefore at $d<4$ these renormalizations are
introduced not for the purpose of eliminating the divergences, which are absent,
but in order to introduce more suitable coordinate $\lambda$ on
$U(\mathcal A_{*G})$ instead of $\lambda_0$. Although for $d<4$ (unlike the case
of $d=4$) the bare coupling $\lambda_0$ can be used as local coordinate on this
unstable manifold, we will see soon that it is not too useful for the purpose of
studying the fixed point $\mathcal A_{*WF}$.

As the fixed point $\mathcal A_{*WF}$ is expected to be located on the critical
surface, it suffices to study the massless $\varphi^4$ theory which is defined by
the condition \eqref{eq:masslesscond},
\[
\Gamma^{(2)}(p^2)\big|_{p^2=0}=0\,.
\]
As before, choose some normalization scale $\mu$ and assume the following simple
normalization conditions
\begin{equation}\label{eq:norm21}
\frac{\dd}{\dd p^2}\,\Gamma^{(2)}(p^2)\Big|_{p^2=\mu^2}=1\,,\qquad
\Gamma^{(4)}(p_1,\dots,p_4)\Big|_{[(p_i+p_j)^2=4\mu^2]}=\lambda\,.
\end{equation}
The first of them fixes the normalization of $\varphi$ while the second gives the
definition of the renormalized coupling constant $\lambda$.

At $d<4$ the Callan-Symanzik equation for the massless theory is
\begin{equation}\label{eq:CS21}
\left[x\frac\partial{\partial x}+N\,D(\lambda)+\beta(\lambda)\,\frac\partial{\partial\lambda}\right]\bigl\langle\varphi(x_1)\cdots\varphi(x_N)\bigr\rangle_\lambda=0\,;
\end{equation}
its form reflects the general properties of the RG we sketched the last time,
which did not depend on $d$. However for generic $d$
\[
D(\lambda)=\frac{d-2}2+\gamma(\lambda)\,,\qquad\gamma(0)=0\,,
\]
because canonical dimension of $\varphi$ is $(d-2)/2$. Our nearest goal is to
determine the RG functions $\beta(\lambda)$ and $\gamma(\lambda)$ for $d<4$. The
first one is particularly important because the RG flow along the massless part of
$U(\mathcal A_{*G})$ (the trajectory $C$ in figure~\ref{fig:wfflow}) is described
by the equation
\[
\frac{\dd}{\dd l}\,\lambda=-\beta(\lambda)\,.
\]
We can calculate the correlation functions perturbatively, order by order in
$\lambda$, and then find $\beta$ and $\gamma$ from the CS equation \eqref{eq:CS21}
itself, as we did before in the case $d=4$. It is convenient to write the
renormalized coupling constant as
\[
\lambda=\mu^{4-d}\,\bar\lambda\,,
\]
where now $\bar\lambda$ is dimensionless parameter. Then by dimensional counting
the $N$-point correlation functions have the form
\[
\bigl\langle\varphi(x_1)\cdots\varphi(x_N)\bigr\rangle_\lambda=\mu^{N(d-2)/2}\,C^{(N)}(\mu x_1,\dots,\mu x_N|\bar\lambda)
\]
and therefore the CS equation can be rewritten as
\begin{equation}\label{eq:CSbar}
\left[\mu\frac\partial{\partial\mu}+N\,\bar\gamma(\bar\lambda)+\bar\beta(\bar\lambda)\,\frac\partial{\partial\bar\lambda}\right]\bigl\langle\varphi(x_1)\cdots\varphi(x_N)\bigr\rangle_{\bar\lambda}=0\,,
\end{equation}
where $\bar\beta$ and $\bar\gamma$ are the dimensionless functions of the
dimensionless parameter $\bar\lambda$ defined as
\[
\beta(\lambda)=\mu^\epsilon\,\bar\beta(\lambda/\mu^\epsilon)\,,\qquad\gamma(\lambda)=\bar\gamma(\lambda/\mu^\epsilon)\,,
\]
and the subscript $\bar\lambda$ indicates that the correlation function is
expressed through $\bar\lambda$ instead of the dimensional $\lambda$. In what
follows I will omit all these bars and use $\lambda$ for $\bar\lambda$, $\beta$
for $\bar\beta$, etc.; anyhow, these quantities coincide if $\mu$ is taken to be
the unit of mass. The equation \eqref{eq:CSbar} can also be rewritten in terms of
the momentum-space proper vertices,
\begin{equation}\label{eq:CSGamma21}
\left[\mu\frac\partial{\partial\mu}-N\,\gamma(\lambda)+\beta(\lambda)\,\frac\partial{\partial\lambda}\right]\Gamma^{(N)}(p_1,\dots,p_N)=0\,.
\end{equation}

We can now repeat our calculations of the leading perturbative contributions to
the proper vertices, which are close to identical to the calculations within the
dimensional regularization method, and obtain
\[
\Gamma^{(2)}(p^2)=p^2+O(\lambda^2)\,,
\]
and
\begin{multline}\label{eq:G4d}
\Gamma^{(4)}(p_1,\dots,p_4)=\mu^{4-d}\Bigl\{\lambda-\frac12\,\frac{\lambda^2}{(4\pi)^{d/2}}\,\frac{\Gamma(2-d/2)\,\Gamma^2(d/2-1)}{\Gamma(d-2)}\\
\times\left[\left(\frac{p_{12}^2}{\mu^2}\right)^{(d-4)/2}+\left(\frac{p_{13}^2}{\mu^2}\right)^{(d-4)/2}+\left(\frac{p_{14}^2}{\mu^2}\right)^{(d-4)/2}-3\,\bigl(2^{d-4}\bigr)\right]+O(\lambda^3)\Bigr\}\,.
\end{multline}
In the last equation the counterterm contribution is fixed by the normalization
condition \eqref{eq:norm21}, i.e.
\[
\Gamma^{(4)}(p_1,\dots,p_4)\big|_{[(p_i+p_j)^2=4\mu^2]}=\mu^{4-d}\,\lambda\,.
\]
Let me remind you that in dimensional regularization of $d=4$ theory the
counterterm contribution was introduced to cancel the $d=4$ singularity of the
Gamma-function $\Gamma(2-d/2)\sim2/(4-d)$. For $d<4$ the loop diagram contributions
are finite even without any counterterm. We can expand in terms of the bare
coupling constant $\lambda_0$ instead of the renormalized one; in this expansion
the counterterm contribution in \eqref{eq:G4d} would be absent. It is instructive
to remember however that when $\epsilon=4-d$ is small the coefficients of this
expansion in terms of the bare coupling constant are large,
\begin{equation}\label{eq:G4bare}
\Gamma^{(4)}(p_1,\dots,p_4)\sim\lambda_0-\frac1\epsilon\,\frac{\lambda_0^2}{(4\pi)^2}\Bigl((p_{12}^2)^{-\epsilon/2}+(p_{13}^2)^{-\epsilon/2}+(p_{14}^2)^{-\epsilon/2}\Bigr)+\cdots\,.
\end{equation}
At $d=4-\epsilon<4$ the ``coordinate transformation'' from $\lambda_0$ to $\lambda$
is nonsingular in some neighborhood of 0, but at small $\epsilon$ it is ``almost
singular''.\footnote{The radius of convergence of the series $\lambda_0(\lambda)$
is $\sim\epsilon$.}

Substituting \eqref{eq:G4d} into \eqref{eq:CSGamma21} we find
$\gamma(\lambda)=O(\lambda^2)$ and
\[
\beta(\lambda)=(d-4)\,\lambda+3\cdot2^{d-4}\,\frac{\lambda^2}{(4\pi)^{d/2}}\,\frac{\Gamma(3-d/2)\,\Gamma^2(d/2-1)}{\Gamma(d-2)}+O(\lambda^3)\,.
\]
The term $(d-4)\lambda$ appeared because of the factor $\mu^{4-d}$; it represents
the canonical dimension of the coupling constant $\lambda$ in the action
\eqref{eq:Aren21}.

Let us assume now that $d=4-\epsilon$ where $\epsilon$ is small,
\[
\epsilon\ll1\,.
\]
Then
\[
\beta(\lambda)=-\epsilon\,\lambda+3\,\frac{\lambda^2}{(4\pi)^2}+O(\lambda^3)\,.
\]
The linear term $-\epsilon\lambda$ brings in an important difference from the case
$d=4$. For small $\epsilon$ the RG flow equation
\[
\frac{\dd}{\dd l}\,\lambda=-\beta(\lambda)
\]
has nontrivial fixed point
\[
\lambda_*=\frac{(4\pi)^2}3\,\epsilon\,.
\]
The RG flow looks like this:
\[
\diac[0.6]{\draw[line width=0.8pt] (-2.5,0) -- (6.5,0) node[right] {$\lambda$};
  \draw[arr] (-0.8,0) -- (-2.2,0); \draw[arr] (0.8,0) -- (2.4,0); \draw[arr] (5.8,0) -- (4.6,0);
  \fill (0,0) circle (2.4pt) node[above=2pt] {$0$}; \fill (3.6,0) circle (2.4pt) node[above=2pt] {$\lambda_*$};}
\]
This fixed point was obtained in the leading order of the renormalized
perturbation theory. Note however that for small $\epsilon$ the fixed point value
$\lambda_*$ is small. The higher order terms can shift its position by terms
$\sim\epsilon^2$,
\begin{equation}\label{eq:lambdastar}
\lambda_*=\frac{(4\pi)^2}3\,\epsilon+O(\epsilon^2)\,,
\end{equation}
but for sufficiently small $\epsilon$ these corrections cannot destroy the fixed
point. The fixed point \eqref{eq:lambdastar} is the Wilson-Fisher fixed point.

The above result was obtained using renormalized perturbation theory for $d<4$
$\varphi^4$ theory. As I mentioned, for $d<4$ the unrenormalized perturbation
theory, i.e.\ the expansion in terms of the bare coupling constant $\lambda_0$, is
(essentially) finite. It is instructive to use this example to explain the
relation between renormalized and unrenormalized perturbation theories.

If we repeat our calculation of the $\beta$-function using unrenormalized
perturbation theory, i.e.\ with the expression \eqref{eq:G4bare} for the
four-point vertex, we would obtain from \eqref{eq:CSGamma21}
\begin{equation}\label{eq:beta0}
\beta_0(\lambda_0)=-\epsilon\,\lambda_0\,,
\end{equation}
with no fixed point visible. Higher order corrections do not change this result;
in fact, it is easy to see that in unrenormalized perturbation theory, without the
counterterm contributions, the formula \eqref{eq:beta0} is exact. If one uses
$\lambda_0$ as the coordinate on $U(\mathcal A_{*G})$, the RG function $\beta$
reduces to its linear part. The bare coupling constant $\lambda_0$ is the
``canonical coordinate'', whose existence is guaranteed by the Poincar\'e-Dulac
theorem.

Indeed, in the case $\epsilon\ll1$ the RG flow equation for small values of the
renormalized parameter $\lambda$ has the form
\begin{equation}\label{eq:logistic}
\frac{(4\pi)^2}3\,\frac{\dd}{\dd l}\,\lambda=\lambda\,(\lambda_*-\lambda)\,.
\end{equation}
It is not difficult to check that the change of variables
\begin{equation}\label{eq:lambda0map}
\lambda=\frac{\lambda_0}{1+\frac{\lambda_0}{\lambda_*}}
\end{equation}
linearizes this flow equation,
\begin{equation}\label{eq:lambda0flow}
\frac{\dd}{\dd l}\,\lambda_0=\epsilon\,\lambda_0\,.
\end{equation}
Note that the transformation \eqref{eq:lambda0map}, or equivalently
\[
\frac1\lambda=\frac1{\lambda_0}+\frac1{\lambda_*}\,,
\]
maps the interval $0<\lambda<\lambda_*$ to the whole positive half-axis of
$\lambda_0$,
\[
\lambda\in[0,\lambda_*)\quad\leftrightarrow\quad\lambda_0\in[0,\infty)\,,
\]
with the fixed point $\lambda_*$ mapped to the point $\lambda_0=\infty$; for this
reason it is invisible in \eqref{eq:lambda0flow}. Note also that if one expands
$\lambda$ in terms of $\lambda_0$ according to \eqref{eq:lambda0map}, for small
$\epsilon$ the coefficients of this expansion become large,
\[
\lambda=\lambda_0-\frac{\lambda_0^2}{\lambda_*}+\cdots\,,
\]
which explains the large coefficients in \eqref{eq:G4bare}. For $\epsilon=0$
(i.e.\ $d=4$) the interval $[0,\lambda_*)$ shrinks to zero, the perturbation
becomes marginal, and the ``canonical coordinate'' $\lambda_0$ ceases to exist.

If the renormalized parameter $\lambda$ is used, for small $\epsilon$ the flow
between $\mathcal A_{*G}$ and $\mathcal A_{*WF}$, as well as the fixed point
$\mathcal A_{*WF}$ itself, remains at small $\lambda\sim\epsilon$, within the
domain where the perturbative results are reliable. Therefore for small
$\epsilon$ we can use the renormalized perturbation theory to calculate all
characteristics of the Wilson-Fisher fixed point. For instance, to the leading
order in $\epsilon$ the flow equation for the mass parameter $M^2$ has the same
form as at $d=4$,
\begin{equation}\label{eq:M2flow21}
\frac{\dd}{\dd l}\,M^2=\left(2-\frac\lambda{(4\pi)^2}+O(\epsilon^2)\right)\,M^2\,.
\end{equation}
Together with \eqref{eq:logistic} this equation yields the pattern of the flow
shown in figure~\ref{fig:wfflow}.

\subsection*{\texorpdfstring{$\epsilon$}{epsilon}-expansion}

The main assumption of $\epsilon$-expansion is that the general features of the
RG flow observed at small $\epsilon$ remain qualitatively the same at
$\epsilon=1$, i.e.\ at $d=3$. Under this assumption, the characteristics of the
Wilson-Fisher fixed point can be calculated perturbatively, as the power series in
$\epsilon$. Although at $d=3$ the expansion parameter $\epsilon$ is not small,
this approach yields remarkably accurate results for the critical exponents in
three dimensions.

To find the critical exponents one must calculate the dimensions of various
relevant fields at the Wilson-Fisher fixed point. Note that $\Phi_0=\varphi^2$ is
the only symmetric relevant field at $\mathcal A_{*WF}$. The dimensions of the
fields $\Phi_0$ and $\varphi$ at this point can be calculated as
\[
D_0=D_{\varphi^2}=d-2+\gamma_{\varphi^2}(\lambda_*)\,,\qquad
D_\varphi=\frac{d-2}2+\gamma(\lambda_*)\,.
\]
In the leading order in $\epsilon$ we can just borrow the functions $\gamma$ from
the $d=4$ calculation,
\[
\gamma_{\varphi^2}(\lambda)=\frac\lambda{(4\pi)^2}+O(\lambda^2)\,,\qquad
\gamma(\lambda)=\frac1{12}\,\frac{\lambda^2}{(4\pi)^4}+O(\lambda^3)\,,
\]
and use the value \eqref{eq:lambdastar} for $\lambda_*$. This yields for the
critical exponents
\[
\frac1\nu=d-D_0=2-\frac\epsilon3+O(\epsilon^2)\,,\qquad
\alpha=2-\frac d{d-D_0}=\frac\epsilon6+O(\epsilon^2)\,,
\]
and similar results for other exponents. For $\epsilon=1$ this gives $\nu=0.60$
and $\alpha=0.16$ for the correlation length and heat capacity exponents, to be
compared to $\nu=0.63$ and $\alpha=0.12$ found in experiment.\footnote{The
definition is $R_c\sim|T_c-T|^{-\nu}$, $C\sim|T_c-T|^{-\alpha}$.}

\subsection*{\texorpdfstring{$O(N)$ $\varphi^4$}{O(N) phi**4} theory}

Let us discuss somewhat more general $\varphi^4$ theory. Assume that the scalar
field has $N$ components $\bm{\varphi}=(\varphi_1,\dots,\varphi_N)$, and the
$\varphi^4$ interaction respects the $O(N)$ transformations of the vector
$\bm{\varphi}$. We have
\[
\mathcal A=\int \dd^dx\,\left[\frac12\,(\partial\bm{\varphi})^2+\frac{M^2}2\,\bm{\varphi}^2+\frac\lambda{4!}\,(\bm{\varphi}^2)^2\right]
\]
Although, as before, we use the notation $M^2/2$ for the coefficient in front of
the term $\bm{\varphi}^2$, this coefficient is allowed to take negative values.

At small $d=4-\epsilon$ with small $\epsilon$ the pattern of the RG flow of the
$O(N)$ $\varphi^4$ theory in the vicinity of its Gaussian fixed point is
essentially the same as in the one-component case we discussed
previously\footnote{Explicit calculation (Problem~\ref{hw:8}) yields, in the
normalization $\frac\lambda{4!}(\bm{\varphi}^2)^2$ used here,
$\beta(\lambda)=-\epsilon\lambda+\frac{N+8}{48\pi^2}\,\lambda^2+\cdots$,
$\gamma_{\varphi^2}(\lambda)=\frac{N+2}{48\pi^2}\,\lambda+\cdots$.}
(figure~\ref{fig:onflow}).
\begin{figure}[ht]
\centering\figfile{onflow}
\caption{RG flow of the $O(N)$ $\varphi^4$ theory near $d=4$. The axis $M^2=0$ is
the critical surface; above it the flow goes up, below it down.}
\label{fig:onflow}
\end{figure}

There is the $O(N)$ version of the Wilson-Fisher fixed point
$\mathcal A_{*,WF}\in\Sigma_c$ at $M^2=0$, $\lambda=\lambda_*\sim\epsilon$. In the
space of $O(N)$ invariant interactions, it has one-dimensional unstable manifold
$U(\mathcal A_{*,WF})$.

In figure~\ref{fig:onflow} the axis $M^2=0$ is the critical surface $\Sigma_c$.
Above it, the flow generally goes upward, while everything below it flows
downward. One can try to figure out the ultimate fate of the flow at positive and
negative $M^2$ on physical grounds. The domain $M^2>0$ corresponds to the
high-temperature phase of the ferromagnetic system. The pattern of the flow shows
that $M^2$ increases with $l$. At large positive $M^2$ the term
$\frac{M^2}2\bm{\varphi}^2$ suppresses fluctuations of $\bm{\varphi}$. Therefore
at $M^2>0$ one expects to have attractive non-critical fixed point
$P_{\text{high}}$, corresponding to the probability distribution
\[
P_{\text{high}}[\bm{\varphi}]\sim\prod_x\delta\bigl(\bm{\varphi}(x)\bigr)\,.
\]
The domain of negative $M^2$ is associated with the low-temperature phase. Here
under the RG flow $M^2$ remains negative while increasing in the absolute value. In
this domain the potential term
$\frac{M^2}2\bm{\varphi}^2+\frac\lambda{4!}(\bm{\varphi}^2)^2$, as the
function of $\bm{\varphi}^2$, has the minimum away from zero
(figure~\ref{fig:mexican}).
\begin{figure}[ht]
\centering\figfile{mexican}
\caption{The potential $V$ as a function of $|\bm{\varphi}|$ for $M^2<0$.}
\label{fig:mexican}
\end{figure}

In terms of $\bm{\varphi}$ we have manifold of degenerate minima, the sphere
$S^{N-1}$, of the radius
\[
|\bm{\varphi}|=\rho:=\sqrt{\frac{6\,|M^2|}\lambda}\,.
\]
Points of the sphere indicate possible directions of $\bm{\varphi}$ at the minimum.
It is easy to see that if $|M^2|$ is large the minimum in figure~\ref{fig:mexican}
becomes deep and narrow. Fluctuations of the radial part, $|\bm{\varphi}|$, are
then suppressed. On the other hand there is no suppression for the fluctuations of
the angular part. Writing
\[
\bm{\varphi}=\rho\,\mathbf n\,,\qquad|\mathbf n|=1\,.
\]
If we completely ignore the radial fluctuations, i.e.\ set $|\bm{\varphi}|=\rho$
everywhere, the action takes simple form
\[
\mathcal A=\int \dd^dx\,\left[\frac1{2g}\,(\partial\mathbf n)^2\right]
\qquad\text{with}\qquad
g=\frac1{\rho^2}=\frac\lambda{6\,|M^2|}\,.
\]
As $|M^2|$ increases with the RG flow, $g$ becomes small.\footnote{This simple
classical analysis is valid only in the limit $\lambda\to0$. More generally,
fluctuations of $|\bm{\varphi}|$ can be taken into account, along with the
fluctuations of the angular part, by renormalized perturbation theory. Some
details about this theory, the so called linear sigma model, can be found in
Chapter 11 of PS.} Due to the non-linear constraint $\mathbf n^2=1$ the above
action leads to a theory more complicated than a free massless field. It is named
the non-linear sigma model.

Let us analyze physical content of the nonlinear sigma-model. Classically, the
most favorable field configurations correspond to minima of the action, which in
this case are constant ($x$-independent) fields
\[
\mathbf n(x)=\mathbf n_0\,.
\]
By the $O(N)$ symmetry we can choose
\[
\mathbf n_0=(0,0,\dots,0,1)
\]
with no loss of generality. To take into account fluctuations around this minimum
one can set
\[
\mathbf n(x)=\bigl(\pi^1(x),\pi^2(x),\dots,\pi^{N-1}(x),\sigma(x)\bigr)\equiv\bigl(\pi^k(x),\sigma(x)\bigr)\qquad(k=1,2,\dots,N-1)
\]
with
\[
\sigma(x)=\sqrt{1-(\pi^k(x))^2}
\]
The action takes the form
\[
\mathcal A=\int \dd^dx\,\frac1{2g}\left[(\partial_\mu\pi^k)^2+\frac{(\pi^k\partial_\mu\pi^k)^2}{1-(\pi^k)^2}\right]
\]
Ignoring first the quartic and higher terms in this action, we observe $N-1$
massless scalar fields $\pi^k(x)$. These are the \emph{Goldstone bosons},
associated with the spontaneous breakdown of the $O(N)$ symmetry, down to
$O(N-1)$ (generalities about the spontaneous symmetry breaking are in Chapter 11
of PS).

Rescaling
\[
\pi^k=\sqrt g\,\phi^k
\]
one can write the above action as
\[
\mathcal A=\int \dd^dx\,\left[\frac12\,(\partial\phi^k)^2+\sum_{n=1}^\infty g^n\,O_n(x)\right]\,,
\qquad\text{where}\qquad
O_n=\frac12\,(\phi^k\partial_\mu\phi^k)^2\,(\phi^k\phi^k)^{n-1}\,.
\]
If $d>2$, it is correct to think of it as $N-1$-component of massless free field
(the first term), with some composite fields added. The dimensions of the fields
$O_n$ are
\[
O_n:\qquad D_n=2+(n+1)(d-2)=d+n\,(d-2)\,.
\]
At $d>2$ all these fields are irrelevant.

By these arguments, at $d>2$ the long-distance physics in the low-temperature
domain $M^2<0$ is controlled by the action
\[
\mathcal A_{*,GS}=\int\frac12\,(\partial\phi^k)^2\dd^dx
\]
which is essentially a Gaussian fixed point of $N-1$ massless bosons. I will call
it the ``Goldstone point'', to distinguish it from the original Gaussian fixed
point $\mathcal A_{*,G}$ (which has $N$ massless scalar fields).

Returning to the pattern of the RG flow on $U(\mathcal A_{*,G})$, we now anticipate
the global structure shown in figure~\ref{fig:onglobal}.
\begin{figure}[ht]
\centering\figfile{onglobal}
\caption{Anticipated global RG flow on $U(\mathcal A_{*,G})$ for $N>1$: the flow
above the critical surface ends at $P_{\text{high}}$, and the flow below it at the
critical Goldstone point $\mathcal A_{*,GS}$.}
\label{fig:onglobal}
\end{figure}
For $N>1$ there is the \emph{critical} fixed point $\mathcal A_{*,GS}$, which
attracts the trajectories in the low-temperature domain $M^2<0$.

At $d=2$ we have $D_n=2=d$, that is all the above fields $O_n$ become marginal.
Obviously, something interesting is happening in this case. Let us analyze it
separately.

\begin{hwproblem}{8}\label{hw:8}
Consider the same $O(N)$ $(\bm{\varphi})^4$ theory as in the HW
Problem~\ref{hw:6}, but this time in $d=4-\epsilon$ space, with $\epsilon\ll1$.

Find the Wilson-Fisher fixed point $\lambda_*$ to the leading order in $\epsilon$,
and evaluate the anomalous dimensions of the fields $(\bm{\varphi})^2$ and
$\bm{\varphi}$ to the order $\sim\epsilon$, at the WF fixed point. Find the critical
exponents $\nu$ in $R_c\sim|T_c-T|^{-\nu}$, and $\alpha$ in
$C\sim|T_c-T|^{-\alpha}$.
\end{hwproblem}

% Source pp. 227-238 (Lectures19-25.pdf).
\section{Nonlinear Sigma Model at \texorpdfstring{$d=2+\epsilon$}{d=2+epsilon}}

Now I would like to discuss how the RG transformation acts on the action of the
$\sigma$-model at $d=2$. In this case the coupling constant $g$ in
\[
\mathcal A=\int d^2x\,\frac1{2g}\,(\partial\mathbf n)^2
\]
is dimensionless. I assume that $g$ is small, so that the fluctuations of
$\mathbf n(x)$ are small, since at small $g$ the exponential factor
$\exp(-\mathcal A)$ strongly suppresses configurations having large gradients
$\partial\mathbf n$. I remind you that the field $\mathbf n(x)$ satisfies the
non-linear constraint $\mathbf n^2(x)=1$.

In principle, it is not difficult to develop renormalized perturbation theory in
the coupling constant $g$, derive the CS equation and find the series expansion of
associated flow function $\beta(g)$. This calculation is displayed in \S\,13.3 of
PS. I will discuss the derivation based explicitly on the idea of integrating out
the fast modes (also described in \S\,13.3 of PS), because it has some novel
features as compared to what we have done before.

First, we introduce cutoff,
\[
\mathbf n(x)=\int_{|k|<\Lambda}\mathbf n_k\,e^{ikx}\,\frac{d^2k}{(2\pi)^2}\,.
\]
Note that now the modes $\mathbf n_k$ are not all independent, but must obey
complicated constraint $\int_q\mathbf n_q\mathbf n_{k-q}=\delta^{(2)}(k)$, because
of the $\sigma$-model constraint $\mathbf n^2(x)=1$.

Next, we want to split $\mathbf n(x)$ into its ``slow'' and ``fast'' components.
The above nonlinear constraint makes this task less straightforward than we had it
before in usual scalar theories. Physically transparent way to make the splitting
is as follows. We define the ``slow'' component $\mathbf n_0(x)$ as the field of
unit length,
\begin{equation}\label{eq:n0unit}
\mathbf n_0^2(x)=1\,,
\end{equation}
whose Fourier transform is restricted to $|k|<\Lambda/L$. To make further
arguments transparent, I will assume that $L\gg1$; we will see that in this case
the dominating contribution to the RG transformed action comes from the
integration over the fast modes with $\Lambda/L\ll|k|\ll\Lambda$. More precisely,
I will assume that $L\gg1$, but $g$ is so small that
\[
g\log L\ll1\,.
\]
At each spatial point $x$, one can choose a system of vectors $\mathbf e_i(x)$,
$i=1,2,\dots,N-1$, such that $N$ vectors $(\mathbf n_0(x),\mathbf e_i(x))$ form a
basis. It is convenient to take orthonormal basis,
\begin{equation}\label{eq:ebasis}
\bigl(\mathbf n_0(x)\,\mathbf e_i(x)\bigr)=0\,,\qquad\bigl(\mathbf e_i(x)\,\mathbf e_k(x)\bigr)=\delta_{ik}\,.
\end{equation}
Since $\mathbf e_i(x)$ only need to be orthogonal to the ``slow'' field
$\mathbf n_0(x)$, one can always choose them to be ``slow'' fields as well,
containing only Fourier modes with $|k|$ of the order of $\Lambda/L$ or smaller.
The full field $\mathbf n(x)$ can be written as
\begin{equation}\label{eq:nsplit}
\mathbf n(x)=\mathbf n_0(x)\,\sqrt{1-(\phi_i)^2(x)}+\mathbf e^i(x)\,\phi_i(x)\,.
\end{equation}
This expression automatically satisfies the constraint $\mathbf n^2(x)=1$. $N-1$
fields $\phi_i(x)$ represent the fast modes; they contain Fourier components with
$|k|$ between $\Lambda/L$ and $\Lambda$. It is important that when $g$ is small we
can assume $\phi_i(x)$ are small.

In order to compute the action, we need spatial derivatives of $\mathbf n(x)$.
These involve the derivatives of the vectors $\mathbf n_0(x)$ and $\mathbf e_i(x)$,
as well as of the fields $\phi_i$. One can expand them in the basic vectors as
\begin{align}
\partial_\mu\mathbf n_0(x)&=B_\mu^i(x)\,\mathbf e_i(x)\,,\label{eq:dn0}\\
\partial_\mu\mathbf e_i(x)&=A_\mu^{ik}(x)\,\mathbf e_k(x)-B_\mu^i(x)\,\mathbf n_0(x)\,,\label{eq:de}
\end{align}
where
\begin{align}
B_\mu^i(x)&=(\mathbf e_i\,\partial_\mu\mathbf n_0)=-(\mathbf n_0\,\partial_\mu\mathbf e_i)\,,\label{eq:Bdef}\\
A_\mu^{ik}(x)&=(\mathbf e_k\,\partial_\mu\mathbf e_i)=-(\mathbf e_i\,\partial_\mu\mathbf e_k)\,.\label{eq:Adef}
\end{align}
Note that the fields $A_\mu^{ik}$ and $B_\mu^i$ are ``slow''. In terms of these
variables the action is
\[
\mathcal A=\frac1{2g}\int d^2x\,\Bigl[\Bigl(\partial_\mu\sqrt{1-\phi^2}-B_\mu^i\partial_\mu\phi_i\Bigr)^2
+\Bigl(\partial_\mu\phi_i-A_\mu^{ik}\phi_k+B_\mu^i\sqrt{1-\phi^2}\Bigr)^2\Bigr]\,.
\]
For small $g$ one can neglect all terms beyond the terms quadratic in $\phi$,
\[
\mathcal A=\mathcal A[\mathbf n_0]+\tilde{\mathcal A}[\mathbf n_0,\phi_i]\,,
\]
where the first term comes from
\[
\frac1{2g}\int d^2x\,B_\mu^iB_\mu^i=\frac1{2g}\int d^2x\,(\partial_\mu\mathbf n_0)^2=\mathcal A[\mathbf n_0]\,,
\]
and
\[
\tilde{\mathcal A}=\frac1{2g}\int d^2x\,\Bigl[\bigl(\partial_\mu\phi_i-A_\mu^{ik}\phi_k\bigr)^2+B_\mu^iB_\mu^k\bigl(\phi_i\phi_k-\delta_{ik}\phi^2\bigr)\Bigr]\,.
\]
Our goal now is to integrate out the fast modes $\phi$. We treat the term
$(1/2g)\int(\partial\phi)^2$ as the unperturbed action, yielding the propagator
for the $\phi$-field
\[
i\ \dia{\draw[prop] (-1.6,0) -- (1.6,0);}\ k=\frac{g\,\delta_{ik}}{p^2}\,.
\]
The remaining terms are treated as perturbations, giving rise to vertices of
three types:
\begin{align*}
\diac[0.45]{\draw[prop,densely dashed] (-2,0) -- (0,0); \draw[prop] (0,0) -- (1.3,1.3) node[above right] {\small$j$}; \draw[prop] (0,0) -- (1.3,-1.3) node[below right] {\small$k$};
  \node[vtx,minimum size=4.5pt] at (0,0) {}; \node[left] at (-2,0) {\small$A$};
  \draw[arr] (0.95,1.35) -- (0.45,0.85); \node[font=\tiny] at (0.4,1.3) {$p_1$};
  \draw[arr] (0.95,-1.35) -- (0.45,-0.85); \node[font=\tiny] at (0.4,-1.3) {$p_2$};}
&=\frac{i\,A_\mu^{jk}}g\,(p_1-p_2)_\mu\,,\\[4pt]
\diac[0.45]{\draw[prop,densely dashed] (-2,0) -- (0,0); \draw[prop] (0,0) -- (1.3,1.3) node[above right] {\small$j$}; \draw[prop] (0,0) -- (1.3,-1.3) node[below right] {\small$k$};
  \node[vtx,minimum size=4.5pt] at (0,0) {}; \node[left] at (-2,0) {\small$A^2$};}
&=-\frac{A_\mu^{ij}A_\mu^{ik}}g\,,\\[4pt]
\diac[0.45]{\draw[prop,densely dashed] (-2,0) -- (0,0); \draw[prop] (0,0) -- (1.3,1.3) node[above right] {\small$j$}; \draw[prop] (0,0) -- (1.3,-1.3) node[below right] {\small$k$};
  \node[vtx,minimum size=4.5pt] at (0,0) {}; \node[left] at (-2,0) {\small$B^2$};}
&=\frac{\delta_{jk}\,B_\mu^iB_\mu^i-B_\mu^jB_\mu^k}g\,.
\end{align*}
As usual
\[
\mathcal A_1=\mathcal A-\text{sum of all connected diagrams}\,,
\]
with the dashed lines taken as the external legs.

The leading order in $g$ comes from the one-loop diagrams. The simplest (and the
most important) of those is
\begin{align*}
\dia[0.45]{\draw[prop,densely dashed] (-2,0) -- (0,0); \draw[prop] (0.9,0) circle (0.9); \node[vtx,minimum size=4.5pt] at (0,0) {}; \node[left] at (-2,0) {\small$B^2$};}
&=\frac{\delta_{jk}\,B_\mu^iB_\mu^i-B_\mu^jB_\mu^k}{2g}\int\frac{d^2p}{(2\pi)^2}\,\frac{g\,\delta_{jk}}{p^2}\\
&=\frac{(N-2)\,(B_\mu^i)^2}2\,\frac1{4\pi}\int_{\Lambda^2/L^2}^{\Lambda^2}\frac{dp^2}{p^2}=\frac{N-2}{4\pi}\,(\partial_\mu\mathbf n_0)^2\,\log L\,.
\end{align*}
Thus, this diagram brings in a term identical in the form to the original action,
i.e.\ it renormalizes the coupling parameter $g$.

The diagram
\[
\dia[0.45]{\draw[prop,densely dashed] (-2.6,0) -- (-1.2,0) (1.2,0) -- (2.6,0); \draw[prop] (0,0) ellipse (1.2 and 0.75);
  \node[vtx] at (-1.2,0) {}; \node[vtx] at (1.2,0) {}; \node[left] at (-2.6,0) {\small$B^2$}; \node[right] at (2.6,0) {\small$B^2$};}
\]
generates term $\sim B^4$ (with the coefficient $\sim1/\Lambda^2$); the field
$B^4$ has dimension $4>d$ and hence it is an irrelevant field; it can be neglected
in the leading approximation. Likewise, all the diagrams of the form
\[
\diac[0.45]{\draw[prop] (0,0) circle (1.1);
  \foreach \a in {180,120,60,0} {\node[vtx] at (\a:1.1) {}; \draw[prop,densely dashed] (\a:1.1) -- (\a:2.2); \node at (\a:2.6) {\scriptsize$B^2$};}
  \draw[densely dotted] (200:1.5) arc[start angle=200,end angle=340,radius=1.5];}
\]
can be ignored, for the same reason.

Let us take a look at the diagrams with $A$-vertices. Consider the following ones
\[
\dia[0.45]{\draw[prop,densely dashed] (-2,0) -- (0,0); \draw[prop] (0.9,0) circle (0.9); \node[vtx,minimum size=4.5pt] at (0,0) {}; \node[left] at (-2,0) {\small$A^2$};}\,,\qquad
\dia[0.45]{\draw[prop,densely dashed] (-2.6,0) -- (-1.2,0) (1.2,0) -- (2.6,0); \draw[prop] (0,0) ellipse (1.2 and 0.75);
  \node[vtx] at (-1.2,0) {}; \node[vtx] at (1.2,0) {}; \node[left] at (-2.6,0) {\small$A$}; \node[right] at (2.6,0) {\small$A$};}\,.
\]
Both these diagrams show logarithmic behavior,
\[
\int\frac{d^2k}{k^2}\sim\log L\,,\qquad\text{and}\qquad\int\frac{d^2k\,k^\mu k^\nu}{k^2\,k^2}\sim\log L\,.
\]
However, it is not difficult to show that these two logarithmic contributions in
fact cancel each other. One can do it by a direct calculation, but it is easy to
prove this statement on the basis of certain symmetry.

The $\log L$ of each of these two diagrams is of the form
\[
\log(L)\,A_\mu^{ij}A_\mu^{ij}\,,
\]
where, as before,
\[
A_\mu^{ij}=\mathbf e^j\,\partial_\mu\mathbf e^i\,.
\]
Recall that $\mathbf e^i(x)$ are arbitrary vectors tangent to $S^{N-1}$ at
$\mathbf n_0(x)$. They satisfy
\[
\mathbf e^i\mathbf n_0=0\,,\qquad\mathbf e^i\mathbf e^j=\delta^{ij}\,,
\]
but otherwise arbitrary. One can always choose them differently,
\[
\tilde{\mathbf e}^i(x)=\Omega_{ij}(x)\,\mathbf e^j(x)\,,
\]
where $\Omega_{ij}(x)$ is orthogonal $(N-1)\times(N-1)$ matrix,
\[
\Omega(x)\,\Omega^t(x)=I\,.
\]
$\Omega(x)$ is ``slow'' (for $\mathbf e^i(x)$ were assumed to be slow), but
otherwise arbitrary function of $x$. Under these local orthogonal transformations,
$A_\mu^{ij}(x)$ transforms as an $O(N-1)$ gauge potential,
\[
\tilde A_\mu(x)=\Omega(x)\,A_\mu(x)\,\Omega^t(x)+\Omega(x)\,\partial_\mu\Omega^t(x)\,,
\]
where $A_\mu(x)$ stands for $(N-1)\times(N-1)$ anti-symmetric matrix
$A_\mu^{ij}(x)$.

As this gauge transformation describes just the freedom in the choice of
variables, the action $\mathcal A_1[\mathbf n_0]$ must be gauge invariant. It is
easy to check that the combination
\[
A_\mu^{ij}A_\mu^{ij}=\tr\bigl(A_\mu A_\mu^t\bigr)
\]
is not gauge invariant, and thus it can not appear in $\mathcal A_1$. The simplest
gauge invariant local combination of course is
\[
\tr\bigl(F_{\mu\nu}F_{\mu\nu}\bigr)\,,\qquad\text{with}\qquad
F_{\mu\nu}=\partial_\mu A_\nu-\partial_\nu A_\mu+[A_\mu,A_\nu]\,.
\]
However it involves extra derivative, and thus corresponds to an irrelevant term.

Therefore, in this approximation the only interesting contribution comes from the
above loop diagram with single $B^2$ vertex. We have
\begin{equation}\label{eq:A1sigma}
\mathcal A_1=\int\Bigl(\frac1{2g}-\frac{N-2}{4\pi}\,\log L\Bigr)\,(\partial_\mu\mathbf n_0)^2\,d^2x\,.
\end{equation}
Thus, the step 1 of the RG transformation leads to
\[
\frac1g\ \to\ \frac1{g_1}=\frac1g-\frac{N-2}{2\pi}\,\log L\,.
\]
The step 2 is nearly trivial. If one changes
\[
\mathbf n_0(x)=\mathbf n(x/L)
\]
the action does not change. Therefore
\[
RG_l:\qquad\frac1g\ \to\ \frac1g-\frac{N-2}{2\pi}\,l\,.
\]
In other words, the beta-function is
\[
\beta(g)=-\frac{N-2}{2\pi}\,g^2+O(g^3)\,,
\]
and for small $g$ the RG flow equation
\[
\frac{dg}{dl}=\frac{N-2}{2\pi}\,g^2
\]
shows that $g(l)$ \emph{increases} towards large distances; at $d=2$ the Goldstone
fixed point becomes unstable.

The above flow equation determines (in the leading order in $g$) the behavior of
the action under the RG transformation. One might be interested in the behavior
of correlation functions of, say, the field $\mathbf n(x)$. To that end one needs
then to find the RG function $\gamma_{\mathbf n}(g)$ in the CS equation
\[
\Bigl(x\frac\partial{\partial x}+n\,\gamma_{\mathbf n}(g)+\beta(g)\,\frac\partial{\partial g}\Bigr)\bigl\langle\mathbf n(x_1)\cdots\mathbf n(x_n)\bigr\rangle_c=0\,.
\]
It is defined in terms of the action of the RG transformation on the field
$\mathbf n(x)$.

Consider any correlation function involving $\mathbf n(x)$,
\[
\bigl\langle\mathbf n(x)\cdots\bigr\rangle\,,
\]
and assume again that the momentum $p$ associated with $x$ is $\ll\Lambda/L$.
Then, splitting as before
\[
\mathbf n=\mathbf n_0\,\sqrt{1-\phi^2}+\mathbf e^i\,\phi_i\simeq\mathbf n_0\,\Bigl(1-\frac12\,\phi^2\Bigr)+\mathbf e^i\,\phi_i
\]
one needs to find expectation value of this combination over the ensemble of the
fast modes $\phi$. The last term can be ignored (its expectation value is zero),
and we have
\[
\mathbf n\ \to\ \mathbf n_0\,\Bigl(1-\frac12\,\langle\phi_i^2\rangle\Bigr):=Z_{\mathbf n}^{-\frac12}(L)\,\mathbf n_0\,,
\qquad\text{where}\qquad
Z_{\mathbf n}^{-\frac12}(L)=1-\frac12\,\langle\phi_i^2\rangle=1-\frac{N-1}2\,\frac g{2\pi}\,l\,.
\]
Since at $d=2$ $\mathbf n$ has zero canonical dimension, we then have
\[
\gamma_{\mathbf n}(g)=D_{\mathbf n}(g)=-\frac d{dl}\,Z_{\mathbf n}^{-\frac12}(l)\Big|_{l=0}=\frac{N-1}{4\pi}\,g\,.
\]
In the above calculation we have assumed that $gl\ll1$; that is why the
higher-order terms in $\phi$ were neglected. But, as usual, we can extend it to
larger $l$. Generally, we have
\[
RG_l\,\mathbf n=Z_{\mathbf n}^{-\frac12}(l)\,\mathbf n\,,
\]
where $Z_{\mathbf n}^{-\frac12}(l)$ is a solution of the equation
\[
-\frac d{dl}\,Z_{\mathbf n}^{-\frac12}(l)=\gamma_{\mathbf n}\bigl(g(l)\bigr)\,Z_{\mathbf n}^{-\frac12}(l)\,,
\]
and $g(l)$ is the integral curve of the equation
\[
\frac{dg(l)}{dl}=-\beta\bigl(g(l)\bigr)\,.
\]
In our case
\[
\frac1{g(l)}=\frac1{g(0)}-\frac{N-2}{2\pi}\,l\,.
\]
Recall that the ``flow'' of $Z_{\mathbf n}^{-\frac12}(l)$ determines the scale
dependence of the correlation functions
\[
\bigl\langle\mathbf n(x_1/L)\cdots\mathbf n(x_n/L)\bigr\rangle_{g(l)}=Z_{\mathbf n}^{-\frac n2}(l)\,\bigl\langle\mathbf n(x_1)\cdots\mathbf n(x_n)\bigr\rangle_{g(0)}\,,
\]
where $l=\log L$.

\begin{exercise}
For the $O(N)$ nonlinear sigma model,

I. At $d=2$
\begin{enumerate}[label=\alph*)]
\item Assuming $g(l)\ll1$, find $Z_{\mathbf n}(l)$.
\item Derive the following $|x|\to0$ asymptotic of the two-point correlation
function
\[
\bigl\langle\mathbf n(x)\mathbf n(0)\bigr\rangle\sim\Bigl[\log\Bigl(\frac1{x^2}\Bigr)\Bigr]^{\frac{N-1}{N-2}}\,.
\]
Explain how this result is in agreement with the exact solution of the model at
$N\to\infty$ (PS, \S13.3, pages 463--465).
\end{enumerate}
II. Find the flow equation for $g(l)$ and $Z_{\mathbf n}(l)$ at $d=2+\epsilon$,
with $\epsilon\ll1$. Again, determine the $|x|\to0$ behavior of the two-point
correlation function.
\end{exercise}

Coming back to the flow equation ($d=2$)
\[
\frac{dg}{dl}=\frac{N-2}{2\pi}\,g^2+\dots
\]
one observes that (at $N>2$) the flow of the coupling $g(l)$ looks like this:
\[
\diac[0.6]{\draw[line width=0.8pt] (0,0) -- (6.5,0) node[right] {$g$}; \draw[arr] (2.5,0) -- (3.6,0);
  \fill (0,0) circle (2.4pt) node[above=2pt] {$0$};}
\]
At large distances $g$ increases, and the interaction becomes strong. Our weak
coupling analysis is no longer valid (in particular the higher order terms in
$\beta(g)$ can not be neglected). On the other hand, if we follow the flow of
$g(l)$ backward, to $l\to-\infty$,
\[
g(l)=\Bigl(\Const-\frac{N-2}{2\pi}\,l\Bigr)^{-1}\ \to\ \frac{2\pi}{N-2}\,(-l)^{-1}\,,
\]
the interaction of the Goldstone modes becomes weaker, and the short-distance
behavior is controlled by a free theory. For this reason, this type of RG flow is
called \emph{asymptotic freedom}. Note that the situation is opposite to that in
$\varphi^4$ theory at $d=4$, where instead $\lambda\to0$ at long scales, whereas
perturbation theory breaks down at short distances.

The asymptotic freedom allows one to establish existence of local field theory;
indeed it shows that at short scales the fluctuations are suppressed, and the
cutoff can be removed with no violation of locality. On the other hand, at longer
scales the fluctuations become stronger. It is not possible to describe these
large scale fluctuations by perturbation theory. Nonetheless, one can make
intuitive picture based on this pattern of the RG flow. If one concentrates
attention on very short scales, the vectors $\mathbf n(x)$ are nearly parallel,
because at $l\to-\infty$ $g(l)$ is small, thus suppressing the gradients of
$\mathbf n(x)$. However, if one looks at larger scales, the directions of
$\mathbf n(x)$ at the points separated by distance $\sim R_c$ are much less
correlated, and if the separation is even much larger, there is no correlation
(figure~\ref{fig:sigmascales}).
\begin{figure}[ht]
\centering\figfile{sigmascales}
\caption{At short scales the vectors $\mathbf n(x)$ are nearly parallel (left).
At scales much larger than $R_c$ their directions are uncorrelated (right).}
\label{fig:sigmascales}
\end{figure}

This picture suggests that at $d=2$ the phase with spontaneously broken $O(N)$
symmetry does not exist. In other words, $d=2$ magnet has no ferromagnetic phase
with spontaneous magnetization. Let me stress that the perturbative analysis of
the RG flow does not prove this statement, which requires much more involved
arguments. Note also that this picture is expected to apply only if
$N>2$.\footnote{At $N=2$ the $g^2$ term in the $\beta$-function vanishes. In fact,
it is easy to see that all higher terms vanish as well. At $N=2$ the direction
$\mathbf n(x)$ can be parameterized by a single angle $\theta(x)$,
$\mathbf n(x)=\bigl(\cos\theta(x),\sin\theta(x)\bigr)$, and formally
$\mathcal A=\frac1{2g}\int(\partial_\mu\theta)^2\,d^2x$, which is just the action
of free massless theory. The action of RG on a free theory is trivial. However,
beyond the perturbation theory, the $N=2$ sigma model has more complex structure.
This is because the field $\theta(x)$ is an angular variable,
$\theta(x)\sim\theta(x)+2\pi N(x)$, where $N(x)$ is integer-valued. As the result,
there are topologically nontrivial field configurations (``vortices'', or
``instantons''). Their interaction leads to remarkable phase transition at finite
$g=g_c$, the so called Berezinsky-Kosterlitz-Thouless transition (See e.g.\
Kosterlitz, Thouless, ``Ordering, metastability and phase transitions in
two-dimensional systems'', Journal of Physics C: Solid State Physics, Vol.~6 pages
1181--1203 (1973)).}

What is different at $d>2$? This question can be studied in details in formal case
of $d=2+\epsilon$ with $\epsilon\ll1$. In this case we can repeat our calculation
of the RG flow at small $g$. The step 1 (elimination of the fast part) yields
essentially the same result,
\begin{equation}\label{eq:A1sigmaeps}
\mathcal A_1=\int\Bigl(\frac1{2g}-\frac{N-2}{4\pi}\,\Lambda^\epsilon\,\log L\Bigr)\,(\partial_\mu\mathbf n_0)^2\,d^dx\,.
\end{equation}
However, the step 2 (scale transformation) is no longer exactly trivial. The change
of variables
\[
\mathbf n_0(x)=\mathbf n(x/L)
\]
leads to
\[
\int(\partial_\mu\mathbf n_0)^2\,d^dx=L^{d-2}\int(\partial_\mu\mathbf n)^2\,d^dx\,.
\]
Therefore
\[
\mathcal A'=\int L^\epsilon\,\Bigl(\frac1{2g}-\frac{N-2}{4\pi}\,\Lambda^\epsilon\,\log L\Bigr)\,(\partial_\mu\mathbf n_0)^2\,d^dx\,.
\]
Since
\[
L^\epsilon\approx1+\epsilon\log L
\]
one finds that under the RG transformation
\[
\frac1g\ \to\ \frac1g+\Bigl(\frac\epsilon g-\frac{N-2}{2\pi}\Bigr)\,\log L\,,
\]
where I have ignored the factor $\Lambda^\epsilon\approx1$.\footnote{One can get
rid of this factor by changing to a dimensionless coupling $g\to\Lambda^{-\epsilon}g$.}
The associated RG flow equation is
\[
\frac{dg}{dl}=-\epsilon\,g+\frac{N-2}{2\pi}\,g^2=\frac{N-2}{2\pi}\,g\,(g-g_*)\,,
\qquad\text{where}\qquad
g_*=\frac{2\pi\,\epsilon}{N-2}\,.
\]
Due to the small linear term, the flow now looks different at small $g\sim\epsilon$,
\[
\diac[0.6]{\draw[line width=0.8pt] (0,0) -- (6.5,0) node[right] {$g$};
  \draw[arr] (2.0,0) -- (0.9,0); \draw[arr] (3.8,0) -- (5.0,0);
  \node at (0,0) {$*$}; \node[above=2pt] at (0,0) {GS}; \node at (2.9,0) {$*$}; \node[above=2pt] at (2.9,0) {$g_*$};}
\]
Nontrivial new fixed point at $g=g_*$ appears, and if $g<g_*$ the theory flows
towards the Goldstone fixed point at large scale.

It is natural to identify the new fixed point $g_*$ with the Wilson-Fisher fixed
point at $d=2+\epsilon$. Indeed, general pattern of RG flow in $O(N)$ $\varphi^4$
theory suggested by our previous analysis was as shown in
figure~\ref{fig:wfgs}.
\begin{figure}[ht]
\centering\figfile{wfgs}
\caption{The flow from $\mathcal A_{*,WF}$ to the Goldstone point
$\mathcal A_{*,GS}$; the deviation $g-g_*$ is the sigma-model coordinate along it.}
\label{fig:wfgs}
\end{figure}

Here $M^2$ and $\lambda$ are two couplings in $\varphi^4$ theory. Closely below
$d=4$ the Wilson-Fisher fixed point appears close to $\mathcal A_{*,G}$, while the
Goldstone fixed point is located far away from it. When $d$ goes down from 4
$\mathcal A_{*,G}$ departs from $\mathcal A_{*,WF}$, but $\mathcal A_{*,GS}$ gets
closer (figure~\ref{fig:wfgs2}).
\begin{figure}[ht]
\centering\figfile{wfgs2}
\caption{The same flow at $d=3.99$ (left) and $d=2.01$ (right). As $d$ decreases
the Wilson-Fisher point moves away from the Gaussian point and towards the
Goldstone point.}
\label{fig:wfgs2}
\end{figure}

We can identify the coordinate along the trajectory going from
$\mathcal A_{*,WF}$ to $\mathcal A_{*,GS}$ with the coupling $g$ of the
non-linear sigma model, with the origin $g=0$ at $\mathcal A_{*,GS}$; then
$\mathcal A_{*,WF}$ appears at some $g_*$. At $d=2+\epsilon$ with small $\epsilon$
$\mathcal A_{*,GS}$ approaches $\mathcal A_{*,WF}$ ($g_*\sim\epsilon$). At $d=2$ the
fixed points $\mathcal A_{*,GS}$ and $\mathcal A_{*,WF}$ merge, and the flow of $g$
(at positive $g$) becomes marginally relevant. This pattern is confirmed by exact
solution of the theory at $N\to\infty$ (partly described in \S13.3 of PS).

\begin{hwproblem}{9}\label{hw:9}
For the $O(N)$ sigma model,
\[
\mathcal A[\mathbf n(x)]=\int d^2x\,\frac1{2g}\,\bigl(\partial\mathbf n(x)\bigr)^2\,.
\]
Assume $N>2$.

I. At $d=2$
\begin{hwparts}
\item Assuming $g(l)\ll1$, find $Z_{\mathbf n}(l)$ (see Section~22 for
definition).
\item Determine the $|x|\to0$ behaviour of the two-point correlation function
$\langle\mathbf n(x)\mathbf n(0)\rangle$.
\item \emph{Optional.} Explain how the result is in agreement with the exact
solution of the model at $N\to\infty$ (PS, \S13.3, pages 463--465).
\end{hwparts}
\emph{Direction:} Find $Z_{\mathbf n}(l)$ from
$-\frac d{dl}Z_{\mathbf n}^{-1/2}(l)=\gamma_{\mathbf n}(g(l))\,Z_{\mathbf n}^{-1/2}(l)$.
Use the equation
\[
\bigl\langle\mathbf n(x_1/L)\,\mathbf n(x_2/L)\bigr\rangle_{g(l)}=Z_{\mathbf n}^{-1}(l)\,\bigl\langle\mathbf n(x_1)\,\mathbf n(x_2)\bigr\rangle_{g(0)}\,,
\]
and remember that this correlation function depends on the coordinates through the
separation $|x_1-x_2|$ only.

II. At $d=2+\epsilon$ with $\epsilon\ll1$, find the flow equation for
$Z_{\mathbf n}(l)$. Again, determine the $|x|\to0$ behavior of the two-point
correlation function.
\end{hwproblem}

% Source pp. 239-253 (Lectures19-25.pdf).
\section{RG and Energy-Momentum Tensor}

Renormalization Group describes the way the scale transformations are realized in
QFT. The scale transformations
\[
x^\mu\to\tilde x^\mu=x^\mu/L
\]
constitute particular class of general coordinate transformations
\[
x^\mu\to\tilde x^\mu=f^\mu(x)\,,
\]
where $f^\mu$ are arbitrary smooth functions. Variations of the action
functional\footnote{The sign minus in this formula corresponds to the relation
$-\mathcal A=iS=i\int dt\,d\xx\,\Lag(x)=\int dx^4\,d\xx\,\Lag(x)=-\int dx^4\,d\xx\,\Lag_E(x)$,
where $\Lag(x)$ is the usual Lagrangian density with $t$ replaced by $-ix_4$, see
Section~7.}
\begin{equation}\label{eq:Aminus}
\mathcal A=-\int d^dx\,\Lag(x)
\end{equation}
under these coordinate transformations in a local field theory is generally
described in terms of its energy-momentum tensor $T^{\mu\nu}(x)$. Therefore this
field, and in particular its trace
\[
\Theta(x):=T^\mu{}_\mu(x)\,,
\]
plays important role in RG theory. In fact, all basic equations of RG theory in
local QFT can be expressed as certain properties of this field.

In the first lecture of this course we have defined the Energy-Momentum tensor (or
EM tensor, for shortness), in a classical field theory, as the Noether's current
associated with the translational symmetry $x^\mu\to x^\mu+a^\mu$,
\[
T^\mu{}_\nu=\frac{\partial\Lag}{\partial(\partial_\mu\varphi)}\,\partial_\nu\varphi-\delta^\mu_\nu\,\Lag\,,
\]
which satisfies the continuity equation
\begin{equation}\label{eq:Tcont}
\partial_\mu T^\mu{}_\nu=0
\end{equation}
as a consequence of the classical equations of motion. Generally, there is an
important ambiguity in the definition of the local EM tensor. One can redefine
$T^\mu{}_\nu$ by adding a total derivative
\[
T^\mu{}_\nu\to T^\mu{}_\nu+\partial_\lambda\psi^{\mu\lambda}{}_\nu\,,\qquad\psi^{\mu\lambda}{}_\nu=-\psi^{\lambda\mu}{}_\nu\,.
\]
The new tensor still satisfies the continuity equation \eqref{eq:Tcont}. By
appropriate choice of $\psi$ the modified EM tensor can always be made symmetric,
$T^{\mu\nu}=T^{\nu\mu}$.

More general way to define the EM tensor is to consider the field theory in the
space with nontrivial metric
\begin{equation}\label{eq:metric}
ds^2=g_{\mu\nu}(x)\,dx^\mu dx^\nu\,.
\end{equation}
In this case the action
\[
\mathcal A[\varphi,g]=-\int d^dx\,\sqrt g\,\Lag(\varphi,\partial_\mu\varphi,g_{\mu\nu})
\]
becomes a functional of the metric tensor $g$ also.

For instance, for generic scalar field theory, in the presence of the metric
background \eqref{eq:metric}, the simplest action is
\[
\mathcal A=\int\sqrt g\,\left[\frac12\,g^{\mu\nu}\,\partial_\mu\varphi\,\partial_\nu\varphi+V(\varphi)\right]\,.
\]
This form is called the ``minimal coupling to gravity''.

In any case, the field $T^{\mu\nu}(x)$ is defined in terms of the variation of the
action with respect to the metric (sign!),
\begin{equation}\label{eq:Tdef}
\delta\mathcal A=-\frac12\int d^dx\,\sqrt g\,T^{\mu\nu}(x)\,\delta g_{\mu\nu}(x)\,.
\end{equation}
Note that $T^{\mu\nu}$ defined this way is automatically symmetric,
\[
T^{\mu\nu}(x)=T^{\nu\mu}(x)\,,
\]
as $g_{\mu\nu}$ is.

Note that for flat space this definition is still ambiguous, as one could couple
the theory to the external metric in many different ways. For example, in the
scalar field theory one could add the ``non-minimal'' terms like
\[
\int d^dx\,\sqrt g\,R\,W(\varphi)\,,\qquad\text{or}\qquad\int d^dx\,\sqrt g\,R_{[\mu\lambda][\nu\rho]}\,W^{[\mu\lambda][\nu\rho]}(\varphi)
\]
where $R$ is the curvature, and $W$ is arbitrary. These terms disappear in the flat
space, however they affect the variation \eqref{eq:Tdef} even when the metric is
set back to the flat Euclidean one. It is possible to show that all such
``non-minimal'' modifications of the action in curved space lead to modifications
of the EM tensor of the form
\begin{equation}\label{eq:Tambig}
T^{\mu\nu}\to T^{\mu\nu}+\partial_\lambda\partial_\rho\,Y^{[\mu\lambda][\nu\rho]}\,,
\end{equation}
where $Y$ is antisymmetric in $[\mu\lambda]$ and in $[\nu\rho]$ but symmetric
w.r.t.\ the interchange of these pairs (this is exactly the symmetry of the
Riemann's curvature tensor $R_{[\mu\lambda][\nu\rho]}$).

This is intrinsic ambiguity which cannot be resolved unless we specify exactly how
the system behaves in curved space. Note that if $Y$ decays sufficiently fast at
$x\to\infty$, such modifications do not affect the notions of energy and momentum,
\[
E=\int d^{d-1}\xx\,T^{00}\,,\qquad P^i=\int d^{d-1}\xx\,T^{0i}\,,
\]
as they add total derivatives to the integrands. If we are interested just in the
theory in flat metric, fixing this ambiguity is a matter of convenience.

Consider now quantum field theory. The correlation functions are defined through
the functional integrals
\begin{equation}\label{eq:corrg}
\bigl\langle\mathcal O_1(x_1)\cdots\mathcal O_n(x_n)\bigr\rangle=Z^{-1}\int[D\varphi]\,\mathcal O_1(x_1)\cdots\mathcal O_n(x_n)\,e^{-\mathcal A[\varphi,g]}\,,
\end{equation}
where for the moment I consider the system in external gravitational field
$g_{\mu\nu}(x)$. Here $\mathcal O(x)$ are some local fields, which can be the field
$\varphi(x)$ itself, or some composite fields built from $\varphi$ and its
derivatives at the point $x$. In the presence of the external metric these
composite fields may depend on this metric in a local way, i.e.\ they may involve
$g_{\mu\nu}(x)$ and its derivatives at the point $x$.

Consider a small variation around the flat metric,
\[
g_{\mu\nu}(x)=\delta_{\mu\nu}+\delta g_{\mu\nu}(x)\,.
\]
For generic variation the correlation functions change in a nontrivial way because
the theory indeed depends on the background metric. However, the change is simple
if we take $\delta g_{\mu\nu}(x)$ corresponding to pure coordinate transformation
\begin{equation}\label{eq:coordtr}
x^\mu\to x^\mu+\varepsilon^\mu(x)
\end{equation}
with infinitesimal $\varepsilon^\mu(x)$; in this case
\[
-\delta g_{\mu\nu}(x)=\partial_\mu\varepsilon_\nu(x)+\partial_\nu\varepsilon_\mu(x)\,.
\]
The coordinate transformation is just a relabeling of the degrees of freedom
$\varphi(x)$, so it amounts to a change of variables in the functional integral
\eqref{eq:corrg}. Therefore the numerical value of this integral does not change.
This leads to the identity
\begin{multline}\label{eq:Ward}
0=\sum_{i=1}^n\bigl\langle\mathcal O_1(x_1)\cdots\delta_\varepsilon\mathcal O_i(x_i)\cdots\mathcal O_n(x_n)\bigr\rangle\\
+\int d^dx\,\partial_\mu\varepsilon_\nu(x)\,\bigl\langle T^{\mu\nu}(x)\,\mathcal O_1(x_1)\cdots\mathcal O_n(x_n)\bigr\rangle\,,
\end{multline}
where $\delta_\varepsilon\mathcal O_i(x)$ stand for local variations of the fields
$\mathcal O_i(x_i)$ under the coordinate transformation \eqref{eq:coordtr}. The
identities of this type are generally known as the \emph{Ward identities}.

What can be said about the variations $\delta_\varepsilon\mathcal O_i(x)$ on
general grounds? These again are local fields which depend on $\varepsilon^\mu(x)$
\begin{itemize}
\item[(i)] linearly, i.e.\
$\delta_{\varepsilon_1+\varepsilon_2}\mathcal O=\delta_{\varepsilon_1}\mathcal O+\delta_{\varepsilon_2}\mathcal O$,
and
\item[(ii)] locally, i.e.\ $\delta_\varepsilon\mathcal O(x)$ depends on
$\varepsilon^\mu(x)$ and its derivatives taken at the point $x$.
\end{itemize}
In fact, the part of $\delta_\varepsilon\mathcal O(x)$ which contains just
$\varepsilon^\mu(x)$, with no derivatives, is completely fixed,
\[
\delta_\varepsilon\mathcal O(x)=\varepsilon^\mu(x)\,\partial_\mu\mathcal O(x)+\text{terms with }\partial_\nu\varepsilon^\mu(x)\text{ and higher derivatives}\,,
\]
because for constant $\varepsilon^\mu(x)=\varepsilon^\mu$ we have
\[
\mathcal O(x)\to\mathcal O(x+\varepsilon)=\mathcal O(x)+\varepsilon^\mu\,\partial_\mu\mathcal O(x)\,.
\]
Moreover, the part which contains the first derivatives $\partial_\nu\varepsilon^\mu(x)$
is also severely restricted. For instance, if $\varphi(x)$ is a scalar field, and we
assume that it transforms independently of other fields, we can write
\[
\delta_\varepsilon\varphi(x)=\varepsilon^\mu(x)\,\partial_\mu\varphi(x)+\frac Dd\,(\partial_\mu\varepsilon^\mu)\,\varphi(x)\,,
\]
where $d$ is the dimensionality of the space, and the parameter $D$ is related to
the dimension of $\varphi$. This generalizes to more complicated scalar composite
fields $\mathcal O_\alpha$ as follows
\[
\delta_\varepsilon\mathcal O_\alpha(x)=\varepsilon^\mu(x)\,\partial_\mu\mathcal O_\alpha(x)+\frac{D_\alpha{}^\beta}d\,(\partial_\mu\varepsilon^\mu)\,\mathcal O_\beta(x)+\text{terms with }\partial_\nu\partial_\lambda\varepsilon^\mu\,,\ \text{etc.}
\]
For tensor fields the terms involving local rotations
\[
\omega_{\mu\nu}(x)=\partial_\mu\varepsilon_\nu(x)-\partial_\nu\varepsilon_\mu(x)
\]
will also appear.

Going back to the Ward Identity \eqref{eq:Ward}, we observe that integral there can
be transformed by parts, hence we have
\begin{multline}\label{eq:Ward2}
0=\sum_{i=1}^n\bigl\langle\mathcal O_1(x_1)\cdots\delta_\varepsilon\mathcal O_i(x_i)\cdots\mathcal O_n(x_n)\bigr\rangle\\
-\int d^dx\,\varepsilon_\nu(x)\,\bigl\langle\partial_\mu T^{\mu\nu}(x)\,\mathcal O_1(x_1)\cdots\mathcal O_n(x_n)\bigr\rangle\,.
\end{multline}
As the first term here depends only on $\varepsilon(x)$ (and its derivatives) at the
points $x_1,\dots,x_n$ only, we conclude that
\[
\bigl\langle\partial_\mu T^{\mu\nu}(x)\,\mathcal O_1(x_1)\cdots\mathcal O_n(x_n)\bigr\rangle=0\qquad\text{for}\quad x\neq x_1,\dots,x_n\,,
\]
i.e.\ in quantum field theory the continuity equation \eqref{eq:Tcont} is satisfied
in the weak sense,
\[
\partial_\mu T^{\mu\nu}(x)\simeq0\,,
\]
exactly as the equations of motion do.

Now let us consider particular case of the coordinate transformations, the scale
transformations (also called the dilations)
\[
x^\mu\to x^\mu+\varepsilon\,x^\mu\,.
\]
In this case
\[
\partial_\mu\varepsilon_\nu(x)=\delta_{\mu\nu}\,\varepsilon\,,
\]
while all higher derivatives of $\varepsilon(x)$ vanish, and so for scalar fields
$\mathcal O_\alpha$ we have
\[
\delta_\varepsilon\mathcal O_\alpha(x)=\varepsilon\,(\mathcal D\mathcal O)_\alpha(x)\,,
\]
where by definition
\begin{equation}\label{eq:DO}
(\mathcal D\mathcal O)_\alpha(x)=x^\mu\frac\partial{\partial x^\mu}\,\mathcal O_\alpha(x)+D_\alpha{}^\beta\,\mathcal O_\beta(x)\,.
\end{equation}
Then the Ward identity \eqref{eq:Ward} reduces to
\begin{equation}\label{eq:WardTheta}
\sum_{i=1}^n\bigl\langle\mathcal O_{\alpha_1}(x_1)\cdots(\mathcal D\mathcal O)_{\alpha_i}(x_i)\cdots\mathcal O_{\alpha_n}(x_n)\bigr\rangle
=-\int d^dx\,\bigl\langle\Theta(x)\,\mathcal O_{\alpha_1}(x_1)\cdots\mathcal O_{\alpha_n}(x_n)\bigr\rangle
\end{equation}
where
\[
\Theta(x)=T^\mu{}_\mu(x)
\]
is the trace component of the EM tensor.

Now let us assume that our QFT depends on some number $N$ of parameters (the
``coupling constants'') $\lambda^k$, $k=1,\dots,N$. These parameters enter the
action $\mathcal A$. The derivatives
\[
\Phi_k(x)=-\frac\partial{\partial\lambda^k}\,\Lag(x)
\]
are in general some scalar composite fields. Then\footnote{Of course, in general
the fields $\mathcal O_\alpha(x)$ may depend on the couplings $\lambda$. If such
dependence is present, the derivatives of $\mathcal O_\alpha$ would add to the
matrix $D_\alpha{}^\beta$, but otherwise would not change our argument. Here I ignore
this contribution for simplicity.}
\[
\frac\partial{\partial\lambda^k}\bigl\langle\mathcal O_{\alpha_1}(x_1)\cdots\mathcal O_{\alpha_n}(x_n)\bigr\rangle
=-\int d^dx\,\bigl\langle\Phi_k(x)\,\mathcal O_{\alpha_1}(x_1)\cdots\mathcal O_{\alpha_n}(x_n)\bigr\rangle
\]
If we assume that $\Theta(x)$ can be expanded in terms of $\Phi_k(x)$,
\begin{equation}\label{eq:ThetaBeta}
\Theta(x)=-\sum_{k=1}^N\beta^k(\lambda)\,\Phi_k(x)\,,
\end{equation}
one arrives again at the Callan-Symanzik equation, i.e.
\begin{multline}\label{eq:CSTheta}
\sum_{i=1}^n\bigl\langle\mathcal O_{\alpha_1}(x_1)\cdots(\mathcal D\mathcal O)_{\alpha_i}(x_i)\cdots\mathcal O_{\alpha_n}(x_n)\bigr\rangle\\
+\sum_{k=1}^N\beta^k(\lambda)\,\frac\partial{\partial\lambda^k}\bigl\langle\mathcal O_{\alpha_1}(x_1)\cdots\mathcal O_{\alpha_n}(x_n)\bigr\rangle=0\,,
\end{multline}
where the coefficients $D_\alpha{}^\beta(\lambda)$ in the operator $\mathcal D$ in
\eqref{eq:DO} can of course depend on the parameters $\lambda$. The possibility to
expand the field $\Theta(x)$ in terms of finitely many ``basic'' fields
$\Phi_k(x)$ is the statement of renormalizability of the field theory.

The energy-momentum tensor $T^{\mu\nu}(x)$ is special composite field. As we know,
in interacting field theory composite fields in general may require infinite
renormalizations. However, the above derivation of the Ward identity
\eqref{eq:Ward} suggests that in the case of the EM tensor no infinite
renormalization is needed. Indeed, the correlation functions in \eqref{eq:Ward}
are understood as the correlation functions of fully renormalized fields
$\mathcal O_\alpha$, they remain finite when the regularization is removed.
Therefore the EM tensor defined through the variation \eqref{eq:Tdef} must have
finite correlation functions without additional renormalization. It is possible to
check this statement in renormalized perturbation theory in $d=4$ $\varphi^4$
theory.

\subsection*{Conformal invariance at the fixed-point}

If $\lambda^k=\lambda^k_*$ such that
\[
\beta^k(\lambda_*)=0\,,
\]
the trace of the EM tensor vanishes in view of \eqref{eq:ThetaBeta}
\begin{equation}\label{eq:Theta0}
\Theta(x)=0\,.
\end{equation}

\begin{remark}
The above statement ignores important subtlety. The derivation of the
Callan-Symanzik equation \eqref{eq:CSTheta} relies in fact on the integrated version
of \eqref{eq:ThetaBeta}. Locally, Eq.~\eqref{eq:ThetaBeta} may hold only up to total
derivatives, i.e.\ the terms of the form $\partial_\mu K^\mu(x)$, with $K^\mu(x)$
being some local field. Thus, vanishing of the beta-functions only implies that
$\Theta(x)=\partial_\mu K^\mu(x)$. Unless $K^\mu$ is itself a total derivative
(i.e.\ $\Theta=\partial_\mu\partial_\nu L$, in which case the ambiguity
\eqref{eq:Tambig} can be used to redefine $T_{\mu\nu}$ to eliminate its trace),
presence of such a term would spoil the argument for conformal invariance. In some
cases (in particular, in $d=2$ unitary quantum field theories) this possibility can
be excluded. Generally it is not excluded, and therefore vanishing of the
beta-functions does not guarantee conformal invariance. Moreover, in general,
counterexamples can be found. For more on this subtlety, see \emph{J.~Polchinski,
Renormalization and effective Lagrangians, Nucl.\ Phys.\ B231, 269 (1984)}. In the
following discussion I assume that the fixed point in question admits the EM tensor
which satisfies \eqref{eq:Theta0}.
\end{remark}

The correlation functions of this fixed-point theory enjoy scale invariance with
the dimensions of the fields $\mathcal O_\alpha$ determined by the eigenvalues of
the matrix
\[
D_\alpha{}^\beta=D_\alpha{}^\beta(\lambda_*)\,.
\]
In fact, vanishing of the trace component $\Theta(x)$ in the fixed-point theory
guarantees much bigger symmetry. Namely, let $v^\mu(x)$ be any vector field which
satisfies the equation
\begin{equation}\label{eq:CKV}
\partial_\mu v_\nu(x)+\partial_\nu v_\mu(x)=\rho(x)\,\delta_{\mu\nu}=\frac2d\,\delta_{\mu\nu}\,\bigl(\partial_\lambda v^\lambda(x)\bigr)\,.
\end{equation}
Consider the coordinate transformation
\begin{equation}\label{eq:confeps}
x^\mu\to x^\mu+\varepsilon^\mu(x)\qquad\text{with}\qquad\varepsilon^\mu(x)=\varepsilon\,v^\mu(x)\,,
\end{equation}
where $\varepsilon$ is infinitesimal parameter. For such $\varepsilon^\mu$
\[
\partial_\mu\varepsilon_\nu\,T^{\mu\nu}=\frac\varepsilon2\,\rho(x)\,\Theta(x)=0\,,
\]
and the second term in the Ward identity \eqref{eq:Ward} disappears, i.e.
\begin{equation}\label{eq:WardConf}
\sum_{i=1}^n\bigl\langle\mathcal O_1(x_1)\cdots\delta_\varepsilon\mathcal O_i(x_i)\cdots\mathcal O_n(x_n)\bigr\rangle=0\,.
\end{equation}
We obtain the closed equation
\begin{equation}\label{eq:WardConf2}
\sum_{i=1}^n\bigl\langle\Phi_{\alpha_1}(x_1)\cdots\mathcal D_v\Phi_{\alpha_i}(x_i)\cdots\Phi_{\alpha_n}(x_n)\bigr\rangle=0\,,
\end{equation}
where
\[
\mathcal D_v\Phi_\alpha(x)=\left(v^\mu(x)\,\frac\partial{\partial x^\mu}+\frac{D_\alpha}d\,\bigl(\partial_\mu v^\mu(x)\bigr)\right)\,\Phi_\alpha(x)\,,
\]
and $\Phi_\alpha(x)$ are the fields which diagonalize the matrix $D_\alpha{}^\beta$,
and $D_\alpha$ are corresponding eigenvalues.

The coordinate transformations \eqref{eq:confeps} which satisfy the equation
\eqref{eq:CKV} are called the \emph{conformal transformations}. It is possible to
find all solutions of this equation. The situation turns out to be rather different
in the cases $d>2$ and $d=2$.

If $d>2$ the general solution has the form\footnote{The Lie algebra of these vector
fields is isomorphic to $so(d+1,1)$, the algebra of infinitesimal Lorentz
transformations of $d+2$-dimensional space-time.}
\[
v^\mu(x)=a^\mu+\omega^{\mu\nu}x^\nu+l\,x^\mu+\bigl(b^\mu x^2-2x^\mu(bx)\bigr)\,,
\]
where $\omega^{\mu\nu}=-\omega^{\nu\mu}$. The parameters $a,\omega,l,b$ are
interpreted as follows:
\begin{center}
\begin{tabular}{lll}
$a^\mu$ & -- & translations\\
$\omega^{\mu\nu}$ & -- & rotations\\
$l$ & -- & dilations\\
$b^\mu$ & -- & special conformal transformations
\end{tabular}
\end{center}
As the first three are very familiar, I will present only the finite version of the
special conformal transformations,
\[
x^\mu\to\tilde x^\mu=\frac{x^\mu+B^\mu x^2}{1+2(Bx)+B^2x^2}\,,
\]
where $B^\mu$ is an arbitrary vector.

Invariance of the fixed-point theory w.r.t.\ the special conformal transformations
provides further limitations on the correlation functions. Consider the Ward
identity \eqref{eq:WardConf2} with
\[
v^\mu(x)=b^\mu x^2-2x^\mu(bx)\,.
\]
In this case
\[
\mathcal D_v\Phi_\alpha(x)=\left[\bigl(b^\mu x^2-2x^\mu(bx)\bigr)\frac\partial{\partial x^\mu}-2D_\alpha\,(bx)\right]\,\Phi_\alpha(x)\,,
\]
and \eqref{eq:WardConf2} becomes nontrivial closed differential equation for the
correlation functions. In fact, this equation completely determines one-, two- and
three-point correlation functions, up to overall constant factors (this is left as
an exercise).

If $d=2$ the equation
\begin{equation}\label{eq:CKV2}
\partial_\mu v_\nu+\partial_\nu v_\mu-\delta_{\mu\nu}\,(\partial_\lambda v^\lambda)=0
\end{equation}
has infinitely many solutions. This is very easy to see by going to the complex
coordinates in $d=2$ space,
\[
z=x^1+ix^2\,,\qquad\bar z=x^1-ix^2\,.
\]
Then \eqref{eq:CKV2} leads to two equations
\[
\partial_{\bar z}v^z=0\,,\qquad\partial_zv^{\bar z}=0\,,
\]
which shows that $v^z(z)$ ($v^{\bar z}(\bar z)$) can be arbitrary holomorphic
(anti-holomorphic) function. This rich symmetry has important consequences for
the structure of $d=2$ conformal field theories.

Scale invariance of the fixed-point theory allows one to establish some important
properties of QFT associated with its unstable manifold.

\subsection*{Scale invariance and density of states}

Important conclusion can be made about its density of states. The density of
states is directly related to the equilibrium state entropy $S(E)$ of the system.
The form of $S(E)$ in a scale invariant quantum field theory can be established
directly by general arguments. However, I will use this opportunity to display some
general features of QFT at finite temperature.

The functional integral of $d$-dimensional Euclidean quantum field theory
\[
\langle\dots\rangle=Z^{-1}\int D[\varphi]\,(\dots)\,e^{-\mathcal A[\varphi]}
\]
has formal interpretation as expectation value over a thermal equilibrium ensemble
in \emph{classical} statistical mechanics in $d$ spatial dimensions. We have used
this interpretation many times already. Here I would like to discuss Euclidean
functional integral representation of \emph{quantum} statistical mechanics of
fields in $d$ dimensional space-time (i.e.\ $d-1$ dimensional space).

Recall that the density matrix of thermal equilibrium state in quantum statistics
has the general form\footnote{The normalization factor $Z(\beta)$ here is different
from the factor $Z$ in e.g.\ Eq.~\eqref{eq:corrdef}. $Z(\beta)$ goes to $Z$ in the
low temperature limit $\beta\to\infty$ (see below).}
\begin{equation}\label{eq:density}
\hat P=Z^{-1}(\beta)\,e^{-\beta\hat H}\,,
\end{equation}
where $\beta=(kT)^{-1}$. The exponential here coincides with the time evolution
operator at pure imaginary value of time $t=-i\hbar\beta$ (I set $\hbar=1$
henceforth). This form is still valid if one is dealing with the equilibrium state
of quantum fields. The imaginary-time evolution operator has representation as the
functional integral,
\begin{equation}\label{eq:PIbeta}
\bigl\langle\varphi_f(\xx)\big|e^{-\beta\hat H}\big|\varphi_i(\xx)\bigr\rangle=\int D[\varphi(\xx,\tau)]\Big|_{\substack{\varphi(\xx,0)=\varphi_i(\xx)\\\varphi(\xx,\beta)=\varphi_f(\xx)}}\,e^{-\mathcal A[\varphi]}
\end{equation}
Here $\xx$ are $d-1$ dimensional ``spatial'' coordinates, and $\tau$ is the
imaginary time (so that $x=(\xx,\tau)$). The integration is over the functions
$\varphi(x)=\varphi(\xx,\tau)$ in the slab $0\le\tau\le\beta$, with the boundary
conditions
\[
\varphi(\xx,0)=\varphi_i(\xx)\,,\qquad\varphi(\xx,\beta)=\varphi_f(\xx)\,.
\]
The partition sum of the quantum system at the finite temperature is the trace
\begin{equation}\label{eq:Zbeta}
Z(\beta)=\tr\bigl(e^{-\beta\hat H}\bigr)\,,
\end{equation}
i.e.\ the sum of its diagonal matrix elements
\begin{equation}\label{eq:Ztrace}
\int D[\varphi(\xx)]\,\bigl\langle\varphi(\xx)\big|e^{-\beta\hat H}\big|\varphi(\xx)\bigr\rangle\,.
\end{equation}
If the representation \eqref{eq:PIbeta} is used, the integration over
$\varphi(\xx)$ in \eqref{eq:Ztrace} combines with the integral $D[\varphi(x)]$ in
\eqref{eq:PIbeta} to produce integral over all fields $\varphi(\xx,\tau)$, with
$0\le\tau\le\beta$, satisfying the periodic boundary condition
\begin{equation}\label{eq:periodic}
\varphi(\xx,0)=\varphi(\xx,\beta)\,.
\end{equation}
Thus, we have
\begin{equation}\label{eq:Zcyl}
Z(\beta)=\int D[\varphi(\xx,\tau)]\big|_{\varphi(\xx,0)=\varphi(\xx,\beta)}\,e^{-\mathcal A[\varphi]}\,.
\end{equation}
Geometrically, this is the functional integral over the functions $\varphi(x)$ on a
$d$-dimensional cylinder $\R^{d-1}\times S^1$, with $\tau$ being the coordinate on
the circle of circumference $\beta$, $\tau\sim\tau+\beta$. Of course this is nothing
but the Matsubara representation specialized to the quantum field theory at hand.

Now, let us change interpretation of the functional integral \eqref{eq:Zcyl}. We
can take one of the coordinates on $\R^{d-1}$, say $x_1$, and claim it to be the
imaginary time. The coordinate $\tau$ will be treated as one of the spatial
coordinates. In this interpretation the imaginary time spans infinite range,
$-\infty<x_1<\infty$, while the ``space'' $\R^{d-2}\times S^1$ is the $d-1$
dimensional cylinder, with the coordinates $(\xx_\parallel,x_\perp)$. Here
$\xx_\parallel=(x_2,\dots,x_{d-1})$, and $x_\perp=\tau$. The latter is
compactified,
\[
x_\perp\sim x_\perp+\beta\,.
\]
Since now the imaginary time spans large (in the limit, infinite) range,
\eqref{eq:Zcyl} is dominated by the contribution from the ground state,
\begin{equation}\label{eq:ZE0}
Z(\beta)\sim e^{-L\,E_0(\beta)}\,,
\end{equation}
where $L\to\infty$ is the range of $x_1$, and $E_0(\beta)$ is the ground state
energy of the quantum field in the spatial geometry $\R^{d-2}\times S^1$.

This geometric setting is similar to that we discussed in the context of the
Casimir effect, with $\beta$ replacing the separation between the plates $a$. The
difference is that in the Casimir case we had a fixed boundary condition
$\varphi=0$ at the plates, which now is replaced by the periodic boundary condition
\eqref{eq:periodic}. For the free massless field it is easy to repeat the
calculation of the Casimir energy (I recommend it as an Exercise), with the result
\begin{equation}\label{eq:E0beta}
E_0(\beta)=-\frac{C\,V_\parallel}{\beta^{d-1}}\,,
\end{equation}
where $C$ is some dimensionless ($d$-dependent) constant, and
\[
V_\parallel=\int d^{d-2}\xx_\parallel
\]
is the volume of the $\R^{d-2}$ component of the space. In fact, apart from the
value of the constant, the form \eqref{eq:E0beta} is dictated by scale invariance
of the theory. Indeed, the ground state energy is an extensive quantity, which at
large $V_\parallel$ must be proportional to this volume. Since the theory has no
intrinsic scale, the $\beta$-dependence then follows from the dimensional analysis.
This simple argument shows that the form \eqref{eq:E0beta} applies to any
scale-invariant field theory, id est fixed point field theory. The constant $C$ of
course would be different for different fixed points.

Returning to original interpretation (the quantum field at finite temperature
$kT=\beta^{-1}$), and using \eqref{eq:ZE0} we have for the free energy of a
scale-invariant field theory
\begin{equation}\label{eq:Fbeta}
\beta F(\beta)=-\log\bigl(Z(\beta)\bigr)=-\frac{C\,V_{d-1}}{\beta^{d-1}}\,,
\end{equation}
where $V_{d-1}=V_\parallel\times L$ is the volume of the space. Of course, one could
have established the form \eqref{eq:Fbeta} directly, by dimensional analysis (the
free energy $F$, being extensive characteristic, must be $\sim V_{d-1}$, whereas
the quantity \eqref{eq:Fbeta} is dimensionless), but I believe our detour through
the Casimir energy was instructive.

The finite temperature partition function is closely related to the energy density
of states. We have
\[
Z(\beta)=\int_0^\infty e^{-\beta E}\,\rho(E)\,dE\,,
\]
where $\rho(E)\,\Delta E$ is the number of states in the energy interval between
$E$ and $E+\Delta E$. Mathematically, $Z(\beta)$ is the Laplace transform of
$\rho(E)$. The inverse transform is
\[
\rho(E)=\int_{-i\infty}^{+i\infty}e^{\beta E}\,Z(\beta)\,\frac{d\beta}{2\pi i}
=\int_{-i\infty}^{+i\infty}e^{\beta E}\,e^{\frac{C\,V_{d-1}}{\beta^{d-1}}}\,\frac{d\beta}{2\pi i}\,.
\]
At large $V_{d-1}$ the integral is dominated by the saddle point at
\[
\beta=\beta_0(E):=\left(\frac{C\,(d-1)\,V_{d-1}}E\right)^{1/d}\,,
\]
and we find
\begin{equation}\label{eq:rhoE}
\rho(E)\sim e^{S(E)}\,,
\end{equation}
where
\begin{equation}\label{eq:SE}
S(E)=\tilde C\,V_{d-1}^{1/d}\,E^{1-1/d}
\end{equation}
where $\tilde C$ is another constant, related to $C$ in
\eqref{eq:E0beta}.\footnote{$\tilde C=\frac d{d-1}\,\bigl(C\,(d-1)\bigr)^{1/d}$.}

By the definition, the quantity $S(E)$ in \eqref{eq:rhoE} coincides with the
equilibrium state entropy of the system at the temperature $kT=1/\beta_0(E)$. The
equation \eqref{eq:SE} expresses the universal scale invariant form of the entropy.
Indeed, the entropy is an extensive property which must have the form
$V_{d-1}\,\sigma$, where the entropy density $\sigma=\sigma(\varepsilon)$ can depend
on the energy density $\varepsilon:=E/V_{d-1}$. On the other hand $S(E)$, being a
dimensionless quantity, can only depend on the dimensionless combination
$V_{d-1}E^{d-1}$. Eq.~\eqref{eq:SE} is the only form that satisfies these
conditions. Thus, the asymptotic density of states \eqref{eq:rhoE}, \eqref{eq:SE}
can be deduced directly, from the scale invariance of the fixed point theory.
Nonetheless, our previous discussion is instructive as it displays some basics of
the finite-temperature formalism of QFT.

Note that in any dimensionality $S(E)$ grows slower than the exponential $\sim E$.
This observation is useful in comparing to the string theories, whose typical
energy density is the exponential.\footnote{Stringy density of states
$S(E)\simeq E/kT_H$ has manifestation in the so called \emph{Hagedorn transition}
at $T=T_H$.}

It is straightforward to argue that the energy density \eqref{eq:rhoE},
\eqref{eq:SE} is valid asymptotically, at large $E$, in any quantum field theory
originating at certain critical fixed point (i.e.\ associated with
$U(\mathcal A_*)$). And if we adopt the Wilsonian idea of quantum field theories as
RG flows on the unstable manifolds of fixed points, the above equation \eqref{eq:SE}
becomes general limitation on the asymptotic density of states in quantum field
theory.

Let me make here a remark on the status of this idea. In our discussion of the
Renormalization Group we (following Wilson) found it natural to identify local field
theories with the special class of RG flows in $\Sigma$, the flows that can be
``integrated backward'' without limit. A critical fixed point $\mathcal A_*$ gives
simple mechanism for generating such flows. There can be several fixed points, or
even continuous manifold of fixed points.\footnote{About limit cycles,
$C$-theorems, etc.} In any case, the above conclusion about the high energy density
of states, which follows from the scale invariance of the fixed point theory,
remains valid.

However, another, somewhat more intricate behavior is conceivable. Let me briefly
discuss it, just in the way of entertainment. To start with, let us imagine a slice
of the critical surface passing through a fixed point $\mathcal A_*$. All RG
trajectories in some domain around $\mathcal A_*$ flow to $\mathcal A_*$ at
$l\to+\infty$, see figure~\ref{fig:nestedfp}. Let's also assume that there is
another fixed point $\mathcal A'_*$ lying in the same critical surface, such that
$\mathcal A_*\in U(\mathcal A'_*)$. Furthermore, assume there is yet another fixed
point $\mathcal A''_*\in U(\mathcal A'_*)$.
\begin{figure}[ht]
\centering\figfile{nestedfp}
\caption{An RG trajectory starting at $\mathcal A''_*$, passing close to
$\mathcal A'_*$, and ending at $\mathcal A_*$.}
\label{fig:nestedfp}
\end{figure}

This figure shows an RG trajectory which originates at $\mathcal A''_*$, then passes
very close to $\mathcal A'_*$, finally converging to $\mathcal A_*$. Starting from
any point on this trajectory, one can ``integrate backward'' without limit, i.e.\
it defines local QFT. While the large scale properties of this QFT are controlled by
the ``infrared'' fixed point $\mathcal A_*$, its behavior at shorter scales is more
intricate. When traced backwards, the RG trajectory stays long ``time'' near the
fixed point $\mathcal A'_*$, before departing towards $\mathcal A''_*$ at
$l\to-\infty$. The associated QFT has its asymptotic short distance behavior
determined by the ``ultraviolet'' fixed point $\mathcal A''_*$, but in some wide
range of intermediate scales it has scaling properties associated with
$\mathcal A'_*$.

One can imagine yet more complicated sequence of nested fixed points, giving rise to
QFT with many different intermediate regimes corresponding to the trajectory passing
close to consecutive fixed points. Ultimately, then, a flow is conceivable which
passes successively near infinitely many fixed points. While such trajectory can be
integrated backward with no limit, its short distance behavior is not controlled by
any given fixed point.

% Source pp. 254-263 (Lectures19-25.pdf).
\section{Operator Formalism vs Functional Integral in QFT}

So far we have discussed QFT mostly from the perspective of (Euclidean space-time)
functional integral. There are aspects which are better understood within the
formalism based on the notions of the Hilbert space of states which supports the
algebra of local field operators. In the previous discussion I made occasional
references to statements that follow from the operator formalism, mostly from the
positivity of norms in the Hilbert space, as well as from the so called spectral
condition. Here I would like to recollect the basics of that formalism, and
clarify its relation to the functional integral.

\subsection*{Operator formalism}

As was stated before (Section~7), the correlation functions defined by the
functional integrals
\begin{equation}\label{eq:corrPhi}
\bigl\langle O_1(x_1)\cdots O_N(x_N)\bigr\rangle:=Z^{-1}\int[D\Phi]\,O_1(x_1)\cdots O_N(x_N)\,e^{-\mathcal A[\Phi]}
\end{equation}
can be interpreted as the vacuum-vacuum matrix element of the ``T-ordered'' product
\begin{equation}\label{eq:Tordered}
\bigl\langle O_1(x_1)\cdots O_N(x_N)\bigr\rangle=\bra0T_\tau\bigl(\hat O_1(x_1)\cdots\hat O_N(x_N)\bigr)\ket0
\end{equation}
of the Heisenberg field operators $\hat O(x)$. The latter notions involve
nominating one of the coordinates to the role of the ``Euclidean time'' $\tau$, so
that $x=(\xx,\tau)$. Usually, the Heisenberg picture is developed in the real time
$t$ which relates to $\tau$ by the analytic continuation to $\tau=it$, but in most
of the development below this specification is not necessary; we generally will
regard $\tau$ as a complex number. Then
\begin{equation}\label{eq:Heis}
\hat O(\xx,\tau)=e^{\hat H\tau}\,\hat O(\xx)\,e^{-\hat H\tau}\,.
\end{equation}
By construction, the field operators satisfy Heisenberg ``equation of motion''
\begin{equation}\label{eq:HeisEOM}
\frac\partial{\partial\tau}\,\hat O(\xx,\tau)=[\hat H,\hat O(\xx,\tau)]
\end{equation}
In fact, most of the calculations can be done at real $\tau$ (``imaginary time'').
In that case the ``time ordering'' $T_\tau$ means that the operators in
\eqref{eq:Tordered} are ordered according to decreasing values of $\tau$: the
operators with greater $\tau$ are placed to the left.

We have derived this relation in quantum mechanics (Section~5), and then assumed
that it can be extended to QFT. In fact, such extension is not straightforward,
and involves many subtleties mostly related to regularization and renormalization
(our brief inspection of the Wilson's RG gives the idea of the size of the
problem). Here I would like to revisit the relation \eqref{eq:Tordered} from more
general perspective. Starting from the functional integral definition of the
correlation functions, we will try to recover basic notions of the operator
formalism. Eq.~\eqref{eq:Heis} involves the Hamiltonian, and we would like to see
how the Hamiltonian, or more generally, the Energy-Momentum $P^\mu$, can be
introduced in terms of the correlation functions.

\subsection*{Correlation Functions and Energy-Momentum}

Consider the correlation function
\[
\bigl\langle T^{\mu\nu}(x)\,O_1(x_1)\cdots O_N(x_N)\bigr\rangle
\]
involving the energy-momentum tensor field $T^{\mu\nu}(x)$. As we have seen in the
previous Section, it satisfies the Ward Identity \eqref{eq:Ward2}, in particular
\[
\bigl\langle\partial_\mu T^{\mu\nu}(x)\,O_1(x_1)\cdots O_N(x_N)\bigr\rangle=0\qquad\text{for}\quad x\neq x_1,\dots,x_N\,,
\]
As a consequence, the integral
\[
\int_\Sigma d\Sigma_\mu(x)\,\bigl\langle T^{\mu\nu}(x)\,O_1(x_1)\cdots O_N(x_N)\bigr\rangle
\]
taken over a $(d-1)$-hypersurface $\Sigma$ does not change under topologically
trivial deformations of $\Sigma$. The hypersurface $\Sigma$ divides the space
$\R^d$ into two parts, and correspondingly splits the set of insertion points
$x_1,\dots,x_N$ into two subsets, according to their locations within the parts.
Topologically trivial are those deformations which do not change this splitting,
or, in other words the deformations which do not make the hypersurface ``cross
over'' any of the insertion points. For the sake of our present discussion it is
convenient to choose $\Sigma$ to be an ``equal-time slice''. As before, we
nominate one of the coordinates in $\R^d$ to be the (imaginary) ``time'' $\tau$,
i.e.\ $x=(\xx,\tau)$, and define $\Sigma(\tau)$ as the union of points with given
$\tau$. Define\footnote{The factor $i$ is added to this notation to agree with the
conventional notation $P^\mu$ for the energy-momentum in real Minkowski space time.
Upon continuation $\tau=it$ the components of the energy-momentum tensor transform
so that $T^{\tau k}=iT^{0k}$, $k=1,2,\dots,d-1$, and $T^{\tau\tau}=-T^{00}$.}
\begin{equation}\label{eq:Pslice}
iP^\mu(\tau):=\int_{\Sigma(\tau)}d\Sigma_\nu(x)\,T^{\mu\nu}(x)=\int_{\Sigma(\tau)}d^{d-1}\xx\,T^{\mu\tau}(\xx,\tau)\,,
\end{equation}
and assume that
\begin{equation}\label{eq:tauorder}
\tau_1>\tau_2>\dots>\tau_M>\tau>\tau_{M+1}>\dots>\tau_N\,,
\end{equation}
as illustrated in figure~\ref{fig:slice}.
\begin{figure}[ht]
\centering\figfile{slice}
\caption{The equal-time slice $\Sigma(\tau)$ separates the insertion points with
$\tau_i>\tau$ from those with $\tau_i<\tau$.}
\label{fig:slice}
\end{figure}
Then the correlation function
\begin{equation}\label{eq:Pcorr}
\bigl\langle P^\mu(\tau)\,O_1(x_1)\cdots O_N(x_N)\bigr\rangle
\end{equation}
does not change if one changes the value of $\tau$, as long as the order
\eqref{eq:tauorder} is preserved. We therefore drop the argument $\tau$ writing in
the notation for $P^\mu:=P^\mu(\tau)$, but adopt the convention that the position of
the insertion $P^\mu$ in the correlation function corresponds to the order
\eqref{eq:tauorder}, i.e.\ \eqref{eq:Pcorr} is written as
\begin{equation}\label{eq:Pinsert}
\bigl\langle O_1(x_1)\cdots O_M(x_M)\,P^\mu\,O_{M+1}(x_{M+1})\cdots O_N(x_N)\bigr\rangle
\end{equation}
With this convention the position of the insertion $P^\mu$ represents the
topological class of the surface $\Sigma$ in \eqref{eq:Pslice} relative to the
points $x_i$.

Now that the position of $P^\mu$ in the correlation function \eqref{eq:Pcorr}
matters, one may ask about the ``commutators'' of $P^\mu$ with the local field
insertions. We have formally (see figure~\ref{fig:commslab})
\begin{multline*}
iP^\mu O(\xx,\tau)-O(\xx,\tau)\,iP^\mu\\
=\int d^{d-1}\yy\,T^{\tau\mu}(\yy,\tau+\Delta\tau)\,O(\xx,\tau)-\int d^{d-1}\yy\,T^{\tau\mu}(\yy,\tau-\Delta\tau)\,O(\xx,\tau)\,,
\end{multline*}
\begin{figure}[ht]
\centering\figfile{commslab}
\caption{The commutator with $P^\mu$ is the difference of the slices at
$\tau+\Delta\tau$ and $\tau-\Delta\tau$, which can be closed into a small surface
around $x$.}
\label{fig:commslab}
\end{figure}
or, by closing the hypersurface
\begin{equation}\label{eq:Pcomm}
i[P^\mu,O(x)]=\int_{\Sigma_x}d\Sigma_\nu(y)\,T^{\mu\nu}(y)\,O(x)
\end{equation}
where $\Sigma_x$ is an arbitrary small closed hypersurface wrapped around the point
$x$. Again, this relation is understood as being valid for the insertions in the
correlation functions, with the above ordering convention.

The integral in the r.h.s.\ can be calculated as follows. Consider again the Ward
identity \eqref{eq:Ward},
\begin{equation}\label{eq:Ward3}
\sum_{i=1}^n\bigl\langle\mathcal O_1(x_1)\cdots\delta_\varepsilon\mathcal O_i(x_i)\cdots\mathcal O_n(x_n)\bigr\rangle
=-\int d^dx\,\partial_\mu\varepsilon_\nu(x)\,\bigl\langle T^{\mu\nu}(x)\,\mathcal O_1(x_1)\cdots\mathcal O_n(x_n)\bigr\rangle\,.
\end{equation}
where $\delta_\varepsilon O(x)$ stand for the variations of the fields $O(x)$ under
the coordinate transformation $x^\mu\to x^\mu+\varepsilon^\mu(x)$. The integration
in the r.h.s.\ is over the whole space $D:=\R^d$. Let us split this domain into
several pieces, as follows
\[
D:=\R^d=\sum_{i=1}^ND_i+D_{\text{out}}\,,
\]
where $D_i$ is a small domain containing the point $x_i$, and $D_{\text{out}}$ is
the remaining part of $D$ left after removing all $D_i$\footnote{In conventional
set theory symbolics the decomposition reads
$D_{\text{out}}:=D\setminus(\cup_iD_i)$.} (figure~\ref{fig:Di}).
\begin{figure}[ht]
\centering\figfile{Di}
\caption{Left: a small domain $D_i$ around the insertion point $x_i$. Right: the
normal element $d\Sigma_\mu$ of its boundary $\partial D_i$.}
\label{fig:Di}
\end{figure}
The integral splits accordingly, as
\[
\int d^dx=\sum_{i=1}^N\int_{D_i}d^dx+\int_{D_{\text{out}}}d^dx
\]
For all $x\in D_{\text{out}}$ we have
\[
\partial_\mu T^{\mu\nu}(x)=0\qquad\text{for}\quad x\in D_{\text{out}}\,,
\]
and hence the integral over $D_{\text{out}}$ can be transformed via the Stokes
theorem, as
\[
\int_{D_{\text{out}}}d^dx\,\partial_\mu\bigl\langle\bigl(\varepsilon_\nu(x)\,T^{\mu\nu}(x)\bigr)\,O_1(x_1)\cdots O_N(x_N)\bigr\rangle=-\sum_{i=1}^NI_i\,,
\]
where $I_i$ stand for the surface integrals over the boundaries of $D_i$,
\[
I_i=\int_{\partial D_i}d\Sigma_\mu(x)\,\varepsilon_\nu(x)\,\bigl\langle T^{\mu\nu}(x)\cdots O_i(x_i)\cdots\bigr\rangle
\]
and $d\Sigma_\mu(x)$ is the normal element of the boundary $\partial D_i$. Hence,
the r.h.s.\ of the Eq.~\eqref{eq:Ward3} splits into the sum
\begin{equation}\label{eq:splitsum}
\sum_{i=1}^N\Bigl(I_i-\int_{D_i}d^dx\,\partial_\mu\varepsilon_\nu(x)\,\bigl\langle T^{\mu\nu}(x)\cdots\bigr\rangle\Bigr)\,.
\end{equation}
By construction, each of the terms here is independent of the shape of the domain
$D_i$, which can be made arbitrary small. This implies that each of these terms
depends on $\varepsilon(x)$ in a local way, i.e.\ the term associated with $D_i$
depends on $\varepsilon(x_i)$, $\partial\varepsilon(x_i)$, $\dots$. Comparing
\eqref{eq:splitsum} with the l.h.s.\ of \eqref{eq:Ward3} we conclude that the
equality is satisfied term by term in the sum over $i$. This leads to the following
representation of the variations $\delta_\varepsilon O(x)$
\begin{equation}\label{eq:deltaO}
\delta_\varepsilon O(x)=\int_{\partial D_x}d\Sigma_\mu(y)\,\varepsilon_\nu(y)\,T^{\mu\nu}(y)\,O(x)-\int_{D_x}d^dy\,\partial_\mu\varepsilon_\nu(y)\,T^{\mu\nu}(y)\,O(x)\,.
\end{equation}
The relation is understood as valid for the correlation functions. Let me note again
that the r.h.s.\ is in fact independent of the choice of the (small) domain $D_x$.

As was already mentioned, the variation $\delta_\varepsilon O(x)$ admits expansion in
the derivatives of $\varepsilon$ taken at the point $x$, with the leading term
\[
\delta_\varepsilon O(x)=\varepsilon^\mu(x)\,\partial_\mu O(x)+\text{terms involving }\partial\varepsilon(x),\dots
\]
In particular, for a constant $\varepsilon^\mu(x)=\varepsilon^\mu$ the terms with
the derivatives do not contribute. On the other hand, in this case the second term
in \eqref{eq:deltaO} vanishes, and we obtain
\[
\partial^\mu O(x)=\int_{\partial D_x}d\Sigma_\nu(y)\,T^{\mu\nu}(y)\,O(x)
\]
The r.h.s.\ here coincides with the expression in \eqref{eq:Pcomm}. Finally, we
find
\[
[P^\mu,O(x)]=-i\,\partial^\mu O(x)\,.
\]
Taking $\mu$ to be the ``time'' index $\tau$, and recalling that $P^\tau=iP^0=iH$,
one recovers \eqref{eq:HeisEOM}.\footnote{Up to the overall sign. With the sign
conventions used in this and the previous Section one finds
$\partial_\tau O=-[H,O]$; the sign in \eqref{eq:HeisEOM} is recovered by reversing
the overall sign of $T^{\mu\nu}$ in \eqref{eq:Tdef}. The notes flag this sign
without resolving it.} The ``spatial'' components of this equation
\[
i\nabla O(\xx,\tau)=[\PPv,O(\xx,\tau)]
\]
reproduce the well known (see e.g.\ Chapter 2 of PS) representation of the
Heisenberg field operators
\[
\hat O(\xx,t)=e^{-i\PPv\cdot\xx+iHt}\,\hat O(\mathbf0,0)\,e^{i\PPv\cdot\xx-iHt}
\]
which I chose here to write directly in Minkowski space-time.

\subsection*{Space of States\protect\footnote{The ideas displayed in this section
are borrowed from the construction in K.~Osterwalder, R.~Schrader, ``Axioms for
Euclidean Greens functions,'' Comm.\ Math.\ Phys.\ \textbf{31} (1973), 83--112;
\textbf{42} (1975), 281--305.}}

The functional integral is particularly suitable for dealing with the correlation
functions \eqref{eq:corrPhi}. Calculations in the previous paragraphs suggest a
natural way to construct the space of states $\Ham$ and an algebra of operators, in
terms of the Euclidean space correlation functions, as follows.

Let us again split the whole space $\R^d$, with the coordinates
$x^\mu=(\xx,\tau)$, into two half-spaces
\[
\R^d=\R^d_+\cup\R^d_-
\]
separated by the ``equal time'' slice, the hypersurface $\tau=0$. Namely
\[
\R^d_+:\ (\xx,\tau)\,,\ \tau\ge0\,,\qquad\R^d_-:\ (\xx,\tau)\,,\ \tau\le0\,,
\]
We associate states with the product of field insertions
$\varphi(x_1)\cdots\varphi(x_M)$ at the points $x_1,x_2,\dots,x_M\in\R^d_-$,
\begin{equation}\label{eq:states}
\varphi(x_1)\cdots\varphi(x_M)\ \to\ \ket{\varphi(x_1)\cdots\varphi(x_M)}\,.
\end{equation}
At this point, the above equation just provides notation for the states. We define
$\Ham$ as the linear envelope of the above states
\[
\Ham=\operatorname{Span}\bigl\{\ket{\varphi(x_1)\cdots\varphi(x_M)}:\ x_i\in\R^d_-\bigr\}
\]
The intuition behind this construction comes from the simple minded idea of $\Ham$
as the space of functionals
\[
\Ham=\text{Space of }\Psi[\varphi(\xx)]
\]
of the field configurations $\varphi(\xx)$, $\xx\in\R^{d-1}$. Let
$x_i=(\xx_i,\tau_i)\in\R^d_-$, and consider the functional integral
\begin{equation}\label{eq:Psiint}
\int_{\varphi(\xx,0)=\varphi(\xx)}[D\varphi(x)]\,\varphi(x_1)\cdots\varphi(x_M)\,e^{-\mathcal A[\varphi(x)]}
\end{equation}
over all field configurations $\varphi(x)$ in the half-space $x\in\R^d_-$. This
integral defines (up to overall normalization) the functional
\[
\Psi_{\varphi(x_1)\cdots\varphi(x_M)}[\varphi(\xx)]
\]
which may be regarded as the Schroedinger wave function representing the state
\eqref{eq:states}. In particular, the integral \eqref{eq:Psiint} with no insertions
represents the vacuum state,
\[
\Psi_0[\varphi(\xx)]=\int_{\varphi(\xx,0)=\varphi(\xx)}[D\varphi(x)]\,e^{-\mathcal A[\varphi(x)]}\ \to\ \ket0\in\Ham\,.
\]
The problem with this representation is in the definition of the scalar products.
The simple minded definition
\begin{equation}\label{eq:naiveprod}
\braket{\Psi_2}{\Psi_1}:=\int[D\varphi(\xx)]\,\Psi_2^*[\varphi(\xx)]\,\Psi_1[\varphi(\xx)]
\end{equation}
has all required formal properties, but involves the functional integration which,
as we already know, requires regularization and renormalization. In this context
the problem of renormalization looks much harder than that in the functional
integral \eqref{eq:corrPhi} because the functionals $\Psi[\varphi(\xx)]$ do not have
(quasi) local structure, and generally are far more complicated than
$e^{-\mathcal A}$.

Instead, one defines the scalar product of the states \eqref{eq:states} as
\begin{equation}\label{eq:OSprod}
\bigl\langle\varphi(y_1)\cdots\varphi(y_N)\big|\varphi(x_1)\cdots\varphi(x_M)\bigr\rangle
:=\bigl\langle\varphi(\bar y_1)\cdots\varphi(\bar y_N)\,\varphi(x_1)\cdots\varphi(x_M)\bigr\rangle
\end{equation}
Here the ``mirror image'' points $\bar x$ are defined as follows:
\begin{equation}\label{eq:reflection}
\text{``Reflection'':}\qquad x=(\xx,\tau)\ \leftrightarrow\ \bar x=(\xx,-\tau)
\end{equation}
Then, if all the points $x_1,\dots,x_M,y_1,\dots,y_N$ in \eqref{eq:OSprod} are in
$\R^d_-$, the ``mirror'' points $\bar y_1,\dots,\bar y_N$ are all in $\R^d_+$. The
configuration of the points in the correlation function in the r.h.s.\ of
\eqref{eq:OSprod} is exactly as in \eqref{eq:tauorder}, see figure~\ref{fig:slice}
above.

The most important property of the scalar product is its \emph{positivity},
\[
\braket\Psi\Psi>0
\]
for any $\ket\Psi\in\Ham$, $\ket\Psi\neq0$. For the formal representation
\eqref{eq:naiveprod} the positivity is obvious: Eq.~\eqref{eq:naiveprod} with
$\Psi_1=\Psi_2$ is explicitly positive. For the scalar product defined in terms of
the correlation functions this property is postulated as the \textbf{reflection
positivity}
\[
\sum_{N,M}\int_{x_i,y_k\in\R^d_-}\bigl\langle\varphi(\bar y_1)\cdots\varphi(\bar y_N)\,\varphi(x_1)\cdots\varphi(x_M)\bigr\rangle\,f_N^*(y_1,\dots,y_N)\,f_M(x_1,\dots,x_M)\ge0
\]
for any sequence of functions $f_N(x_1,\dots,x_N)$, which expresses positivity of the
norm of the state
\[
\ket\Psi=\sum_N\int_{x_i\in\R^d_-}f_N(x_1,\dots,x_N)\,\ket{\varphi(x_1)\cdots\varphi(x_N)}\,.
\]
More generally, we can associate states with general products of local fields
$O_a(x)$,
\begin{equation}\label{eq:Ostates}
O_{a_1}(x_1)\cdots O_{a_N}(x_N)\ \to\ \ket{O_{a_1}(x_1)\cdots O_{a_N}(x_N)}\in\Ham\,,\qquad x_1,\dots,x_N\in\R^d_-
\end{equation}
according to the same idea
\[
\Psi_{O_1(x_1)\cdots O_N(x_N)}[\varphi(\xx)]=\int_{\varphi(\xx,0)=\varphi(\xx)}[D\varphi(x)]\,O_1(x_1)\cdots O_N(x_N)\,e^{-\mathcal A[\varphi(x)]}
\]
with the scalar product related to their correlation functions similarly,
\[
\bigl\langle O_{b_1}(y_1)\cdots O_{b_M}(y_M)\big|O_{a_1}(x_1)\cdots O_{a_N}(x_N)\bigr\rangle
:=\bigl\langle\bar O_{b_1}(\bar y_1)\cdots\bar O_{b_M}(\bar y_M)\,O_{a_1}(x_1)\cdots O_{a_N}(x_N)\bigr\rangle
\]
where again $x_1,\dots,x_N;y_1,\dots,y_M\in\R^d_-$, and the points
$\bar y_i\in\R^d_+$ are the ``mirror images'' of $y_i$, as in
Eq.~\eqref{eq:reflection}. However, when the fields $O_i(x_i)$ are not all scalars,
the fields $\bar O_b$ are related to $O_b$ by certain ``reflection''. To get an idea
assume that the field $O(x)$ is real-valued vector with the components $O^\mu(x)$.
Then, if $\mu=1,2,\dots,d-1$ represent ``spatial'' components,
\[
\bar O^\mu(\bar x)=O^\mu(\bar x)\,,\qquad\bar O^\tau(\bar x)=-O^\tau(\bar x)\,,
\]
i.e.\ the sign of the ``temporal'' component $O^\tau$ gets flipped upon reflection.
For example,
\[
\overline{\partial_\tau\varphi}(x)=-\partial_\tau\varphi(x)
\]
so that
\[
\bigl\langle\partial_\nu\varphi(y)\big|\partial_\mu\varphi(x)\bigr\rangle=\pm\bigl\langle\partial_\nu\varphi(\bar y)\,\partial_\mu\varphi(x)\bigr\rangle
\]
where the sign $+$ applies when both $\mu$ and $\nu$ are ``spatial'' or both are
``temporal'', otherwise the sign is $-$. Similar rule applies to higher rank tensors.
For complex valued fields the ``reflection'' involves complex conjugation. The
correlation functions of composite fields $O_a$ obey then the suitably generalized
reflection positivity conditions, which follow from the positivity of the norms of
the states \eqref{eq:Ostates}.

As an elementary application, let us show that in a scale invariant (i.e.\ RG fixed
point) theory the mass dimensions $D_\alpha$ of all nontrivial fields
$\Phi_\alpha(x)$ are positive (we denote $\Phi_\alpha(x)$ the scalar fields that
diagonalize the matrix of anomalous dimensions, see Section~18). From the scale
invariance
\[
\bigl\langle\Phi_\alpha(x)\Phi_\alpha(y)\bigr\rangle=\frac{C_\alpha}{|x-y|^{2D_\alpha}}
\]
where $|x-y|$ is the Euclidean distance between the points $x$ and $y$. Let us take
$\tau>0$, so that $x=(\xx,-\tau/2)\in\R^d_-$. Set $y=\bar x=(\xx,\tau/2)$. We have
\[
\braket{\Phi_\alpha(x)}{\Phi_\alpha(x)}=\frac{C_\alpha}{(\tau^2)^{D_\alpha}}\,.
\]
The positivity requires $C_\alpha>0$; one can set $C_\alpha=1$ by adjusting
normalization of $\Phi_\alpha$. Consider the norm
\begin{equation}\label{eq:normdPhi}
\Bigl\langle\frac\partial{\partial x^1}\Phi_\alpha(\xx,-\tau/2)\Big|\frac\partial{\partial x^1}\Phi_\alpha(\xx,-\tau/2)\Bigr\rangle
\end{equation}
where $x^1$ is one of the spatial coordinates in $\xx=(x^1,x^2,\dots,x^{d-1})$. This
norm is expressed as the second derivative of the correlation function
\[
\frac{\partial^2}{\partial x^1\,\partial y^1}\bigl\langle\Phi_\alpha(\xx,\tau/2)\,\Phi_\alpha(\yy,-\tau/2)\bigr\rangle
\]
evaluated at $\yy=\xx$. Elementary calculation yields for \eqref{eq:normdPhi}
\[
\frac{2D_\alpha}{(\tau^2)^{D_\alpha+1}}
\]
which must be positive, hence $D_\alpha\ge0$.

The case $D=0$ is instructive. If the field $\Phi$ has mass dimension $D=0$ it
follows that
\[
\Bigl|\frac\partial{\partial x^1}\Phi(x)\Bigr\rangle=0\,.
\]
It is then easy to prove that any correlation function involving the derivative of
$\Phi(x)$ vanishes
\[
\bigl\langle\partial_\mu\Phi(x)\,O_1(x_1)\cdots O_N(x_N)\bigr\rangle=0\,,\qquad\Rightarrow\qquad\partial_\mu\Phi(x)\simeq0\,.
\]
Then one can show (using cluster factorization, see below) that
\[
\bigl\langle\Phi(x)\,O_1(x_1)\cdots O_N(x_N)\bigr\rangle=\bigl\langle O_1(x_1)\cdots O_N(x_N)\bigr\rangle
\]
i.e.\ the only field with zero mass dimension is the identity operator.

% Source pp. 264-272 (Lectures19-25.pdf).
\section{Energy-Momentum Operators}

The expression \eqref{eq:Pcomm} can be now promoted to the definition of the
operators $\hat P^\mu$ acting in the space $\Ham$. Specifically, for any
$x_1,\dots,x_N,y_1,\dots,y_M$ in $\R^d_-$ define
\begin{equation}\label{eq:Pmatrix}
\bigl\langle\varphi(y_1)\cdots\varphi(y_M)\big|\hat P^\mu\big|\varphi(x_1)\cdots\varphi(x_N)\bigr\rangle
=\bigl\langle\varphi(\bar y_1)\cdots\varphi(\bar y_M)\,P^\mu\,\varphi(x_1)\cdots\varphi(x_N)\bigr\rangle
\end{equation}
where the correlation function in the r.h.s.\ involves the integral of the
energy-momentum tensor over the equal-time slice $\tau=0$,
\begin{equation}\label{eq:Pdef}
P^\mu(\tau):=-i\int_\Sigma \dd\Sigma_\nu(x)\,T^{\mu\nu}(x)=-i\int_{\Sigma(\tau=0)}\dd^{d-1}\xx\,T^{\mu\tau}(\xx,\tau=0)\,,
\end{equation}
It is straightforward to check, using the definition of the inner product, that the
``spatial'' components $\PPv=(P^1,\dots,P^{d-1})$ are Hermitian, while $P^\tau$ is
anti-Hermitian. We therefore define $d$ Hermitian operators $(\hat{\PPv},\hat H)$ as
\begin{align}
\hat P^k&:=-i\int_{\Sigma(\tau=0)}\dd^{d-1}\xx\,T^{k\tau}(\xx,\tau)\,,\label{eq:Pk}\\
\hat H&=-\int_{\Sigma(\tau=0)}\dd^{d-1}\xx\,T^{\tau\tau}(\xx,\tau)\,.\label{eq:Hdef}
\end{align}
Obviously, $\hat H$ coincides with $\hat P^0:=-i\hat P^\tau$, the real-time
component of $\hat P^\mu$.

A remark is in order. If we are dealing with the field theory in the infinite
space-time $\R^d$, the matrix elements of these operators diverge because of the
contribution of generally nonzero expectation value
\[
\bigl\langle T^{\mu\nu}(x)\bigr\rangle=g^{\mu\nu}\,\varepsilon
\]
where $g^{\mu\nu}=\delta^{\mu\nu}$ is the Euclidean metric tensor, and
$\varepsilon$ is the vacuum energy density. It is convenient to get rid of this
divergence by defining reduced energy-momentum tensor
\[
T^{\mu\nu}(x)\to T^{\mu\nu}(x)-g^{\mu\nu}\,\varepsilon
\]
which has zero expectation value. This can be interpreted as adding the
``cosmological term'' to the action
\[
\mathcal A\to\mathcal A+\varepsilon\int \dd^dx\,\sqrt{g(x)}
\]
For the fixed metric background such term has no effect on the content of the
theory, but brings in the above subtraction when $T^{\mu\nu}(x)$ is defined as the
response to the variation of the background metric,
\[
\delta\mathcal A=-\frac12\int \dd^dx\,\sqrt g\,T^{\mu\nu}(x)\,\delta g_{\mu\nu}(x)\,.
\]
In what follows we adopt this convention, assuming that $T^{\mu\nu}(x)$ is chosen in
such a way that
\begin{equation}\label{eq:Tvev0}
\bigl\langle T^{\mu\nu}(x)\bigr\rangle=0\,.
\end{equation}

With this convention, it is easy to prove that the vacuum $\ket0$ is the eigenstate
of $\hat P^\mu$ with the eigenvalue 0, i.e.
\begin{equation}\label{eq:vacP}
\hat H\ket0=0\,,\qquad\hat{\PPv}\ket0=0\,.
\end{equation}
Indeed, consider the element
\begin{equation}\label{eq:vacelement}
\bigl\langle O_1(x_1)\cdots O_N(x_N)\big|\hat P^\mu\big|0\bigr\rangle=\bigl\langle O_1(\bar x_1)\cdots O_N(\bar x_N)\,P^\mu\bigr\rangle
\end{equation}
where $P^\mu$ is the surface integral \eqref{eq:Pdef}. As all insertions
$\bar x_1,\dots,\bar x_N$ are in $\R^d_+$, the hypersurface $\Sigma$ involved in the
definition \eqref{eq:Pmatrix} can be deformed away to the domain $\tau\to-\infty$,
where the integral vanishes. As the result, all matrix elements
\eqref{eq:vacelement} are equal to zero. Eq.~\eqref{eq:vacP} follows.

Similar calculation shows that the operators $P_\mu$ act on the states
$\ket{O_1(x_1)\cdots O_N(x_N)}$ as follows
\[
\hat P_\mu\ket{O_1(x_1)\cdots O_N(x_N)}=i\left[\sum_{k=1}^N\frac\partial{\partial x_k^\mu}\right]\ket{O_1(x_1)\cdots O_N(x_N)}
\]
Indeed, this time deforming $\Sigma$ to the $\tau\to-\infty$ domain leaves behind
integrals around the insertion points $x_k$, which, according to
\[
\partial^\mu O(x)=\int_{\partial D_x}\dd\Sigma_\nu(y)\,T^{\mu\nu}(y)\,O(x)
\]
leaves behind the derivatives in the r.h.s. As usual, the infinitesimal
translations generated by $\hat P_\mu$ can be exponentiated to finite translations.
In terms of the operators \eqref{eq:Pk}, \eqref{eq:Hdef} we have
\begin{equation}\label{eq:finitetr}
\begin{aligned}
e^{-i\hat{\PPv}\cdot\mathbf X}\ket{O_1(\xx_1,\tau_1)\cdots O_N(\xx_N,\tau_N)}&=\ket{O_1(\xx_1+\mathbf X,\tau_1)\cdots O_N(\xx_N+\mathbf X,\tau_N)}\\
e^{-\hat HL}\ket{O_1(\xx_1,\tau_1)\cdots O_N(\xx_N,\tau_N)}&=\ket{O_1(\xx_1,\tau_1-L)\cdots O_N(\xx_N,\tau_N-L)}
\end{aligned}
\end{equation}

\subsection*{Stationary States. Energy Positivity}

The operators $\hat P^\mu$ commute, and thus can be simultaneously diagonalized in
$\Ham$. It is convenient to discuss the spectrum of the Hermitian operators
$\hat P^\mu$, $\mu=0,\dots,d-1$ with $\hat P^0:=-i\hat P^\tau=\hat H$, which, of
course, coincide with the conventional Energy-Momentum operators in the Minkowski
space-time. Let $\ket\alpha$ be an eigenstate
\[
\hat P^\mu\ket\alpha=p^\mu_\alpha\ket\alpha\,.
\]
We will make few assumptions about the spectrum of these operators. One is about
the spectrum of $\hat H$. We assume that the eigenvalues $E_\alpha$ of $\hat H$
satisfy $E_\alpha\ge0$, with the only zero eigenvalue $E_0=0$ being associated with
the vacuum state, Eq.~\eqref{eq:vacP}. This assumption implies
\begin{equation}\label{eq:projector}
\lim_{L\to\infty}e^{-L\hat H}=\ket0\bra0
\end{equation}
There are several simple consequences from this condition alone. One is the
so-called cluster factorization property of the correlation functions, which states
that if a collection of $N$ points consists of two groups with very large separation
between the groups, correlations between the fields in the different groups tend to
zero, that is
\begin{multline}\label{eq:cluster}
\lim_{L\to\infty}\bigl\langle O_1(x_1)\cdots O_M(x_M)\,O_{M+1}(x_{M+1}+L)\cdots O_N(x_N+L)\bigr\rangle\\
=\bigl\langle O_1(x_1)\cdots O_M(x_M)\bigr\rangle\,\bigl\langle O_{M+1}(x_{M+1})\cdots O_N(x_N)\bigr\rangle
\end{multline}
To prove, take the direction of $L$ to be the ``time'', then the l.h.s.\ can be
represented as
\[
\bigl\langle O_1(x_1)\cdots O_M(x_M)\big|e^{-L\hat H}\big|O_{M+1}(x_{M+1})\cdots O_N(x_N)\bigr\rangle
\]
and \eqref{eq:cluster} follows from \eqref{eq:projector}.

Less obvious consequence concerns analytic properties of the correlation functions
\[
\bigl\langle O_1(\xx_1,\tau_1)\cdots O_N(\xx_N,\tau_N)\bigr\rangle
\]
at complex values of $\tau_1,\dots,\tau_N$. We have briefly studied this question at
the beginning of the course, for the case of the Klein-Gordon propagator
(Section~4). In fact, our arguments there depend only on the condition of the
positivity of the energy spectrum, which was explicit in the KG theory. Therefore,
the analysis can be repeated word for word in general case. Let me restate the
result for the two point function
\[
\bigl\langle O_1(\xx,\tau)\,O_2(\mathbf0,0)\bigr\rangle\,.
\]
In the complex plane of $t=-i\tau$ we have the picture of
figure~\ref{fig:tcontours2}.
\begin{figure}[ht]
\centering\figfile{tcontours2}
\caption{The two-point function in the complex $t$-plane. Above the light-cone
cuts (along $C_+$) it gives $\bra0\hat O_2(\mathbf0,0)\,\hat O_1(\xx,t)\ket0$,
below them (along $C_-$) it gives $\bra0\hat O_1(\xx,t)\,\hat O_2(\mathbf0,0)\ket0$.}
\label{fig:tcontours2}
\end{figure}
Here
\[
\hat O(\xx,t)=e^{-i\PPv\cdot\xx+iHt}\,\hat O(\mathbf0,0)\,e^{i\PPv\cdot\xx-iHt}
\]
are conventional Heisenberg field operators. The discontinuity across the branch
cuts represents the commutators which lie within the future and past light cones
(``local commutativity'').

\subsection*{Spectrum of Energy-Momentum}

Some properties of the spectrum of the operator $\hat P^\mu=(\hat H,\hat{\PPv})$ can
be established on general grounds. Let $\mathcal P$ be the spectrum of this
operator, i.e.\ the set of its eigenvalues
\[
p^\mu\in\mathcal P\ \Rightarrow\ \exists\ \ket\alpha\in\Ham\ \text{such that}\ p^\mu=p^\mu_\alpha:\quad\hat P^\mu\ket\alpha=p^\mu_\alpha\ket\alpha\,.
\]
Lorentz invariance suggests that
\[
p^\mu\in\mathcal P\ \Rightarrow\ \Lambda^\mu{}_\nu\,p^\nu\in\mathcal P\,,
\]
where $\Lambda$ is a Lorentz transformation. This means that any $p^\mu\in\mathcal P$
such that $p^\mu p_\mu=-m^2$ enters $\mathcal P$ together with the entire hyperboloid
\[
\mathcal P_{m^2}=\bigl\{p^\mu:\ -p^\mu p_\mu=(p^0)^2-\pp^2=m^2\bigr\}\,.
\]
It follows then from the positivity of the energy that $\mathcal P$ lies entirely
within the future light cone $p^\mu p_\mu\le0$, $p^0\ge0$.

Another important property of $\mathcal P$ is its additivity, i.e.
\[
p^\mu_\alpha,\,p^\mu_\beta\in\mathcal P\ \Rightarrow\ p^\mu_\alpha+p^\mu_\beta\in\mathcal P\,.
\]
Intuitively, if we think of the state $\ket\alpha$ carrying the momentum
$p^\mu_\alpha$ as a local excitation over the ground state, it is possible to add
another excitation, with the momentum $p^\mu_\beta$ arbitrarily far away from the
first one. At large separations the excitations do not interact, and their energies
and momenta add up. Formal derivation using the cluster property will be considered
later.

The additivity implies that if $\mathcal P$ contains $\mathcal P_{m^2}$, it also
contains all $p^\mu$ of the form
\[
p^\mu=p_1^\mu+\dots+p_N^\mu\,;\qquad p_i^\mu\in\mathcal P_{m^2}\,.
\]
The totality of such vectors $p^\mu$ will be denoted $\bar{\mathcal P}_{m^2}$,
\[
\bar{\mathcal P}_{m^2}=\bigcup_N\ \bigcup_{p_1,\dots,p_N}\bigl\{p^\mu=p_1^\mu+\dots+p_N^\mu;\ p_i^\mu\in\mathcal P_{m^2}\bigr\}\,.
\]
The set $\bar{\mathcal P}_{m^2}$ looks like figure~\ref{fig:spectrum} of
Section~4.

Note that $\bar{\mathcal P}_{m^2}$ is exactly the spectrum of $\hat P^\mu$ in the
Klein-Gordon field theory, where we have found that
\[
\hat P^\mu=\int\frac{\dd^3\pp}{(2\pi)^3}\,p^\mu\,a^\dagger_\pp\,a_\pp\,,
\]
where $p^\mu=(\omega_\pp,\pp)$, $\omega_\pp=\sqrt{\pp^2+m^2}$, and $a_\pp$,
$a^\dagger_\pp$ are particle creation and annihilation operators. There these
states are interpreted as $N$-particle states of free bosons.

The spectrum $\bar{\mathcal P}_{m^2}$ in some sense is the minimal Lorentz-invariant
set compatible with the above additivity condition. If this spectrum appears in an
interacting QFT we say that the theory contains just one massive particle of the
mass $m$. The states associated with the hyperboloid $\mathcal P_{m^2}$ are the
one-particle states, and the states with $-p^\mu p_\mu\ge4m^2$ are interpreted as the
scattering states of two or more interacting particles.

In general case the spectrum can be more complicated (even in a scalar QFT like we
consider here). It can involve more than one kind of stable particles, and hence
would contain several hyperboloids, say $\mathcal P_{m_1^2}$ and
$\mathcal P_{m_2^2}$, and correspondingly its continuous part would look more
complicated, as in figure~\ref{fig:spectrum2}.
\begin{figure}[ht]
\centering\figfile{spectrum2}
\caption{Energy-momentum spectrum of a theory with two stable particles of masses
$m_1<m_2$: two mass hyperboloids, and a continuum of scattering states starting at
$\min(2m_1,m_1+m_2)=2m_1$, with further thresholds at $m_1+m_2$, $2m_2$, $3m_1$,
\ldots{} inside it.}
\label{fig:spectrum2}
\end{figure}

In any case, if we form the operator
\[
-\hat P^\mu\hat P_\mu=\hat H^2-\hat{\PPv}^2\,,
\]
the discrete components of its spectrum are associated with the one-particle
states, and the continuous spectrum incorporates all the multi-particle scattering
states.

``Asymptotic completeness'': $\mathcal P=\cup_a\bar{\mathcal P}_{m_a^2}$.

Lorentz invariance of the spectrum has certain implications for the general
structure of the correlation functions in QFT. As the simplest example, consider the
two-point correlation function
\[
\bigl\langle\varphi(\xx,\tau)\,\varphi(\mathbf0,0)\bigr\rangle=\bra0T_\tau\bigl(\hat\varphi(\xx,\tau)\,\hat\varphi(\mathbf0,0)\bigr)\ket0\,,
\]
where $(\xx,\tau)$ are the Euclidean space-time coordinates. Assuming $\tau>0$ we can
write it as
\[
\sum_\alpha\bra0\hat\varphi(\xx,\tau)\ket\alpha\bra\alpha\hat\varphi(\mathbf0,0)\ket0\,.
\]
Using the fact that
\[
\hat\varphi(\xx,\tau)=e^{\hat H\tau-i\hat{\PPv}\cdot\xx}\,\hat\varphi(\mathbf0,0)\,e^{-\hat H\tau+i\hat{\PPv}\cdot\xx}\,.
\]
we can write the above sum as
\[
\sum_\alpha\bra0\hat\varphi(\mathbf0,0)\ket\alpha\bra\alpha\hat\varphi(\mathbf0,0)\ket0\,e^{-E_\alpha\tau+i\PPv_\alpha\cdot\xx}\,,
\]
where $E_\alpha$, $\PPv_\alpha$ are the energy and momentum of the state
$\ket\alpha$. In fact, we have continuously many states, so that
\[
\sum_\alpha\to\int \dd\mu(\alpha)\,.
\]
where $\alpha$ denotes complete set of quantum numbers characterizing the state.
Define
\begin{equation}\label{eq:rhodef}
\rho(q)=\int \dd\mu(\alpha)\,\bigl|\bra0\hat\varphi(0)\ket\alpha\bigr|^2\,(2\pi)^4\,\delta^{(4)}(q^\mu-p^\mu_\alpha)
\end{equation}
Then, assuming $\tau>0$, we have
\begin{equation}\label{eq:spectralrep}
\bigl\langle\varphi(\xx,\tau)\,\varphi(\mathbf0,0)\bigr\rangle=\int \dd\mu(\alpha)\,\bra0\hat\varphi(\xx,\tau)\,\hat\varphi(\mathbf0,0)\ket0
=\int\frac{\dd^4q}{(2\pi)^4}\,\rho(q)\,e^{-q^0\tau+i\qq\cdot\xx}
\end{equation}
Evidently,
\[
\rho(q)=0\quad\text{if}\quad q^\mu\notin\mathcal P\,,\qquad\rho(q)\ge0\quad\text{if}\quad q^\mu\in\mathcal P\,,
\]
so that $\rho(q)\ge0$ everywhere.

$\rho(q)$ vanishes if $q^0<0$ (Hamiltonian is assumed to be positive definite), and
at $q^0>0$ it depends on $-q^2=(q^0)^2-\qq^2$ (Lorentz invariance); I will write
\[
\rho=\theta(q^0)\,\rho(-q^2)\,.
\]
Then, for $\tau>0$ we have
\begin{align*}
\bigl\langle\varphi(\xx,\tau)\,\varphi(\mathbf0,0)\bigr\rangle
&=\int_0^\infty \dd\mu^2\,\rho(\mu^2)\int\frac{\dd^4q}{(2\pi)^4}\,\delta\bigl((q^0)^2-\qq^2-\mu^2\bigr)\,\theta(q^0)\,e^{-q^0\tau+i\qq\cdot\xx}\\
&=\int_0^\infty\frac{\dd\mu^2}{2\pi}\,\rho(\mu^2)\int\frac{\dd^3\qq}{(2\pi)^3}\,\frac1{2\omega_\qq}\,e^{-\omega_\qq\tau+i\qq\cdot\xx}\,,
\end{align*}
where
\[
\omega_\qq=\sqrt{\qq^2+\mu^2}\,.
\]
Proceeding similarly, the case $\tau<0$ we obtain the same expression with
\[
e^{-\omega_\qq\tau}\to e^{\omega_\qq\tau}
\]
Both cases $\tau>0$ and $\tau<0$ can be put together by using
\[
\frac1{2\omega_q}\,e^{-\omega_q|\tau|}=\int\frac{\dd q_4}{2\pi}\,\frac{e^{iq_4\tau}}{q_4^2+\qq^2+\mu^2}
\]
Finally
\begin{equation}\label{eq:KL}
\bigl\langle\varphi(\xx,\tau)\,\varphi(\mathbf0,0)\bigr\rangle=\int_0^\infty\frac{\dd\mu^2}{2\pi}\,\rho(\mu^2)\,D(|x|,\mu^2)\,,
\end{equation}
where $|x|=\sqrt{\tau^2+\xx^2}$, and $D(|x|,\mu^2)$ is the Klein-Gordon propagator
with the mass $\mu$,
\[
D(|x|,\mu^2)=\int\frac{\dd^4k}{(2\pi)^4}\,\frac{e^{i\kk\cdot\xx+ik_4\tau}}{k_4^2+\kk^2+\mu^2}\,.
\]
Eq.~\eqref{eq:KL} is known as the \emph{K\"all\'en-Lehmann representation}.

For $\bar{\mathcal P}_{m^2}$
\[
\rho(\mu^2)=(2\pi)\,Z\,\delta(\mu^2-m^2)+\tilde\rho(\mu^2)
\]
where $\tilde\rho(\mu^2)=0$ at $\mu^2<4m^2$,
\begin{equation}\label{eq:Zres}
Z^{1/2}=\bra0\hat\varphi(0)\ket\pp
\end{equation}
\begin{equation}\label{eq:KLmom}
\tilde W^{(2)}(p)=\frac Z{p^2+m^2}+\int_{4m^2}^\infty\frac{\tilde\rho(\mu^2)}{p^2+\mu^2}\,\frac{\dd\mu^2}{2\pi}
\end{equation}
which is another form of the K\"all\'en-Lehmann representation.

Note the pole at $p^2=-m^2$ explicit in this formula, appearing entirely from the
contributions of the one-particle states in the intermediate-state decomposition.
The residue $Z$ is positive (as it is the square of the absolute value of the matrix
element \eqref{eq:Zres}), but its value depends on the normalization of the field
$\varphi$. That normalization was subject to one of the normalization conditions in
our calculations in the perturbation theory. Thus, under the normalization
condition
\begin{equation}\label{eq:SIrecall}
\tilde\Gamma^{(2)}(p^2)=p^2+m^2+O\bigl((p^2+m^2)^2\bigr)
\end{equation}
of the scheme SI in the $\varphi^4$ theory (see Eqs.~\eqref{eq:SIa},
\eqref{eq:SIW}) we just choose $Z=1$. Let us also note that $\tilde\rho(\mu^2)$ is
non-negative at all $\mu^2$, as the consequence of the positivity of the norms in
the Hilbert space, see Eq.~\eqref{eq:rhodef}.

At $|x|\to\infty$ we have
\[
D(|x|,\mu^2)\simeq\frac{\mu^{1/2}}{|x|^{3/2}}\,e^{-\mu|x|}\,,
\]
and the form \eqref{eq:KL} shows that at large $|x|$ the asymptotic of
$\langle\varphi(x)\varphi(0)\rangle$ is always dominated by the one-particle
contribution, provided the matrix element \eqref{eq:Zres} is not zero.



\end{document}
